\documentclass[booklanguage=brazilian]{byrnebook}
%\usepackage{lua-visual-debug}

% Define missing angle style functions and constants for SOLID_SECTOR, ARC_SECTOR, etc.
% Note: SOLID_SECTOR, ARC_SECTOR, etc. are used as string literals in the code,
% but MetaPost needs the corresponding functions to exist.
\def\mpPre{u:=1cm;
numeric SOLID_LINE; SOLID_LINE := 0;
numeric DASHED_LINE; DASHED_LINE := 1;
numeric REGULAR_WIDTH; REGULAR_WIDTH := 0;
numeric THIN_WIDTH; THIN_WIDTH := 1;
vardef byAngleMSOLID_SECTOR (expr angleArc, angleColor)(suffix angleOptionalColors) =
    save p;
    path p;
    p := angleArc scaled (angleScale*angleSize);
    image(
        fill ((0, 0)--p--cycle) withcolor angleColor;
    )
enddef;
vardef byAngleMARC_SECTOR (expr angleArc, angleColor)(suffix angleOptionalColors) =
    save p;
    path p;
    p := angleArc scaled ((angleScale*angleSize) - lineWidthThin);
    image(
        draw p withpen pencircle scaled lineWidthThin withcolor angleColor;
    )
enddef;
vardef byAngleMDASHED_ARC_SECTOR (expr angleArc, angleColor)(suffix angleOptionalColors) =
    save p;
    path p;
    p := angleArc scaled ((angleScale*angleSize) - lineWidthThin);
    image(
        draw p withpen pencircle scaled lineWidthThin withcolor angleColor dashed byDashPattern(arclength(p), 1);
    )
enddef;
}

\begin{document}
\newcommand{\triangleABC}{\ensuremath{\triangle ABC}}
\newcommand{\triangleABD}{\ensuremath{\triangle ABD}}
\newcommand{\triangleABG}{\ensuremath{\triangle ABG}}
\newcommand{\triangleADE}{\ensuremath{\triangle ADE}}
\newcommand{\triangleAEF}{\ensuremath{\triangle AEF}}
\newcommand{\triangleBDE}{\ensuremath{\triangle BDE}}
\newcommand{\triangleCDE}{\ensuremath{\triangle CDE}}
\newcommand{\triangleDEF}{\ensuremath{\triangle DEF}}
\newcommand{\triangleFEB}{\ensuremath{\triangle FEB}}
\newcommand{\triangleKGH}{\ensuremath{\triangle KGH}}

\thispagestyle{empty}

\begin{center}

\Large \uppercase{Os primeiros seis livros dos}

\LARGE \uppercase{Elementos de Euclides}

\vskip 0.5\baselineskip

\normalsize \uppercase{nos quais são utilizados diagramas e símbolos coloridos em vez de letras para maior facilidade dos aprendizes}

\vskip 0.75\baselineskip

\Large \uppercase{Por Oliver Byrne}

%{\uppercase{SURVEYOR OF HER MAJESTY'S SETTLEMENTS IN THE FALKLAND ISLANDS AND AUTHOR OF NUMEROUS MATHEMATICAL WORKS}}

\defineNewPicture{
textLabels := false;
scaleFactor := 7/6;
angleScale := 4/3;
pair A, B, C, D, E, F, G, H, I, J, K, L, M, d[];
A := (0, 0);
B := A shifted (-7/10u, -8/7u);
C = whatever[A, A shifted ((A-B) rotated 90)] = whatever[B, B shifted dir(0)];
d1 := (B-A) rotated -90;
D := A shifted d1;
E := B shifted d1;
d2 := (A-C) rotated -90;
F := C shifted d2;
G := A shifted d2;
d3 := (C-B) rotated -90;
H := B shifted d3;
I := C shifted d3;
J = whatever[A, A shifted dir(90)];
J = whatever[B, C];
K = whatever[A, A shifted dir(90)];
K = whatever[H, I];
L = whatever[B, F];
L = whatever[A, C];
M = whatever[A, I];
M = whatever[B, C];
draw byPolygon(A,B,E,D)(byblack);
draw byPolygon(L,A,G,F)(byred);
draw byPolygon(C,L,F)(byred);
draw byPolygon(J,M,I,K)(byyellow);
draw byPolygon (M,C,I)(byyellow);
draw byPolygon(B,J,K,H)(byblue);
byAngleDefine.FCA(F, C, A, byyellow, SOLID_SECTOR);
byAngleDefine.BCI(B, C, I, byblue, SOLID_SECTOR);
byAngleDefine.ACB(A, C, B, byblack, SOLID_SECTOR);
draw byNamedAngleResized(FCA, BCI, ACB);
draw byLineFull(A, K, byred, 1, 0)(I, I, 1, 1, -1);
draw byLineFull(B, F, byblack, 0, 0)(G, G, 1, 1, -1);
draw byLineFull(A, I, byblack, 0, 0)(K, K, 1, 1, 1);
byLineDefine(C, F, byblue, DASHED_LINE, REGULAR_WIDTH);
byLineDefine(C, I, byblack, DASHED_LINE, REGULAR_WIDTH);
draw byNamedLineSeq(0)(CF,CI);
byLineDefine(A, B, byyellow, SOLID_LINE, REGULAR_WIDTH);
byLineDefine(B, C, byred, SOLID_LINE, REGULAR_WIDTH);
byLineDefine(C, A, byblue, SOLID_LINE, REGULAR_WIDTH);
draw byNamedLineSeq(1)(AB,BC,CA);
byLineDefine.CAb(C, A, byblack, SOLID_LINE, REGULAR_WIDTH);
byLineStylize (M, M, 1, 0, -1) (CAb);
byLineDefine.AMb(A, M, byblack, SOLID_LINE, REGULAR_WIDTH);
byLineStylize (C, C, 0, 1, -1) (AMb);
byLineDefine.BCb(B, C, byblack, SOLID_LINE, REGULAR_WIDTH);
byLineStylize (L, L, 0, 1, -1) (BCb);
byLineDefine.BLb(L, B, byblack, SOLID_LINE, REGULAR_WIDTH);
byLineStylize (C, C, 1, 0, -1) (BLb);
draw byLabelsOnPolygon(B, E, D, A, G, F, C, I, K, H)(ALL_LABELS, -1);
draw byLabelsOnPolygon(A, J, C)(OMIT_FIRST_LABEL+OMIT_LAST_LABEL, 1);
}

\vfill\vfill

~\hfill\drawCurrentPicture\hfill~

\vfill\vfill\vfill

\large github.com/jemmybutton
\vskip 0.25\baselineskip

\Large 2025 ed.\,0.8-latex-alpha
\vskip \baselineskip

\ccbysa
\vskip 0.25\baselineskip

\footnotesize Esta versão de \emph{Os primeiros seis livros dos Elementos de Euclides} de Oliver Byrne foi feita por Slyusarev Sergey e é distribuída sob a licença CC-BY-SA~4.0

\vskip -\baselineskip

\normalsize
\end{center}
\pagebreak

\part*{Introdução}

\charspacing{-2}{\regularLettrine{A}{s} artes e ciências tornaram-se tão extensas, que facilitar sua aquisição é de tanta importância quanto estender seus limites. A ilustração, se não reduz o tempo de estudo, ao menos o torna mais agradável. Esta Obra tem um objetivo maior do que mera ilustração; não introduzimos cores com o propósito de entretenimento, ou para divertir \emph{por certas combinações de tom e forma}, % don't know where this is from
mas para auxiliar a mente em suas pesquisas pela verdade, aumentar as facilidades de introdução e difundir conhecimento permanente. Se quiséssemos autoridades para provar a importância e utilidade da geometria, poderíamos citar todo filósofo desde os dias de Platão. Entre os gregos, na antiguidade, assim como na escola de Pestalozzi e outras em tempos recentes, a geometria foi adotada como a melhor ginástica da mente. De fato, os Elementos de Euclides tornaram-se, por consentimento comum, a base da ciência matemática em todo o globo civilizado. Mas isso não parecerá extraordinário, se considerarmos que esta ciência sublime não apenas é melhor calculada do que qualquer outra para despertar o espírito de investigação, elevar a mente e fortalecer as faculdades de raciocínio, mas também forma a melhor introdução à maioria das vocações úteis e importantes da vida humana. Aritmética, agrimensura, hidrostática, pneumática, óptica, astronomia física, \&c.\ são todas dependentes das proposições da geometria.}

\charspacing{-1}{Muito, no entanto, depende da primeira comunicação de qualquer ciência a um aprendiz, embora os melhores e mais fáceis métodos raramente sejam adotados. Proposições são colocadas diante de um estudante, que embora tenha compreensão suficiente, é informado apenas o suficiente sobre elas ao entrar no próprio limiar da ciência, como se lhe fosse dada uma predisposição muito desfavorável ao seu futuro estudo deste assunto delicioso; ou \enquote{as formalidades e parafernália do rigor são tão ostensivamente apresentadas, que quase escondem a realidade. Repetições infinitas e perplexas, que não conferem maior exatidão ao raciocínio, tornam as demonstrações envolvidas e obscuras, e ocultam da vista do estudante a consecução da evidência.} % quotation seems to be from The First Six Books of the Elements of Euclid by Dionysius Lardner https://books.google.ru/books?id=YnkAAAAAMAAJ
Assim, uma aversão é criada na mente do aluno, e um assunto tão calculado para melhorar os poderes de raciocínio e dar o hábito de pensamento próximo, é degradado por um curso seco e rígido de instrução em um exercício desinteressante da memória. Despertar a curiosidade e acordar os poderes indolentes e adormecidos de mentes mais jovens deve ser o objetivo de todo professor; mas onde exemplos de excelência estão faltando, as tentativas de alcançá-la são poucas, enquanto a eminência excita atenção e produz imitação. O objetivo desta Obra é introduzir um método de ensino de geometria, que tem sido muito aprovado por muitos homens científicos neste país, assim como na França e na América. O plano aqui adotado apela fortemente ao olho, o mais sensível e o mais abrangente de nossos órgãos externos, e sua preeminência para imprimir seu assunto na mente é apoiada pela máxima incontestável expressa nas bem conhecidas palavras de Horácio:—}

\begin{center}
\emph{Segnius irritant animos demissa per aurem\\
Quam quae sunt oculis subjecta fidelibus}\\
\vskip 0.5\baselineskip
Uma impressão mais fraca através do ouvido é feita,\\
Do que aquilo que é transmitido pelo olho fiel.
\end{center}

\charspacing{-1}{Toda linguagem consiste de sinais representativos, e aqueles sinais são os melhores que efetuam seus propósitos com a maior precisão e rapidez. Tais, para todos os propósitos comuns, são os sinais audíveis chamados palavras, que ainda são considerados como audíveis, sejam dirigidos imediatamente ao ouvido, ou através do meio de letras ao olho. Diagramas geométricos não são sinais, mas os materiais da ciência geométrica, cujo objeto é mostrar as quantidades relativas de suas partes por um processo de raciocínio chamado Demonstração. Este raciocínio tem sido geralmente conduzido por palavras, letras e diagramas pretos ou não coloridos; mas como o uso de símbolos, sinais e diagramas coloridos nas artes e ciências lineares torna o processo de raciocínio mais preciso e a obtenção mais expedita, eles foram neste caso correspondentemente adotados.}

Tal é a rapidez deste modo sedutor de comunicar conhecimento, que os Elementos de Euclides podem ser adquiridos em menos de um terço do tempo usualmente empregado, e a retenção pela memória é muito mais permanente; estes fatos foram verificados por numerosos experimentos feitos pelo inventor, e vários outros que adotaram seus planos. Os detalhes dos quais são poucos e óbvios; as letras anexadas a pontos, linhas ou outras partes de um diagrama são de fato apenas nomes arbitrários, e as representam na demonstração; em vez destas, as partes sendo diferentemente coloridas, são feitas para nomear a si mesmas, pois suas formas em cores correspondentes as representam na demonstração.

Para dar uma melhor ideia deste sistema, e das vantagens obtidas por sua adoção, tomemos um triângulo retângulo, e expressemos algumas de suas propriedades tanto por cores quanto pelo método geralmente empregado.

\defineNewPicture{
	pair A, B, C;
	B := (0, 0);
	A := B shifted (dir(-145)*3u);
	C = whatever[A, A shifted (1,0)] = whatever[B, B shifted dir(-145+90)];
		byAngleDefine.ABC(A, B, C, byyellow, SOLID_SECTOR);
		byAngleDefine.BCA(B, C, A, byblue, SOLID_SECTOR);
		byAngleDefine.CAB(C, A, B, byred, SOLID_SECTOR);
		draw byNamedAngleResized(ABC, BCA, CAB);
		byLineDefine(A, B, byblue, SOLID_LINE, REGULAR_WIDTH);
		byLineDefine(B, C, byred, SOLID_LINE, REGULAR_WIDTH);
		byLineDefine(C, A, byyellow, SOLID_LINE, REGULAR_WIDTH);
		draw byNamedLineSeq(0)(AB,BC,CA);
		label.top(btex B etex, B);
		label.rt(btex C etex, C);
		label.lft(btex A etex, A);
	angleScale := 4/5;
}

\begin{center}
\drawCurrentPictureInMargin \emph{Algumas das propriedades do triângulo retângulo ABC, expressas pelo método geralmente empregado:}
\end{center}

\vskip 0.5\baselineskip

\begin{enumerate}
\item O ângulo BAC, juntamente com os ângulos BCA e ABC são iguais a dois ângulos retos, ou duas vezes o ângulo ABC.
\item O ângulo CAB adicionado ao ângulo ACB será igual ao ângulo ABC.
\item O ângulo ABC é maior do que qualquer dos ângulos BAC ou BCA.
\item O ângulo BCA ou o ângulo CAB é menor do que o ângulo ABC.
\item Se do ângulo ABC for tomado o ângulo BAC, o resto será igual ao ângulo ACB.
\item O quadrado de AC é igual à soma dos quadrados de AB e BC.
\end{enumerate}

\vskip 0.5\baselineskip

\begin{center}
\emph{As mesmas propriedades expressas colorindo as diferentes partes:}
\end{center}

\vskip 0.5\baselineskip

\begin{enumerate}
\item $\drawAngle{A} + \drawAngle{B} + \drawAngle{C} = 2 \drawAngle{B} = \drawTwoRightAngles$. \\ Isto é, o ângulo vermelho adicionado ao ângulo amarelo adicionado ao ângulo azul, igual a duas vezes o ângulo amarelo, igual a dois ângulos retos.
\item $\drawAngle{A} + \drawAngle{C} = \drawAngle{B}$. \\ Ou em palavras, o ângulo vermelho adicionado ao ângulo azul, igual ao ângulo amarelo.
\item $\drawAngle{B} > \drawAngle{A} \mbox{ ou } > \drawAngle{C}$. \\ O ângulo amarelo é maior do que qualquer dos ângulos vermelho ou azul.
\item $\drawAngle{A} \mbox{ ou } \drawAngle{C} < \drawAngle{B}$. \\ Qualquer dos ângulos vermelho ou azul é menor do que o ângulo amarelo.
\item $\drawAngle{B} - \drawAngle{C} = \drawAngle{A}$. \\ Em outros termos, o ângulo amarelo diminuído do ângulo azul igual ao ângulo vermelho. %\mbox{ minus }
\item $\drawUnitLine{CA}^2 = \drawUnitLine{AB}^2 + \drawUnitLine{BC}^2$. \\ Isto é, o quadrado da linha amarela é igual à soma dos quadrados das linhas azul e vermelha.
\end{enumerate}

Em demonstrações orais ganhamos com as cores esta importante vantagem, o olho e o ouvido podem ser abordados ao mesmo tempo, de modo que para ensinar geometria, e outras artes e ciências lineares, em classes, o sistema é o melhor já proposto, isso é aparente pelos exemplos dados.

\charspacing{-2}{Donde é evidente que uma referência do texto ao diagrama é mais rápida e segura, dando as formas e cores das partes, ou nomeando as partes e suas cores, do que nomeando as partes e letras no diagrama. Além da simplicidade superior, este sistema é igualmente notável pela concentração, e exclui totalmente a prática prejudicial embora prevalente de permitir que o estudante cometa a demonstração à memória; até que a razão, o fato e a prova façam apenas impressões do entendimento.}

Novamente, ao dar palestras sobre os princípios ou propriedades de figuras, se mencionarmos a cor da parte ou partes referidas, como ao dizer, o ângulo vermelho, a linha azul, ou linhas, \&c, a parte ou partes assim nomeadas serão imediatamente vistas por toda a classe no mesmo instante; não é assim se dissermos o ângulo ABC, o triângulo PFQ, a figura EGKt, e assim por diante; pois as letras devem ser rastreadas uma por uma antes que os estudantes organizem em suas mentes a magnitude particular referida, o que frequentemente causa confusão e erro, bem como perda de tempo. Também se as partes que são dadas como iguais, tiverem as mesmas cores em qualquer diagrama, a mente não se desviará do objeto diante dela; isto é, tal arranjo apresenta uma demonstração ocular das partes a serem provadas iguais, e o aprendiz retém os dados ao longo de todo o raciocínio. Mas quaisquer que sejam as vantagens do plano presente, se não for substituído, sempre pode ser feito um auxiliar poderoso para os outros métodos, para o propósito de introdução, ou de uma reminiscência mais rápida, ou de retenção mais permanente pela memória.

\charspacing{-2.5}{A experiência de todos que formaram sistemas para impressionar fatos no entendimento, concordam em provar que representações coloridas, como imagens, gravuras, diagramas, \&c.\ são mais facilmente fixadas na mente do que meras sentenças não marcadas por qualquer peculiaridade. Curioso como possa parecer, poetas parecem estar cientes deste fato mais do que matemáticos; muitos poetas modernos aludem a este sistema visível de comunicar conhecimento, um deles assim se expressou:}

\vskip 0.5\baselineskip

\begin{center} % Once again, the same verse by Horace, this time translated by Isaac Watts https://archive.org/stream/improvementofmin00wattuoft#page/198/mode/2up
Sons que abordam o ouvido se perdem e morrem\\
Em uma hora curta, mas estes que atingem o olho,\\
Vivem muito tempo na mente, a visão fiel\\
Grava o conhecimento com um feixe de luz.
% I'd put something like this instead:
% For man loves knowledge, and the beams of Truth
% More welcome touch his understanding's eye,
% Than all the blandishments of sound his ear,
% Than all of taste his tongue...
% from The Pleasures Of Imagination by Mark Akenside (1721–1770) https://archive.org/stream/pleasuresofimagi00aken#page/50/mode/2up
\end{center}

\vskip 0.5\baselineskip

Isto talvez possa ser considerado a única melhoria que a geometria plana recebeu desde os dias de Euclides, e se houvesse algum geômetra de nota antes daquele tempo, o sucesso de Euclides eclipsou completamente sua memória, e até causou que todas as coisas boas desse tipo fossem atribuídas a ele; como \AE sop entre os escritores de Fábulas. Também pode ser digno de nota, como diagramas tangíveis fornecem o único meio através do qual a geometria e outras artes lineares podem ser ensinadas aos cegos, o sistema visível não é menos adaptado às exigências dos surdos e mudos.

\charspacing{-1}{Deve-se ter cuidado para mostrar que a cor não tem nada a ver com as linhas, ângulos ou magnitudes, exceto meramente para nomeá-las. Uma linha matemática, que é comprimento sem largura, não pode possuir cor, ainda assim a junção de duas cores no mesmo plano dá uma boa ideia do que se entende por uma linha matemática; lembrem-se que estamos falando familiarmente, tal junção deve ser entendida e não a cor, quando dizemos a linha preta, a linha vermelha ou linhas, \&c.}

Cores e diagramas coloridos podem a princípio parecer um método desajeitado para transmitir noções adequadas das propriedades e partes de figuras e magnitudes matemáticas, no entanto descobrir-se-á que oferecem um meio mais refinado e extenso do que qualquer que tenha sido proposto até agora.

Definiremos aqui um ponto, uma linha e uma superfície, e demonstraremos uma proposição para mostrar a verdade desta afirmação.

Um ponto é aquilo que tem posição, mas não magnitude; ou um ponto é apenas posição, abstraído da consideração de comprimento, largura e espessura. Talvez a seguinte descrição seja melhor calculada para explicar a natureza do ponto matemático àqueles que não adquiriram a ideia, do que a definição acima especulosa.

\defineNewPicture{
	angleScale := 2;
	pair O, A, B, C;
	O := (0, 0);
	A := dir(0) scaled 3u;
	B := dir(120) scaled 3u;
	C := dir(240) scaled 3u;
		draw byAngle(A, O, B, byred, SOLID_SECTOR);
		draw byAngle(B, O, C, byblue, SOLID_SECTOR);
		draw byAngle(C, O, A, byyellow, SOLID_SECTOR);
}
Seja três cores \drawCurrentPictureInMargin se encontrarem e cobrirem uma porção do papel, onde se encontram não é azul, nem é amarelo, nem é vermelho, pois ocupa nenhuma porção do plano, pois se ocupasse, pertenceria à parte azul, vermelha ou amarela; ainda assim existe, e tem posição sem magnitude, de modo que com um pouco de reflexão, esta junção de três cores em um plano, dá uma boa ideia de um ponto matemático.

Uma linha é comprimento sem largura. Com a assistência de cores, quase da mesma maneira que antes, uma ideia de uma linha pode ser assim dada:—

\defineNewPicture{
	pair A, B, C, D, E, F;
	A := (0, 0);
	B := (5/2u, ypart(A));
	C := (xpart(A), -2u);
	D := (xpart(B), ypart(C));
	E := 1/2[A, C];
	F := 1/2[B, D];
		draw byPolygon(A,B,F,E)(byred);
		draw byPolygon(C,D,F,E)(byblue);
}
\drawCurrentPictureInMargin
Seja duas cores se encontrarem e cobrirem uma porção de papel; onde se encontram não é vermelho, nem azul; portanto a junção não ocupa nenhuma porção do plano, e portanto não pode ter largura, mas apenas comprimento: disso podemos prontamente formar uma ideia do que se entende por uma linha matemática. Para fins de ilustração, uma cor diferente da do papel, ou plano sobre o qual é traçada, teria sido suficiente; portanto, no futuro, se dissermos a linha vermelha, a linha azul ou linhas, \&c.\ devem ser entendidas como as junções com o plano sobre o qual são traçadas.

\defineNewPicture{
	pair A', A'', A''', B', B'', B''', C', C'', C''', D', D'', D''', d[];
	d1 := (3/2u, 0);
	d2 := (-3/4u, -2/3u);
	d3 := (0, -3/2u);
	A' := (0, 0);
	B' := A' shifted d1;
	C' := A' shifted d2;
	D' := C' shifted d1;
	A'' := A' shifted d3;
	B'' := B' shifted d3;
	C'' := C' shifted d3;
	D'' := D' shifted d3;
	A''' := A'' shifted d3;
	B''' := B'' shifted d3;
	C''' := C'' shifted d3;
	D''' := D'' shifted d3;
		draw byPolygon(A',B',B'',A'',C'',C')(byred);
		draw byPolygon(A'',B'',D'',C'')(byblue);
		draw byPolygon(C'',D'',B'',B''',D''',C''')(byyellow);
		draw byLine(A''', B''', white, SOLID_LINE, 2);
		draw byLine(A''', C''', white, SOLID_LINE, 2);
		draw byLine(A''', A', white, SOLID_LINE, 2);
		draw byLine(D', C', white, SOLID_LINE, THIN_WIDTH);
		draw byLine(D', B', white, SOLID_LINE, THIN_WIDTH);
		draw byLine(D', D''', white, SOLID_LINE, THIN_WIDTH);
		label.lft(btex P etex, C');
		label.lft(btex R etex, C'');
		label.rt(btex S etex, B'');
		label.rt(btex Q etex, B''');
}
\drawCurrentPictureInMargin
Superfície é aquilo que tem comprimento e largura sem espessura.

Quando consideramos um corpo sólido (PQ), percebemos imediatamente que tem três dimensões, a saber:— comprimento, largura e espessura; suponha que uma parte deste sólido (PS) seja vermelha, e a outra parte (QR) amarela, e que as cores sejam distintas sem se misturarem, a superfície azul (RS) que separa estas partes, ou que é a mesma coisa, aquilo que divide o sólido sem perda de material, deve ser sem espessura, e possui apenas comprimento e largura; isso claramente aparece do raciocínio, similar ao que acabamos de empregar ao definir, ou melhor, descrever um ponto e uma linha.

A proposição que selecionamos para elucidar a maneira pela qual os princípios são aplicados, é a quinta do primeiro Livro.

\defineNewPicture[1/4]{
angleScale := 5/6;
pair A, B, C, D, E;
A := (0, 0);
B := A shifted (u, -2u);
C := B xscaled -1;
D := 9/5[A,B];
E := 9/5[A,C];
byAngleDefine(B, A, C, byblack, SOLID_SECTOR);
byAngleDefine(A, B, C, byblue, SOLID_SECTOR);
byAngleDefine(B, C, A, byblue, SOLID_SECTOR);
byAngleDefine(C, B, E, byyellow, SOLID_SECTOR);
byAngleDefine(D, C, B, byyellow, SOLID_SECTOR);
byAngleDefine(B, D, C, byred, SOLID_SECTOR);
byAngleDefine(C, E, B, byred, SOLID_SECTOR);
byAngleDefine(E, B, D, byblack, ARC_SECTOR);
byAngleDefine(D, C, E, byblack, ARC_SECTOR);
draw byNamedAngleResized(BAC,ABC,BCA,CBE,DCB,BDC,CEB);
byLineDefine(B, D, byyellow, SOLID_LINE, REGULAR_WIDTH);
byLineDefine(C, E, byyellow, SOLID_LINE, REGULAR_WIDTH);
byLineDefine(B, E, byblue, SOLID_LINE, REGULAR_WIDTH);
byLineDefine(C, D, byblue, SOLID_LINE, REGULAR_WIDTH);
byLineDefine(A, B, byred, SOLID_LINE, REGULAR_WIDTH);
byLineDefine(A, C, byred, SOLID_LINE, REGULAR_WIDTH);
byLineDefine(B, C, byblack, SOLID_LINE, REGULAR_WIDTH);
draw byNamedLineSeq(0)(CD,noLine,BC,noLine,BE,CE,AC,AB,BD);
label.top(btex A etex, A);
label.lft(btex C etex, C);
label.rt(btex B etex, B);
label.lft(btex E etex, E);
label.rt(btex D etex, D);
}
Em um triângulo isósceles ABC, os ângulos internos na base ABC, ACB são iguais, e quando os lados AB, AC são prolongados, os ângulos externos na base BCE, CBD também são iguais.

\begin{center}
\drawCurrentPictureInMargin Prolongue \drawUnitLine{AB} e \drawUnitLine{AC},\\
faça $\drawUnitLine{BD} = \drawUnitLine{CE}$, trace \drawUnitLine{BE} e \drawUnitLine{CD}.\\
Em
\drawFromCurrentPicture{
draw byNamedAngle(BAC);
startAutoLabeling;
draw byNamedLineSeq(0)(BE,CE,AC,AB);
stopAutoLabeling;
}
e
\drawFromCurrentPicture{
draw byNamedAngle(BAC);
startAutoLabeling;
draw byNamedLineSeq(0)(BD,CD,AC,AB);
stopAutoLabeling;
}\\
temos $\drawUnitLine{AB,BD} = \drawUnitLine{AC,CE}$,\\
\drawAngle{BAC} comum e $\drawUnitLine{AB} = \drawUnitLine{AC}$:\\
$\therefore \drawAngle{BCA,DCB} = \drawAngle{ABC,CBE}$, $\drawUnitLine{BE} = \drawUnitLine{CD}$\\
e $\drawAngle{CEB} = \drawAngle{BDC}$ \byref{prop:I.IV}.\\
Novamente em \drawLine{BC,BE,CE} e \drawLine{BC,BD,CD},\\
$\drawUnitLine{BD} = \drawUnitLine{CE}$, $\drawAngle{CEB} = \drawAngle{BDC}$\\
e $\drawUnitLine{BE} = \drawUnitLine{CD}$;\\
$\therefore \drawAngle{DCE,DCB} = \drawAngle{EBD,CBE}$
e $\drawAngle{DCB} = \drawAngle{CBE}$ \byref{prop:I.IV}\\
Mas $\drawAngle{BCA,DCB} = \drawAngle{ABC,CBE}$, $\therefore \drawAngle{BCA} = \drawAngle{ABC}$.
\end{center}

\qedNB

\begin{center}
\emph{Anexando Letras ao Diagrama.}
\end{center}

Seja os lados iguais AB e AC serem prolongados através das extremidades BC, do terceiro lado, e na parte prolongada BD de qualquer um, seja assumido qualquer ponto D, e do outro seja cortado AE igual a AD \byref{prop:I.III}. Seja os pontos E e D, assim tomados nos lados prolongados, serem conectados por linhas retas DC e BE com as extremidades alternadas do terceiro lado do triângulo.

Nos triângulos DAC e EAB os lados DA e AC são respectivamente iguais a EA e AB, e o ângulo incluído A é comum a ambos os triângulos. Portanto \byref{prop:I.IV} a linha DC é igual a BE, o ângulo ADC ao ângulo AEB, e o ângulo ACD ao ângulo ABE; se das linhas iguais AD e AE forem tomados os lados iguais AB e AC, os restos BD e CE serão iguais. Portanto nos triângulos BDC e CEB, os lados BD e DC são respectivamente iguais a CE e EB, e os ângulos D e E incluídos por esses lados também são iguais. Portanto \byref{prop:I.IV} os ângulos DBC e ECB, que são aqueles incluídos pelo terceiro lado BC e os prolongamentos dos lados iguais AB e AC são iguais. Também os ângulos DCB e EBC são iguais se esses iguais forem tomados dos ângulos DCA e EBA antes provados iguais, os restos, que são os ângulos ABC e ACB opostos aos lados iguais, serão iguais.

\emph{Portanto em um triângulo isósceles,} \&c.

\qedNB

\charspacing{-2}{Nosso objetivo neste lugar sendo introduzir o sistema ao invés de ensinar qualquer conjunto particular de proposições, portanto selecionamos o acima mencionado fora do curso regular. Para escolas e outros lugares públicos de instrução, gizes tingidos servirão para descrever os diagramas, \&c.\ para uso privado lápis coloridos serão encontrados muito convenientes.}

Estamos felizes em descobrir que os Elementos da Matemática agora formam uma parte considerável de toda educação feminina sólida, portanto chamamos a atenção daqueles interessados ou engajados na educação de senhoras para este modo muito atraente de comunicar conhecimento, e para o trabalho subsequente para seu futuro desenvolvimento.

\charspacing{-2.5}{Concluiremos por enquanto observando, como os sentidos da visão e audição podem ser tão forçosa e instantaneamente abordados igualmente com mil quanto com um, \emph{o milhão} poderia ser ensinado geometria e outros ramos da matemática com grande facilidade, isso avançaria o propósito da educação mais do que qualquer coisa \emph{poderia} ser nomeada, pois ensinaria o povo como pensar, e não o que pensar; é neste particular que o grande erro da educação se origina.}

\chapter*{Esclarecimentos}

A geometria tem como seus principais objetos a exposição e explicação das propriedades de \emph{figura}, e figura é definida como a relação que subsiste entre os limites do espaço. Espaço ou magnitude é de três tipos, \emph{linear}, \emph{superficial}, e \emph{sólido}.

\defineNewPicture{
	pair A, B, C;
	numeric s;
	s := 3/2u;
	A := (0, s);
	B := (1/2s, 0);
	C := B xscaled -1;
		draw byAngle.A(B, A, C, byyellow, SOLID_SECTOR);
		byLineDefine(B, A, byblue, SOLID_LINE, REGULAR_WIDTH);
		byLineDefine(C, A, byred, SOLID_LINE, REGULAR_WIDTH);
		draw byNamedLineSeq(0)(CA,BA);
		label.urt(btex A etex, A);
}\drawCurrentPictureInMargin
Ângulos podem propriamente ser considerados como uma quarta espécie de magnitude. Magnitude angular evidentemente consiste de partes, e deve portanto ser admitida como uma espécie de quantidade. O estudante não deve supor que a magnitude de um ângulo é afetada pelo comprimento das linhas retas que o incluem e cuja divergência mútua é a medida. O \emph{vértice} de um ângulo é o ponto onde os \emph{lados} do ângulo se encontram, como A.

\defineNewPicture{
	pair B, C, D, E, F, G, H;
	numeric s;
	s := 5/4u;
	C := (0, 0);
	B := dir(0)*s;
	D := dir(50)*s;
	E := dir(-30)*s;
	F := E scaled -1;
	G := D scaled -1;
	H := B scaled -1;
	angleScale := 4/3;
		draw byAngle(E, C, B, byyellow, SOLID_SECTOR);
		draw byAngle(B, C, D, byblack, SOLID_SECTOR);
		draw byAngle(D, C, F, byblue, SOLID_SECTOR);
		draw byAngle(F, C, H, byred, SOLID_SECTOR);
		draw byAngle(H, C, G, byyellow, ARC_SECTOR);
		draw byAngle(G, C, E, byblue, ARC_SECTOR);
		draw byLine(B, H, byblue, SOLID_LINE, REGULAR_WIDTH);
		draw byLine(D, G, byred, SOLID_LINE, REGULAR_WIDTH);
		draw byLine(E, F, byblack, SOLID_LINE, REGULAR_WIDTH);
		label.bot(btex C etex, C shifted (0, -3pt));
		label.bot(btex B etex, B);
		label.lrt(btex D etex, D);
		label.llft(btex F etex, F);
		label.bot(btex H etex, H);
		label.lrt(btex G etex, G);
		label.llft(btex E etex, E);
}
\drawCurrentPictureInMargin
\charspacing{-2}{Um ângulo é frequentemente designado por uma única letra quando seus lados são as únicas linhas que se encontram juntas em seu vértice. Assim as linhas vermelha e azul formam o ângulo amarelo, que em outros sistemas seria chamado ângulo A. Mas quando mais de duas linhas se encontram no mesmo ponto, era necessário por métodos anteriores, para evitar confusão, empregar três letras para designar um ângulo sobre esse ponto, a letra que marcava o vértice do ângulo sendo sempre colocada no meio. Assim as linhas preta e vermelha se encontrando juntas em C, formam o ângulo azul, e tem sido usualmente denominado o ângulo FCD ou DCF. As linhas FC e CD são os lados do ângulo; o ponto C é seu vértice. Da mesma maneira o ângulo preto seria designado o ângulo DCB ou BCD. Os ângulos vermelho e azul adicionados juntos, ou o ângulo HCF adicionado a FCD, fazem o ângulo HCD; e assim de outros ângulos.}

Quando os lados de um ângulo são prolongados além de seu vértice, os ângulos feitos por elas em ambos os lados do vértice são ditos serem \emph{verticalmente opostos} um ao outro: assim os ângulos vermelho e amarelo são ditos serem ângulos verticalmente opostos.

\charspacing{-2}{\emph{Sobreposição} é o processo pelo qual uma magnitude pode ser concebida como sendo colocada sobre outra, de modo a exatamente cobri-la, ou de modo que cada parte de cada uma coincida exatamente.}

Uma linha é dita ser \emph{prolongada}, quando é estendida, prolongada, ou tem seu comprimento aumentado, e o aumento de comprimento que recebe é chamado \emph{parte prolongada}, ou seu \emph{prolongamento}.

\charspacing{-2.5}{O comprimento inteiro da linha ou linhas que encerram uma figura, é chamado seu \emph{perímetro}. Os primeiros seis livros de Euclides tratam apenas de figuras planas. Uma reta traçada do centro de um círculo até sua circunferência chama-se \emph{raio}. Aquele lado de um triângulo retângulo, que é oposto ao ângulo reto, é chamado a \emph{hipotenusa}. Um oblongo é definido no segundo livro, e chamado um \emph{retângulo}. Todas as linhas que são consideradas nos primeiros seis livros dos Elementos são supostas estar no mesmo plano.}

\charspacing{-3}{A \emph{régua} e o \emph{compasso} são os únicos instrumentos, cujo uso é permitido em Euclides, ou Geometria plana. Declarar esta restrição é o objeto dos \emph{postulados}.}

Os \emph{Axiomas} da geometria são certas proposições gerais, cuja verdade é tomada como auto-evidente e incapaz de ser estabelecida por demonstração.

\emph{Proposições} são aqueles resultados que são obtidos em geometria por um processo de raciocínio. Há duas espécies de proposições em geometria, \emph{problemas} e \emph{teoremas}.

Um \emph{Problema} é uma proposição na qual algo é proposto a ser feito; como uma linha a ser traçada sob algumas condições dadas, um círculo a ser descrito, alguma figura a ser construída, \&c.

A \emph{solução} do problema consiste em mostrar como a coisa requerida pode ser feita com o auxílio da régua e do compasso.

A \emph{demonstração} consiste em provar que o processo indicado na solução atinge o fim requerido.

Um \emph{Teorema} é uma proposição na qual a verdade de algum princípio é afirmada. Este princípio deve ser deduzido dos axiomas e definições, ou outras verdades previamente e independentemente estabelecidas. Mostrar isto é o objeto da demonstração.

Um \emph{Problema} é análogo a um postulado.

Um \emph{Teorema} se assemelha a um axioma.

Um \emph{Postulado} é um problema, cuja solução é assumida.

Um \emph{Axioma} é um teorema, cuja verdade é concedida sem demonstração.

Um \emph{Corolário} é uma inferência deduzida imediatamente de uma proposição.

Um \emph{Escólio} é uma nota ou observação sobre uma proposição não contendo uma inferência de importância suficiente para dar-lhe o nome de \emph{corolário}.

\charspacing{-2}{Um \emph{Lema} é uma proposição meramente introduzida para o propósito de estabelecer alguma proposição mais importante.}

\chapter*{Símbolos e abreviações}

\symb{$\therefore$}
expressa a palavra \emph{portanto}.

\symb{$\because$}
 expressa a palavra \emph{porque}.

\symb{$=$}
 expressa a palavra \emph{igual}. Este sinal de igualdade pode ser lido \emph{igual a}, ou \emph{é igual a}, ou \emph{são iguais a}; mas a discrepância em relação à introdução dos verbos auxiliares \emph{é}, \emph{são}, \&c.\ não pode afetar o rigor geométrico.

\symb{$\neq$}
 significa o mesmo que se as palavras \emph{‘não igual’} fossem escritas.

\symb{$>$}
 significa \emph{maior que}.

\symb{$<$}
 significa \emph{menor que}.

\symb{$\ngtr$}
 significa \emph{não maior que}.

\symb{$\nless$}
 significa \emph{não menor que}.

\symb{$+$}
 é lido \emph{mais}, o sinal de adição; quando interposto entre duas ou mais magnitudes, significa sua soma.

\symb{$-$}
 é lido \emph{menos}, significa subtração; e quando colocado entre duas quantidades denota que a última é tomada da primeira.

\symb{$\times$}
 este sinal expressa o produto de dois ou mais números quando colocado entre eles em aritmética e álgebra; mas em geometria é geralmente usado para expressar um \emph{retângulo}, quando colocado entre \enquote{duas linhas retas que contêm um de seus ângulos retos.} Um \emph{retângulo} também pode ser representado colocando um ponto entre dois de seus lados contíguos.

\symb{$:\ ::\ :$}
 expressa uma \emph{analogia} ou \emph{proporção}; assim se A, B, C e D representam quatro magnitudes, e A tem para B a mesma razão que C tem para D, a proporção é assim brevemente escrita

$A : B :: C : D$, $A : B = C : D$, ou $\dfrac{A}{B} = \dfrac{C}{D}$.

Esta igualdade ou mesmice de razão é lida,

como A está para B, assim C está para D;

ou A está para B, como C está para D.

\symb{$\parallel$}
 significa \emph{paralelo a}.

\symb{$\perp$}
 significa \emph{perpendicular a}.

\defineNewPicture{
	pair A, B, C, D;
	numeric s;
	s := 3/2u;
	A := (0, 0);
	B := dir(0)*s;
	C := dir(50)*s;
	D := dir(90)*s;
	byAngleDefine(B, A, C, byblack, ARC_SECTOR);
	byAngleDefine(B, A, D, byblack, ARC_SECTOR);
	byPointLabelRemove(A,B,C,D);
}

\symb{\drawAngle{BAC}}
 significa \emph{ângulo}.

\symb{\drawAngle{BAD}}
 significa \emph{ângulo reto}.

\symb{\drawTwoRightAngles}
 significa \emph{dois ângulos retos}.

\defineNewPicture{
	pair A, B, C, D;
	A := (0, -1/4u);
	B := (u, 0);
	C := (-u, 0);
	D := (0, u);
		byLineDefine (A, D, byblack, SOLID_LINE, REGULAR_WIDTH);
		byLineDefine (B, D, byblack, SOLID_LINE, REGULAR_WIDTH);
		byLineDefine (C, D, byblack, SOLID_LINE, REGULAR_WIDTH);
	byPointLabelRemove(A,D);
}

\symb{\drawFromCurrentPicture{
draw byNamedLine(AD);
draw byNamedLineSeq(0)(BD,CD);
}
or
\drawFromCurrentPicture{
draw byNamedLineSeq(0)(AD,BD);
}}
brevemente designa um \emph{ponto}.

O quadrado descrito sobre uma linha é concisamente escrito assim, $\drawUnitLine{AD}^2$.

Da mesma maneira duas vezes o quadrado de, é expresso por $2 \cdot \drawUnitLine{AD}^2$.

\symb{\indefstr}
 significa \emph{definição}.

\symb{\inpoststr}
 significa \emph{postulado}.

\symb{\inaxstr}
 significa \emph{axioma}.

\symb{hip.}
 significa \emph{hipótese}. Pode ser necessário aqui observar, que \emph{hipótese} é a condição assumida ou tomada como garantida. Assim, a hipótese da proposição dada na Introdução, é que o triângulo é isósceles, ou que seus lados são iguais.

\symb{\conststr}
 significa \emph{construção}. A \emph{construção} é a mudança feita na figura original, traçando linhas, fazendo ângulos, descrevendo círculos, \&c.\ para adaptá-la ao argumento da demonstração ou à solução do problema. As condições sob as quais essas mudanças são feitas, são tão indiscutíveis quanto aquelas contidas na hipótese. Por exemplo, se fizermos um ângulo igual a um ângulo dado, esses dois ângulos são iguais por construção.

\symb{\qedstr}
significa \emph{Quod erat demonstrandum}. O que era para ser demonstrado.

\part{Livro I}

\chapter*{Definições}

\startdefinition{}\label{def:I.I}
\begin{center}
Um \emph{ponto} é aquilo que não tem partes.
\end{center}

\startdefinition{}\label{def:I.II}
\begin{center}
Uma \emph{linha} é comprimento sem largura.
\end{center}

\startdefinition{}\label{def:I.III}
\begin{center}
As extremidades de uma linha são pontos.
\end{center}

\startdefinition{}\label{def:I.IV}
\begin{center}
A linha reta é aquela que se estende uniformemente entre suas extremidades.
\end{center}

\startdefinition{}\label{def:I.V}
\begin{center}
Uma superfície é aquilo que tem apenas comprimento e largura.
\end{center}

\startdefinition{}\label{def:I.VI}
\begin{center}
As extremidades de uma superfície são linhas.
\end{center}

\startdefinition{}\label{def:I.VII}
\begin{center}
Uma superfície plana é aquela que se estende uniformemente entre suas extremidades.
\end{center}

\startdefinition{}\label{def:I.VIII}
\begin{center}
Um ângulo plano é a inclinação de duas linhas uma em relação à outra, em um plano, que se encontram juntas, mas não estão na mesma direção.
\end{center}

\defineNewPicture{
	pair A, B, C;
	A := (0, 0);
	B := (2/3u, 2/3u);
	C := (u, ypart(A));
		byAngleDefine(B, A, C, byyellow, SOLID_SECTOR);
		draw byNamedAngleResized();
		byLineDefine(A, B, byblue, SOLID_LINE, REGULAR_WIDTH);
		byLineDefine(A, C, byred, SOLID_LINE, REGULAR_WIDTH);
		draw byNamedLineSeq(0)(AC,AB);
}
\startdefinition{}\label{def:I.IX}
\begin{center}
\drawCurrentPictureInMargin[inside] Um ângulo retilíneo plano é a inclinação de duas linhas retas uma em relação à outra, que se encontram juntas, mas não estão na mesma linha reta.
\end{center}

\defineNewPicture{
	pair A, B, C, D;
	A := (0, 0);
	B := (u, 0);
	C := (0, 2/3u);
	D := (-u, 0);
		byAngleDefine(B, A, C, byblack, ARC_SECTOR);
		byAngleDefine(D, A, C, byblack, ARC_SECTOR);
		draw byNamedAngleResized();
		draw byLine(D, B, byblack, SOLID_LINE, REGULAR_WIDTH);
		draw byLine(A, C, byblack, SOLID_LINE, REGULAR_WIDTH);
}
\startdefinition{}\label{def:I.X}
\begin{center}
\drawCurrentPictureInMargin[inside] Quando uma linha reta apoiada sobre outra linha reta torna os ângulos adjacentes iguais, cada um desses ângulos é chamado um \emph{ângulo reto}, e cada uma dessas linhas é dita ser \emph{perpendicular} uma à outra.
\end{center}

\defineNewPicture{
	pair A, B, C;
	A := (0, 0);
	B := (-2/3u, 2/3u);
	C := (u, ypart(A));
		byAngleDefine(B, A, C, byred, SOLID_SECTOR);
		draw byNamedAngleResized();
		byLineDefine(A, B, byyellow, SOLID_LINE, REGULAR_WIDTH);
		byLineDefine(A, C, byblue, SOLID_LINE, REGULAR_WIDTH);
		draw byNamedLineSeq(0)(AC,AB);
}
\startdefinition{}\label{def:I.XI}
\begin{center}
\drawCurrentPictureInMargin[inside] Um ângulo obtuso é um ângulo maior que um ângulo reto
\end{center}

\defineNewPicture{
	pair A, B, C;
	A := (0, 0);
	B := (2/3u, 2/3u);
	C := (u, ypart(A));
		byAngleDefine(B, A, C, byblue, SOLID_SECTOR);
		draw byNamedAngleResized();
		byLineDefine(A, B, byyellow, SOLID_LINE, REGULAR_WIDTH);
		byLineDefine(A, C, byred, SOLID_LINE, REGULAR_WIDTH);
		draw byNamedLineSeq(0)(AC,AB);
}
\startdefinition{}\label{def:I.XII}
\begin{center}
\drawCurrentPictureInMargin[inside] Um ângulo agudo é um ângulo menor que um ângulo reto.
\end{center}

\startdefinition{}\label{def:I.XIII}
\begin{center}
Um termo ou limite é a extremidade de qualquer coisa.
\end{center}

\startdefinition{}\label{def:I.XIV}
\begin{center}
Uma figura é uma superfície encerrada em todos os lados por uma linha ou linhas.
\end{center}

\defineNewPicture{
	pair O, A, B, C, D, E;
	numeric r;
	r := 1/2u;
	O := (0, 0);
	A := dir(0) scaled r;
	B := dir(60) scaled r;
	C := dir(130) scaled r;
	D := dir(180) scaled r;
	E := dir(-60) scaled r;
		draw byLine(O, B)(byblack, SOLID_LINE, REGULAR_WIDTH);
		draw byLine(O, C)(byred, SOLID_LINE, REGULAR_WIDTH);
		draw byLine(O, E)(byyellow, SOLID_LINE, REGULAR_WIDTH);
		draw byLine(D, A)(byblue, SOLID_LINE, REGULAR_WIDTH);
		draw byCircleR(O, r, byred, 0, 0, 0);
}
\startdefinition{}\label{def:I.XV}
\begin{center}
\drawCurrentPictureInMargin[inside] Um círculo é uma figura plana, limitada por uma linha contínua, chamada sua circunferência ou periferia; e tendo um certo ponto dentro dele, do qual todas as linhas retas traçadas até sua circunferência são iguais.
\end{center}

\startdefinition{}\label{def:I.XVI}
\begin{center}
Este ponto, do qual as retas iguais são traçadas, é chamado centro do círculo.
\end{center}

\defineNewPicture{
	pair O, A, B;
	numeric r;
	r := 1/2u;
	O := (0, 0);
	A := dir(0) scaled r;
	B := dir(180) scaled r;
		draw byLine(A, B)(byyellow, SOLID_LINE, REGULAR_WIDTH);
		draw byCircleR(O, r, byred, 0, 0, 0);
}
\startdefinition{}\label{def:I.XVII}
\begin{center}
\drawCurrentPictureInMargin[inside] Um diâmetro de um círculo é uma reta traçada através do centro, terminada em ambos os sentidos na circunferência.
\end{center}

\defineNewPicture{
	pair O, A, B;
	numeric r;
	r := 1/2u;
	O := (0, 0);
	A := dir(0) scaled r;
	B := dir(180) scaled r;
		draw byLine(A, B)(byblue, SOLID_LINE, REGULAR_WIDTH);
		draw byArc(O, A, B)(r, byyellow, 0, 0, 0, 0);
		draw byArc(O, B, A)(r, byyellow, 1, 0, 0, 0);
}
\startdefinition{}\label{def:I.XVIII}
\begin{center}
\drawCurrentPictureInMargin[inside] Um semicírculo é a figura contida pelo diâmetro, e a parte do círculo cortada pelo diâmetro.
\end{center}

\defineNewPicture{
	pair O, A, B;
	path P;
	numeric r;
	r := 1/2u;
	P := fullcircle scaled 2r;
	O := (0, 0);
	A := point 1 of P;
	B := point 3 of P;
		draw byLine(A, B)(byred, SOLID_LINE, REGULAR_WIDTH);
		draw byArc(O, A, B)(r, byblue, 0, 0, 0, 0);
		draw byArc(O, B, A)(r, byblue, 1, 0, 0, 0);
}
\startdefinition{}\label{def:I.XIX}
\begin{center}
\drawCurrentPictureInMargin[inside] Um segmento de um círculo é uma figura contida por uma linha reta e a parte da circunferência que ela corta.
\end{center}

\startdefinition{}\label{def:I.XX}
\begin{center}
Uma figura contida apenas por linhas retas, é chamada uma figura retilínea.
\end{center}

\startdefinition{}\label{def:I.XXI}
\begin{center}
Um triângulo é uma figura retilínea incluída por três lados.
\end{center}

\defineNewPicture{
	pair A, B, C, D;
	A := (0, 0);
	B := (u, 1/2u);
	C := (-1/2u, -4/3u);
	D := (4/3u, -u);
		draw byLine(C, B)(byred, SOLID_LINE, REGULAR_WIDTH);
		draw byLine(A, D)(byblue, SOLID_LINE, REGULAR_WIDTH);
		byLineDefine(A, B, byyellow, SOLID_LINE, REGULAR_WIDTH);
		byLineDefine(A, C, byyellow, SOLID_LINE, REGULAR_WIDTH);
		byLineDefine(B, D, byyellow, SOLID_LINE, REGULAR_WIDTH);
		byLineDefine(C, D, byblack, SOLID_LINE, REGULAR_WIDTH);
		draw byNamedLineSeq(0)(BD,CD,AC,AB);
		draw byLabelsOnPolygon(A, B, D, C)(ALL_LABELS, 0);
}
\startdefinition{}\label{def:I.XXII}
\begin{center}
\drawCurrentPictureInMargin[inside] Uma figura quadrilátera é aquela que é limitada por quatro lados. As linhas retas \drawUnitLine{AD} e \drawUnitLine{CB} conectando os vértices dos ângulos opostos de uma figura quadrilátera, são chamadas suas diagonais.
\end{center}

\startdefinition{}\label{def:I.XXIII}
\begin{center}
Um polígono é uma figura retilínea limitada por mais de quatro lados.
\end{center}

\defineNewPicture{
	pair A, B, C;
	A := dir(-30) scaled 1/2u;
	B := dir(-150) scaled 1/2u;
	C := dir(90) scaled 1/2u;
		byLineDefine(A, B, byblue, SOLID_LINE, REGULAR_WIDTH);
		byLineDefine(B, C, byred, SOLID_LINE, REGULAR_WIDTH);
		byLineDefine(C, A, byyellow, SOLID_LINE, REGULAR_WIDTH);
		draw byNamedLineSeq(0)(AB,BC,CA);
}
\startdefinition{}\label{def:I.XXIV}
\begin{center}
\drawCurrentPictureInMargin[inside] Um triângulo cujos lados são iguais, é dito ser equilátero.
\end{center}

\defineNewPicture{
	pair A, B, C;
	A := dir(-60) scaled 1/2u;
	B := dir(-120) scaled 1/2u;
	C := dir(90) scaled 1/2u;
		byLineDefine(A, B, byblue, SOLID_LINE, REGULAR_WIDTH);
		byLineDefine(B, C, byred, SOLID_LINE, REGULAR_WIDTH);
		byLineDefine(C, A, byred, SOLID_LINE, REGULAR_WIDTH);
		draw byNamedLineSeq(0)(AB,BC,CA);
}
\startdefinition{}\label{def:I.XXV}
\begin{center}
\drawCurrentPictureInMargin[inside] Um triângulo que tem apenas dois lados iguais é chamado um triângulo isósceles.
\end{center}

\startdefinition{}\label{def:I.XXVI}
\begin{center}
Um triângulo escaleno é aquele que não tem dois lados iguais.
\end{center}

\defineNewPicture{
	pair A, B, C;
	A := (0, 0);
	B := (-u, 0);
	C := (0, 3/4u);
		byLineDefine(A, B, byred, SOLID_LINE, REGULAR_WIDTH);
		byLineDefine(B, C, byyellow, SOLID_LINE, REGULAR_WIDTH);
		byLineDefine(C, A, byblue, SOLID_LINE, REGULAR_WIDTH);
		draw byNamedLineSeq(0)(AB,BC,CA);
}
\startdefinition{}\label{def:I.XXVII}
\begin{center}
\drawCurrentPictureInMargin[inside]Um triângulo retângulo é aquele que tem um ângulo reto.
\end{center}

\vskip \baselineskip

\defineNewPicture{
	pair A, B, C;
	A := (-1/4u, 0);
	B := (-u, 0);
	C := (0, 3/4u);
		byLineDefine(A, B, byred, SOLID_LINE, REGULAR_WIDTH);
		byLineDefine(B, C, byblue, SOLID_LINE, REGULAR_WIDTH);
		byLineDefine(C, A, byyellow, SOLID_LINE, REGULAR_WIDTH);
		draw byNamedLineSeq(0)(AB,BC,CA);
}
\startdefinition{}\label{def:I.XXVIII}
\begin{center}
\drawCurrentPictureInMargin[inside] Um triângulo obtusângulo é aquele que tem um ângulo obtuso.
\end{center}

\defineNewPicture{
	pair A, B, C;
	A := (0, 0);
	B := (-u, 0);
	C := (-1/4u, 3/4u);
		byLineDefine(A, B, byblue, SOLID_LINE, REGULAR_WIDTH);
		byLineDefine(B, C, byyellow, SOLID_LINE, REGULAR_WIDTH);
		byLineDefine(C, A, byred, SOLID_LINE, REGULAR_WIDTH);
		draw byNamedLineSeq(0)(AB,BC,CA);
}
\startdefinition{}\label{def:I.XXIX}
\begin{center}
\drawCurrentPictureInMargin[inside] Um triângulo acutângulo é aquele que tem três ângulos agudos.
\end{center}

\defineNewPicture{
	pair A, B, C, D;
	numeric s;
	s := u;
	A := (0, 0);
	B := (s, 0);
	C := (0, s);
	D := (s, s);
		byLineDefine(A, B, byred, SOLID_LINE, REGULAR_WIDTH);
		byLineDefine(A, C, byblue, SOLID_LINE, REGULAR_WIDTH);
		byLineDefine(B, D, byyellow, SOLID_LINE, REGULAR_WIDTH);
		byLineDefine(C, D, byblack, SOLID_LINE, REGULAR_WIDTH);
		draw byNamedLineSeq(0)(AB,AC,CD,BD);
}
\startdefinition{}\label{def:I.XXX}
\begin{center}
\drawCurrentPictureInMargin[inside] De figuras de quatro lados, um quadrado é aquele que tem todos os seus lados iguais, e todos os seus ângulos são ângulos retos.
\end{center}

\defineNewPicture{
	pair A, B, C, D;
	numeric s;
	s := u;
	A := (0, 0);
	B := (s, 0);
	C := A shifted (dir(80) scaled s);
	D := B shifted (dir(80) scaled s);
		byLineDefine(A, B, byred, SOLID_LINE, REGULAR_WIDTH);
		byLineDefine(A, C, byblue, SOLID_LINE, REGULAR_WIDTH);
		byLineDefine(B, D, byyellow, SOLID_LINE, REGULAR_WIDTH);
		byLineDefine(C, D, byblack, SOLID_LINE, REGULAR_WIDTH);
		draw byNamedLineSeq(0)(AB,AC,CD,BD);
}
\startdefinition{}\label{def:I.XXXI}
\begin{center}
\drawCurrentPictureInMargin[inside] Um losango é aquele que tem todos os seus lados iguais, mas seus ângulos não são ângulos retos.
\end{center}

\defineNewPicture{
	pair A,B,C,D;
	numeric s;
	s := u;
	A := (0, 0);
	B := (4/3s, 0);
	C := (0, 3/4s);
	D := (4/3s, 3/4s);
		byLineDefine(A, B, byblue, SOLID_LINE, REGULAR_WIDTH);
		byLineDefine(A, C, byred, SOLID_LINE, REGULAR_WIDTH);
		byLineDefine(B, D, byred, SOLID_LINE, REGULAR_WIDTH);
		byLineDefine(C, D, byblue, SOLID_LINE, REGULAR_WIDTH);
		draw byNamedLineSeq(0)(AB,AC,CD,BD);
}
\startdefinition{}\label{def:I.XXXII}
\begin{center}
\drawCurrentPictureInMargin[inside] Um oblongo é aquele que tem todos os seus ângulos retos, mas não tem todos os seus lados iguais.
\end{center}

\defineNewPicture{
	pair A, B, C, D;
	numeric s;
	s := u;
	A := (0, 0);
	B := (s, 0);
	C := (1/4s, 3/4s);
	D := (s + 1/4s, 3/4s);
		byLineDefine(A, B, byblue, SOLID_LINE, REGULAR_WIDTH);
		byLineDefine(A, C, byred, SOLID_LINE, REGULAR_WIDTH);
		byLineDefine(B, D, byred, SOLID_LINE, REGULAR_WIDTH);
		byLineDefine(C, D, byblue, SOLID_LINE, REGULAR_WIDTH);
		draw byNamedLineSeq(0)(AB,AC,CD,BD);
}
\startdefinition{}\label{def:I.XXXIII}
\begin{center}
\drawCurrentPictureInMargin[inside] Um romboide é aquele que tem seus lados opostos iguais um ao outro, mas seus lados não são todos iguais nem seus ângulos são retos.
\end{center}

\startdefinition{}\label{def:I.XXXIV}
\begin{center}
Todas as outras figuras quadriláteras são chamadas trapézios.
\end{center}

\defineNewPicture{
	pair A, B, C, D;
	numeric s;
	s := u;
	A := (0, 0);
	B := (4/3s, 0);
	C := (0, 1/2s);
	D := (4/3s, 1/2s);
		draw byLine(A, B, byred, SOLID_LINE, REGULAR_WIDTH);
		draw byLine(C, D, byyellow, SOLID_LINE, REGULAR_WIDTH);
}
\startdefinition{}\label{def:I.XXXV}
\begin{center}
\drawCurrentPictureInMargin[inside] Linhas retas paralelas são aquelas que estão no mesmo plano, e que sendo prolongadas continuamente em ambas as direções nunca se encontrariam.
\end{center}

\chapter*{Postulados}

\startpostulate{}\label{post:I.I}
Seja concedido que se trace uma reta de qualquer ponto a qualquer outro ponto.

\startpostulate{}\label{post:I.II}
Seja concedido que uma linha reta finita possa ser prolongada a qualquer comprimento em uma linha reta.

\startpostulate{}\label{post:I.III}
Seja concedido que um círculo possa ser descrito com qualquer centro a qualquer distância desse centro.

\chapter*{Axiomas}

\startaxiom{}\label{ax:I.I}
Magnitudes que são iguais à mesma são iguais entre si.

\startaxiom{}\label{ax:I.II}
Se iguais forem adicionados a iguais as somas serão iguais.

\startaxiom{}\label{ax:I.III}
Se iguais forem retirados de iguais os restos serão iguais.

\startaxiom{}\label{ax:I.IV}
Se iguais forem adicionados a desiguais as somas serão desiguais.

\startaxiom{}\label{ax:I.V}
Se iguais forem retirados de desiguais os restos serão desiguais.

\startaxiom{}\label{ax:I.VI}
Os dobros da mesma ou de magnitudes iguais são iguais.

\startaxiom{}\label{ax:I.VII}
As metades da mesma ou de magnitudes iguais são iguais.

\startaxiom{}\label{ax:I.VIII}
As magnitudes que coincidem uma com a outra, ou exatamente preenchem o mesmo espaço, são iguais.

\startaxiom{}\label{ax:I.IX}
O todo é maior que sua parte.

\startaxiom{}\label{ax:I.X}
Duas linhas retas não podem incluir um espaço.

\startaxiom{}\label{ax:I.XI}
Todos os ângulos retos são iguais.

\startaxiom{}\label{ax:I.XII}
\defineNewPicture{
	pair A, B, C, D, E, F, G, H;
	numeric s;
	s := 3/2u;
	A := (0, 0);
	B := (4/3s, 0);
	C := (0, s);
	D := (4/3s, s);
	E := (1/3s, 8/6s);
	F := (xpart(E), -2/6s);
	G = whatever[A, B] = whatever[E, F];
	H = whatever[C, D] = whatever[E, F];
		byAngleDefine(B, G, E, byred, SOLID_SECTOR);
		byAngleDefine(D, H, F, byyellow, SOLID_SECTOR);
		draw byNamedAngleResized();
		draw byLine(A, B, byblue, SOLID_LINE, REGULAR_WIDTH);
		draw byLine(C, D, byred, SOLID_LINE, REGULAR_WIDTH);
		draw byLine(E, F, byblack, SOLID_LINE, REGULAR_WIDTH);
		draw byLabelLine(0)(AB, CD, EF);
		draw byLabelsOnPolygon(E, H, D)(OMIT_FIRST_LABEL+OMIT_LAST_LABEL, 0);
		draw byLabelsOnPolygon(B, G, F)(OMIT_FIRST_LABEL+OMIT_LAST_LABEL, 0);
}
\drawCurrentPictureInMargin[inside]
Se duas linhas retas $\left(\vcenter{\nointerlineskip\hbox{\drawUnitLine{AB}}\nointerlineskip\hbox{\drawUnitLine{CD}}}\right)$ encontram uma terceira linha reta (\drawUnitLine{EF}) de modo a fazer os dois ângulos interiores (\drawAngle{H} e \drawAngle{G}) no mesmo lado menores que dois ângulos retos, essas duas linhas retas se encontrarão se forem prolongadas naquele lado em que os ângulos são menores que dois ângulos retos.

O décimo segundo axioma pode ser expresso de qualquer das seguintes maneiras:
\begin{enumerate}
\item Duas linhas retas divergentes não podem ser ambas paralelas à mesma linha reta.
\item Se uma linha reta intersectar uma das duas linhas retas paralelas, também deve intersectar a outra.
\item Apenas uma linha reta pode ser traçada através de um ponto dado, paralela a uma linha reta dada.
\end{enumerate}

\chapter*{Proposições}

\startproblem{Prop. I. Prob.}\label{prop:I.I}

\defineNewPicture[1/2]{
	pair A, B, C;
	path P[];
	numeric r;
	r := 3/2u;
	A := (0, 0);
	B := (r, 0);
	P1 := fullcircle scaled 2r;
	P2 := fullcircle scaled 2r shifted B;
	C := P1 intersectionpoint P2;
		byLineDefine(A, B, byblack, SOLID_LINE, REGULAR_WIDTH);
		byLineDefine(B, C, byred, SOLID_LINE, REGULAR_WIDTH);
		byLineDefine(C, A, byyellow, SOLID_LINE, REGULAR_WIDTH);
		draw byNamedLineSeq(-1)(AB,CA,BC);
		draw byCircle.A(A, B, byblue, 0, 0, 1/2);
		draw byCircle.B(B, A, byred, 0, 0, 1/2);
		draw byLabelsOnPolygon(A, C, B)(ALL_LABELS, 1);
}
\drawCurrentPictureInMargin
\problem{S}{obre}{uma linha reta finita dada (\drawUnitLine{AB}) descrever um triângulo equilátero.}

\begin{center}
Descreva \offsetPicture{15pt}{0pt}{\drawFromCurrentPicture{
draw byNamedLine(AB);
draw byNamedCircle(A);
draw byLabelLineEnd(A, B, 0);
draw byLabelLineEnd(B, A, 1);
}} e \offsetPicture{15pt}{0pt}{\drawFromCurrentPicture{
draw byNamedLine(AB);
draw byNamedCircle(B);
draw byLabelLineEnd(A, B, 1);
draw byLabelLineEnd(B, A, 0);
}}
\byref{post:I.III};\\
trace \drawUnitLine{CA} e \drawUnitLine{BC} \byref{post:I.I}.\\
Então \drawLine[bottom][triangleABC]{AB,CA,BC} será equilátero.

Pois $\drawUnitLine{AB} = \drawUnitLine{CA}$ \byref{def:I.XV};\\
e $\drawUnitLine{AB} = \drawUnitLine{BC}$ \byref{def:I.XV};\\
$\therefore \drawUnitLine{CA} = \drawUnitLine{BC}$ \byref{ax:I.I};\\
e portanto \triangleABC\ é o triângulo equilátero requerido.
\end{center}

\qed

\startproblem{Prop. II. Prob.}\label{prop:I.II}

\defineNewPicture{
pair A, B, C, D, E, F;
path P[];
numeric r[];
A := (0, 0);
B := (-3/5u, -3/5u);
C := (-2u, -1/3u);
r1 := abs(A-B);
D := (fullcircle scaled 2r1 shifted A) intersectionpoint (fullcircle scaled 2r1 shifted B);
r2 := abs(B-C);
r3 := r1 + r2;
P1 := fullcircle scaled 2r2 shifted B;
P2 := fullcircle scaled 2r3 shifted D;
E := (D -- 10[D, B]) intersectionpoint P1;
F := (D -- 10[D, A]) intersectionpoint P2;
byLineDefine(A, B, byblack, DASHED_LINE, REGULAR_WIDTH);
byLineDefine(B, C, byblack, SOLID_LINE, REGULAR_WIDTH);
byLineDefine(B, D, byred, SOLID_LINE, REGULAR_WIDTH);
byLineDefine(D, A, byred, SOLID_LINE, REGULAR_WIDTH); % improvement: change style of either BD or DA
byLineDefine(B, E, byyellow, SOLID_LINE, REGULAR_WIDTH);
byLineDefine(A, F, byblue, SOLID_LINE, REGULAR_WIDTH);
draw byNamedLineSeq(0)(AB,BC);
draw byNamedLineSeq(0)(BE,BD,DA,AF);
draw byCircle.A(D, E, byred, 0, 0, 1/2);
draw byCircle.B(B, C, byblue, 0, 0, -1/2);
draw byLabelsOnPolygon(E, D, A, F)(OMIT_FIRST_LABEL+OMIT_LAST_LABEL, 1);
draw byLabelsOnPolygon(E, B, C)(OMIT_FIRST_LABEL+OMIT_LAST_LABEL, 1);
draw byLabelsOnCircle(C)(B);
draw byLabelsOnCircle(E, F)(A);
}
\drawCurrentPictureInMargin
\problem{D}{e}{um ponto dado (\drawFromCurrentPicture[middle][pointA]{
startGlobalRotation(-lineAngle.DA);
draw byNamedPointLines(A,"AB");
stopGlobalRotation;
}), traçar uma linha reta igual a uma linha reta dada (\drawUnitLine{BC}).}

\begin{center}
Trace \drawUnitLine{AB} \byref{post:I.I}, descreva \drawFromCurrentPicture[bottom]{
startAutoLabeling;
startTempScale(scaleFactor*3);
startGlobalRotation(180-lineAngle.AB);
draw byNamedLineSeq(0)(AB,BD,DA);
stopGlobalRotation;
stopTempScale;
stopAutoLabeling;
} \byref{prop:I.I},\\
prolongue \drawUnitLine{BD} \byref{post:I.II},\\
descreva
\drawFromCurrentPicture{
draw byNamedLine (BC);
draw byNamedCircle(B);
draw byLabelLineEnd(B, C, 0);
draw byLabelLineEnd(C, B, 0);
}
\byref{post:I.III}, e
\drawFromCurrentPicture{
draw byNamedLineSeq(0)(BD, BE);
draw byNamedCircle(A);
draw byLabelLineEnd(D, E, 0);
draw byLabelLineEnd(E, D, 1);
}
\byref{post:I.III};\\
prolongue \drawUnitLine{DA} \byref{post:I.II},\\
então \drawUnitLine{AF} é a linha requerida.

Pois $\drawUnitLine{BE,BD} = \drawUnitLine{DA,AF}$ \byref{def:I.XV},\\
e $\drawUnitLine{BD} = \drawUnitLine{DA}$ \byref{\constref},\\
$\therefore \drawUnitLine{BE} = \drawUnitLine{AF}$ \byref{post:I.III},\\
mas \byref{def:I.XV} $\drawUnitLine{BC} = \drawUnitLine{BE} = \drawUnitLine{AF}$;

$\therefore \drawUnitLine{AF}$ traçada do ponto dado (\pointA), é igual à linha dada \drawUnitLine{BC}.
\end{center}

\qed

\startproblem{Prop. III. Prob.}\label{prop:I.III}
\defineNewPicture{
pair A, B, C, D, E, F;
path P;
numeric r;
A := (0, 0);
r := 7/4u;
B := A shifted (r, 0);
C := A shifted (4/3r, 0);
D := A shifted dir(30)*r;
E := A shifted (7/6r, -1/6r);
F := A shifted (7/6r, -7/6r);
byLineDefine(A, B, byblack, SOLID_LINE, REGULAR_WIDTH);
byLineDefine(B, C, byblack, DASHED_LINE, REGULAR_WIDTH);
byLineDefine(A, D, byred, SOLID_LINE, REGULAR_WIDTH);
draw byNamedLineSeq(0)(BC,AB,AD);
draw byLine(E, F, byblue, SOLID_LINE, REGULAR_WIDTH);
draw byCircle.A(A, D, byblue, 0, 0, 0);
draw byLabelsOnPolygon(B, A, D)(OMIT_FIRST_LABEL+OMIT_LAST_LABEL, 1);
draw byLabelLineEnd(D, A, 0);
draw byLabelLineEnd(C, A, 0);
draw byLabelPoint(B, angle(B-A) + 45, 2);
draw byLabelsOnPolygon(E, F)(ALL_LABELS, 0);
}
\drawCurrentPictureInMargin
\problem{D}{a}{maior (\drawUnitLine{AB,BC}) de duas linhas retas dadas, cortar uma parte igual à menor (\drawUnitLine{EF}).}

\begin{center}
Trace $\drawUnitLine{AD} = \drawUnitLine{EF}$ \byref{prop:I.II};\\
descreva
\drawFromCurrentPicture{
draw byNamedLine (AD);
draw byNamedCircle(A);
draw byLabelLineEnd(D, A, 0);
draw byLabelLineEnd(A, D, 0);
} \byref{post:I.III},\\
então $\drawUnitLine{EF} = \drawUnitLine{AB}$

Pois $\drawUnitLine{AD} = \drawUnitLine{AB}$ \byref{def:I.XV},\\
e $\drawUnitLine{EF} = \drawUnitLine{AD}$ \byref{\constref};

$\therefore \drawUnitLine{EF} = \drawUnitLine{AB}$ \byref{ax:I.I}.
\end{center}

\qed

\starttheorem{Prop. IV. Theor.}\label{prop:I.IV}
\defineNewPicture[1/5]{
pair A, B, C, D, E, F, d;
A := (0, 0);
B := A shifted (-5/2u, -7/2u);
C := A shifted (1/2u, -3u);
d := (0, -4u);
D := A shifted d;
E := B shifted d;
F := C shifted d;
byAngleDefine(B, A, C, byyellow, SOLID_SECTOR);
byAngleDefine(A, B, C, byblue, SOLID_SECTOR);
byAngleDefine(B, C, A, byred, SOLID_SECTOR);
draw byNamedAngleResized(BAC, ABC, BCA);
byLineDefine(A, B, byred, SOLID_LINE, REGULAR_WIDTH);
byLineDefine(B, C, byblack, SOLID_LINE, REGULAR_WIDTH);
byLineDefine(C, A, byblue, SOLID_LINE, REGULAR_WIDTH);
draw byNamedLineSeq(0)(CA,BC,AB);
byAngleDefine(E, D, F, byyellow, SOLID_SECTOR);
byAngleDefine(D, E, F, byblue, SOLID_SECTOR);
byAngleDefine(E, F, D, byred, SOLID_SECTOR);
draw byNamedAngleResized(EDF, DEF, EFD);
byLineDefine(D, E, byred, SOLID_LINE, THIN_WIDTH);
byLineDefine(E, F, byblack, SOLID_LINE, THIN_WIDTH);
byLineDefine(F, D, byblue, SOLID_LINE, THIN_WIDTH);
draw byNamedLineSeq(0)(FD,EF,DE);
draw byLabelsOnPolygon(F, E, D)(ALL_LABELS, 0);
draw byLabelsOnPolygon(B, A, C)(ALL_LABELS, 1);
}
\drawCurrentPictureInMargin
\problem[4]{S}{e}{dois triângulos têm dois lados de um respectivamente iguais a dois lados do outro, (\drawUnitLine{AB} a \drawUnitLine{DE} e \drawUnitLine{CA} a \drawUnitLine{FD}) e os ângulos (\drawAngle{A} e \drawAngle{D}) contidos por esses lados iguais também iguais; então suas bases ou seus lados (\drawUnitLine{BC} e \drawUnitLine{EF}) também são iguais: e os ângulos restantes opostos a lados iguais são respectivamente iguais ($\drawAngle{B} = \drawAngle{E}$ e $\drawAngle{C} = \drawAngle{F}$): e os triângulos são iguais em todos os aspectos.}

Sejam dois triângulos concebidos, de modo a serem colocados de tal forma que o vértice de um dos ângulos iguais, \drawAngle{A} ou \drawAngle{D}; caia sobre o do outro, e \drawUnitLine{AB} coincida com \drawUnitLine{DE}, então \drawUnitLine{CA} coincidirá com \drawUnitLine{FD} se aplicado: consequentemente \drawUnitLine{BC} coincidirá com \drawUnitLine{EF}, ou duas linhas retas encerrarão um espaço, o que é impossível \byref{ax:I.X}, portanto $\drawUnitLine{BC} = \drawUnitLine{EF}$, $\drawAngle{B} = \drawAngle{E}$ e $\drawAngle{C} = \drawAngle{F}$, e como os triângulos \drawLine{CA,BC,AB} e \drawLine{FD,EF,DE} coincidem, quando aplicados, eles são iguais em todos os aspectos.

\qed

\starttheorem{Prop. V. Theor.}\label{prop:I.V}
\defineNewPicture{
pair A, B, C, D, E;
A := (0, 0);
B := A shifted (u, -2u);
C := B xscaled -1;
D := 9/5[A,B];
E := 9/5[A,C];
byAngleDefine(B, A, C, byblack, SOLID_SECTOR);
byAngleDefine(A, B, C, byblue, SOLID_SECTOR);
byAngleDefine(B, C, A, byblue, SOLID_SECTOR);
byAngleDefine(C, B, E, byyellow, SOLID_SECTOR);
byAngleDefine(D, C, B, byyellow, SOLID_SECTOR);
byAngleDefine(B, D, C, byred, SOLID_SECTOR);
byAngleDefine(C, E, B, byred, SOLID_SECTOR);
byAngleDefine(E, B, D, byblack, ARC_SECTOR);
byAngleDefine(D, C, E, byblack, ARC_SECTOR);
draw byNamedAngleResized(BAC,ABC,BCA,CBE,DCB,BDC,CEB);
byLineDefine(B, D, byyellow, SOLID_LINE, REGULAR_WIDTH);
byLineDefine(C, E, byyellow, SOLID_LINE, REGULAR_WIDTH);
byLineDefine(B, E, byblue, SOLID_LINE, REGULAR_WIDTH);
byLineDefine(C, D, byblue, SOLID_LINE, REGULAR_WIDTH);
byLineDefine(A, B, byred, SOLID_LINE, REGULAR_WIDTH);
byLineDefine(A, C, byred, SOLID_LINE, REGULAR_WIDTH);
byLineDefine(B, C, byblack, SOLID_LINE, REGULAR_WIDTH);
draw byNamedLineSeq(0)(CD,noLine,BC,noLine,BE,CE,AC,AB,BD);
draw byLabelsOnPolygon(E, C, A, B, D, C, B)(ALL_LABELS, 0);
}
\drawCurrentPictureInMargin
\problem[4]{E}{m}{qualquer triângulo isósceles \drawLine[bottom]{BC,AC,AB} se os lados iguais forem prolongados, os ângulos externos na base são iguais, e os ângulos internos na base também são iguais.}

\begin{center}
Prolongue \drawUnitLine{AB} e \drawUnitLine{AC} \byref{post:I.II},\\
tome $\drawUnitLine{BD} = \drawUnitLine{CE}$ \byref{prop:I.III};\\
trace \drawUnitLine{BE} e \drawUnitLine{CD}.

Então em
\drawFromCurrentPicture{
startAutoLabeling;
draw byNamedAngle(BAC);
draw byNamedLineSeq(0)(BE,CE,AC,AB);
stopAutoLabeling;
}
e
\drawFromCurrentPicture{
startAutoLabeling;
draw byNamedAngle(BAC);
draw byNamedLineSeq(0)(BD,CD,AC,AB);
stopAutoLabeling;
}\\
temos $\drawUnitLine{AB,BD} = \drawUnitLine{AC,CE}$ \byref{\constref},\\
\drawAngle{BAC} comum a ambos,\\
e $\drawUnitLine{AB} = \drawUnitLine{AC}$ \byref{\hypref}\\
$\therefore \drawAngle{BCA,DCB} = \drawAngle{ABC,CBE}$, $\drawUnitLine{BE} = \drawUnitLine{CD}$ e $\drawAngle{CEB} = \drawAngle{BDC}$ \byref{prop:I.IV}.

Novamente em \drawLine{BE,CE,BC} e \drawLine{BD,CD,BC}\\
temos $\drawUnitLine{BD} = \drawUnitLine{CE}$, $\drawAngle{CEB} = \drawAngle{BDC}$ e $\drawUnitLine{BE} = \drawUnitLine{CD}$,\\
$\therefore \drawAngle{DCE,DCB} = \drawAngle{EBD,CBE}$ e $\drawAngle{DCB} = \drawAngle{CBE}$ \byref{prop:I.IV}\\
mas $\drawAngle{BCA,DCB} = \drawAngle{ABC,CBE}$, $\therefore \drawAngle{BCA} = \drawAngle{ABC}$ \byref{ax:I.III}.
\end{center}

\qed

\starttheorem{Prop VI. Theor.}\label{prop:I.VI}
\defineNewPicture[1/4]{
pair A, B, C, D;
A := (0, 0);
B := A shifted (7/2u, 0);
D := 1/2[A,B] shifted (0, 3u);
C := 2/3[A, D];
byAngleDefine(B, A, D, byyellow, SOLID_SECTOR);
byAngleDefine(A, B, D, byblack, SOLID_SECTOR);
draw byNamedAngleResized();
byLineDefine(B, C, byyellow, SOLID_LINE, REGULAR_WIDTH);
byLineDefine(A, B, byred, SOLID_LINE, REGULAR_WIDTH);
byLineDefine(B, D, byblue, SOLID_LINE, REGULAR_WIDTH);
byLineDefine(C, A, byblack, SOLID_LINE, REGULAR_WIDTH);
byLineDefine(C, D, byblack, DASHED_LINE, REGULAR_WIDTH);
draw byNamedLine(BC);
draw byNamedLineSeq(0)(CA,CD,BD,AB);
draw byLabelsOnPolygon(A, C, D, B)(ALL_LABELS, 0);
}
\drawCurrentPictureInMargin
\problem{E}{m}{qualquer triângulo (\drawLine[bottom][triangleABD]{CA,CD,BD,AB}) se dois ângulos (\drawAngle{A} e \drawAngle{B}) são iguais, os lados (\drawUnitLine{CA,CD} e \drawUnitLine{BD}) opostos a eles também são iguais.}

Pois se os lados não forem iguais, seja um deles \drawUnitLine{CA,CD} maior que o outro \drawUnitLine{BD}, e dele corte $\drawUnitLine{CA} = \drawUnitLine{BD}$ \byref{prop:I.III}, trace \drawUnitLine{BC}.

\begin{center}
Então em \drawLine[bottom]{BC,AB,CA} e \triangleABD,\\
$\drawUnitLine{CA} = \drawUnitLine{BD}$ \byref{\constref},\\
$\drawAngle{A} = \drawAngle{B}$ \byref{\hypref}\\
e \drawUnitLine{AB} comum,\\
$\therefore$ os triângulos são iguais \byref{prop:I.IV}\\
uma parte igual ao todo, o que é absurdo;\\
$\therefore$ nenhum dos lados \drawUnitLine{CA,CD} ou \drawUnitLine{BD} é maior que o outro,\\
portanto eles são iguais.
\end{center}

\qed

\starttheorem{Prop VII. Theor.}\label{prop:I.VII}
\defineNewPicture{
pair A, B, C, D, E, F, G, H;
A := (0, 0);
B := A shifted (4u, 0);
C := A shifted (u, 3u);
D := C shifted (7/4u, 0);
E := 1/2[C, D] yscaled -0.7;
F := E shifted (0, -2u);
G := 5/4[A, E];
H := 5/4[A, F];
byAngleDefine.C(B, C, A, byblack, SOLID_SECTOR);
byAngleDefine(D, C, B, byred, SOLID_SECTOR);
byAngleDefine.D(A, D, B, byyellow, SOLID_SECTOR);
byAngleDefine(C, D, A, byblue, SOLID_SECTOR);
byAngleDefine(B, F, H, byblack, SOLID_SECTOR);
byAngleDefine(B, F, E, byred, SOLID_SECTOR);
byAngleDefine(B, E, G, byyellow, SOLID_SECTOR);
byAngleDefine(G, E, F, byblue, SOLID_SECTOR);
draw byNamedAngleResized();
draw byLine(C, D, byblack, DASHED_LINE, REGULAR_WIDTH);
draw byLine(E, F, byblack, DASHED_LINE, REGULAR_WIDTH);
draw byLine(A, B, byblack, SOLID_LINE, REGULAR_WIDTH);
byLineDefine(B, C, byblue, SOLID_LINE, REGULAR_WIDTH);
byLineDefine(C, A, byred, SOLID_LINE, REGULAR_WIDTH);
byLineDefine(B, D, byblue, SOLID_LINE, REGULAR_WIDTH);
byLineDefine(D, A, byred, SOLID_LINE, REGULAR_WIDTH);
byLineDefine(B, E, byblue, SOLID_LINE, REGULAR_WIDTH);
byLineDefine(E, A, byred, SOLID_LINE, REGULAR_WIDTH);
byLineDefine(B, F, byblue, SOLID_LINE, REGULAR_WIDTH);
byLineDefine(F, A, byred, SOLID_LINE, REGULAR_WIDTH);
byLineDefine(E, G, byred, DASHED_LINE, REGULAR_WIDTH);
byLineDefine(F, H, byred, DASHED_LINE, REGULAR_WIDTH);
draw byNamedLine(EG,FH);
draw byNamedLineSeq(0)(BC,CA,EA,BE);
draw byNamedLineSeq(0)(BD,DA,FA,BF);
byPointLabelDefine(F, "C");
byPointLabelDefine(E, "D");
draw byLabelsOnPolygon(F, A, C, D, B, F, noPoint)(OMIT_FIRST_LABEL+OMIT_LAST_LABEL, 0);
draw byLabelsOnPolygon(A, E, B)(OMIT_FIRST_LABEL+OMIT_LAST_LABEL, 0);
draw byLabelsOnPolygon(H, F, A)(OMIT_FIRST_LABEL+OMIT_LAST_LABEL, 0);
}
\drawCurrentPictureInMargin
\problem{S}{obre}{a mesma base (\drawUnitLine{AB}), e do mesmo lado dela não podem existir dois triângulos tendo seus lados conterminantes (\drawUnitLine{CA} e \drawUnitLine{DA}, \drawUnitLine{BC} e \drawUnitLine{BD}) em ambas as extremidades da base, iguais entre si.}

Quando dois triângulos estão sobre a mesma base, e do mesmo lado dela, o vértice de um cairá ou fora do outro triângulo, ou dentro dele; ou, por fim, em um de seus lados.

Se for possível, sejam os dois triângulos construídos de modo que
$\left\{
	\begin{aligned}
		\drawUnitLine{CA}&=\drawUnitLine{DA}\\
		\drawUnitLine{BC}&=\drawUnitLine{BD}\\
	\end{aligned}
	\right\}$
, então trace \drawUnitLine{CD} e,

\begin{center}
$\drawAngle{C,DCB} = \drawAngle{CDA}$ \byref{prop:I.V}

$\therefore \drawAngle{DCB} < \drawAngle{CDA}$ e

$\left.
	\begin{aligned}
		\therefore\drawAngle{DCB} &< \drawAngle{CDA,D}\\
		\mbox{mas \byref{prop:I.V}} \drawAngle{DCB} &= \drawAngle{CDA,D}\\
	\end{aligned}
	\right\}\mbox{o que é absurdo,}$
\end{center}

\noindent portanto os dois triângulos não podem ter seus lados conterminantes iguais em ambas as extremidades da base.

\qed

\starttheorem{Prop VIII. Theor.}\label{prop:I.VIII}
\defineNewPicture{
pair A, B, C, D, E, F, d;
A := (0, 0);
B := A shifted (-u, -4u);
C := A shifted (3/2u, -3u);
d := (0, -9/2u);
D := A shifted d;
E := B shifted d;
F := C shifted d;
byAngleDefine(F, D, E, byblack, SOLID_SECTOR);
byAngleDefine(C, A, B, byblack, SOLID_SECTOR);
draw byNamedAngleResized();
byLineDefine(A, B, byred, SOLID_LINE, REGULAR_WIDTH);
byLineDefine(B, C, byblack, SOLID_LINE, REGULAR_WIDTH);
byLineDefine(C, A, byblue, SOLID_LINE, REGULAR_WIDTH);
byLineDefine(D, E, byred, SOLID_LINE, THIN_WIDTH);
byLineDefine(E, F, byblack, SOLID_LINE, THIN_WIDTH);
byLineDefine(F, D, byblue, SOLID_LINE, THIN_WIDTH);
draw byNamedLineSeq(0)(CA,BC,AB);
draw byNamedLineSeq(0)(FD,EF,DE);
draw byLabelsOnPolygon(C, B, A)(ALL_LABELS, 0);
draw byLabelsOnPolygon(F, E, D)(ALL_LABELS, 0);
}
\drawCurrentPictureInMargin
\problem{S}{e}{dois triângulos têm dois lados de um respectivamente iguais a dois lados do outro ($\drawUnitLine{CA} = \drawUnitLine{FD}$ e $\drawUnitLine{AB} = \drawUnitLine{DE}$) e também suas bases ($\drawUnitLine{BC} = \drawUnitLine{EF}$), iguais; então os ângulos
(\drawFromCurrentPicture{
startAutoLabeling;
startGlobalRotation(-getAttribute("angle","Direction","A"));
draw byNamedAngleWithDummySides(A);
stopGlobalRotation;
stopAutoLabeling;
} e
\drawFromCurrentPicture{
startAutoLabeling;
startGlobalRotation(-getAttribute("angle","Direction","D"));
draw byNamedAngleWithDummySides(D);
stopGlobalRotation;
stopAutoLabeling;
}) contidos por seus lados iguais também são iguais.}

Se as bases iguais \drawUnitLine{BC} e \drawUnitLine{EF} forem concebidas como colocadas uma sobre a outra, de modo que os triângulos fiquem do mesmo lado delas, e que os lados iguais \drawUnitLine{AB} e \drawUnitLine{DE}, \drawUnitLine{CA} e \drawUnitLine{FD} sejam conterminantes, o vértice de um deve cair sobre o vértice do outro; pois supor que não sejam coincidentes contradiria a última proposição.

Portanto os lados \drawUnitLine{AB} e \drawUnitLine{CA}, sendo coincidentes com \drawUnitLine{DE} e \drawUnitLine{FD}, $\therefore \drawAngle{A} = \drawAngle{D}$.

\qed

\startproblem{Prop IX. Prob.}\label{prop:I.IX}
\defineNewPicture{
pair A, B, C, D, E, F;
A := (0, 5/3u);
B := (-4/3u, 0);
C := B xscaled -1;
D = whatever[B, B shifted ((C-B) rotated -60)] = whatever[C, C shifted ((B-C) rotated 60)];
E := 5/4[A, B];
F := 5/4[A, C];
byAngleDefine(B, A, D, byblue, SOLID_SECTOR);
byAngleDefine(C, A, D, byyellow, SOLID_SECTOR);
draw byNamedAngleResized();
byLineDefine(B, C, byyellow, SOLID_LINE, REGULAR_WIDTH);
byLineDefine(A, D, byblack, SOLID_LINE, REGULAR_WIDTH);
byLineDefine(D, B, byblue, SOLID_LINE, REGULAR_WIDTH);
byLineDefine(C, D, byblue, SOLID_LINE, REGULAR_WIDTH);
byLineDefine(A, B, byred, SOLID_LINE, REGULAR_WIDTH);
byLineDefine(C, A, byred, SOLID_LINE, REGULAR_WIDTH);
byLineDefine(B, E, byred, DASHED_LINE, REGULAR_WIDTH);
byLineDefine(C, F, byred, DASHED_LINE, REGULAR_WIDTH);
draw byNamedLine(BC,AD);
draw byNamedLineSeq(0)(DB,CD);
draw byNamedLineSeq(0)(BE,AB,CA,CF);
draw byLabelsOnPolygon(D, B, A, C)(ALL_LABELS, 0);
}
\drawCurrentPictureInMargin
\problem{B}{issectar}{um ângulo retilíneo dado (\drawAngle{BAD,CAD}).}

\begin{center}
Tome $\drawUnitLine{AB} = \drawUnitLine{CA}$ \byref{prop:I.III}

trace \drawUnitLine{BC}, sobre o qual descreva \drawLine{CD,DB,BC} \byref{prop:I.I},\\
trace \drawUnitLine{AD}.

$\because \drawUnitLine{AB} = \drawUnitLine{CA}$ \byref{\constref}\\ %Because
e \drawUnitLine{AD} comum aos dois triângulos\\  % what triangles?
e $\drawUnitLine{CD} = \drawUnitLine{DB}$ \byref{\constref},

$\therefore \drawAngle{BAD} = \drawAngle{CAD}$ \byref{prop:I.VIII}.
\end{center}

\qed

\startproblem{Prop X. Prob.}\label{prop:I.X}
\defineNewPicture{
pair A, B, C, D;
A := (0, 3u);
B := (-ypart(A)/sqrt(3), 0);
C := B xscaled -1;
D := 1/2[B, C];
byAngleDefine(B, A, D, byblue, SOLID_SECTOR);
byAngleDefine(C, A, D, byyellow, SOLID_SECTOR);
draw byNamedAngleResized();
draw byLine(A, D, byred, SOLID_LINE, REGULAR_WIDTH);
byLineDefine(D, B, byblack, SOLID_LINE, REGULAR_WIDTH);
byLineDefine(C, D, byblack, DASHED_LINE, REGULAR_WIDTH);
byLineDefine(A, B, byyellow, SOLID_LINE, REGULAR_WIDTH);
byLineDefine(C, A, byblue, SOLID_LINE, REGULAR_WIDTH);
draw byNamedLineSeq(0)(AB,CA,CD,DB);
draw byLabelsOnPolygon(B, A, C, D)(ALL_LABELS, 0);
}
\drawCurrentPictureInMargin
\problem{B}{issectar}{uma linha reta finita dada (\drawUnitLine{DB,CD}).}

\begin{center}
Construa \drawLine[bottom]{AB,CA,CD,DB} \byref{prop:I.I},\\
trace \drawUnitLine{AD}, fazendo $\drawAngle{BAD} = \drawAngle{CAD}$ \byref{prop:I.IX}.

Então $\drawUnitLine{BD} = \drawUnitLine{DC}$ por \byref{prop:I.IV},

pois~$\drawUnitLine{AB} = \drawUnitLine{AC}$ \byref{\constref} $\drawAngle{BAD} = \drawAngle{CAD}$\\
e \drawUnitLine{AD} comum aos dois triângulos. % what triangles? obviously DBA and DAC

Portanto a linha dada é bissectada.
\end{center}

\qed

\startproblem{Prop XI. Prob.}\label{prop:I.XI}
\defineNewPicture{
pair A, B, C, D, E, F;
A := (0, 5/2u);
B := (-ypart(A)/sqrt(3), 0);
C := B xscaled -1;
D := 1/2[B, C];
E := 3/2[D, B];
F := 3/2[D, C];
byAngleDefine(A, D, B, byred, SOLID_SECTOR);
byAngleDefine(C, D, A, byblue, SOLID_SECTOR);
draw byNamedAngleResized();
draw byLine(A, D, byyellow, SOLID_LINE, REGULAR_WIDTH);
byLineDefine(A, B, byblue, SOLID_LINE, REGULAR_WIDTH);
byLineDefine(C, A, byblue, SOLID_LINE, REGULAR_WIDTH);
draw byNamedLineSeq(0)(AB,CA);
byLineDefine(D, B, byblack, SOLID_LINE, REGULAR_WIDTH);
byLineDefine(B, E, byblack, DASHED_LINE, REGULAR_WIDTH);
byLineDefine(C, D, byred, SOLID_LINE, REGULAR_WIDTH);
byLineDefine(F, C, byred, DASHED_LINE, REGULAR_WIDTH);
draw byNamedLineSeq(0)(BE,DB,CD,FC);
draw byLabelsOnPolygon(F, C, D, B, E)(OMIT_FIRST_LABEL+OMIT_LAST_LABEL, 0);
draw byLabelsOnPolygon(B, A, C)(OMIT_FIRST_LABEL+OMIT_LAST_LABEL, 0);
}
\drawCurrentPictureInMargin
\problem{D}{e}{um ponto dado (\drawPointL[middle][AD]{D}), em uma linha reta dada (\drawUnitLine[3/2cm]{BD,DC}), traçar uma perpendicular.}

\begin{center}
Tome qualquer ponto (\drawPointL[middle][CA]{C}) na linha dada,\\
corte $\drawUnitLine{DB} = \drawUnitLine{CD}$ \byref{prop:I.III},\\
construa \drawLine[bottom]{AB,CA,CD,DB} \byref{prop:I.I},\\
trace \drawUnitLine{AD} e ela será perpendicular à linha dada.

Pois $\drawUnitLine{AB} = \drawUnitLine{CA}$ \byref{\constref}\\
$\drawUnitLine{CD} = \drawUnitLine{DB}$ \byref{\constref}\\
e \drawUnitLine{AD} comum aos dois triângulos. % what triangles? obviously DBA and DAC

Portanto $\drawAngle{ADB} = \drawAngle{CDA}$ \byref{prop:I.VIII}

$\therefore \drawUnitLine{AD} \perp \drawUnitLine{DB,CD}$ \byref{def:I.X}.
\end{center}

\qed

\startproblem{Prop XII. Prob.}\label{prop:I.XII}
\defineNewPicture{
pair A, B, C, D, E, F;
path c;
numeric r, a[];
A := (0, 2u);
B := (-7/4u, 0);
C := B xscaled -1;
D := 1/2[B, C];
E := 4/3[D, B];
F := 4/3[D, C];
r := abs(A-B);
c := fullcircle scaled 2r shifted A;
a1 := xpart(c intersectiontimes (F--1/2[B, C]));
a2 := xpart(c intersectiontimes (E--1/2[B, C]));
byAngleDefine(A, D, B, byyellow, SOLID_SECTOR);
byAngleDefine(C, D, A, byblue, SOLID_SECTOR);
draw byNamedAngleResized();
draw byLine(A, D, byred, SOLID_LINE, REGULAR_WIDTH);
byLineDefine(A, B, byblue, SOLID_LINE, REGULAR_WIDTH);
byLineDefine(C, A, byblue, SOLID_LINE, REGULAR_WIDTH);
draw byNamedLineSeq(0)(AB,CA);
draw byArc.O(A, B, C)(r, byred, 0, 0, 0, 0);
draw byArcBE.Ol(A, a2-1/4, a2, r, byred, 1, 0, 0, 0);
draw byArcBE.Or(A, a1, a1+1/4, r, byred, 1, 0, 0, 0);
byLineDefine(D, B, byblack, SOLID_LINE, REGULAR_WIDTH);
byLineDefine(B, E, byblack, DASHED_LINE, REGULAR_WIDTH);
byLineDefine(C, D, byyellow, SOLID_LINE, REGULAR_WIDTH);
byLineDefine(F, C, byyellow, DASHED_LINE, REGULAR_WIDTH);
draw byNamedLineSeq(1)(BE, DB, CD,FC);
draw byLabelsOnPolygon(B, A, C)(OMIT_FIRST_LABEL+OMIT_LAST_LABEL, 0);
draw byLabelLineEnd(B, A, 0);
draw byLabelLineEnd(D, A, 0);
draw byLabelLineEnd(C, A, 0);
}
\drawCurrentPictureInMargin
\problem{D}{esenhar}{uma linha reta perpendicular a uma linha reta indefinida dada (\drawUnitLine[1.2cm]{DB,CD}) de um ponto dado (\drawPointL{A}) fora dela.}

\begin{center}
Com o ponto dado \drawPointL{A} como centro, de um lado da linha, e qualquer distância \drawUnitLine{DB} capaz de se estender ao outro lado, descreva \drawArc{O}. % improvement: the distance mentioned here seems not to be in the original drawing, it should either be drawn, or not be referenced graphically as DB

Faça $\drawUnitLine{DB} = \drawUnitLine{CD}$ \byref{prop:I.X},\\
trace \drawUnitLine{AB}, \drawUnitLine{CA} e \drawUnitLine{AD},\\
então $\drawUnitLine{AD} \perp \drawUnitLine{DB,CD}$.

Pois \byref{prop:I.VIII} como $\drawUnitLine{DB} = \drawUnitLine{CD}$ \byref{\constref},\\
\drawUnitLine{AD} comum a ambos,\\ % to both what? obviously DBA and DAC
e $\drawUnitLine{AB} = \drawUnitLine{CA}$ \byref{def:I.XV},

$\therefore \drawAngle{ADB} = \drawAngle{CDA}$, e

$\therefore \drawUnitLine{AD} \perp \drawUnitLine{DB,CD}$ \byref{def:I.X}.
\end{center}

\qed

\starttheorem{Prop XIII. Theor.}\label{prop:I.XIII}
\defineNewPicture{
pair A, B, C, D, E;
A := (0, 5/2u);
B := (-7/4u, 0);
C := B xscaled -1;
D := (xpart(A), ypart(B));
E := (2/3xpart(C), 2/3ypart(A));
byAngleDefine(A, D, B, byyellow, SOLID_SECTOR);
byAngleDefine(E, D, A, byred, SOLID_SECTOR);
byAngleDefine(C, D, E, byblue, SOLID_SECTOR);
draw byNamedAngleResized();
draw byLine(A, D, byblack, SOLID_LINE, REGULAR_WIDTH);
draw byLine(E, D, byyellow, SOLID_LINE, REGULAR_WIDTH);
draw byLine(B, C, byred, SOLID_LINE, REGULAR_WIDTH);
draw byLabelsOnPolygon(C, D, B, noPoint)(ALL_LABELS, 0);
draw byLabelLineEnd(E, D, 0);
draw byLabelLineEnd(A, D, 0);
}
\drawCurrentPictureInMargin
\problem{Q}{uando}{uma linha reta (\drawUnitLine{ED}) apoiada sobre outra linha reta (\drawUnitLine{BC}) forma ângulos com ela; eles são ou dois ângulos retos ou juntos iguais a dois ângulos retos.}

\begin{center}
Se \drawUnitLine{ED} for $\perp$ a \drawUnitLine{BC} então,\\
\drawAngle{ADB,EDA} e $\drawAngle{CDE} = \drawTwoRightAngles$ \byref{def:I.X}, % improvement: def. 7 in the original is likely to be a mistake

mas se \drawUnitLine{ED} não for $\perp$ a \drawUnitLine{BC},\\
trace $\drawUnitLine{AD} \perp \drawUnitLine{BC}$ \byref{prop:I.XI};\\
$\drawAngle{ADB} +\drawAngle{CDE,EDA} = \drawTwoRightAngles$ \byref{\constref},\\
$\drawAngle{ADB} = \drawAngle{CDE,EDA} = \drawAngle{EDA} + \drawAngle{CDE}$

$\therefore \drawAngle{ADB} + \drawAngle{CDE,EDA} = \drawAngle{ADB} + \drawAngle{EDA} + \drawAngle{CDE}$ \byref{ax:I.II}

$= \drawAngle{ADB,EDA} + \drawAngle{CDE} = \drawTwoRightAngles$.
\end{center}

\qed

\starttheorem{Prop XIV. Theor.}\label{prop:I.XIV}
\defineNewPicture[1/4]{
pair A, B, C, D, E;
A := (u, 5/2u);
B := (-7/4u, 0);
C := B xscaled -1;
D := (0, 0);
E := (xpart(C), -1/2ypart(A));
byAngleDefine(B, D, A, byyellow, SOLID_SECTOR);
byAngleDefine(C, D, A, byblue, SOLID_SECTOR);
byAngleDefine(E, D, C, byred, SOLID_SECTOR);
draw byNamedAngleResized();
draw byLine(A, D, byred, SOLID_LINE, REGULAR_WIDTH);
draw byLine(E, D, byyellow, SOLID_LINE, REGULAR_WIDTH);
draw byLine(B, D, byblue, SOLID_LINE, REGULAR_WIDTH);
draw byLine(C, D, byblack, SOLID_LINE, REGULAR_WIDTH);
draw byLabelsOnPolygon(E, D, B, noPoint)(ALL_LABELS, 0);
draw byLabelsOnPolygon(C, B, noPoint)(OMIT_LAST_LABEL, 0);
draw byLabelLineEnd(A, D, 0);
}
\drawCurrentPictureInMargin
\problem{S}{e}{duas linhas retas (\drawUnitLine{BD} e \drawUnitLine{DC}), encontrando uma terceira linha reta (\drawUnitLine{AD}), no mesmo ponto, e em lados opostos dela, formam com ela ângulos adjacentes (\offsetPicture{0pt}{15pt}{\drawAngle{BDA}} e \drawAngle{CDA}) iguais a dois ângulos retos; essas linhas retas estão em uma linha reta contínua.}

\begin{center}
Pois, se possível seja \drawUnitLine{ED}, e não \drawUnitLine{DC},\\
a continuação de \drawUnitLine{BD},\\
então $\drawAngle{BDA} + \drawAngle{CDA,EDC} = \drawTwoRightAngles$

mas pela hipótese $\drawAngle{BDA} + \drawAngle{CDA} = \drawTwoRightAngles$

$\therefore\drawAngle{CDA,EDC} = \drawAngle{CDA}$ \byref{ax:I.III}; o que é absurdo \byref{ax:I.IX}.

$\therefore \drawUnitLine{ED}$ não é a continuação de \drawUnitLine{BD}, e o mesmo pode ser demonstrado de qualquer outra linha reta exceto \drawUnitLine{DC}, $\therefore \drawUnitLine{DC}$ é a continuação de \drawUnitLine{BD}.
\end{center}

\qed

\starttheorem{Prop XV. Theor.}\label{prop:I.XV}
\defineNewPicture{
pair A, B, C, D, E;
A := (7/4u, 3/2u);
B := A scaled -1;
C := A xscaled -1;
D := C scaled -1;
E := (A--B) intersectionpoint (C--D);
byAngleDefine(B, E, C, byyellow, SOLID_SECTOR);
byAngleDefine(C, E, A, byred, SOLID_SECTOR);
byAngleDefine(A, E, D, byblack, SOLID_SECTOR);
byAngleDefine(D, E, B, byblue, SOLID_SECTOR);
draw byNamedAngleResized();
draw byLine(A, B, byred, SOLID_LINE, REGULAR_WIDTH);
draw byLine(C, D, byblack, SOLID_LINE, REGULAR_WIDTH);
draw byLabelsOnPolygon(C, E, A, noPoint)(ALL_LABELS, 0);
draw byLabelPoint(B, lineAngle.AB + 90, 1);
draw byLabelPoint(D, lineAngle.CD - 90, 1);
}
\drawCurrentPictureInMargin
\problem{S}{e}{duas linhas retas (\drawUnitLine{AB} e \drawUnitLine{CD}) se intersectam, os ângulos opostos pelo vértice \drawAngle{BEC} e \drawAngle{AED}, \drawAngle{CEA} e \drawAngle{DEB} são iguais.}

\begin{center}
$\drawAngle{BEC} + \drawAngle{CEA} = \drawTwoRightAngles$ % Here and in the next line should be a link to prop I.XIII

$\drawAngle{AED} + \drawAngle{CEA} = \drawTwoRightAngles$ %

$\therefore \drawAngle{BEC} = \drawAngle{AED}$. % And this is according to \byref{ax:I.III}

Da mesma forma pode ser mostrado que\\
$\drawAngle{CEA} = \drawAngle{DEB}$.
\end{center}

\qed

\starttheorem{Prop XVI. Theor.}\label{prop:I.XVI}
\defineNewPicture[1/4]{
pair A, B, C, D, E, F, G;
A := (0, 0);
B := A shifted (u, 7/2u);
C := A shifted (3u, 0);
D := B shifted (3u, 0);
E = whatever[A, D] = whatever[B, C];
F := (xpart(D), ypart(A));
G := 4/3[B, C];
byAngleDefine(B, A, C, byblue, SOLID_SECTOR);
byAngleDefine(C, B, A, byblack, SOLID_SECTOR);
byAngleDefine(A, E, B, byyellow, SOLID_SECTOR);
byAngleDefine(D, E, C, byyellow, SOLID_SECTOR);
byAngleDefine(E, C, D, byblack, SOLID_SECTOR);
byAngleDefine(G, C, A, byred, SOLID_SECTOR);
byAngleDefine(D, C, F, byblack, ARC_SECTOR);
draw byNamedAngleResized();
byLineDefine(C, F, byblack, DASHED_LINE, REGULAR_WIDTH);
byLineDefine(C, G, byblack, SOLID_LINE, REGULAR_WIDTH);
byLineDefine(B, E, byblue, SOLID_LINE, REGULAR_WIDTH);
byLineDefine(E, C, byblue, DASHED_LINE, REGULAR_WIDTH);
byLineDefine(A, E, byred, SOLID_LINE, REGULAR_WIDTH);
byLineDefine(E, D, byred, DASHED_LINE, REGULAR_WIDTH);
byLineDefine(A, B, byyellow, DASHED_LINE, REGULAR_WIDTH);
byLineDefine(A, C, byblack, SOLID_LINE, REGULAR_WIDTH);
byLineDefine(C, D, byyellow, SOLID_LINE, REGULAR_WIDTH);
draw byNamedLineSeq(0)(AE,ED,CD);
draw byNamedLineSeq(0)(EC,CG,noLine,CF,AC,AB,BE);
draw byLabelsOnPolygon(F, A, B, E, D, C)(OMIT_FIRST_LABEL+OMIT_LAST_LABEL, 0);
draw byLabelsOnPolygon(F, C, G, noPoint)(ALL_LABELS, 0);
}
\drawCurrentPictureInMargin
\problem{S}{e}{um lado de um triângulo (\drawLine[bottom]{BE,EC,AC,AB}) é prolongado, o ângulo externo (\drawFromCurrentPicture[middle][anglesECDpDCF]{
startAutoLabeling;
draw byNamedAngleSides(ECD,DCF)(CF);
stopAutoLabeling;
}) é maior que qualquer dos ângulos internos remotos (\drawAngle{B} ou \drawAngle{A}).
}

\begin{center}
Faça $\drawUnitLine{BE} = \drawUnitLine{EC}$ \byref{prop:I.X};\\
Trace \drawUnitLine{AE} e prolongue até $\drawUnitLine{ED} = \drawUnitLine{AE}$;\\
trace \drawUnitLine{CD}.

Em \drawLine{BE,AE,AB} e \drawLine{EC,ED,CD};\\
$\drawUnitLine{BE} = \drawUnitLine{EC}$, $\drawAngle{AEB} = \drawAngle{DEC}$ e $\drawUnitLine{AE} = \drawUnitLine{ED}$ \byref{\constref,prop:I.XV},\\
$\therefore \drawAngle{B} = \drawAngle{ECD}$ \byref{prop:I.IV},\\
$\therefore \anglesECDpDCF\ > \drawAngle{B}$.

Da mesma forma pode ser mostrado que, se \drawUnitLine{BC} for prolongado, $\drawAngle{GCA} > \drawAngle{A}$ \\
e portanto \anglesECDpDCF\ que é $= \drawAngle{GCA}$ é $> \drawAngle{A}$.
\end{center}

\qed

\starttheorem{Prop XVII. Theor.}\label{prop:I.XVII}
\defineNewPicture{
pair A, B, C, D;
A := (0, 0);
B := A shifted (2u, 5/2u);
C := A shifted (5/2u, 0);
D := C shifted (3/4u, 0);
byAngleDefine(B, A, C, byblue, SOLID_SECTOR);
byAngleDefine(A, B, C, byblack, SOLID_SECTOR);
byAngleDefine(A, C, B, byred, SOLID_SECTOR);
byAngleDefine(B, C, D, byyellow, SOLID_SECTOR);
draw byNamedAngleResized();
byLineDefine(A, B, byred, SOLID_LINE, REGULAR_WIDTH);
byLineDefine(B, C, byblue, SOLID_LINE, REGULAR_WIDTH);
byLineDefine(A, C, byblack, SOLID_LINE, REGULAR_WIDTH);
byLineDefine(C, D, byblack, SOLID_LINE, REGULAR_WIDTH);
draw byNamedLineSeq(0)(noLine,BC,AB,AC,CD);
draw byLabelsOnPolygon(D, C, A, B)(ALL_LABELS, 0);
}
\drawCurrentPictureInMargin
\problem{Q}{uaisquer}{dois ângulos de um triângulo \drawLine[bottom]{AB,BC,AC} juntos são menores que dois ângulos retos.}

\begin{center}
Prolongue \drawUnitLine{AC}, então\\
$\drawAngle{ACB} + \drawAngle{BCD} = \drawTwoRightAngles$.

Mas $\drawAngle{BCD} > \drawAngle{A}$ \byref{prop:I.XVI}

$\therefore \drawAngle{ACB} + \drawAngle{A} < \drawTwoRightAngles$,
\end{center}

\noindent e da mesma forma pode ser mostrado que quaisquer outros dois ângulos do triângulo tomados juntos são menores que dois ângulos retos.

\qed

\starttheorem{Prop XVIII. Theor.}\label{prop:I.XVIII}
\defineNewPicture[1/4]{
pair A, B, C, D;
numeric a;
A := (0, 0);
B := A shifted (5/2u, 2u);
C := B shifted (-3/2u, 2u);
D := C shifted (unitvector(A-C) scaled abs(B-C));
a := angle(B-D);
forsuffixes i=A, B, C, D:
i := i rotated -a;
endfor;
byAngleDefine(C, D, B, byblue, SOLID_SECTOR);
byAngleDefine(D, B, C, byblack, SOLID_SECTOR);
byAngleDefine(A, B, D, byred, SOLID_SECTOR);
byAngleDefine(B, A, D, byyellow, SOLID_SECTOR);
draw byNamedAngleResized();
draw byLine(D, B, byyellow, SOLID_LINE, REGULAR_WIDTH);
byLineDefine(D, C, byred, SOLID_LINE, REGULAR_WIDTH);
byLineDefine(B, C, byblue, SOLID_LINE, REGULAR_WIDTH);
byLineDefine(B, A, byblack, SOLID_LINE, REGULAR_WIDTH);
byLineDefine(A, D, byred, DASHED_LINE, REGULAR_WIDTH);
draw byNamedLineSeq(0)(DC,BC,BA,AD);
draw byLabelsOnPolygon(D, C, B, A)(ALL_LABELS, 0);
}
\drawCurrentPictureInMargin
\problem[3]{E}{m}{qualquer triângulo \drawLine{DC,BC,BA,AD} se um lado \drawUnitLine{AD,DC} for maior que outro \drawUnitLine{BC}, o ângulo oposto ao lado maior é maior que o ângulo oposto ao menor. I.\ e.\ $\drawAngle{DBC,ABD} > \drawAngle{A}$.}

\begin{center}
Faça $\drawUnitLine{DC} = \drawUnitLine{BC}$ \byref{prop:I.III}, trace \drawUnitLine{DB}.

Então $\drawAngle{D} = \drawAngle{DBC}$ \byref{prop:I.V};

mas $\drawAngle{D} > \drawAngle{A}$ \byref{prop:I.XVI}

$\therefore \drawAngle{DBC} > \drawAngle{A}$ e muito mais\\
é $\drawAngle{DBC,ABD} > \drawAngle{A}$.
\end{center}

\qed

\starttheorem{Prop XIX. Theor.}\label{prop:I.XIX}
\defineNewPicture[1/4]{
pair A, B, C;
A := (0, 0);
B := A shifted (7/2u, 0);
C := A shifted (u, 3u);
byAngleDefine(C, A, B, byblue, SOLID_SECTOR);
byAngleDefine(A, B, C, byred, SOLID_SECTOR);
draw byNamedAngleResized();
byLineDefine(A, B, byblack, SOLID_LINE, REGULAR_WIDTH);
byLineDefine(B, C, byblue, SOLID_LINE, REGULAR_WIDTH);
byLineDefine(C, A, byred, SOLID_LINE, REGULAR_WIDTH);
draw byNamedLineSeq(0)(CA,BC,AB);
draw byLabelsOnPolygon(B, A, C)(ALL_LABELS, 0);
}
\drawCurrentPictureInMargin
\problem{S}{e}{em qualquer triângulo \drawLine[bottom]{CA,BC,AB} um ângulo \drawAngle{A} for maior que outro \drawAngle{B} o lado \drawUnitLine{BC} que é oposto ao ângulo maior, é maior que o lado \drawUnitLine{CA} oposto ao menor.}

\begin{center}
Se \drawUnitLine{BC} não for maior que \drawUnitLine{CA} então deve\\
$\drawUnitLine{BC} =$ ou $< \drawUnitLine{CA}$.

Se $\drawUnitLine{BC} = \drawUnitLine{CA}$ então\\
$\drawAngle{A} = \drawAngle{B}$ \byref{prop:I.V};\\
o que é contrário à hipótese.

\drawUnitLine{BC} não é menor que \drawUnitLine{CA}; pois se fosse,\\
$\drawAngle{A} < \drawAngle{B}$ \byref{prop:I.XVIII}\\
o que é contrário à hipótese:

$\therefore \drawUnitLine{BC} > \drawUnitLine{CA}$.
\end{center}

\qed

\starttheorem{Prop XX. Theor.}\label{prop:I.XX}
\defineNewPicture{
pair A, B, C, D;
A := (0, 0);
B := A shifted (7/2u, 0);
D := A shifted (4/3u, 3/2u);
C := ((fullcircle scaled 2arclength(D--B)) shifted D) intersectionpoint (D--10[A, D]);
byAngleDefine(B, C, A, byred, SOLID_SECTOR);
byAngleDefine(C, B, D, byblue, SOLID_SECTOR);
byAngleDefine(D, B, A, byyellow, SOLID_SECTOR);
draw byNamedAngleResized();
byLineDefine(B, D, byred, SOLID_LINE, REGULAR_WIDTH);
byLineDefine(A, B, byblack, SOLID_LINE, REGULAR_WIDTH);
byLineDefine(B, C, byyellow, SOLID_LINE, REGULAR_WIDTH);
byLineDefine(C, D, byblue, DASHED_LINE, REGULAR_WIDTH);
byLineDefine(D, A, byblue, SOLID_LINE, REGULAR_WIDTH);
draw byNamedLineSeq(0)(BD);
draw byNamedLineSeq(0)(DA,CD,BC,AB);
draw byLabelsOnPolygon(D, C, B, A)(ALL_LABELS, 0);
}
\drawCurrentPictureInMargin
\problem{Q}{uaisquer}{dois lados \drawUnitLine{DA} e \drawUnitLine{BD} de um triângulo \drawLine[bottom]{DA,BD,AB} tomados juntos são maiores que o terceiro lado (\drawUnitLine{AB}).}

\begin{center}
Prolongue \drawUnitLine{DA}, e\\
faça $\drawUnitLine{CD} = \drawUnitLine{BD}$ \byref{prop:I.III};\\
trace \drawUnitLine{BC}.

Então $\because \drawUnitLine{CD} = \drawUnitLine{BD}$ \byref{\constref},\\ %because
$\drawAngle{CBD} = \drawAngle{C}$ \byref{prop:I.V}

$\therefore \drawAngle{CBD,DBA} > \drawAngle{C}$ \byref{ax:I.IX}

$\therefore \drawUnitLine{DA} + \drawUnitLine{CD} > \drawUnitLine{AB}$ \byref{prop:I.XIX}

e $\therefore \drawUnitLine{DA} + \drawUnitLine{BD} > \drawUnitLine{AB}$.
\end{center}

\qed

\starttheorem{Prop XXI. Theor.}\label{prop:I.XXI}
\defineNewPicture[1/5]{
pair A, B, C, D, E;
A := (0, 0);
B := A shifted (7/2u, 0);
C := A shifted (3u, 4u);
D := 1/2[1/2[A, B], C];
E = whatever[A, D] = whatever[B, C];
byAngleDefine(B, D, A, byred, SOLID_SECTOR);
byAngleDefine(B, E, D, byblue, SOLID_SECTOR);
byAngleDefine(B, C, A, byyellow, SOLID_SECTOR);
draw byNamedAngleResized();
byLineDefine(B, D, byyellow, SOLID_LINE, REGULAR_WIDTH);
byLineDefine(A, D, byblack, SOLID_LINE, REGULAR_WIDTH);
byLineDefine(D, E, byblack, DASHED_LINE, REGULAR_WIDTH);
byLineDefine(A, B, byblue, DASHED_LINE, REGULAR_WIDTH);
byLineDefine(B, E, byred, DASHED_LINE, REGULAR_WIDTH);
byLineDefine(E, C, byred, SOLID_LINE, REGULAR_WIDTH);
byLineDefine(C, A, byblue, SOLID_LINE, REGULAR_WIDTH);
draw byNamedLine(BD);
draw byNamedLineSeq(0)(AD,DE);
draw byNamedLineSeq(0)(CA,EC,BE,AB);
draw byLabelsOnPolygon(A, C, E, B)(ALL_LABELS, 0);
draw byLabelsOnPolygon(A, D, E)(OMIT_FIRST_LABEL+OMIT_LAST_LABEL, 0);
}
\drawCurrentPictureInMargin
\problem[2]{S}{e}{de qualquer ponto (\drawPointL[middle][DE]{D}) dentro de um triângulo \drawLine[bottom]{CA,EC,BE,AB} linhas retas forem traçadas até as extremidades de um lado (\drawSizedLine{AB}), essas linhas devem ser juntas menores que os outros dois lados, mas devem conter um ângulo maior.}

\begin{center}
Prolongue \drawSizedLine{AD},\\
$\drawSizedLine{CA} + \drawSizedLine{EC} > \drawSizedLine{AD,DE}$ \byref{prop:I.XX},\\
adicione \drawSizedLine{BE} a cada um,\\
$\drawSizedLine{CA} + \drawSizedLine{EC,BE} > \drawSizedLine{AD,DE} + \drawSizedLine{BE}$ \byref{ax:I.IV}

da mesma forma pode ser mostrado que\\
$\drawSizedLine{AD,DE} + \drawSizedLine{BE} > \drawSizedLine{AD} + \drawSizedLine{BD}$,\\
$\therefore \drawSizedLine{CA} + \drawSizedLine{EC,BE} > \drawSizedLine{AD} + \drawSizedLine{BD}$,\\
o que devia ser provado.

Novamente $\drawAngle{E} > \drawAngle{C}$ \byref{prop:I.XVI},\\
e também $\drawAngle{D} > \drawAngle{E}$ \byref{prop:I.XVI},

$\therefore \drawAngle{D} > \drawAngle{C}$.
\end{center}

\qed

\startproblem{Prop XXII. Prob.}\label{prop:I.XXII}
\defineNewPicture[1/2]{
numeric r[], d;
pair A, B, C, D, E, LI, LII, LIII, LIV, LV, LVI;
path q[];
r1 := 2u;
r2 := 4/3u;
r3 := (1/2)*(r1+r2);
d := 1/3u;
A := (0, 0);
B := A shifted (r3, 0);
q1 := (fullcircle scaled 2r1) shifted A;
q2 := (fullcircle scaled 2r2) shifted B;
C := q1 intersectionpoint q2;
D := point 11/2 of q1;
E := point 3/4 of q2;
LI := (xpart(point 0 of q2), ypart(point 6 of q1) - 1/2d);
LII := LI shifted (-r3, 0);
LIII := LI shifted (0, -d);
LIV := LIII shifted (-r2, 0);
LV := LIII shifted (0, -d);
LVI := LV shifted (-r1, 0);
draw byCircle.A(A, D, byblue, 0, 0, 0);
byLineDefine(A, D, byblue, SOLID_LINE, REGULAR_WIDTH);
byLineDefine(B, E, byred, SOLID_LINE, REGULAR_WIDTH);
byLineDefine(A, B, byblack, SOLID_LINE, REGULAR_WIDTH);
byLineDefine(B, C, byyellow, SOLID_LINE, REGULAR_WIDTH);
byLineDefine(C, A, byyellow, DASHED_LINE, REGULAR_WIDTH);
draw byNamedLineSeq(0)(BC,CA);
draw byNamedLineSeq(0)(AD,AB,BE);
draw byLineWithName (LII, LI, byblack, 1, 0)(L');
draw byLineWithName (LIV, LIII, byred, 1, 0)(L'');
draw byLineWithName (LVI, LV, byblue, 1, 0)(L''');
draw byCircle.B(B, E, byred, 0, 0, 0);
draw byLabelsOnPolygon(D, A, C)(OMIT_FIRST_LABEL+OMIT_LAST_LABEL, 0);
draw byLabelsOnPolygon(E, B, A)(OMIT_FIRST_LABEL+OMIT_LAST_LABEL, 0);
draw byLabelsOnPolygon(A, C, B)(OMIT_FIRST_LABEL+OMIT_LAST_LABEL, 0);
draw byLabelsOnCircle(D)(A);
draw byLabelsOnCircle(E)(B);
draw byLabelLine(0)(L', L'', L''');
}
\drawCurrentPictureInMargin
\problem[4]{D}{adas}{três linhas retas $\left\{\vcenter{
\nointerlineskip\hbox{\drawSizedLine{L'}}
\nointerlineskip\hbox{\drawSizedLine{L''}}
\nointerlineskip\hbox{\drawSizedLine{L'''}}}\right.$
a soma de quaisquer duas maior que a terceira, construir um triângulo cujos lados sejam respectivamente iguais às linhas dadas.}

\begin{center}
Assuma $\drawSizedLine{AB} = \drawSizedLine{L'}$ \byref{prop:I.III}.

$\left.
	\begin{aligned}
		\mbox{Trace } \drawSizedLine{BE} &= \drawSizedLine{L''}\\
		\mbox{e } \drawSizedLine{AD} &= \drawSizedLine{L'''}\\
	\end{aligned}
	\right\}\mbox{\byref{prop:I.II}.}$

Com \drawSizedLine{AD} e \drawSizedLine{BE} como raios, descreva
\drawFromCurrentPicture{
draw byNamedLine(AD); draw byNamedCircle(A);
draw byLabelLineEnd(A, D, 0);
draw byLabelLineEnd(D, A, 0);
} e
\offsetPicture{12pt}{0pt}{\drawFromCurrentPicture{
draw byNamedLine(BE); draw byNamedCircle(B);
draw byLabelLineEnd(B, E, 0);
draw byLabelLineEnd(E, B, 0);
}} \byref{post:I.III};\\
trace \drawSizedLine{CA} e \drawSizedLine{BC},\\
então \drawLine[bottom]{CA,BC,AB} será o triângulo requerido.

$\left.
	\begin{aligned}
		\mbox{Pois } \drawSizedLine{AB} &= \drawSizedLine{L'} \mbox{,} \\
		\drawSizedLine{BC} &= \drawSizedLine{BE} = \drawSizedLine{L''} \mbox{,} \\
		\mbox{e } \drawSizedLine{CA} &= \drawSizedLine{AD} = \drawSizedLine{L'''} \mbox{.} \\
	\end{aligned}
	\right\}\mbox{\byref{\constref}}$
\end{center}

\qed

\startproblem{Prop XXIII. Prob.}\label{prop:I.XXIII}
\defineNewPicture{
pair A, B, C, D, E, F, G, H, J, d;
A := (0, 0);
B := A shifted (7/2u, 0);
C := A shifted (3u, 11/5u);
D := 5/4[A, B];
E := 7/6[A, C];
d := (0, -3u);
F := A shifted d;
G := B shifted d;
H := C shifted d;
J := D shifted d;
byAngleDefine(B, A, C, byred, SOLID_SECTOR);
byAngleDefine(G, F, H, byblue, SOLID_SECTOR);
draw byNamedAngleResized();
byLineDefine(B, D, byblack, DASHED_LINE, THIN_WIDTH);
byLineDefine(C, E, byblue, DASHED_LINE, THIN_WIDTH);
byLineDefine(A, B, byblack, SOLID_LINE, THIN_WIDTH);
byLineDefine(C, A, byblue, SOLID_LINE, THIN_WIDTH);
draw byLine(B, C, byred, SOLID_LINE, THIN_WIDTH);
draw byNamedLineSeq(0)(CE,CA,AB,BD);
byLineDefine(G, J, byblack, DASHED_LINE, REGULAR_WIDTH);
byLineDefine(F, G, byblack, SOLID_LINE, REGULAR_WIDTH);
byLineDefine(G, H, byred, SOLID_LINE, REGULAR_WIDTH);
byLineDefine(H, F, byyellow, SOLID_LINE, REGULAR_WIDTH);
draw byNamedLineSeq(0)(noLine,GH,HF,FG,GJ);
draw byLabelsOnPolygon(D, B, A, C, E)(OMIT_FIRST_LABEL+OMIT_LAST_LABEL, 0);
draw byLabelsOnPolygon(noPoint,J, G, F, H, G)(OMIT_FIRST_LABEL+OMIT_LAST_LABEL, 0);
}
\drawCurrentPictureInMargin
\problem{E}{m}{um ponto dado (\drawPointL{F}) em uma linha reta dada (\drawUnitLine{FG,GJ}), fazer um ângulo igual a um ângulo retilíneo dado (\drawAngle{A}).}

Trace \drawUnitLine{BC} entre quaisquer dois pontos nas pernas do ângulo dado.

\begin{center}
Construa \drawLine[bottom]{HF,GH,FG} \byref{prop:I.XXII}\\
de modo que $\drawUnitLine{FG} = \drawUnitLine{AB}$, $\drawUnitLine{HF} = \drawUnitLine{CA}$\\
e $\drawUnitLine{GH} = \drawUnitLine{BC}$.

Então $\drawAngle{A} = \drawAngle{F}$ \byref{prop:I.VIII}.
\end{center}

\qed

\starttheorem{Prop XXIV. Theor.}\label{prop:I.XXIV}
\defineNewPicture{
byAngleMacroName[3] := "byAngleMArcs";
vardef byAngleMArcs (expr angleArc, angleColor)(suffix angleOptionalColors) =
    save p;
    path p;
    p := angleArc scaled ((angleSize*angleScale) - lineWidthThin);
    image(
        for i := 1 step -1/floor((angleSize*angleScale)/2mm) until (2mm/(angleSize*angleScale)):
            draw (p scaled i) withpen pencircle scaled lineWidthThin withcolor angleColor;
        endfor;
    )
enddef;
pair A, B, C, D, E, F, G, d;
A := (0, 0);
B := A shifted (u, -5/2u);
C := A shifted (-u, -7/2u);
D := ((fullcircle scaled 2abs(C-A)) shifted A) intersectionpoint (B--B shifted (-2abs(C-A), 0));
d := (0, -4u);
E := A shifted d;
F := B shifted d;
G := C shifted d;
byAngleDefine(B, A, C, byblack, DASHED_ARC_SECTOR);
byAngleDefine(C, A, D, byblack, 3);
byAngleDefine(F, E, G, byblack, DASHED_ARC_SECTOR);
byAngleDefine(B, D, A, byblue, SOLID_SECTOR);
byAngleDefine(C, D, B, byred, SOLID_SECTOR);
byAngleDefine(D, C, A, byyellow, SOLID_SECTOR);
byAngleDefine(A, C, B, byblack, SOLID_SECTOR);
draw byNamedAngleResized();
byLineDefine(A, B, byblue, SOLID_LINE, REGULAR_WIDTH);
byLineDefine(B, C, byblack, DASHED_LINE, REGULAR_WIDTH);
draw byLine(C, A, byred, SOLID_LINE, REGULAR_WIDTH);
draw byLine(B, D, byblack, SOLID_LINE, REGULAR_WIDTH);
byLineDefine(A, D, byred, DASHED_LINE, REGULAR_WIDTH);
byLineDefine(C, D, byblue, DASHED_LINE, REGULAR_WIDTH);
draw byNamedLineSeq(0)(AB,AD,CD,BC);
byLineDefine(E, F, byblue, SOLID_LINE, THIN_WIDTH);
byLineDefine(F, G, byyellow, SOLID_LINE, THIN_WIDTH);
byLineDefine(G, E, byred, SOLID_LINE, THIN_WIDTH);
draw byNamedLineSeq(0)(EF,FG,GE);
draw byLabelsOnPolygon(D, A, B, C)(ALL_LABELS, 0);
draw byLabelsOnPolygon(G, E, F)(ALL_LABELS, 0);
}
\drawCurrentPictureInMargin
\problem{S}{e}{dois triângulos têm dois lados de um respectivamente iguais a dois lados do outro (\drawUnitLine{AB} a \drawUnitLine{EF} e \drawUnitLine{AD} a \drawUnitLine{GE}), e se um dos ângulos
(\drawFromCurrentPicture[bottom]{
startAutoLabeling;
draw byNamedAngleSides(BAC,CAD)(AB,CA,AD);
stopAutoLabeling;
})
contidos pelos lados iguais for maior que o outro
(\drawFromCurrentPicture[bottom][angleFEG]{
startAutoLabeling;
draw byNamedAngleSides(FEG)(EF, GE);
stopAutoLabeling;
}),
o lado (\drawUnitLine{DB}) que é oposto ao ângulo maior é maior que o lado (\drawUnitLine{FG}) que é oposto ao ângulo menor.
}

\begin{center}
Faça $\drawFromCurrentPicture[bottom][angleBAC]{
startAutoLabeling;
draw byNamedAngleSides(BAC)(AB, CA);
stopAutoLabeling;
} = \angleFEG$ \byref{prop:I.XXIII},\\
e $\drawUnitLine{CA} = \drawUnitLine{GE}$ \byref{prop:I.III},\\
trace \drawUnitLine{CD} e \drawUnitLine{BC}.

$\because \drawUnitLine{CA} = \drawUnitLine{AD}$ \byref{ax:I.I,\hypref,\constref}\\ %because
$\therefore \drawAngle{BDA,CDB} = \drawAngle{DCA}$ \byref{prop:I.V}\\
mas $\drawAngle{CDB} < \drawAngle{DCA}$,\\
e $\therefore \drawAngle{CDB} < \drawAngle{DCA,ACB}$,

$\therefore \drawUnitLine{DB} > \drawUnitLine{BC}$ \byref{prop:I.XIX}

mas $\drawUnitLine{BC} = \drawUnitLine{FG}$ \byref{prop:I.IV}

$\therefore \drawUnitLine{DB} > \drawUnitLine{FG}$.
\end{center}

\qed

\starttheorem{Prop XXV. Theor.}\label{prop:I.XXV}
\defineNewPicture{
pair A, B, C, D, E, F, d;
A := (0, 0);
B := A shifted (u, -3u);
C := A shifted (-7/4u, -4u);
d := (0, -9/2u);
D := A shifted d;
E := ((B shifted -A) rotated -10) shifted d;
F := C shifted d;
byAngleDefine(B, A, C, byyellow, SOLID_SECTOR);
byAngleDefine(E, D, F, byblack, SOLID_SECTOR);
draw byNamedAngleResized();
byLineDefine(A, B, byblue, SOLID_LINE, REGULAR_WIDTH);
byLineDefine(B, C, byblack, SOLID_LINE, REGULAR_WIDTH);
byLineDefine(C, A, byred, SOLID_LINE, REGULAR_WIDTH);
draw byNamedLineSeq(0)(AB,BC,CA);
byLineDefine(D, E, byblue, SOLID_LINE, THIN_WIDTH);
byLineDefine(E, F, byyellow, SOLID_LINE, THIN_WIDTH);
byLineDefine(F, D, byred, SOLID_LINE, THIN_WIDTH);
draw byNamedLineSeq(0)(DE,EF,FD);
draw byLabelsOnPolygon(A, B, C)(ALL_LABELS, 0);
draw byLabelsOnPolygon(D, E, F)(ALL_LABELS, 0);
}
\drawCurrentPictureInMargin
\problem{S}{e}{dois triângulos têm dois lados (\drawUnitLine[0.7cm]{AB} e \drawUnitLine[0.7cm]{CA}) de um respectivamente iguais a dois lados (\drawUnitLine{DE} e \drawUnitLine{FD}) do outro, mas suas bases desiguais, o ângulo subtendido pela base maior (\drawUnitLine{BC}) de um, deve ser maior que o ângulo subtendido pela base menor (\drawUnitLine{EF}) do outro.}

\begin{center}
$\drawAngle{A} =\mbox{, } > \mbox{ ou } < \drawAngle{D}$

\drawAngle{A} não é igual a \drawAngle{D}\\
pois se $\drawAngle{A} = \drawAngle{D}$ então $\drawUnitLine{CB} = \drawUnitLine{FE}$ \byref{prop:I.IV}\\
o que é contrário à hipótese;

\drawAngle{A} não é menor que \drawAngle{D}\\
pois se $\drawAngle{A} < \drawAngle{D}$\\
então $\drawUnitLine{CB} < \drawUnitLine{FE}$ \byref{prop:I.XXIV},\\
o que também é contrário à hipótese:

$\therefore \drawAngle{A} > \drawAngle{D}$.
\end{center}

\qed

\starttheorem{Prop XXVI. Theor.}\label{prop:I.XXVI}
\defineNewPicture{
pair A, B, C, D, E, F, G, d;
A := (0, 0);
B := A shifted (3u, 0);
C := A shifted (5/2u, 3u);
d := (0, -4u);
D := A shifted d;
E := B shifted d;
F := C shifted d;
G := 3/4[D, F];
byAngleDefine(B, A, C, byyellow, SOLID_SECTOR);
byAngleDefine(C, B, A, byred, SOLID_SECTOR);
byAngleDefine(A, C, B, byblack, ARC_SECTOR);
draw byNamedAngleResized(BAC, CBA, ACB);
byLineDefine(A, B, byblue, SOLID_LINE, REGULAR_WIDTH);
byLineDefine(B, C, byblack, SOLID_LINE, REGULAR_WIDTH);
byLineDefine(C, A, byred, SOLID_LINE, REGULAR_WIDTH);
draw byNamedLineSeq(0)(CA,BC,AB);
byAngleDefine(E, D, F, byyellow, SOLID_SECTOR);
byAngleDefine(G, E, D, byblack, SOLID_SECTOR);
byAngleDefine(F, E, G, byblue, SOLID_SECTOR);
byAngleDefine(D, F, E, byblack, ARC_SECTOR);
draw byNamedAngleResized(EDF, GED, FEG, DFE);
draw byLine(E, G, byyellow, SOLID_LINE, THIN_WIDTH);
byLineDefine(D, E, byblue, SOLID_LINE, THIN_WIDTH);
byLineDefine(E, F, byblack, SOLID_LINE, THIN_WIDTH);
byLineDefine(F, G, byred, DASHED_LINE, THIN_WIDTH);
byLineDefine(G, D, byred, SOLID_LINE, THIN_WIDTH);
draw byNamedLineSeq(0)(GD,FG,EF,DE);
draw byLabelsOnPolygon(C, B, A)(ALL_LABELS, 0);
draw byLabelsOnPolygon(D, G, F, E)(ALL_LABELS, 0);
}
\problem{S}{e}{dois triângulos têm dois ângulos de um respectivamente iguais a dois ângulos do outro ($\drawAngle{A} = \drawAngle{D}$ e $\drawAngle{B} = \drawAngle{GED,FEG}$), e um lado de um igual a um lado do outro similarmente posicionado em relação aos ângulos iguais, os lados e ângulos restantes são respectivamente iguais entre si.}
\drawCurrentPictureInMargin
\startsubproposition{Case I.}
\begin{center}
Sejam \drawUnitLine{AB} e \drawUnitLine{DE} que estão entre os ângulos iguais iguais,\\
então $\drawUnitLine{CA} = \drawUnitLine{GD,FG}$.

Pois se for possível, seja um deles \drawUnitLine{GD,FG} maior que o outro;\\
faça $\drawUnitLine{CA} = \drawUnitLine{GD}$, trace \drawUnitLine{EG}.

Em \drawLine[bottom]{CA,BC,AB} e
\drawLine[bottom]{GD,EG,DE} temos\\
$\drawUnitLine{CA} = \drawUnitLine{GD}$, $\drawAngle{A} = \drawAngle{D}$, $\drawUnitLine{AB} = \drawUnitLine{DE}$;\\
$\therefore \drawAngle{B} = \drawAngle{GED}$ (pr. 4.)\\
mas $\drawAngle{B} = \drawAngle{GED,FEG}$ \byref{\hypref}

e portanto $\drawAngle{GED} = \drawAngle{GED,FEG}$, o que é absurdo;\\
portanto nenhum dos lados \drawUnitLine{CA} e \drawUnitLine{GD,FG} é maior que o outro;\\
e $\therefore$ eles são iguais;

$\therefore \drawUnitLine{BC} = \drawUnitLine{EF}$, e $\drawAngle{C} = \drawAngle{F}$, \byref{prop:I.IV}.
\end{center}

\vfill\pagebreak

\defineNewPicture{
pair A, B, C, D, E, F, G, d;
d := (0, -4u);
A := (0, 0);
B := A shifted (3u, 0);
C := A shifted (1/2u, 3u);
D := A shifted d;
E := B shifted d;
F := C shifted d;
G := 3/4[D, E];
byAngleDefine(B, A, C, byyellow, SOLID_SECTOR);
byAngleDefine(C, B, A, byred, SOLID_SECTOR);
draw byNamedAngleResized(BAC, CBA);
byLineDefine(A, B, byblue, SOLID_LINE, REGULAR_WIDTH);
byLineDefine(B, C, byblack, SOLID_LINE, REGULAR_WIDTH);
byLineDefine(C, A, byred, SOLID_LINE, REGULAR_WIDTH);
draw byNamedLineSeq(0)(CA,AB,BC);
byAngleDefine(F, D, E, byyellow, SOLID_SECTOR);
byAngleDefine(F, G, D, byblack, SOLID_SECTOR);
byAngleDefine(F, E, D, byred, SOLID_SECTOR);
draw byNamedAngleResized(FDE, FGD, FED);
draw byLine(F, G, byyellow, SOLID_LINE, THIN_WIDTH);
byLineDefine(D, G, byblue, SOLID_LINE, THIN_WIDTH);
byLineDefine(G, E, byblue, DASHED_LINE, THIN_WIDTH);
byLineDefine(E, F, byblack, SOLID_LINE, THIN_WIDTH);
byLineDefine(F, D, byred, SOLID_LINE, THIN_WIDTH);
draw byNamedLineSeq(0)(FD,EF,GE,DG);
draw byLabelsOnPolygon(C, B, A)(ALL_LABELS, 0);
draw byLabelsOnPolygon(D, F, E, G)(ALL_LABELS, 0);
}
\drawCurrentPictureInMargin
\startsubproposition{Case II.}
\begin{center}
Novamente, seja $\drawUnitLine{CA} = \drawUnitLine{FD}$, que estão opostos aos ângulos iguais \drawAngle{B} e \drawAngle{E}.\\
Se for possível, seja $\drawUnitLine{DG,GE} > \drawUnitLine{AB}$, então tome $\drawUnitLine{DG} = \drawUnitLine{AB}$, trace \drawUnitLine{FG}.

Então em \drawLine[bottom]{CA,BC,AB} e \drawLine[bottom]{FD,FG,DG} temos $\drawUnitLine{CA} = \drawUnitLine{FD}$, $\drawUnitLine{AB} = \drawUnitLine{DG}$ e $\drawAngle{A} = \drawAngle{D}$,

$\therefore \drawAngle{B} = \drawAngle{G}$ \byref{prop:I.IV}\\
mas $\drawAngle{B} = \drawAngle{E}$ \byref{\hypref}

$\therefore \drawAngle{G} = \drawAngle{E}$ o que é absurdo \byref{prop:I.XVI}.

Consequentemente, nenhum dos lados \drawUnitLine{AB} ou \drawUnitLine{DG,GE} é maior que o outro, portanto eles devem ser iguais. Segue-se (por \byref{prop:I.IV}) que os triângulos são iguais em todos os aspectos.
\end{center}

\qed

\starttheorem{Prop XXVII. Theor.}\label{prop:I.XXVII}
\defineNewPicture{
pair A, B, C, D, E, F, G, H, I, d;
A := (0, 0);
B := A shifted (8/3u, 0);
d := (0, -7/4u);
C := A shifted d;
D := B shifted d;
E := 1/3[A, B];
F := 2/3[C, D];
G := 3/2[F, E];
H := 3/2[E, F];
I := 1/2[A, C] shifted (-2u, 0);
byAngleDefine(A, E, H, byyellow, SOLID_SECTOR);
byAngleDefine(H, E, B, byred, SOLID_SECTOR);
byAngleDefine(C, F, G, byblue, SOLID_SECTOR);
byAngleDefine(G, F, D, byyellow, SOLID_SECTOR);
draw byNamedAngleResized();
byLineDefine(I, A, byblue, SOLID_LINE, REGULAR_WIDTH);
byLineDefine(A, B, byblue, SOLID_LINE, REGULAR_WIDTH);
byLineDefine(I, C, byred, SOLID_LINE, REGULAR_WIDTH);
byLineDefine(C, D, byred, SOLID_LINE, REGULAR_WIDTH);
draw byNamedLineSeq(0)(CD,IC,IA,AB);
draw byLine(G, H, byblack, SOLID_LINE, REGULAR_WIDTH);
draw byLabelLine(0)(AB, CD, GH);
draw byLabelsOnPolygon(G, E, B)(OMIT_FIRST_LABEL+OMIT_LAST_LABEL, 0);
draw byLabelsOnPolygon(H, F, C)(OMIT_FIRST_LABEL+OMIT_LAST_LABEL, 0);
}
\drawCurrentPictureInMargin
\problem{S}{e}{uma linha reta (\drawUnitLine{GH}) encontrando duas outras linhas retas (\drawUnitLine{CD} e \drawUnitLine{AB}) forma com elas os ângulos alternados (\drawAngle{CFG} e \drawAngle{HEB}; \drawAngle{GFD} e \drawAngle{AEH}) iguais, essas duas linhas retas são paralelas.}

Se \drawUnitLine{CD} não for paralela a \drawUnitLine{AB} elas se encontrarão quando prolongadas.

Se for possível, sejam essas linhas não paralelas, mas se encontrem quando prolongadas; então o ângulo externo \drawAngle{HEB} é maior que \drawAngle{CFG} \byref{prop:I.XVI}, mas eles também são iguais \byref{\hypref}, o que é absurdo: da mesma forma pode ser mostrado que elas não podem se encontrar do outro lado; $\therefore$ elas são paralelas.

\qed

\starttheorem{Prop XXVIII. Theor.}\label{prop:I.XXVIII}
\defineNewPicture{
pair A, B, C, D, E, F, G, H, d;
A := (0, 0);
B := A shifted (7/2u, 0);
d := (0, -3/2u);
C := A shifted d;
D := B shifted d;
E := 9/20[A, B];
F := 11/20[C, D];
G := 7/4[F, E];
H := 4/3[E, F];
byAngleDefine(G, E, A, byblack, SOLID_SECTOR);
byAngleDefine(B, E, G, byyellow, SOLID_SECTOR);
byAngleDefine(A, E, H, byred, SOLID_SECTOR);
byAngleDefine(H, E, B, byblue, SOLID_SECTOR);
byAngleDefine(C, F, G, byblue, SOLID_SECTOR);
byAngleDefine(G, F, D, byred, SOLID_SECTOR);
draw byNamedAngleResized();
draw byLine(A, B, byred, SOLID_LINE, REGULAR_WIDTH);
draw byLine(C, D, byyellow, SOLID_LINE, REGULAR_WIDTH);
draw byLine(G, H, byblack, SOLID_LINE, REGULAR_WIDTH);
draw byLabelLine(0)(AB, CD, GH);
draw byLabelsOnPolygon(G, E, B)(OMIT_FIRST_LABEL+OMIT_LAST_LABEL, 0);
draw byLabelsOnPolygon(H, F, C)(OMIT_FIRST_LABEL+OMIT_LAST_LABEL, 0);
}
\drawCurrentPictureInMargin
\problem{S}{e}{uma linha reta (\drawUnitLine{GH}), cortando duas outras linhas retas (\drawUnitLine{AB} e \drawUnitLine{CD}), torna o externo igual ao interno e ao ângulo oposto, no mesmo lado da linha cortante (ou seja, $\drawAngle{GEA} = \drawAngle{CFG}$ ou $\drawAngle{BEG} = \drawAngle{GFD}$), ou se torna os dois ângulos internos no mesmo lado (\drawAngle{GFD} e \drawAngle{HEB}, ou \drawAngle{CFG} e \drawAngle{AEH}) juntos iguais a dois ângulos retos, essas duas linhas retas são paralelas.}

\begin{center}
Primeiro, se $\drawAngle{GEA} = \drawAngle{CFG}$, então $\drawAngle{GEA} = \drawAngle{HEB}$ \byref{prop:I.XV},\\
$\therefore \drawAngle{CFG} = \drawAngle{HEB} \therefore \drawUnitLine{AB} \parallel \drawUnitLine{CD}$ \byref{prop:I.XXVII}.

Segundo, se $\drawAngle{CFG} + \drawAngle{AEH} = \drawTwoRightAngles$,\\
então $\drawAngle{AEH} + \drawAngle{HEB} = \drawTwoRightAngles$ \byref{prop:I.XIII},\\
$\therefore \drawAngle{CFG} + \drawAngle{AEH} = \drawAngle{AEH} + \drawAngle{HEB}$ \byref{ax:I.III}

$\therefore \drawAngle{CFG} = \drawAngle{HEB}$

$\therefore \drawUnitLine{AB} \parallel \drawUnitLine{CD}$ \byref{prop:I.XXVII}.
\end{center}

\qed

\starttheorem{Prop XXIX. Theor.}\label{prop:I.XXIX}
\defineNewPicture{
pair A, B, C, D, E, F, G, H, I, J, d[];
A := (0, 0);
B := A shifted (7/2u, 0);
d1 := (0, -2u);
C := A shifted d1;
D := B shifted d1;
E := 11/20[A, B];
F := 9/20[C, D];
G := 7/4[F, E];
H := 4/3[E, F];
d2 := (3/2u, 1/2u);
I := E shifted -d2;
J := E shifted d2;
byAngleDefine(I, E, A, byblue, SOLID_SECTOR);
byAngleDefine(H, E, I, byyellow, SOLID_SECTOR);
byAngleDefine(B, E, H, byblack, SOLID_SECTOR);
byAngleDefine(G, E, B, byred, SOLID_SECTOR);
byAngleDefine(G, F, D, byblack, SOLID_SECTOR);
draw byNamedAngleResized();
draw byLine(I, E, byblack, SOLID_LINE, REGULAR_WIDTH);
draw byLine(E, J, byblack, DASHED_LINE, REGULAR_WIDTH);
draw byLine(A, B, byyellow, SOLID_LINE, REGULAR_WIDTH);
draw byLine(C, D, byred, SOLID_LINE, REGULAR_WIDTH);
draw byLine(G, H, byblue, SOLID_LINE, REGULAR_WIDTH);
draw byLabelLine(0)(AB, CD, GH);
draw byLabelsOnPolygon(A, E, G)(OMIT_FIRST_LABEL+OMIT_LAST_LABEL, 0);
draw byLabelsOnPolygon(H, F, C)(OMIT_FIRST_LABEL+OMIT_LAST_LABEL, 0);
draw byLabelPoint(I, lineAngle.IE - 90, 1);
draw byLabelPoint(J, lineAngle.EJ + 90, 1);
}
\drawCurrentPictureInMargin
\problem{U}{ma}{linha reta (\drawUnitLine{GH}) caindo sobre duas linhas retas paralelas (\drawUnitLine{AB} e \drawUnitLine{CD}), torna os ângulos alternados iguais entre si; e também o externo igual ao interno e ao ângulo oposto no mesmo lado; e os dois ângulos internos no mesmo lado juntos iguais a dois ângulos retos.}

Pois se os ângulos alternados \drawAngle{IEA,HEI} e \drawAngle{GFD} não forem iguais, trace \drawUnitLine{IE}, fazendo $\drawAngle{HEI} = \drawAngle{GFD}$ \byref{prop:I.XXIII}.

Portanto $\drawUnitLine{IE,EJ} \parallel \drawUnitLine{CD}$ \byref{prop:I.XXVII} e portanto duas linhas retas que se intersectam são paralelas à mesma linha reta, o que é impossível \byref{ax:I.XII}.

Portanto \drawAngle{IEA,HEI} e \drawAngle{GFD} não são desiguais, isto é, eles são iguais: $\drawAngle{IEA,HEI} = \drawAngle{GEB}$ \byref{prop:I.XV}; $\therefore \drawAngle{GEB} = \drawAngle{GFD}$, o ângulo externo igual ao interno e oposto no mesmo lado: se \drawAngle{BEH} for adicionado a ambos, então $\drawAngle{GFD} + \drawAngle{BEH} = \drawAngle{BEH,GEB} = \drawTwoRightAngles$ \byref{prop:I.XIII}. Isto é dizer, os dois ângulos internos no mesmo lado da linha cortante são iguais a dois ângulos retos.

\qed

\starttheorem{Prop XXX. Theor.}\label{prop:I.XXX}
\defineNewPicture{
pair A, B, C, D, E, F, G, H, I, J, K, d;
A := (0, 0);
B := A shifted (7/2u, 0);
d := (0, -u);
C := A shifted d;
D := B shifted d;
E := C shifted d;
F := D shifted d;
G := 13/20[A, B];
H := 7/20[E, F];
I := (G--H) intersectionpoint (C--D);
J := 3/2[H, G];
K := 5/4[G, H];
byAngleDefine(B, G, J, byyellow, SOLID_SECTOR);
byAngleDefine(D, I, J, byblue, SOLID_SECTOR);
byAngleDefine(F, H, J, byred, SOLID_SECTOR);
draw byNamedAngleResized();
draw byLine(A, B, byred, SOLID_LINE, REGULAR_WIDTH);
draw byLine(C, D, byyellow, SOLID_LINE, REGULAR_WIDTH);
draw byLine(E, F, byblue, SOLID_LINE, REGULAR_WIDTH);
draw byLine(J, K, byblack, SOLID_LINE, REGULAR_WIDTH);
draw byLabelLine(0)(AB, CD, EF, JK);
draw byLabelsOnPolygon(J, G, A)(OMIT_FIRST_LABEL+OMIT_LAST_LABEL, 0);
draw byLabelsOnPolygon(J, I, C)(OMIT_FIRST_LABEL+OMIT_LAST_LABEL, 0);
draw byLabelsOnPolygon(J, H, E)(OMIT_FIRST_LABEL+OMIT_LAST_LABEL, 0);
}
\drawCurrentPictureInMargin
\problem{L}{inhas}{retas (\drawUnitLine{AB} e \drawUnitLine{EF}) que são paralelas à mesma linha reta (\drawUnitLine{CD}), são paralelas entre si.}

\begin{center}
Seja \drawUnitLine{JK} intersectando $\left\{\vcenter{\nointerlineskip\hbox{\drawUnitLine{AB}}\nointerlineskip\hbox{\drawUnitLine{CD}}\nointerlineskip\hbox{\drawUnitLine{EF}}}\right\}$;

Então, $\drawAngle{G} = \drawAngle{I} = \drawAngle{H}$ \byref{prop:I.XXIX},

$\therefore \drawAngle{G} = \drawAngle{H}$

$\therefore \drawUnitLine{AB} \parallel \drawUnitLine{EF}$ \byref{prop:I.XXVIII}. % improvement: byrne references I.XXVII here, which is wrong https://github.com/jemmybutton/byrne-euclid/issues/27#issuecomment-507012558
\end{center}

\qed

\startproblem{Prop XXXI. Prob.}\label{prop:I.XXXI}
\defineNewPicture[1/6]{
pair A, B, C, D, E, F, d;
A := (0, 0);
B := A shifted (7/2u, 0);
d := (0, -3u);
C := A shifted d;
D := B shifted d;
E := 4/5[A, B];
F := 1/5[C, D];
byAngleDefine(F, E, A, byyellow, SOLID_SECTOR);
byAngleDefine(E, F, D, byred, SOLID_SECTOR);
draw byNamedAngleResized();
draw byLine(E, F, byblack, SOLID_LINE, REGULAR_WIDTH);
byLineDefine(A, E, byred, SOLID_LINE, REGULAR_WIDTH);
byLineDefine(E, B, byred, DASHED_LINE, REGULAR_WIDTH);
byLineDefine(C, F, byblue, SOLID_LINE, REGULAR_WIDTH);
byLineDefine(F, D, byblue, SOLID_LINE, REGULAR_WIDTH);
draw byNamedLineSeq(0)(AE, EB);
draw byNamedLineSeq(0)(CF, FD);
draw byLabelsOnPolygon(A, E, B, noPoint, D, F, C, noPoint)(ALL_LABELS, 0);
}
\drawCurrentPictureInMargin
\problem{D}{e}{um ponto dado \drawPointL[middle][EB]{E} traçar uma linha reta paralela a uma linha reta dada (\drawUnitLine{CF,FD}).}

\begin{center}
Trace \drawUnitLine{EF} do ponto \drawPointL[middle][EB]{E} a qualquer ponto  \drawPointL[middle][CF]{F} em \drawUnitLine{CF,FD},

faça $\drawAngle{E} = \drawAngle{F}$ \byref{prop:I.XXIII},

então $\drawUnitLine{AE,EB} \parallel \drawUnitLine{CF,FD}$ \byref{prop:I.XXVII}.
\end{center}

\qed

\starttheorem{Prop XXXII. Theor.}\label{prop:I.XXXII}
\defineNewPicture[1/6]{
pair A, B, C, D, E;
A := (0, 0);
B := A shifted (-1u, -3u);
C := A shifted (5/2u, -3u);
D := 4/3[B, A];
E := A shifted (unitvector(C-B) scaled 3/2u);
byAngleDefine(D, A, E, byred, SOLID_SECTOR);
byAngleDefine(E, A, C, byblack, SOLID_SECTOR);
byAngleDefine(C, A, B, byblue, SOLID_SECTOR);
byAngleDefine(A, B, C, byyellow, SOLID_SECTOR);
byAngleDefine(B, C, A, byblack, SOLID_SECTOR);
draw byNamedAngleResized();
draw byLine(A, E, byblue, SOLID_LINE, REGULAR_WIDTH);
byLineDefine(A, D, byblack, DASHED_LINE, REGULAR_WIDTH);
byLineDefine(A, B, byblack, SOLID_LINE, REGULAR_WIDTH);
byLineDefine(B, C, byred, SOLID_LINE, REGULAR_WIDTH);
byLineDefine(C, A, byyellow, SOLID_LINE, REGULAR_WIDTH);
draw byNamedLineSeq(0)(CA,noLine,AD,AB,BC);
draw byLabelsOnPolygon(C, B, A, D, noPoint, E, A)(ALL_LABELS, 0);
}
\drawCurrentPictureInMargin
\problem[3]{S}{e}{qualquer lado (\drawUnitLine{AB}) de um triângulo for prolongado, o ângulo externo (\drawAngle{DAE,EAC}) é igual à soma dos dois ângulos internos e opostos (\drawAngle{B} e \drawAngle{C}), e os três ângulos internos de qualquer triângulo tomados juntos são iguais a dois ângulos retos.}

\begin{center}
Através do ponto \drawPointL[middle][AD,AE]{A} desenhe\\
$\drawUnitLine{AE} \parallel \drawUnitLine{BC}$ \byref{prop:I.XXXI}.

Então $\left\{
	\begin{aligned}
		\drawAngle{DAE} &= \drawAngle{B}\\
		\drawAngle{EAC} &= \drawAngle{C}\\
	\end{aligned}
	\right\}$ \byref{prop:I.XXIX},

$\therefore \drawAngle{B} + \drawAngle{C} = \drawAngle{DAE,EAC}$ \byref{ax:I.II},

e portanto\\
$\drawAngle{B} + \drawAngle{CAB} + \drawAngle{C} = \drawAngle{DAE,EAC,CAB} = \drawTwoRightAngles$ \byref{prop:I.XIII}.
\end{center}

\qed

\starttheorem{Prop XXXIII. Theor.}\label{prop:I.XXXIII}
\defineNewPicture[1/4]{
pair A, B, C, D, d[];
d1 := (5/2u, 0);
d2 := (-7/8u, -3u);
A := (0, 0);
B := A shifted d1;
C := A shifted d2;
D := C shifted d1;
byAngleDefine(B, A, D, byyellow, SOLID_SECTOR);
byAngleDefine(D, A, C, byred, SOLID_SECTOR);
byAngleDefine(C, D, A, byblack, SOLID_SECTOR);
byAngleDefine(A, D, B, byblue, SOLID_SECTOR);
draw byNamedAngleResized();
draw byLine(A, D, byblack, SOLID_LINE, REGULAR_WIDTH);
byLineDefine(A, B, byred, SOLID_LINE, REGULAR_WIDTH);
byLineDefine(C, D, byred, DASHED_LINE, REGULAR_WIDTH);
byLineDefine(A, C, byblue, SOLID_LINE, REGULAR_WIDTH);
byLineDefine(B, D, byyellow, SOLID_LINE, REGULAR_WIDTH);
draw byNamedLineSeq(0)(AB,BD,CD,AC);
draw byLabelsOnPolygon(A, B, D, C)(ALL_LABELS, 0);
}
\drawCurrentPictureInMargin
\problem{L}{inhas}{retas (\drawUnitLine{AC} e \drawUnitLine{BD}) que unem as extremidades adjacentes de duas linhas retas iguais e paralelas (\drawUnitLine{AB} e \drawUnitLine{CD}), são elas mesmas iguais e paralelas.}

\begin{center}
Trace \drawUnitLine{AD} a diagonal.

$\drawUnitLine{AB} = \drawUnitLine{CD}$ \byref{\hypref}\\
$\drawAngle{BAD} = \drawAngle{CDA}$ \byref{prop:I.XXIX}\\
e \drawUnitLine{AD} comum aos dois triângulos; % what triangles?

$\therefore \drawUnitLine{AC} = \drawUnitLine{BD}$, e $\drawAngle{ADB} = \drawAngle{DAC}$ \byref{prop:I.IV};

e $\therefore \drawUnitLine{AC} \parallel \drawUnitLine{BD}$ \byref{prop:I.XXVII}.
\end{center}

\qed

\starttheorem{Prop XXXIV. Theor.}\label{prop:I.XXXIV}
\defineNewPicture{
pair A, B, C, D, d[];
d1 := (5/2u, 0);
d2 := (-7/8u, -3u);
A := (0, 0);
B := A shifted d1;
C := A shifted d2;
D := C shifted d1;
byAngleDefine(B, A, D, byblue, SOLID_SECTOR);
byAngleDefine(D, A, C, byred, SOLID_SECTOR);
byAngleDefine(C, D, A, byyellow, SOLID_SECTOR);
byAngleDefine(A, D, B, byred, SOLID_SECTOR);
byAngleDefine(A, C, D, byblack, SOLID_SECTOR);
byAngleDefine(D, B, A, byblack, SOLID_SECTOR);
draw byNamedAngleResized();
draw byLine(A, D, byblack, SOLID_LINE, REGULAR_WIDTH);
byLineDefine(A, B, byred, SOLID_LINE, REGULAR_WIDTH);
byLineDefine(C, D, byred, DASHED_LINE, REGULAR_WIDTH);
byLineDefine(A, C, byyellow, SOLID_LINE, REGULAR_WIDTH);
byLineDefine(B, D, byblue, SOLID_LINE, REGULAR_WIDTH);
draw byNamedLineSeq(0)(AB,BD,CD,AC);
draw byLabelsOnPolygon(A, B, D, C)(ALL_LABELS, 0);
}
\drawCurrentPictureInMargin
\problem{O}{s}{lados e ângulos opostos de qualquer paralelogramo são iguais, e a diagonal (\drawUnitLine{AD}) o divide em duas partes iguais.}

\begin{center}
Como $\left\{
	\begin{aligned}
		\drawAngle{BAD} &= \drawAngle{CDA}\\
		\drawAngle{DAC} &= \drawAngle{ADB}\\
	\end{aligned}
	\right\}$ \byref{prop:I.XXIX}

e \drawUnitLine{AD} comum aos dois triângulos.

$\therefore \left\{
	\begin{aligned}
		\drawUnitLine{AB} &= \drawUnitLine{CD}\\
		\drawUnitLine{AC} &= \drawUnitLine{BD}\\
		\drawAngle{B} &= \drawAngle{C}\\
	\end{aligned}
	\right\}$ \byref{prop:I.XXVI}\\
e $\drawAngle{BAD,DAC} = \drawAngle{CDA,ADB}$ \byref{ax:I.II}.
\end{center}

Portanto os lados e ângulos opostos do paralelogramo são iguais: e como os triângulos \drawLine{AD,CD,AC} e \drawLine{AB,BD,AD} são iguais em todos os aspectos \byref{prop:I.IV}, a diagonal divide o paralelogramo em duas partes iguais.

\qed

\starttheorem{Prop XXXV. Theor.}\label{prop:I.XXXV}
\defineNewPicture{
pair A, B, C, D, E, F, G, d[];
d1 := (7/4u, 0);
d2 := (u, -3u);
d3 := (-2u, -3u);
A := (0, 0);
B := A shifted d1;
C := A shifted d2;
D := C shifted d1;
E := C shifted -d3;
F := D shifted -d3;
G := (B--D) intersectionpoint (C--E);
draw byPolygon(A,B,G,C)(byyellow);
draw byPolygon(E,F,D,G)(byyellow);
draw byPolygon(B,E,G)(byyellow);
draw byPolygon(C,D,G)(byblue);
byAngleDefine(E, A, C, byred, SOLID_SECTOR);
byAngleDefine(F, B, D, byblue, SOLID_SECTOR);
byAngleDefine(A, E, C, byblack, SOLID_SECTOR);
byAngleDefine(B, F, D, white, SOLID_SECTOR);
draw byNamedAngleResized();
draw byNamedAngleDummySides(BFD);
byLineDefine(A, C, byblue, SOLID_LINE, REGULAR_WIDTH);
byLineDefine(B, D, byred, SOLID_LINE, REGULAR_WIDTH);
byLineDefine(C, D, byblack, SOLID_LINE, REGULAR_WIDTH);
draw byNamedLineSeq(0)(AC,CD,BD);
draw byLabelsOnPolygon(C, A, B, E, F, D)(ALL_LABELS, 0);
}
\drawCurrentPictureInMargin
\problem{P}{aralelogramos}{sobre a mesma base, e entre as mesmas paralelas, são (em área) iguais.}

\begin{center}
Por causa das paralelas,\\
$\left.
	\begin{aligned}
		\drawAngle{A} &= \drawAngle{B};\\
		\drawAngle{E} &= \drawAngle{F};\\
		\mbox{e } \drawUnitLine{AC} &= \drawUnitLine{BD}\\
	\end{aligned}
	\right\}
	\begin{aligned}
	&\mbox{\byref{prop:I.XXIX}}\\
	&\mbox{\byref{prop:I.XXIX}}\\
	&\mbox{\byref{prop:I.XXXIV}.}\\
	\end{aligned}
	$

Mas
$\drawFromCurrentPicture[middle][polygonABC]{
startAutoLabeling;
draw byNamedPolygon (ABGC, BEG);
stopAutoLabeling;
draw byNamedLine (AC);
}
=
\drawFromCurrentPicture[middle][polygonEFD]{
startAutoLabeling;
draw byNamedPolygon (EFDG, BEG);
stopAutoLabeling;
draw byNamedLine (BD);
}$ \byref{prop:I.VIII}

$\therefore
\drawFromCurrentPicture[middle][polygonAFDC]{
startAutoLabeling;
draw byNamedPolygon (ABGC, EFDG, BEG, CDG);
stopAutoLabeling;
draw byNamedLine (AC);
} - \polygonEFD =
\drawPolygon{ABGC, CDG}$,\\
e $\polygonAFDC - \polygonABC =
\drawPolygon{EFDG, CDG}$;

$\therefore \drawPolygon{ABGC, CDG} = \drawPolygon{EFDG, CDG}$.
\end{center}

\qed

\starttheorem{Prop XXXVI. Theor.}\label{prop:I.XXXVI}
\defineNewPicture{
pair A, B, C, D, E, F, G, H, J, I, d[];
numeric h;
h := 3u;
d1 := (3/2u, 0);
d2 := (2/3u, -h);
d3 := (-8/3u, -h);
d4 := (-1/2u, -h);
A := (0, 0);
B := A shifted d1;
C := A shifted d2;
D := C shifted d1;
E := C shifted -d3;
F := D shifted -d3;
G := E shifted d4;
H := F shifted d4;
I := (B--D) intersectionpoint (C--E);
J := (E--G) intersectionpoint (D--F);
draw byPolygon(A,B,I,C)(byred);
draw byPolygon(C,D,I)(byred);
draw byPolygon(I,D,J,E)(byblue);
draw byPolygon(E,F,J)(byyellow);
draw byPolygon(J,F,H,G)(byyellow);
byLineDefine(C, E, byyellow, SOLID_LINE, REGULAR_WIDTH);
byLineDefine(D, F, byblack, DASHED_LINE, REGULAR_WIDTH);
byLineDefine(C, D, byblack, SOLID_LINE, REGULAR_WIDTH);
byLineDefine(E, F, byred, SOLID_LINE, REGULAR_WIDTH);
draw byNamedLineSeq(0)(CE,EF,DF,CD);
draw byLineFull(G, H, byblue, 0, 0)(E, F, 1, 1, 0);
draw byLabelsOnPolygon(A, B, D, C)(ALL_LABELS, 0);
draw byLabelsOnPolygon(E, F, H, G)(ALL_LABELS, 0);
}
\drawCurrentPictureInMargin
\problem{P}{aralelogramos}{(\drawPolygon[bottom][polygonABDC]{ABIC, CDI}~e~\drawPolygon[bottom][polygonEFHG]{EFJ, JFHG}) sobre bases iguais, e entre as mesmas paralelas, são iguais.}

\begin{center}
Trace \drawUnitLine{CE} e \drawUnitLine{DF}\\
$\drawUnitLine{CD} = \drawUnitLine{GH} = \drawUnitLine{EF}$ por \byref{prop:I.XXXIV,\hypref};

$\therefore \drawUnitLine{CD} = \mbox{ e } \parallel \drawUnitLine{EF}$;

$\therefore \drawUnitLine{CE} = \mbox{ e } \parallel \drawUnitLine{DF}$ \byref{prop:I.XXXIII}

E portanto
\drawPolygon[bottom][polygonCDFE]{CDI, IDJE, EFJ}
é um paralelogramo:

mas $\polygonABDC = \polygonCDFE = \polygonEFHG$ \byref{prop:I.XXXV}

$\therefore \polygonABDC = \polygonEFHG$ \byref{ax:I.I}.
\end{center}

\qed

\starttheorem{Prop XXXVII. Theor.}\label{prop:I.XXXVII}
\defineNewPicture{
pair A, B, C, D, E, F, G, H, I, d[];
d1 := (3/2u, 0);
d2 := (1/2u, -3u);
d3 := (-7/4u, -3u);
A := (0, 0);
B := A shifted d1;
C := A shifted d2;
D := C shifted d1;
E := C shifted -d3;
F := D shifted -d3;
G := (B--D) intersectionpoint (C--E);
H := 11/10[F, A];
I := 11/10[A, F];
draw byPolygon(A,B,C)(byblue);
draw byPolygon(B,C,G)(byyellow);
draw byPolygon(C,D,G)(byyellow);
draw byPolygon(D,G,E)(byblack);
draw byPolygon(E,F,D)(byred);
draw byLine(B, D, byred, SOLID_LINE, REGULAR_WIDTH);
draw byLine(E, C, byblue, SOLID_LINE, REGULAR_WIDTH);
byLineDefine(A, C, byred, DASHED_LINE, REGULAR_WIDTH);
byLineDefine(F, D, byblue, DASHED_LINE, REGULAR_WIDTH);
byLineDefine(C, D, byblack, SOLID_LINE, REGULAR_WIDTH);
draw byNamedLineSeq(0)(AC,CD,FD);
draw byLine(H, I, byblack, DASHED_LINE, REGULAR_WIDTH);
draw byLabelsOnPolygon(H, A, B, E, F, I, noPoint)(ALL_LABELS, 0);
draw byLabelsOnPolygon(F, D, C, A)(OMIT_FIRST_LABEL+OMIT_LAST_LABEL, 0);
}
\drawCurrentPictureInMargin
\problem{T}{riângulos}{
\drawPolygon[bottom][polygonBCD]{BCG, CDG} e~\drawPolygon[bottom][polygonCDE]{DGE, CDG}
sobre a mesma base (\drawUnitLine{CD}) e entre as mesmas paralelas são iguais.
}

\begin{center}
$\left.
	\begin{aligned}
		\mbox{Trace } \drawUnitLine{AC} &\parallel \drawUnitLine{BD}\\
		\drawUnitLine{FD} &\parallel \drawUnitLine{EC}\\
	\end{aligned}
\right\}\mbox{\byref{prop:I.XXXI}}$

Prolongue \drawUnitLine{HI}.

\drawPolygon[bottom][polygonABDC]{ABC, BCG, CDG}
e
\drawPolygon[bottom][polygonEFDC]{DGE, CDG, EFD}
são paralelogramos sobre a mesma base e entre as mesmas paralelas, e portanto iguais. \byref{prop:I.XXXV}

$\therefore \left\{
	\begin{aligned}
		\polygonABDC &= \mbox{ duas vezes } \polygonBCD\\
		\polygonEFDC &= \mbox{ duas vezes } \polygonCDE\\
	\end{aligned}
	\right\}$ \byref{prop:I.XXXIV}
$\therefore \polygonBCD = \polygonCDE$.
\end{center}

\qed

\starttheorem{Prop XXXVIII. Theor.}\label{prop:I.XXXVIII}
\defineNewPicture{
pair A, B, C, D, E, F, G, H, J, I, d[];
numeric h;
h := 5/2u;
d1 := (3/2u, 0);
d2 := (3/4u, -h);
d3 := (7/3u, h);
d4 := (-1/4u, -h);
A := (0, 0);
B := A shifted d1;
C := A shifted d2;
D := C shifted d1;
E := C shifted d3;
F := D shifted d3;
G := E shifted d4;
H := F shifted d4;
I := (xpart(A), ypart(C));
J := (xpart(F), ypart(C));
draw byPolygon(A,B,C)(byyellow);
draw byPolygon(B,C,D)(byred);
draw byPolygon(E,F,H)(byblack);
draw byPolygon(E,G,H)(byblue);
draw byLine(B, D, byblue, SOLID_LINE, REGULAR_WIDTH);
draw byLine(E, G, byred, SOLID_LINE, REGULAR_WIDTH);
byLineDefine(A, C, byblue, DASHED_LINE, REGULAR_WIDTH);
byLineDefine(F, H, byred, DASHED_LINE, REGULAR_WIDTH);
byLineDefine(A, F, byblack, DASHED_LINE, REGULAR_WIDTH);
draw byNamedLineSeq(0)(AC,AF,FH);
draw byLine(I, J, byblack, DASHED_LINE, REGULAR_WIDTH);
draw byLabelsOnPolygon(A, B, E, F, noPoint)(ALL_LABELS, 0);
draw byLabelsOnPolygon(H, G, D, C, noPoint)(ALL_LABELS, 0);
}
\drawCurrentPictureInMargin
\problem{T}{riângulos}{(\drawPolygon[bottom][polygonBCD]{BCD} e \drawPolygon[bottom][polygonEGH]{EGH}) sobre bases iguais e entre as mesmas paralelas são iguais.}

\begin{center}
$\left.
	\begin{aligned}
		\mbox{Trace } \drawUnitLine{AC} &\parallel \drawUnitLine{BD}\\
		\mbox{e } \drawUnitLine{FH} &\parallel \drawUnitLine{EG}\\
	\end{aligned}
	\right\}\mbox{\byref{prop:I.XXXI}}$\\
$\drawPolygon[bottom][polygonABDC]{ABC, BCD} = \drawPolygon[bottom][polygonEFHG]{EFH, EGH}$ \byref{prop:I.XXXVI};

mas $\polygonABDC = \mbox{ duas vezes } \polygonBCD$ \byref{prop:I.XXXIV},\\
e $\polygonEFHG = \mbox{ duas vezes } \polygonEGH$ \byref{prop:I.XXXIV},

$\therefore \polygonBCD = \polygonEGH$ \byref{ax:I.VII}.
\end{center}

\qed

\starttheorem{Prop XXXIX. Theor.}\label{prop:I.XXXIX}
\defineNewPicture{
pair A, B, C, D, E, F, G;
A := (0, 0);
B := A shifted (5/2u, 0);
C := A shifted (3/4u, -5/2u);
D := C shifted (3u, 0);
E = whatever[A, D] = whatever[B, C];
F := 11/8[C, B];
G := 13/8[C, B];
draw byPolygon(A,B,E)(byred);
draw byPolygon(A,E,C)(byyellow);
draw byPolygon(B,E,D)(byblack);
draw byPolygon(E,C,D)(byyellow);
draw byPolygon(F,B,D)(byblue);
byLineDefine(A, F, byred, SOLID_LINE, REGULAR_WIDTH);
byLineDefine(D, F, byyellow, SOLID_LINE, REGULAR_WIDTH);
byLineDefine(A, B, byblue, SOLID_LINE, REGULAR_WIDTH);
byLineDefine(C, D, byblack, SOLID_LINE, REGULAR_WIDTH);
byLineDefine(C, G, byblack, DASHED_LINE, REGULAR_WIDTH);
draw byNamedLineSeq(-4/5)(AB,AF,DF,CD,CG);
draw byLabelsOnPolygon(F, D, C, A)(OMIT_FIRST_LABEL+OMIT_LAST_LABEL, 0);
draw byLabelsOnPolygon(A, B, F)(OMIT_FIRST_LABEL+OMIT_LAST_LABEL, 0);
draw byLabelsOnPolygon(A, F, G, noPoint)(ALL_LABELS, 0);
}
\drawCurrentPictureInMargin
\problem{T}{riângulos}{iguais
\drawPolygon[bottom][polygonADC]{AEC, ECD} e~\drawPolygon[bottom][polygonBDC]{BED, ECD}
sobre a mesma base (\drawUnitLine{CD}) e do mesmo lado dela, estão entre as mesmas paralelas.}

\begin{center}
Se \drawUnitLine{AB}, que une os vértices dos triângulos, não for $\parallel \drawUnitLine{CD}$, trace $\drawUnitLine{AF} \parallel \drawUnitLine{CD}$ \byref{prop:I.XXXI}, encontrando \drawUnitLine{CG}.

Trace \drawUnitLine{DF}.

$\because \drawUnitLine{AF} \parallel \drawUnitLine{CD}$ \byref{\constref}\\ %Because
$\polygonADC =
\drawPolygon[bottom][polygonFDC]{BED, ECD, FBD}$ \byref{prop:I.XXXVII};\\
mas $\polygonADC = \polygonBDC$ \byref{\hypref};

$\therefore \polygonBDC = \polygonFDC$, uma parte igual ao todo, o que é absurdo.

$\therefore \drawUnitLine{AF} \nparallel \drawUnitLine{CD}$; e da mesma forma pode ser demonstrado que nenhuma outra linha exceto \drawUnitLine{AB} é $\parallel \drawUnitLine{CD}$; $\therefore \drawUnitLine{AB} \parallel \drawUnitLine{CD}$.
\end{center}

\qed

\starttheorem{Prop XL. Theor.}\label{prop:I.XL}
\defineNewPicture{
pair A, B, C, D, E, F, G, H, d;
A := (0, 0);
B := A shifted (3/2u, 0);
C := A shifted (-3/2u, -9/4u);
d := (7/4u, 0);
D := C shifted d;
E := B shifted (-2/3u, -9/4u);
F := E shifted d;
G := 5/4[E, B];
H := 2[B, G];
draw byPolygon(A,C,D)(byyellow);
draw byPolygon(B,E,F)(byred);
draw byPolygon(G,B,F)(byblue);
draw byLine(E, H, byblack, DASHED_LINE, REGULAR_WIDTH);
byLineDefine(A, G, byred, SOLID_LINE, REGULAR_WIDTH);
byLineDefine(F, G, byyellow, SOLID_LINE, REGULAR_WIDTH);
byLineDefine(A, B, byblue, SOLID_LINE, REGULAR_WIDTH);
byLineDefine(C, D, byblack, SOLID_LINE, REGULAR_WIDTH);
byLineDefine(E, F, byblack, SOLID_LINE, REGULAR_WIDTH);
byLineDefine(D, E, byblue, DASHED_LINE, REGULAR_WIDTH);
draw byNamedLineSeq(-4/5)(AB,AG,FG,EF,DE,CD);
draw byLabelsOnPolygon(G, F, E, D, C, A)(OMIT_FIRST_LABEL+OMIT_LAST_LABEL, 0);
draw byLabelsOnPolygon(A, B, G)(OMIT_FIRST_LABEL+OMIT_LAST_LABEL, 0);
draw byLabelsOnPolygon(A, G, H, noPoint)(ALL_LABELS, 0);
}
\drawCurrentPictureInMargin
\problem{T}{riângulos}{iguais
(\drawFromCurrentPicture[bottom][polygonACD]{
startAutoLabeling;
draw byNamedPolygon (ACD);
stopAutoLabeling;
draw byNamedLineFull(A, A, 1, 1, 0, 0)(CD);
}
e
\drawFromCurrentPicture[bottom][polygonBEF]{
startAutoLabeling;
draw byNamedPolygon (BEF);
stopAutoLabeling;
draw byNamedLineFull(B, B, 1, 1,  0, 0)(EF);
}) sobre bases iguais, e do mesmo lado, estão entre as mesmas paralelas.}

\begin{center}
Se \drawSizedLine{AB} que une os vértices dos triângulos não for $\parallel \drawSizedLine{CD,DE,EF}$,\\
trace \drawSizedLine{AG} $\parallel \drawSizedLine{CD,DE,EF}$ \byref{prop:I.XXXI},\\
encontrando \drawSizedLine{EH}.\\
Trace \drawSizedLine{FG}.

$\because \drawSizedLine{AG} \parallel \drawSizedLine{CD,DE,EF}$ \byref{\constref}\\ %Because
$\polygonACD =
\drawFromCurrentPicture[bottom][polygonGEF]{
startAutoLabeling;
draw byNamedPolygon (BEF, GBF);
stopAutoLabeling;
draw byNamedLineFull(G, G, 1, 1,  0, 0)(EF);
}$ mas $\polygonACD = \polygonBEF$

$\therefore \polygonBEF = \polygonGEF$, uma parte igual ao todo, o que é absurdo.

$\therefore \drawSizedLine{AG} \nparallel \drawSizedLine{CD,DE,EF}$: e da mesma forma pode ser demonstrado que nenhuma outra linha exceto \drawSizedLine{AB} é $\parallel \drawSizedLine{CD,DE,EF}$: $\therefore \drawSizedLine{AB} \parallel \drawSizedLine{CD,DE,EF}$.
\end{center}

\qed

\starttheorem{Prop XLI. Theor.}\label{prop:I.XLI}
\defineNewPicture{
pair A, B, C, D, E, F, G, d;
A := (0, 0);
d := (2u, 0);
B := A shifted d;
C := B shifted (2u, 0);
D := A shifted (4/3u, -5/2u);
E := D shifted d;
F = whatever[B, E] = whatever[D, C];
G = whatever[A, E] = whatever[D, C];
draw byPolygon(A,B,F,G)(byyellow);
draw byPolygon(G,F,E)(byyellow);
draw byPolygon(A,G,D)(byblue);
draw byPolygon(D,E,G)(byblue);
draw byPolygon(C,F,E)(byred);
draw byLine(A, E, byred, SOLID_LINE, REGULAR_WIDTH);
draw byLineFull(A, C, byblack, 1, 0)(D, E, 1, 1, 0);
draw byLineFull(D, E, byblack, 0, 0)(A, C, 1, 1, 0);
draw byLabelsOnPolygon(E, D, A, B, C)(ALL_LABELS, 0);
}
\drawCurrentPictureInMargin
\problem[3]{S}{e}{um paralelogramo
\drawPolygon[bottom][polygonABED]{ABFG,GFE,AGD,DEG}
e um triângulo
\drawPolygon[bottom][polygonCED]{GFE,DEG,CFE}
estão sobre a mesma base \drawUnitLine{DE} e entre as mesmas paralelas \drawUnitLine{AC} e \drawUnitLine{DE}, o paralelogramo é o dobro do triângulo.}

\begin{center}
Trace \drawUnitLine{AE} a diagonal.

Então
$\drawPolygon[bottom][polygonAED]{AGD,DEG} = \polygonCED$ \byref{prop:I.XXXVII}\\
$\polygonABED = \mbox{ duas vezes } \polygonAED$ \byref{prop:I.XXXIV}

$\therefore \polygonABED = \mbox{ duas vezes } \polygonCED$.
\end{center}

\qed

\startproblem{Prop XLII. Prob.}\label{prop:I.XLII}
\defineNewPicture{
pair A, B, C, D, E, F, G, H, I, J, d[];
A := (0, 0);
d1 := (8/5u, 0);
B := A shifted d1;
C := B shifted (3/2u, 0);
D := A shifted (u, -13/5u);
E := D shifted d1;
F := (B--E) intersectionpoint (D--C);
G := 2[D, E];
d2 := (-xpart(E)+xpart(D)-1/2u, 0);
H := B shifted d2;
I := E shifted d2;
J := D shifted d2;
draw byPolygon(A,B,F,D)(byyellow);
draw byPolygon(D,E,F)(byyellow);
draw byPolygon(C,F,E)(byblue);
draw byPolygon(E,C,G)(byblack);
byAngleDefine(B, E, D, byblue, SOLID_SECTOR);
byAngleDefine(H, I, J, byyellow, SOLID_SECTOR);
setAttribute("angle", "Standalone", "HIJ", 1);
draw byNamedAngleResized();
draw byLine(B, E, byred, SOLID_LINE, REGULAR_WIDTH);
byLineDefine(A, D, byred, DASHED_LINE, REGULAR_WIDTH);
byLineDefine(C, E, byyellow, SOLID_LINE, REGULAR_WIDTH);
byLineDefine(A, C, byblue, SOLID_LINE, REGULAR_WIDTH);
byLineDefine(D, E, byblack, SOLID_LINE, REGULAR_WIDTH);
byLineDefine(E, G, byblack, DASHED_LINE, REGULAR_WIDTH);
draw byNamedLineSeq(0)(CE,AC,AD,DE,EG,noLine);
draw byLabelsOnPolygon(G, E, D, A, B, C)(ALL_LABELS, 0);
startAutoLabeling;
draw byNamedAngleSidesFull(HIJ)();
stopAutoLabeling;
}
\drawCurrentPictureInMargin
\problem[3]{C}{onstruir}{um paralelogramo igual a um triângulo dado
\drawPolygon[bottom][polygonCDG]{DEF,CFE,ECG} e tendo um ângulo igual a um ângulo retilíneo dado \drawAngle{I}.}

\begin{center}
Faça $\drawUnitLine{DE} = \drawUnitLine{EG}$ \byref{prop:I.X}\\
Trace \drawUnitLine{CE}.

Faça $\drawAngle{E} = \drawAngle{I}$ \byref{prop:I.XXIII}\\
Trace $\left\{
	\begin{aligned}
		\drawUnitLine{AD} &\parallel \drawUnitLine{BE}\\
		\drawUnitLine{AC} &\parallel \drawUnitLine{DE}\\
	\end{aligned}
	\right\}$ \byref{prop:I.XXXI}

$\drawPolygon[bottom][polygonABED]{ABFD,DEF}
= \mbox{ duas vezes }
\drawPolygon[bottom][polygonCED]{DEF,CFE}$ \byref{prop:I.XLI}\\
mas $\polygonCED =
\drawPolygon[bottom][polygonDCG]{ECG}$ \byref{prop:I.XXXVIII}

$\therefore \polygonABED = \polygonCDG$.
\end{center}

\qed

\starttheorem{Prop XLIII. Theor.}\label{prop:I.XLIII}
\defineNewPicture[1/2]{
pair A, B, C, D, E, F, G, H, I, d[];
path q[];
d1 := (5/2u, 0);
d2 := (-u, -3u);
A := (0, 0);
B := A shifted d1;
C := A shifted d2;
D := C shifted d1;
E := 2/5[A, D];
q1 := (E shifted d1) -- (E shifted -d1);
q2 := (E shifted d2) -- (E shifted -d2);
F := q1 intersectionpoint (A--C);
G := q1 intersectionpoint (B--D);
H := q2 intersectionpoint (A--B);
I := q2 intersectionpoint (C--D);
draw byPolygon(A,E,H)(byyellow);
draw byPolygon(A,E,F)(byyellow);
draw byPolygon(H,B,G,E)(byblue);
draw byPolygon(F,C,I,E)(byblack);
draw byPolygon(I,D,E)(byred);
draw byPolygon(G,D,E)(byred);
draw byLabelsOnPolygon(C, F, A, H, B, G, D, I)(ALL_LABELS, 0);
draw byLabelsOnPolygon(F, E, H)(OMIT_FIRST_LABEL+OMIT_LAST_LABEL, 0);
}
\drawCurrentPictureInMargin
\problem{O}{s}{complementos
\drawPolygon[bottom][polygonHBGE]{HBGE}
e
\drawPolygon[bottom][polygonFCIE]{FCIE}
dos paralelogramos que estão sobre a diagonal de um paralelogramo são iguais.}

\begin{center}
$\drawPolygon[bottom][polygonADC]{AEF,FCIE,IDE} = \drawPolygon[bottom][polygonABD]{AEH,HBGE,GDE}$ \byref{prop:I.XXXIV}\\
e $\drawPolygon[bottom][polygonAEFpIDE]{AEF,IDE} = \drawPolygon[bottom][polygonAEHpGDE]{AEH,GDE}$ \byref{prop:I.XXXIV}

$\therefore \polygonFCIE = \polygonHBGE$ \byref{ax:I.III}.
\end{center}

\qed

\startproblem{Prop XLIV. Prob.}\label{prop:I.XLIV}
\defineNewPicture{
pair A, B, C, D, E, F, G, H, I, J, K, L, M, N, O, d[];
path q[];
d1 := (3u, 0);
d2 := (-3/2u, -3u);
d3 := (3/2u, 5/2u);
d4 := -d2 +1/2d1;
A := (0, 0);
B := A shifted d1;
C := A shifted d2;
D := C shifted d1;
E := 2/5[C, B];
q1 := (E shifted d1) -- (E shifted -d1);
q2 := (E shifted d2) -- (E shifted -d2);
F := q1 intersectionpoint (A--C);
G := q1 intersectionpoint (B--D);
H := q2 intersectionpoint (A--B);
I := q2 intersectionpoint (C--D);
J := A shifted d3;
K := J shifted (2(xpart(A)-xpart(H)), 0);
L := (xpart(1/3[J, K]), ypart(F)-ypart(A)+ypart(J));
M := A shifted d4;
N := F shifted d4;
O := E shifted d4;
draw byPolygon(J,K,L)(byred);
draw byPolygon(A,H,E,F)(byyellow);
draw byPolygon(E,G,D,I)(byblue);
byAngleDefine(A, F, E, byblue, SOLID_SECTOR);
byAngleDefine(F, E, I, byred, SOLID_SECTOR);
byAngleDefine(E, I, D, byblack, SOLID_SECTOR);
byAngleDefine(M, N, O, byyellow, SOLID_SECTOR);
setAttribute("angle", "Standalone", "MNO", 1);
draw byNamedAngleResized();
draw byLine(B, E, byred, SOLID_LINE, REGULAR_WIDTH);
draw byLine(E, C, byblack, SOLID_LINE, THIN_WIDTH);
byLineDefine(A, F, byred, DASHED_LINE, REGULAR_WIDTH);
byLineDefine(F, C, byblack, SOLID_LINE, THIN_WIDTH);
draw byLine(H, E, byblue, DASHED_LINE, REGULAR_WIDTH);
byLineDefine(B, G, byyellow, SOLID_LINE, REGULAR_WIDTH);
byLineDefine(A, H, byblue, SOLID_LINE, REGULAR_WIDTH);
byLineDefine(H, B, byblack, SOLID_LINE, REGULAR_WIDTH);
byLineDefine(F, E, byblack, DASHED_LINE, REGULAR_WIDTH);
byLineDefine(E, G, byblack, SOLID_LINE, REGULAR_WIDTH);
byLineDefine(C, D, byyellow, DASHED_LINE, REGULAR_WIDTH);
draw byNamedLineSeq(0)(FE,EG,BG,HB,AH,AF,FC,CD);
draw byLabelsOnPolygon(K, J, L)(ALL_LABELS, 0);
draw byLabelsOnPolygon(F, A, H, B, G, D, I, C)(ALL_LABELS, 0);
draw byLabelsOnPolygon(F, E, H)(OMIT_FIRST_LABEL+OMIT_LAST_LABEL, 0);
startAutoLabeling;
draw byNamedAngleSidesFull(MNO)();
stopAutoLabeling;
}
\drawCurrentPictureInMargin
\problem[4]{A}{}{uma linha reta dada (\drawUnitLine{EG}) aplicar um paralelogramo igual a um triângulo dado (\drawPolygon[middle][polygonJKL]{JKL}), e tendo um ângulo igual a um ângulo retilíneo dado (\drawAngle{N}).}

\begin{center}
Faça $\drawPolygon[middle][polygonAHEF]{AHEF} = \polygonJKL$ com $\drawAngle{F} = \drawAngle{N}$ \byref{prop:I.XLII}\\
e tendo um de seus lados \drawUnitLine{FE} conterminante com e na continuação de \drawUnitLine{EG}.

Prolongue \drawUnitLine{AH} até encontrar $\drawUnitLine{BG} \parallel \drawUnitLine{HE}$\\
trace \drawUnitLine{BE} prolongue até encontrar \drawUnitLine{AF} continuado;\\
trace $\drawUnitLine{CD} \parallel \drawUnitLine{FE,EG}$ encontrando \drawUnitLine{BG} prolongado e prolongue \drawUnitLine{HE}.

$\polygonAHEF = \drawPolygon[middle][polygonEGDI]{EGDI}$ \byref{prop:I.XLIII}\\
mas $\polygonAHEF = \polygonJKL$ \byref{\constref}

$\therefore \polygonEGDI = \polygonJKL$;

e $\drawAngle{F} = \drawAngle{E} =\drawAngle{I} = \drawAngle{N}$ \byref{prop:I.XXIX,\constref}. % improvement: proposition 19 in the original seems to be a typo
\end{center}

\qed

\startproblem{Prop XLV. Prob.}\label{prop:I.XLV}
\defineNewPicture{
pair A, B, C, D, E, F, G, H, I, J, K, L, M, N, O, P, d[];
numeric a, h[], b[], s[];
a := 15;
A := (0, 0);
B := A shifted (0, 2u);
C := A shifted (4/3u, u);
D := A shifted (2u, -3/2u);
E := A shifted (-6/5u, -u);
b1 := arclength(B--C);
h1 := distanceToLine(A, B--C);
s1 := (b1 * h1)/2;
b2 := arclength(C--D);
h2 := distanceToLine(A, C--D);
s2 := (b2 * h2)/2;
b3 := arclength(D--E);
h3 := distanceToLine(A, D--E);
s3 := (b3 * h3)/2;
d1 := (0, ypart(D)-u);
d2 := (0, -b3/2) rotated -a;
d3 := (h3*(1/cosd(a)), 0);
d6 := (u, 0);
F := (-u, 0) shifted d1;
G := F shifted d3;
H := F shifted d2;
I := G shifted d2;
d4 := (2*(s2/b3)*(1/cosd(a)), 0);
J := G shifted d4;
K := J shifted d2;
d5 := (2*(s1/b3)*(1/cosd(a)), 0);
L := J shifted d5;
M := L shifted d2;
N := L shifted d6;
O := M shifted d6;
P := K shifted d6;
draw byPolygon(A,B,C)(byred);
draw byPolygon(A,C,D)(byyellow);
draw byPolygon(A,D,E)(byblue);
byLineDefine(A, C, byblue, SOLID_LINE, REGULAR_WIDTH);
byLineDefine(A, D, byred, SOLID_LINE, REGULAR_WIDTH);
draw byNamedLineSeq(0)(AC,AD);
draw byPolygon(F,G,I,H)(byblue);
draw byPolygon(G,J,K,I)(byyellow);
draw byPolygon(J,L,M,K)(byred);
byAngleDefine(G, I, H, byyellow, SOLID_SECTOR);
byAngleDefine(J, K, I, byblack, SOLID_SECTOR);
byAngleDefine(L, M, K, byblue, SOLID_SECTOR);
byAngleDefine(N, O, P, byred, SOLID_SECTOR);
setAttribute("angle", "Standalone", "NOP", 1);
draw byNamedAngleResized();
draw byLine(G, I, byred, SOLID_LINE, REGULAR_WIDTH);
draw byLine(J, K, byblue, SOLID_LINE, REGULAR_WIDTH);
draw byLabelsOnPolygon(A, B, C, D, E)(ALL_LABELS, 0);
draw byLabelsOnPolygon(F, G, J, L, M, K, I, H)(ALL_LABELS, 0);
startAutoLabeling;
draw byNamedAngleSidesFull(NOP)();
stopAutoLabeling;
}
\drawCurrentPictureInMargin
\problem[2]{C}{onstruir}{um paralelogramo igual a uma figura retilínea dada (\drawPolygon[middle][polygonABCDE]{ABC,ACD,ADE}) e tendo um ângulo igual a um ângulo retilíneo dado (\drawAngle{O}).}

\begin{center}
Trace \drawUnitLine{AD} e \drawUnitLine{AC} dividindo a figura retilínea em triângulos.

Construa $\drawPolygon{FGIH} = \drawPolygon{ADE}$\\
tendo $\drawAngle{I} = \drawAngle{O}$ \byref{prop:I.XLII}

a \drawUnitLine{GI} aplique $\drawPolygon{GJKI} = \drawPolygon{ACD}$\\
tendo $\drawAngle{K} = \drawAngle{O}$ \byref{prop:I.XLIV}

a \drawUnitLine{JK} aplique $\drawPolygon{JLMK} = \drawPolygon{ABC}$\\
tendo $\drawAngle{M} = \drawAngle{O}$ \byref{prop:I.XLIV}

$\therefore \drawPolygon[middle][polygonFLMH]{FGIH,GJKI,JLMK} = \polygonABCDE$

e \polygonFLMH é um paralelogramo. \byref{prop:I.XXIX,prop:I.XIV,prop:I.XXX}\\
tendo $\drawAngle{M} = \drawAngle{O}$.
\end{center}

\qed

\startproblem{Prop XLVI. Prob.}\label{prop:I.XLVI}
\defineNewPicture{
pair A, B, C, D;
numeric d;
d := 7/2u;
A := (0, 0);
B := A shifted (d, 0);
C := A shifted (0, -d);
D := A shifted (d, -d);
byAngleDefine(B, A, C, byblack, SOLID_SECTOR);
byAngleDefine(D, B, A, byblue, SOLID_SECTOR);
byAngleDefine(C, D, B, byred, SOLID_SECTOR);
byAngleDefine(A, C, D, byyellow, SOLID_SECTOR);
draw byNamedAngleResized();
byLineDefine(A, B, byred, SOLID_LINE, REGULAR_WIDTH);
byLineDefine(B, D, byyellow, SOLID_LINE, REGULAR_WIDTH);
byLineDefine(D, C, byblack, SOLID_LINE, REGULAR_WIDTH);
byLineDefine(C, A, byblue, SOLID_LINE, REGULAR_WIDTH);
draw byNamedLineSeq(0)(AB,BD,DC,CA);
draw byLabelsOnPolygon(A, B, D, C)(ALL_LABELS, 0);
}
\drawCurrentPictureInMargin
\problem{S}{obre}{uma linha reta dada (\drawUnitLine{DC}) construir um quadrado.}

\begin{center}
Trace $\drawUnitLine{CA} \perp \mbox{ e } = \drawUnitLine{DC}$ \byref{prop:I.XI,prop:I.III}

Trace $\drawUnitLine{AB} \parallel \drawUnitLine{DC}$,\\
e encontrando \drawUnitLine{BD} traçado $\parallel \drawUnitLine{CA}$.

Em
\drawFromCurrentPicture[bottom][polygonABDC]{
startTempAngleScale(angleScale*3/4);
draw byNamedAngle(A,B,C,D);
draw byNamedLineSeq(0)(AB,BD,DC,CA);
draw byLabelsOnPolygon(A, B, D, C)(ALL_LABELS, 0);
stopTempAngleScale;
}
$\drawUnitLine{CA} = \drawUnitLine{DC}$ \byref{\constref}\\
$\drawAngle{C} = \mbox{um ângulo reto}$ \byref{\constref}

$\therefore \drawAngle{D} = \drawAngle{C} = \mbox{um ângulo reto}$ \byref{prop:I.XXIX},\\
e~os~lados e~ângulos restantes devem ser iguais \byref{prop:I.XXXIV}.

E $\therefore \polygonABDC$ é um quadrado \byref{def:I.XXX}. % improvement: in the original definition 27 was wrongly referenced
\end{center}

\qed

\starttheorem{Prop XLVII. Theor.}\label{prop:I.XLVII}
\defineNewPicture[1/2]{
pair A, B, C, D, E, F, G, H, I, J, K, L, M, d[];
%byPointLabelDefine(A, "α");
%byPointLabelDefine(B, "β");
%byPointLabelDefine(C, "γ");
%byPointLabelDefine(D, "δ");
%byPointLabelDefine(E, "ε");
%byPointLabelDefine(F, "ζ");
%byPointLabelDefine(G, "η");
%byPointLabelDefine(H, "θ");
%byPointLabelDefine(I, "ι");
%byPointLabelDefine(J, "κ");
%byPointLabelDefine(K, "λ");
%byPointLabelDefine(L, "μ");
%byPointLabelDefine(M, "ν");
A := (0, 0);
B := A shifted (-7/10u, -8/7u);
C = whatever[A, A shifted ((A-B) rotated 90)] = whatever[B, B shifted dir(0)];
d1 := (B-A) rotated -90;
D := A shifted d1;
E := B shifted d1;
d2 := (A-C) rotated -90;
F := C shifted d2;
G := A shifted d2;
d3 := (C-B) rotated -90;
H := B shifted d3;
I := C shifted d3;
J = whatever[A, A shifted dir(90)];
J = whatever[B, C];
K = whatever[A, A shifted dir(90)];
K = whatever[H, I];
L = whatever[B, F];
L = whatever[A, C];
M = whatever[A, I];
M = whatever[B, C];
draw byPolygon(A,B,E,D)(byblack);
draw byPolygon(L,A,G,F)(byred);
draw byPolygon(C,L,F)(byred);
draw byPolygon(J,M,I,K)(byblue);
draw byPolygon (M,C,I)(byblue);
draw byPolygon(B,J,K,H)(byyellow);
byAngleDefine(F, C, A, byyellow, SOLID_SECTOR);
byAngleDefine(B, C, I, byyellow, SOLID_SECTOR);
byAngleDefine(A, C, B, byblack, SOLID_SECTOR);
draw byNamedAngleResized();
draw byLineFull(A, K, byblack, 1, 0)(I, I, 1, 1, -1);
draw byLineFull(B, F, byblack, 0, 0)(G, G, 1, 1, -1);
draw byLineFull(A, I, byblack, 0, 0)(K, K, 1, 1, 1);
byLineDefine(C, F, byblue, DASHED_LINE, REGULAR_WIDTH);
byLineDefine(C, I, byred, DASHED_LINE, REGULAR_WIDTH);
draw byNamedLineSeq(0)(CF,CI);
byLineDefine(A, B, byyellow, SOLID_LINE, REGULAR_WIDTH);
byLineDefine(B, C, byred, SOLID_LINE, REGULAR_WIDTH);
byLineDefine(C, A, byblue, SOLID_LINE, REGULAR_WIDTH);
draw byNamedLineSeq(1)(AB,BC,CA);
byLineDefine.CAb(C, A, byblack, SOLID_LINE, REGULAR_WIDTH);
byLineStylize (M, M, 1, 0, -1) (CAb);
byLineDefine.AMb(A, M, byblack, SOLID_LINE, REGULAR_WIDTH);
byLineStylize (C, C, 0, 1, -1) (AMb);
byLineDefine.BCb(B, C, byblack, SOLID_LINE, REGULAR_WIDTH);
byLineStylize (L, L, 0, 1, -1) (BCb);
byLineDefine.BLb(L, B, byblack, SOLID_LINE, REGULAR_WIDTH);
byLineStylize (C, C, 1, 0, -1) (BLb);
draw byLabelsOnPolygon(B, E, D, A, G, F, C, I, K, H)(ALL_LABELS, -1);
draw byLabelsOnPolygon(A, J, C)(OMIT_FIRST_LABEL+OMIT_LAST_LABEL, 1);
}
\drawCurrentPictureInMargin
\problem[4]{E}{m}{um triângulo retângulo \drawLine[bottom][triangleABC]{CA,BC,AB} o quadrado sobre a hipotenusa \drawUnitLine{BC} é igual à soma dos quadrados dos lados (\drawUnitLine{CA} e \drawUnitLine{AB}).}

\begin{center}
Sobre \drawUnitLine{BC}, \drawUnitLine{CA}, \drawUnitLine{AB} descreva quadrados, \byref{prop:I.XLVI}

Trace $\drawUnitLine{AK} \parallel \drawUnitLine{CI}$ \byref{prop:I.XXXI}\\
também trace \drawUnitLine{BF} e \drawUnitLine{AI}.\\
$\drawAngle{BCI} = \drawAngle{FCA}$.

A cada um adicione \drawAngle{ACB} $\therefore \drawAngle{BCI,ACB} = \drawAngle{FCA,ACB}$,\\
$\drawUnitLine{BC} = \drawUnitLine{CI}$ e $\drawUnitLine{CA} = \drawUnitLine{CF}$;

$\therefore
\drawFromCurrentPicture[middle][polygonAFC]{
draw byNamedPolygon(MCI);
draw byNamedAngle(ACB);
draw byNamedLine(CAb,AMb);
draw byLabelsOnPolygon(I,A,C)(OMIT_LABELS_AT_STRAIGHT_ANGLES, -1);
}
=
\drawFromCurrentPicture[middle][polygonBLC]{
draw byNamedPolygon(CLF);
draw byNamedAngle(ACB);
draw byNamedLine(BCb,BLb);
draw byLabelsOnPolygon(B,F,C)(OMIT_LABELS_AT_STRAIGHT_ANGLES, -1);
}
$.

Novamente, $\because \drawUnitLine{AB} \parallel \drawUnitLine{CF}$\\ %because
$\drawPolygon[middle][polygonACFG]{LAGF,CLF} = \mbox{ duas vezes } \polygonBLC$,\\
e $\drawPolygon[middle][polygonJMCK]{JMIK,MCI} = \mbox{ duas vezes } \polygonAFC$;

$\therefore \polygonACFG = \polygonJMCK$.

Da mesma forma pode ser mostrado que $\drawPolygon[middle][polygonABED]{ABED} = \drawPolygon[middle][polygonBJKH]{BJKH}$;

portanto $\drawPolygon[middle][polygonABEDpACFG]{ABED,LAGF,CLF} = \drawPolygon[middle][polygonBCIH]{JMIK,MCI,BJKH}$.

\end{center}

\qed

\starttheorem{Prop XLVIII. Theor.}\label{prop:I.XLVIII}
\defineNewPicture{
pair A, B, C, D;
numeric d;
d := 7/4u;
A := (0, 0);
B := A shifted (0, 4u);
C := A shifted (d, 0);
D := A shifted (-d, 0);
byAngleDefine(B, A, C, byred, SOLID_SECTOR);
byAngleDefine(D, A, B, byyellow, SOLID_SECTOR);
draw byNamedAngleResized();
draw byLine(A, B, byblue, SOLID_LINE, REGULAR_WIDTH);
byLineDefine(A, C, byblack, SOLID_LINE, REGULAR_WIDTH);
byLineDefine(A, D, byblack, DASHED_LINE, REGULAR_WIDTH);
byLineDefine(B, C, byred, SOLID_LINE, REGULAR_WIDTH);
byLineDefine(B, D, byred, DASHED_LINE, REGULAR_WIDTH);
draw byNamedLineSeq(0)(AC,AD,BD,BC);
draw byLabelsOnPolygon(D, B, C, A)(ALL_LABELS, 0);
}
\drawCurrentPictureInMargin
\problem{S}{e}{o quadrado de um lado (\drawUnitLine{BC}) de um triângulo é igual aos quadrados dos outros dois lados (\drawUnitLine{AB} e \drawUnitLine{AC}), o ângulo (\drawAngle{BAC}) subtendido por esse lado é um ângulo reto.} % improvement: DAB sould be BAC https://github.com/jemmybutton/byrne-euclid/issues/31#issuecomment-350511743

\begin{center}
Trace $\drawUnitLine{AD} \perp \drawUnitLine{AB}$ e $= \drawUnitLine{AC}$ \byref{prop:I.XI,prop:I.III}\\
e trace \drawUnitLine{BD} também.

Como $\drawUnitLine{AD} = \drawUnitLine{AC}$ \byref{\constref}\\
$\drawUnitLine{AD}^2 = \drawUnitLine{AC}^2$;

$\therefore \drawUnitLine{AD}^2 + \drawUnitLine{AB}^2 = \drawUnitLine{AC}^2 + \drawUnitLine{AB}^2$

mas $\drawUnitLine{AD}^2 + \drawUnitLine{AB}^2 = \drawUnitLine{BD}^2$ \byref{prop:I.XLVII},\\
e $\drawUnitLine{AC}^2 + \drawUnitLine{AB}^2 = \drawUnitLine{BC}^2$ \byref{\hypref}

$\therefore \drawUnitLine{BD}^2 = \drawUnitLine{BC}^2$,

$\therefore \drawUnitLine{BD} = \drawUnitLine{BC}$;

e $\therefore \drawAngle{DAB} = \drawAngle{BAC}$ \byref{prop:I.VIII},

consequentemente \drawAngle{BAC} é um ângulo reto.
\end{center}

\qed

\part{Livro II}

\startdefinition{Definição I}\label{def:II.I}

\defineNewPicture{
pair A, B, C, D;
numeric w, h;
w := 7/2u;
h := 3u;
A := (0, 0);
B := (w, 0);
C := (0, h);
D := (w, h);
draw byPolygon(A,B,D,C)(byblue);
byLineDefine(A, B, byblack, SOLID_LINE, REGULAR_WIDTH);
byLineDefine(A, C, byred, SOLID_LINE, REGULAR_WIDTH);
draw byNamedLineSeq(0)(AB,AC);
draw byLabelsOnPolygon(A, C, D, B)(ALL_LABELS, 0);
}
\drawCurrentPictureInMargin
\problem{U}{m}{retângulo ou um paralelogramo retângulo é dito ser compreendido por quaisquer dois de seus lados adjacentes ou contíguos.}

Assim: diz-se que o paralelogramo retângulo \drawPolygon{ABDC} é compreendido pelos lados \drawUnitLine{AB} e \drawUnitLine{AC}; ou pode ser brevemente designado por $\drawUnitLine{AB} \cdot \drawUnitLine{AC}$.

Se os lados adjacentes forem iguais; i.\ e.\ $\drawUnitLine{AB} = \drawUnitLine{AC}$, então $\drawUnitLine{AB} \cdot \drawUnitLine{AC}$, que é a expressão do retângulo compreendido por \drawUnitLine{AB} e \drawUnitLine{AC}, é um quadrado, e

é igual a $\left\{
	\begin{aligned}
		\drawUnitLine{AB} \cdot \drawUnitLine{AC}&\mbox{ ou } \drawUnitLine{AC}^2\\
		\drawUnitLine{AB} \cdot \drawUnitLine{AC}&\mbox{ ou } \drawUnitLine{AB}^2\\
	\end{aligned}\right.$

\vfill\pagebreak

\startdefinition{Definição II}\label{def:II.II}

\defineNewPicture{
pair A, B, C, D, E, F, G, H, I, d[];
d1 := (3u, 0);
d2 := (1/2u, 5/2u);
A := (0, 0);
B := A shifted d1;
C := A shifted d2;
D := A shifted d1 shifted d2;
E := 2/3[A, B];
F := 2/3[C, D];
G := 2/3[A, C];
H := 2/3[B, D];
I = whatever[E, F] = whatever[G, H];
draw byPolygon(A,E,I,G)(byblue);
draw byPolygon(E,B,H,I)(byyellow);
draw byPolygon(G,I,F,C)(byyellow);
draw byPolygon(I,H,D,F)(byred);
draw byLabelsOnPolygon(C, F, D, H, B, E, A, G)(ALL_LABELS,0);
draw byLabelsOnPolygon(G, I, F)(OMIT_FIRST_LABEL+OMIT_LAST_LABEL,0);
}
\drawCurrentPictureInMargin
\problem[4]{N}{um}{paralelogramo, a figura composta por um dos paralelogramos em torno da diagonal, juntamente com os dois complementos, chama-se gnômon.}

Assim, \drawPolygon{AEIG,EBHI,GIFC} e \drawPolygon{EBHI,GIFC,IHDF} são chamados de gnômons.

\vfill\pagebreak

\starttheorem{Prop. I. Teor.}\label{prop:II.I}
\defineNewPicture[1/4]{
pair B, C, D, E, G, H, K, L, M, N;
numeric w, h;
w := 7/2u;
h := 3u;
G := (0, 0);
H := (w, 0);
B := (0, h);
C := (w, h);
K := 2/5[G, H];
D := 2/5[B, C];
L := 3/4[G, H];
E := 3/4[B, C];
M := G shifted (0, -2/3u);
N := M shifted (h, 0);
draw byPolygon(G,K,D,B)(byyellow);
draw byPolygon(K,L,E,D)(byblue);
draw byPolygon(L,H,C,E)(byred);
draw byLine(K, D, byblack, DASHED_LINE, REGULAR_WIDTH);
draw byLine(L, E, byblack, DASHED_LINE, REGULAR_WIDTH);
byLineDefine(G, B, byblack, SOLID_LINE, REGULAR_WIDTH);
draw byLine.A(M, N, byblack, SOLID_LINE, REGULAR_WIDTH); % improvement: the line wasn't there, but it should've been
byLineDefine(H, C, byblack, DASHED_LINE, REGULAR_WIDTH);
byLineDefine(G, K, byblue, SOLID_LINE, REGULAR_WIDTH);
byLineDefine(K, L, byred, SOLID_LINE, REGULAR_WIDTH);
byLineDefine(L, H, byyellow, SOLID_LINE, REGULAR_WIDTH);
byLineDefine(B, D, byblue, DASHED_LINE, REGULAR_WIDTH);
byLineDefine(D, E, byred, DASHED_LINE, REGULAR_WIDTH);
byLineDefine(E, C, byyellow, DASHED_LINE, REGULAR_WIDTH);
draw byNamedLineSeq(0)(GK,KL,LH,HC,EC,DE,BD,GB);
draw byLabelsOnPolygon(B, D, E, C, H, L, K, G)(ALL_LABELS, 0);
draw byLabelLine(0)(A);
}
\drawCurrentPictureInMargin
\problem[3]{T}{he}{retângulo compreendido por duas linhas retas, uma das quais é dividida em qualquer número de partes,
$\drawProportionalLine{GK,KL,LH} \cdot \drawProportionalLine{A} = \left\{
	\begin{aligned}
		&\drawProportionalLine{A} \cdot \drawProportionalLine{GK}\\
		+&\drawProportionalLine{A} \cdot \drawProportionalLine{KL}\\
		+&\drawProportionalLine{A} \cdot \drawProportionalLine{LH}\\
	\end{aligned}
\right.$\\
é igual à soma dos retângulos compreendido por da linha indivisa e das várias partes da linha dividida.}

\begin{center}
Trace $\drawProportionalLine{GB} \perp \drawProportionalLine{GK,KL,LH} \mbox{ and} = \drawProportionalLine{GB}$ \byref{prop:I.XI,prop:I.III}; completar os paralelogramos, ou seja, % de melhoria: no original \byref{prop:I.II} foi referenciado em vez de \byref{prop:I.XI}

draw $\left\{
	\begin{aligned}
		\drawProportionalLine{BD,DE,EC} &\parallel \drawProportionalLine{GK,KL,LH} \\
		\vcenter{
		\nointerlineskip\hbox{\drawProportionalLine{KD}}
		\nointerlineskip\hbox{\drawProportionalLine{LE}}
		\nointerlineskip\hbox{\drawProportionalLine{HC}}} &\parallel \drawProportionalLine{GB}\\
	\end{aligned}
\right\}$ \byref{prop:I.XXXI}

$\drawPolygon[bottom]{GKDB,KLED,LHCE} =
\drawPolygon[bottom]{GKDB} +
\drawPolygon[bottom]{KLED} +
\drawPolygon[bottom]{LHCE}$

$\drawPolygon[bottom]{GKDB,KLED,LHCE} = \drawProportionalLine{GK,KL,LH} \cdot \drawProportionalLine{GB}$

$\polygonGKDB = \drawProportionalLine{GK} \cdot \drawProportionalLine{GB}$,
$\polygonKLED = \drawProportionalLine{KL} \cdot \drawProportionalLine{GB}$,

$\polygonLHCE = \drawProportionalLine{LH} \cdot \drawProportionalLine{GB}$

$\therefore \drawProportionalLine{GK,KL,LH} \cdot \drawProportionalLine{A} = \drawProportionalLine{GK} \cdot \drawProportionalLine{A} + \drawProportionalLine{KL} \cdot \drawProportionalLine{A} + \drawProportionalLine{LH} \cdot \drawProportionalLine{A}$.
\end{center}

\qed

\starttheorem{Prop. II. Teor.}\label{prop:II.II}
\defineNewPicture{
pair A, B, C, D, E, F;
numeric w;
w := 7/2u;
A := (0, w);
B := (w, w);
C := 2/3[A, B];
D := (0, 0);
E := (w, 0);
F := 2/3[D, E];
draw byPolygon(A,C,F,D)(byred);
draw byPolygon(C,B,E,F)(byyellow);
draw byLine(C, F, byblack, SOLID_LINE, REGULAR_WIDTH);
byLineDefine(A, D, byblack, DASHED_LINE, REGULAR_WIDTH);
byLineDefine(B, E, byblack, DASHED_LINE, REGULAR_WIDTH);
byLineDefine(A, C, byblue, SOLID_LINE, REGULAR_WIDTH);
byLineDefine(C, B, byred, SOLID_LINE, REGULAR_WIDTH);
draw byNamedLineSeq(0)(AD,AC,CB,BE);
draw byLabelsOnPolygon(A, C, B, E, F, D)(ALL_LABELS, 0);
}
\drawCurrentPictureInMargin
\problem{S}{e}{uma linha reta for dividida em duas partes quaisquer \drawProportionalLine{AC,CB}
$\drawProportionalLine{AC,CB}^2 = \left\{
	\begin{aligned}
		& \drawProportionalLine{AC,CB} \cdot \drawProportionalLine{AC}\\
		+ & \drawProportionalLine{AC,CB} \cdot \drawProportionalLine{CB}\\
	\end{aligned}
\right.$
}

\begin{center}
Describe \drawPolygon[bottom][polygonABED]{ACFD,CBEF} \byref{prop:I.XLVI}

Trace \drawProportionalLine{CF} paralelo a \drawProportionalLine{AD} \byref{prop:I.XXXI}

$\polygonABED = \drawProportionalLine{AC,CB}^2$

$\drawPolygon[bottom]{ACFD} = \drawProportionalLine{CF} \cdot \drawProportionalLine{AC} = \drawProportionalLine{AC,CB} \cdot \drawProportionalLine{AC}$

$\drawPolygon[bottom]{CBEF} = \drawProportionalLine{CF} \cdot \drawProportionalLine{CB} = \drawProportionalLine{AC,CB} \cdot \drawProportionalLine{CB}$

$\polygonABED = \drawPolygon[bottom]{ACFD} + \drawPolygon[bottom]{CBEF}$

$\therefore \drawProportionalLine{AC,CB}^2 = \drawProportionalLine{AC,CB} \cdot \drawProportionalLine{AC} + \drawProportionalLine{AC,CB} \cdot \drawProportionalLine{CB}$.
\end{center}

\qed

\starttheorem{Prop. III. Teor.}\label{prop:II.III}
\defineNewPicture{
pair A, B, C, D, E, F;
numeric w, h;
w := -4u;
h := 11/4u;
A := (0, h);
B := (w, h);
C := (w+h, h);
D := (w+h, 0);
E := (w, 0);
F := (0, 0);
draw byPolygon(A,C,D,F)(byyellow);
draw byPolygon(C,B,E,D)(byred);
byLineDefine(D, F, byred, SOLID_LINE, REGULAR_WIDTH);
byLineDefine(B, C, byblue, SOLID_LINE, REGULAR_WIDTH);
byLineDefine(C, D, byblue, SOLID_LINE, REGULAR_WIDTH);
byLineDefine(D, E, byblue, SOLID_LINE, REGULAR_WIDTH);
byLineDefine(E, B, byblue, SOLID_LINE, REGULAR_WIDTH);
draw byNamedLineSeq(0)(CD,noLine,DF,DE,EB,BC);
draw byLabelsOnPolygon(B, C, A, F, D, E)(ALL_LABELS, 0);
}
\drawCurrentPictureInMargin
\problem{S}{e}{uma linha reta for dividida em duas partes quaisquer \drawProportionalLine{DE,DF}
$\drawProportionalLine{DE,DF} \cdot \drawProportionalLine{DE} = \drawProportionalLine{DE}^2 + \drawProportionalLine{DE} \cdot \drawProportionalLine{DF}$, or \\
$\drawProportionalLine{DE,DF} \cdot \drawProportionalLine{DF} = \drawProportionalLine{DF}^2 + \drawProportionalLine{DE} \cdot \drawProportionalLine{DF}$.}

\begin{center}
Describe \drawPolygon[bottom]{CBED} \byref{prop:I.XLVI}

Describe \drawPolygon[bottom]{ACDF} \byref{prop:I.XXXI}

Then $\drawPolygon[bottom][polygonABEF]{ACDF,CBED} = \polygonCBED + \polygonACDF$, but\\
$\polygonABEF = \drawProportionalLine{DE,DF} \cdot \drawProportionalLine{DE}$ and\\
$\polygonCBED = \drawProportionalLine{DE}^2$, $\polygonACDF = \drawProportionalLine{DE} \cdot \drawProportionalLine{DF}$,

$\therefore \drawProportionalLine{DE,DF} \cdot \drawProportionalLine{DE} = \drawProportionalLine{DE}^2 + \drawProportionalLine{DE} \cdot \drawProportionalLine{DF}$.

In a similar manner it may be readily shown that $\drawProportionalLine{DE,DF} \cdot \drawProportionalLine{DF} = \drawProportionalLine{DF}^2 + \drawProportionalLine{DE} \cdot \drawProportionalLine{DF}$.
\end{center}

\qed

\starttheorem{Prop. IV. Teor.}\label{prop:II.IV}
\defineNewPicture[1/4]{
pair A, B, C, D, E, F, G, H, K;
numeric w;
w := 9/2u;
A := (0, w);
B := (w, w);
C :=2/3[A, B];
D := (0, 0);
E := (w, 0);
F := 2/3[D, E];
H := 2/3[D, A];
K := 2/3[E, B];
G = whatever[H, K] = whatever[F, C];
draw byPolygon(A,C,G,H)(byyellow);
draw byPolygon(G,K,E,F)(byyellow);
draw byPolygon(D,H,G)(byblue);
draw byPolygon(G,B,C)(byred);
byAngleDefine(F, D, G, byyellow, SOLID_SECTOR);
byAngleDefine(F, G, D, byred, SOLID_SECTOR);
byAngleDefine(K, G, B, byblack, SOLID_SECTOR);
byAngleDefine(G, B, K, byblue, SOLID_SECTOR);
draw byNamedAngleResized();
draw byLine(H, G, byred, DASHED_LINE, REGULAR_WIDTH);
draw byLine(G, K, byred, SOLID_LINE, REGULAR_WIDTH);
draw byLine(C, G, byblue, DASHED_LINE, REGULAR_WIDTH);
draw byLine(F, G, byblue, SOLID_LINE, REGULAR_WIDTH);
byLineDefine(E, K, byblue, SOLID_LINE, REGULAR_WIDTH);
byLineDefine(F, E, byred, SOLID_LINE, REGULAR_WIDTH);
byLineDefine(D, F, byblue, SOLID_LINE, REGULAR_WIDTH);
byLineDefine(K, B, byred, SOLID_LINE, REGULAR_WIDTH);
byLineDefine(G, D, byblack, SOLID_LINE, REGULAR_WIDTH);
byLineDefine(B, G, byblack, DASHED_LINE, REGULAR_WIDTH);
draw byNamedLineSeq(1)(BG,GD,DF,FE,EK,KB);
byLineDefine(A, D, byblack, SOLID_LINE, REGULAR_WIDTH);
byLineStylize(B, E, 0, 0, -1)(AD);
byLineDefine(B, A, byblack, SOLID_LINE, REGULAR_WIDTH);
byLineStylize(E, D, 0, 0, -1)(BA);
draw byLabelsOnPolygon(A, C, B, K, E, F, D, H)(ALL_LABELS, 0);
draw byLabelsOnPolygon(H, G, C)(OMIT_FIRST_LABEL+OMIT_LAST_LABEL, 0);
}
\drawCurrentPictureInMargin
\problem{S}{e}{uma linha reta for dividida em duas partes quaisquer \drawProportionalLine{DF,FE}
$\drawProportionalLine{DF,FE}^2 = \drawProportionalLine{DF}^2 + \drawProportionalLine{FE}^2 + \mbox{twice} \drawProportionalLine{DF} \cdot \drawProportionalLine{FE}$
}

\begin{center}
Describe \drawLine[bottom][squareABED]{AD,BA,KB,EK,FE,DF} \byref{prop:I.XLVI}

draw \drawProportionalLine{BG,GD} \byref{post:I.I}\\
and $\left\{
	\begin{aligned}
		\drawProportionalLine{FG,CG} &\parallel \drawProportionalLine{EK,KB}\\
		\drawProportionalLine{HG,GK} &\parallel \drawProportionalLine{DF,FE}\\
	\end{aligned}
\right\}$ \byref{prop:I.XXXI}\\
$\drawAngle{GBK} = \drawAngle{FDG}$ \byref{prop:I.V},\\
$\drawAngle{GBK} = \drawAngle{FGD}$ \byref{prop:I.XXIX},

$\therefore \drawAngle{FDG} = \drawAngle{FGD}$

$\therefore$ by \byref{prop:I.VI,prop:I.XXIX,prop:I.XXXIV}\\ $\drawFromCurrentPicture[bottom][squareFGHD]{
draw byNamedPolygon(DHG);
draw byNamedLineFull(G, G, 1, 0,  0, -1)(DF);
draw byNamedLineFull(D, D, 0, 1,  0, -1)(FG);
draw byLabelsOnPolygon(D, H, G, F)(ALL_LABELS, 0);
} \mbox{ is a square } = \drawProportionalLine{DF}^2$.

For the same reasons \drawFromCurrentPicture[bottom][squareKBCG]{
draw byNamedPolygon(GBC);
draw byNamedLineFull(B, B, 1, 0,  0, -1)(GK);
draw byNamedLineFull(G, G, 0, 1,  0, -1)(KB);
draw byLabelsOnPolygon(G, C, B, K)(ALL_LABELS, 0);
} is a square $= \drawProportionalLine{GK}^2$,

$\drawPolygon[bottom]{ACGH} = \drawPolygon[middle]{GKEF} = \drawProportionalLine{DF} \cdot \drawProportionalLine{GK}$ \byref{prop:I.XLIII}

But $\drawFromCurrentPicture[bottom][squareABEDf]{
draw byNamedPolygon(GBC,DHG,ACGH,GKEF);
draw byNamedLineFull(G, G, 1, 0,  0, -1)(DF);
draw byNamedLineFull(D, D, 0, 1,  0, -1)(FG);
draw byNamedLineFull(B, B, 1, 0,  0, -1)(GK);
draw byNamedLineFull(G, G, 0, 1,  0, -1)(KB);
draw byLabelsOnPolygon(A, B, E, D)(ALL_LABELS, 0);
} = \squareFGHD + \drawPolygon{ACGH} + \drawPolygon{GKEF} + \squareKBCG$,

$\therefore \drawProportionalLine{DF,FE}^2 = \drawProportionalLine{DF}^2 + \drawProportionalLine{FE}^2 + \mbox{ twice } \drawProportionalLine{DF} \cdot \drawProportionalLine{FE}$.
\end{center}

\qed

\starttheorem{Prop. V. Teor.}\label{prop:II.V}
\defineNewPicture[1/4]{
pair A, B, C, D, E, F, G, H, K, L, M;
numeric h;
h := 6u;
A := (0, -h);
B := (0, 0);
C := 1/2[A, B];
D := 2/5[B, C];
E := (1/2h, -1/2h);
F := (1/2h, 0);
G := (xpart(F), ypart(D));
H = whatever[D, G] = whatever[B, E];
K := (xpart(H), ypart(A));
L := (xpart(H), ypart(C));
M := (xpart(H), ypart(B));
draw byPolygon(B,D,H,M)(byblue);
draw byPolygon(D,C,L,H)(byyellow);
draw byPolygon(C,L,K,A)(byblack);
draw byPolygon(M,H,G,F)(byyellow);
draw byPolygon(H,L,E,G)(byred);
draw byLine(H, G, byblack, DASHED_LINE, REGULAR_WIDTH);
draw byLine(D, H, byred, SOLID_LINE, REGULAR_WIDTH);
byLineDefine(L, M, byblack, DASHED_LINE, REGULAR_WIDTH);
byLineDefine(C, L, byred, SOLID_LINE, REGULAR_WIDTH);
byLineDefine(B, D, byred, SOLID_LINE, REGULAR_WIDTH);
byLineDefine(D, C, byblue, SOLID_LINE, REGULAR_WIDTH);
byLineDefine(C, A, byyellow, SOLID_LINE, REGULAR_WIDTH);
byLineDefine(E, B, byblack, SOLID_LINE, REGULAR_WIDTH);
byLineDefine(A, K, byred, DASHED_LINE, REGULAR_WIDTH);
byLineDefine(K, L, byyellow, SOLID_LINE, REGULAR_WIDTH);
byLineDefine(L, E, byblue, DASHED_LINE, REGULAR_WIDTH);
draw byNamedLineSeq(1)(BD,DC,CA,AK,KL,LM,CL,LE,EB);
draw byLabelsOnPolygon(A, C, D, B, M, F, G, E, L, K)(ALL_LABELS, 0);
draw byLabelsOnPolygon(M, H, G)(OMIT_FIRST_LABEL+OMIT_LAST_LABEL, 0);
}
\drawCurrentPictureInMargin
\problem{S}{e}{uma linha reta for dividida \drawProportionalLine{BD,DC} \drawProportionalLine{CA} em duas partes iguais e também \drawProportionalLine{BD} \drawProportionalLine{DC,CA}
$\drawProportionalLine{BD} \cdot \drawProportionalLine{DC,CA} + \drawProportionalLine{DC}^2 = \drawProportionalLine{CA}^2 = \drawProportionalLine{BD,DC}^2$.
}

\begin{center}
Describe \drawPolygon[middle][squareCBFE]{BDHM,DCLH,MHGF,HLEG} \byref{prop:I.XLVI},\\
draw \drawProportionalLine{EB}\\
and $\left\{
	\begin{aligned}
		\drawProportionalLine{DH,HG} & \parallel \drawProportionalLine{CL,LE} \\
		\drawProportionalLine{LM,KL} & \parallel \drawProportionalLine{BD,DC,CA} \\
		\drawProportionalLine{AK} & \parallel \drawProportionalLine{CL,LE}\\
	\end{aligned}
\right\}$ \byref{prop:I.XXXI}

$\drawPolygon{CLKA} = \drawPolygon{DCLH,BDHM}$ \byref{prop:I.XXXVI}\\
$\drawPolygon{MHGF} = \drawPolygon{DCLH}$ \byref{prop:I.XLIII}

$\therefore \mbox{\byref{ax:I.II} } \drawPolygon{DCLH,BDHM,MHGF} = \drawPolygon{DCLH,CLKA} = \drawProportionalLine{BD} \cdot \drawProportionalLine{DC,CA}$
but $\drawPolygon{HLEG} = \drawProportionalLine{DC}^2$ \byref{prop:II.IV}

and $\squareCBFE = \drawProportionalLine{BD,DC}^2$ \byref{\constref}

$\therefore \mbox{\byref{ax:I.II} } \squareCBFE = \drawPolygon{DCLH,CLKA,HLEG}$

$\therefore \drawProportionalLine{BD} \cdot \drawProportionalLine{DC,CA} + \drawProportionalLine{DC}^2 = \drawProportionalLine{CA}^2 = \drawProportionalLine{BD,DC}^2$.
\end{center}

\qed

\starttheorem{Prop. VI. Teor.}\label{prop:II.VI}
\defineNewPicture{
pair A, B, C, D, E, F, G, H, K, L, M;
numeric h, s;
h := 5u;
s := 2/5h;
A := (0, h);
B := (0, 0);
C := A shifted (0, -s);
D := A shifted (0, -2s);
E := (-h+s, h-s);
F := (-h+s, 0);
G := (xpart(F), ypart(D));
H = whatever[D, G] = whatever[B, E];
K := (xpart(H), ypart(A));
L := (xpart(H), ypart(C));
M := (xpart(H), ypart(B));
draw byPolygon(B,D,H,M)(byblue);
draw byPolygon(D,C,L,H)(byyellow);
draw byPolygon(C,L,K,A)(byblack);
draw byPolygon(M,H,G,F)(byyellow);
draw byPolygon(H,L,E,G)(byred);
draw byLine(H, G, byblue, DASHED_LINE, REGULAR_WIDTH);
draw byLine(D, H, byred, SOLID_LINE, REGULAR_WIDTH);
byLineDefine(L, M, byblack, DASHED_LINE, REGULAR_WIDTH);
byLineDefine(C, L, byred, SOLID_LINE, REGULAR_WIDTH);
byLineDefine(B, D, byred, SOLID_LINE, REGULAR_WIDTH);
byLineDefine(D, C, byblue, SOLID_LINE, REGULAR_WIDTH);
byLineDefine(C, A, byyellow, SOLID_LINE, REGULAR_WIDTH);
byLineDefine(E, B, byblack, SOLID_LINE, REGULAR_WIDTH);
byLineDefine(A, K, byred, DASHED_LINE, REGULAR_WIDTH);
byLineDefine(K, L, byyellow, SOLID_LINE, REGULAR_WIDTH);
byLineDefine(L, E, byblack, DASHED_LINE, REGULAR_WIDTH);
draw byNamedLineSeq(1)(BD,DC,CA,AK,KL,LM,CL,LE,EB);
draw byLabelsOnPolygon(F, G, E, L, K, A, C, D, B, M)(ALL_LABELS, 0);
draw byLabelsOnPolygon(L, H, D)(OMIT_FIRST_LABEL+OMIT_LAST_LABEL, 1);
}
\drawCurrentPictureInMargin
\problem{S}{e}{uma linha reta for dividida ao meio \drawProportionalLine{DC,CA} e produzida em qualquer ponto \drawProportionalLine{BD,DC,CA}
$\drawProportionalLine{BD,DC,CA} \cdot \drawProportionalLine{BD} + \drawProportionalLine{DC}^2 = \drawProportionalLine{BD,DC}^2$
}

\begin{center}
Descreva \drawPolygon[bottom][squareCBFE]{BDHM,DCLH,MHGF,HLEG} \byref{prop:I.XLVI}, trace \drawProportionalLine{EB}\\
and $\left\{
	\begin{aligned}
		\drawProportionalLine{HG,DH} & \parallel \drawProportionalLine{CL,LE} \\
		\drawProportionalLine{LM,KL} & \parallel \drawProportionalLine{BD,DC,CA} \\
		\drawProportionalLine{AK} & \parallel \drawProportionalLine{CL,LE}\\
	\end{aligned}
\right\}$ \byref{prop:I.XXXI}

$\drawPolygon[bottom]{MHGF} = \drawPolygon[bottom]{DCLH} = \drawPolygon[bottom]{CLKA}$ \byref{prop:I.XXXVI,prop:I.XLIII}

$\therefore \drawPolygon[bottom]{BDHM,DCLH,MHGF} = \drawPolygon[bottom]{BDHM,DCLH,CLKA} = \drawProportionalLine{BD} \cdot \drawProportionalLine{BD,DC,CA}$;\\
but $\drawPolygon[bottom]{HLEG} = \drawProportionalLine{DC}^2$ \byref{prop:II.IV}

$\therefore \squareCBFE = \drawProportionalLine{CB}^2 = \drawPolygon[bottom]{BDHM,DCLH,CLKA,HLEG}$ \byref{\constref,ax:I.II}

$\therefore \drawProportionalLine{BD,DC,CA} \cdot \drawProportionalLine{BD} + \drawProportionalLine{DC}^2 = \drawProportionalLine{BD,DC}^2$.
\end{center}

\qed

\starttheorem{Prop. VII. Teor.}\label{prop:II.VII}
\defineNewPicture{
pair A, B, C, D, E, F, G, H, N;
numeric w;
w := 7/2u;
A := (0, w);
B := (w, w);
C := 3/5[A, B];
D := (0, 0);
E := (w, 0);
N := 3/5[D, E];
F := 3/5[E, B];
G = whatever[D, B] = whatever[N, C];
H := whatever[A, D] = whatever[F, G];
draw byPolygon(D,N,G,H)(byred);
draw byPolygon(N,E,F,G)(byblack);
draw byPolygon(H,G,C,A)(byyellow);
draw byPolygon(G,F,B,C)(byblue);
draw byLine(G, N, byblue, SOLID_LINE, REGULAR_WIDTH);
draw byLine(G, F, byred, SOLID_LINE, REGULAR_WIDTH);
draw byLine(G, H, byblack, DASHED_LINE, REGULAR_WIDTH);
draw byLine(G, C, byblack, DASHED_LINE, REGULAR_WIDTH);
byLineDefine(B, D, byblack, SOLID_LINE, REGULAR_WIDTH);
byLineDefine(D, N, byblue, SOLID_LINE, REGULAR_WIDTH);
byLineDefine(N, E, byred, SOLID_LINE, REGULAR_WIDTH);
byLineDefine(E, B, byyellow, SOLID_LINE, REGULAR_WIDTH);
draw byNamedLineSeq(1)(BD,DN,NE,EB);
draw byLabelsOnPolygon(D, H, A, C, B, F, E, N)(ALL_LABELS, 0);
draw byLabelsOnPolygon(H, G, C)(OMIT_FIRST_LABEL+OMIT_LAST_LABEL, 0);
}
\drawCurrentPictureInMargin
\problem{S}{e}{uma linha reta for dividida em duas partes quaisquer \drawProportionalLine{DN,NE}
$\drawProportionalLine{DN,NE}^2 + \drawProportionalLine{NE}^2 = 2\drawProportionalLine{DN,NE} \cdot \drawProportionalLine{NE} + \drawProportionalLine{DN}^2$
}

\begin{center}
Descreva \drawPolygon[bottom][squareABED]{DNGH,NEFG,HGCA,GFBC} \byref{prop:I.XLVI}, trace \drawProportionalLine{BD} \byref{post:I.I},\\
and $\left\{
	\begin{aligned}
		\drawProportionalLine{GN,GC} &\parallel \drawProportionalLine{EB} \\
		\drawProportionalLine{GH,GF} &\parallel \drawProportionalLine{DN,NE}\\
	\end{aligned}
\right\}$\\
$\drawPolygon[bottom]{HGCA} = \drawPolygon[bottom]{NEFG}$ \byref{prop:I.XLIII},\\
add $\drawPolygon[bottom]{GFBC} = \drawProportionalLine{NE}^2$ to both \byref{prop:II.IV}

$\drawPolygon[bottom]{HGCA,GFBC} = \drawPolygon[bottom]{NEFG,GFBC} = \drawProportionalLine{DN,NE} \cdot \drawProportionalLine{NE}$\\
$\drawPolygon[bottom]{DNGH} = \drawProportionalLine{DN}^2$ \byref{prop:II.IV}\\
$\drawPolygon[bottom]{HGCA,GFBC} + \drawPolygon[bottom]{NEFG,GFBC} + \drawPolygon[bottom]{DNGH} = 2\drawProportionalLine{DN,NE} \cdot \drawProportionalLine{NE} + \drawProportionalLine{DN}^2 = \squareABED + \drawPolygon[bottom]{GFBC}$;

$\drawProportionalLine{DN,NE}^2 + \drawProportionalLine{NE}^2 = 2\drawProportionalLine{DN,NE} \cdot \drawProportionalLine{NE} + \drawProportionalLine{DN}^2$.
\end{center}

\qed

\starttheorem{Prop. VIII. Teor.}\label{prop:II.VIII}
\defineNewPicture{
pair A, B, C, D, E, F, G, H, K, L, M, N, O, P, Q, R;
numeric w, d;
w := 7/2u;
d := u;
A := (0, w + d);
B := (w, w + d);
C := (w - d, w + d);
D := (w + d, w + d);
E := (0, 0);
F := (w + d, 0);
G := (w - d, w);
H := (w - d, 0);
K := (w, w);
L := (w, 0);
M := (0, w);
N := (w + d, w);
O := (0, w - d);
P := (w + d, w - d);
Q := (w - d, w - d);
R := (w, w - d);
draw byLine(C, Q, byblue, DASHED_LINE, REGULAR_WIDTH);
draw byLine(B, R, byblack, DASHED_LINE, REGULAR_WIDTH);
draw byLine(Q, H, byblue, SOLID_LINE, REGULAR_WIDTH);
draw byLine(R, L, byblue, SOLID_LINE, REGULAR_WIDTH);
draw byLine(M, G, byblack, DASHED_LINE, REGULAR_WIDTH);
draw byLine(O, Q, byred, DASHED_LINE, REGULAR_WIDTH);
draw byLine(G, N, byred, SOLID_LINE, REGULAR_WIDTH);
draw byLine(Q, P, byred, SOLID_LINE, REGULAR_WIDTH);
draw byLine(D, E, byblack, SOLID_LINE, REGULAR_WIDTH);
byLineDefine(E, H, byblue, SOLID_LINE, REGULAR_WIDTH);
byLineDefine(H, L, byred, SOLID_LINE, REGULAR_WIDTH);
byLineDefine(L, F, byyellow, SOLID_LINE, REGULAR_WIDTH);
byLineDefine(F, P, byblue, SOLID_LINE, REGULAR_WIDTH);
byLineDefine(P, N, byyellow, SOLID_LINE, REGULAR_WIDTH);
byLineDefine(N, D, byred, SOLID_LINE, REGULAR_WIDTH);
byLineDefine(A, D, byblack, SOLID_LINE, REGULAR_WIDTH);
byLineDefine(E, A, byblack, SOLID_LINE, REGULAR_WIDTH);
draw byNamedLineSeq(0)(EH,HL,LF,FP,PN,ND,AD,EA);
byLineDefine(D, F, byblack, SOLID_LINE, REGULAR_WIDTH);
byLineDefine(B, L, byblack, SOLID_LINE, THIN_WIDTH);
byLineDefine(C, H, byblack, SOLID_LINE, THIN_WIDTH);
byLineDefine(M, N, byblack, SOLID_LINE, THIN_WIDTH);
byLineDefine(O, P, byblack, SOLID_LINE, THIN_WIDTH);
draw byLabelsOnPolygon(A, C, B, D, N, P, F, L, H, E, O, M)(ALL_LABELS, 0);
}
\drawCurrentPictureInMargin
\problem{S}{e}{uma linha reta for dividida em duas partes quaisquer \drawProportionalLine{EH,HL}
$\drawProportionalLine{EH,HL,LF}^2 = 4 \cdot \drawProportionalLine{EH,HL} \cdot \drawProportionalLine{HL} + \drawProportionalLine{EH}^2$
}

\begin{center}
Produce \drawProportionalLine{EH,HL} and make $\drawProportionalLine{LF} = \drawProportionalLine{HL}$

Construct \drawFromCurrentPicture[bottom]{
draw byNamedLine(BL,CH,MN,OP);
draw byNamedLineSeq(0)(LF,HL,EH,EA,AD,DF);
} \byref{prop:I.XLVI};\\
draw \drawProportionalLine{DE},

$\left.
\begin{aligned}
\left.
	\begin{aligned}
		\vcenter{
		\nointerlineskip\hbox{\drawProportionalLine{CQ,QH}}
		\nointerlineskip\hbox{\drawProportionalLine{BR,RL}}
		}
	\end{aligned}\right\} & \parallel \drawProportionalLine{FP,PN,ND}\\
\left.
	\begin{aligned}
		\vcenter{
		\nointerlineskip\hbox{\drawProportionalLine{OQ,QP}}
		\nointerlineskip\hbox{\drawProportionalLine{MG,GN}}
		}
	\end{aligned}
\right\} & \parallel \drawProportionalLine{EH,HL,LF}\\
\end{aligned}
\right\}$ \byref{prop:I.XXXI}

$\drawProportionalLine{EH,HL,LF}^2 = \drawProportionalLine{LF}^2 + \drawProportionalLine{EH,HL}^2 + 2 \cdot \drawProportionalLine{EH,HL} \cdot \drawProportionalLine{LF}$ \byref{prop:II.IV}\\
but $\drawProportionalLine{HL}^2 + \drawProportionalLine{EH,HL}^2 = 2 \cdot \drawProportionalLine{EH,HL} \cdot \drawProportionalLine{HL} + \drawProportionalLine{EH}^2$ \byref{prop:II.VII}

$\therefore \drawProportionalLine{EH,HL,LF}^2 = 4 \cdot \drawProportionalLine{EH,HL} \cdot \drawProportionalLine{HL} + \drawProportionalLine{EH}^2$.
\end{center}

\qed

\starttheorem{Prop. IX. Teor.}\label{prop:II.IX}
\defineNewPicture{
pair A, B, C, D, E, F, G, H;
numeric w;
w := 5u;
A := (-1/2w, 0);
B := (1/2w, 0);
C := (0, 0);
D := (1/5w, 0);
E := (0, 1/2w);
F = whatever[E, B] = (xpart(D), whatever);
G = whatever[E, C] = (whatever, ypart(F));
H := 1/2[E, F];
byAngleDefine(E, A, B, byyellow, SOLID_SECTOR);
byAngleDefine(A, B, E, byblue, SOLID_SECTOR);
byAngleDefine(A, E, C, byyellow, SOLID_SECTOR);
byAngleDefine(C, E, B, byred, SOLID_SECTOR);
byAngleDefine(E, F, G, byred, SOLID_SECTOR);
byAngleDefine(D, F, B, byblack, SOLID_SECTOR);
draw byNamedAngleResized();
draw byLine(A, F, byblack, SOLID_LINE, REGULAR_WIDTH);
draw byLine(F, D, byred, DASHED_LINE, REGULAR_WIDTH);
draw byLine(G, F, byyellow, 0, -1);
byLineDefine(C, G, byblue, SOLID_LINE, REGULAR_WIDTH);
byLineDefine(G, E, byblue, DASHED_LINE, REGULAR_WIDTH);
draw byNamedLineSeq(0)(CG, GE);
byLineDefine(A, C, byblue, SOLID_LINE, REGULAR_WIDTH);
byLineDefine(C, D, byyellow, SOLID_LINE, REGULAR_WIDTH);
byLineDefine(D, B, byred, SOLID_LINE, REGULAR_WIDTH);
byLineDefine(F, B, byyellow, DASHED_LINE, REGULAR_WIDTH);
byLineDefine(H, F, byblack, SOLID_LINE, REGULAR_WIDTH);
byLineDefine(E, H, byblack, 2, REGULAR_WIDTH);
byLineDefine(A, E, byblack, DASHED_LINE, REGULAR_WIDTH);
draw byNamedLineSeq(0)(AC,CD,DB,FB,HF,EH,AE);
draw byLabelsOnPolygon(E, F, B, D, C, A)(ALL_LABELS, 0);
draw byLabelsOnPolygon(C, G, E)(OMIT_FIRST_LABEL+OMIT_LAST_LABEL, 0);
}
\drawCurrentPictureInMargin
\problem{S}{e}{uma linha reta for dividida em duas partes iguais \drawProportionalLine{AC} \drawProportionalLine{CD,DB} e também em duas partes desiguais \drawProportionalLine{AC,CD} \drawProportionalLine{DB}
$\drawProportionalLine{AC,CD}^2 + \drawProportionalLine{DB}^2 = 2 \cdot \drawProportionalLine{AC}^2 + 2 \cdot \drawProportionalLine{CD}^2$
}

\begin{center}
Make $\drawProportionalLine{CG,GE} \perp \mbox{ e } = \drawProportionalLine{AC} \mbox{ ou } \drawProportionalLine{CD,DB}$,\\
Trace \drawProportionalLine{AE} e \drawProportionalLine{EH,HF,FB},\\
$\drawProportionalLine{FD} \parallel \drawProportionalLine{CG,GE}$, $ \drawProportionalLine{GF} \parallel \drawProportionalLine{CD,DB}$ e traçar \drawProportionalLine{AF}.

$\drawAngle{A} = \drawAngle{AEC}$ \byref{prop:I.V} $=$ half a right angle \byref{prop:I.XXXII}\\
$\drawAngle{B} = \drawAngle{DFB}$ \byref{prop:I.V} $=$ half a right angle \byref{prop:I.XXXII}\\
$\therefore \drawAngle{AEC,CEB} =$ a right angle.

$\drawAngle{B} = \drawAngle{CEB} = \drawAngle{EFG} = \drawAngle{DFB}$ \byref{prop:I.V, prop:I.XXIX}.\\
hence $\drawProportionalLine{FD} = \drawProportionalLine{DB}$, $\drawProportionalLine{GE} = \drawProportionalLine{GF} = \drawProportionalLine{CD}$ \byref{prop:I.VI,prop:I.XXXIV}

$\drawProportionalLine{AF}^2 = \left\{
\begin{aligned}
&\drawProportionalLine{AC,CD}^2 + \drawProportionalLine{FD}^2 \mbox{, or } + \drawProportionalLine{DB}^2 \\
&= \left\{
	\begin{aligned}
		&= \drawProportionalLine{AE}^2 + \drawProportionalLine{EH,HF}^2 \\
		&= 2 \cdot \drawProportionalLine{AC}^2 + 2 \cdot \drawProportionalLine{CD}^2\\
	\end{aligned}
	\right. \mbox{ \byref{prop:I.XLVII}}
\end{aligned}\right.$

$\therefore \drawProportionalLine{AC,CD}^2 + \drawProportionalLine{DB}^2 = 2 \cdot \drawProportionalLine{AC}^2 + 2 \cdot \drawProportionalLine{CD}^2$.
\end{center}

\qed

\starttheorem{Prop. X. Teor.}\label{prop:II.X}
\defineNewPicture{
pair A, B, C, D, E, F, G;
numeric w;
w := 7/2u;
A := (0, 0);
B := (w, 0);
C := 1/2[A, B];
D := 4/3[A, B];
E := (w/2, w/2);
F := (xpart(D), ypart(E));
G = whatever[E, B] = whatever[F, D];
byAngleDefine(E, A, C, byblack, SOLID_SECTOR);
byAngleDefine(C, E, A, byyellow, SOLID_SECTOR);
byAngleDefine(B, E, C, byyellow, SOLID_SECTOR);
byAngleDefine(F, E, B, byblue, SOLID_SECTOR);
byAngleDefine(C, B, E, byred, SOLID_SECTOR);
byAngleDefine(D, B, G, byred, SOLID_SECTOR);
byAngleDefine(B, G, D, byblue, SOLID_SECTOR);
draw byNamedAngleResized();
draw byLine(C, E, byred, 0, -1);
draw byLine(A, C, byred, SOLID_LINE, REGULAR_WIDTH);
draw byLine(C, B, byyellow, SOLID_LINE, REGULAR_WIDTH);
draw byLine(B, D, byblue, SOLID_LINE, REGULAR_WIDTH);
byLineDefine(E, B, byblack, SOLID_LINE, REGULAR_WIDTH);
byLineDefine(B, G, byblack, DASHED_LINE, REGULAR_WIDTH);
draw byNamedLineSeq(0)(EB,BG);
byLineDefine(E, F, byyellow, DASHED_LINE, REGULAR_WIDTH);
byLineDefine(F, D, byred, SOLID_LINE, REGULAR_WIDTH);
byLineDefine(D, G, byred, DASHED_LINE, REGULAR_WIDTH);
byLineDefine(A, G, byblack, SOLID_LINE, REGULAR_WIDTH);
byLineDefine(A, E, byblue, DASHED_LINE, REGULAR_WIDTH);
draw byNamedLineSeq(0)(EF,FD,DG,AG,AE);
draw byLabelsOnPolygon(A, E, F, D, G)(ALL_LABELS, 0);
draw byLabelsOnPolygon(G, B, C, A)(OMIT_FIRST_LABEL+OMIT_LAST_LABEL, 0);
}
\drawCurrentPictureInMargin
\problem{S}{e}{uma linha reta \drawProportionalLine{AC,CB} for dividida ao meio e produzida em qualquer ponto \drawProportionalLine{AC,CB,BD}
$\drawProportionalLine{AC,CB,BD}^2 + \drawProportionalLine{BD}^2 = 2 \cdot \drawProportionalLine{CB}^2 + 2 \cdot \drawProportionalLine{CB,BD}^2$
}

\begin{center}
Faça $\drawProportionalLine{CE} \perp \mbox{ e } = \mbox{ to } \drawProportionalLine{AC} \mbox{ ou } \drawProportionalLine{CB}$,\\
traçar \drawProportionalLine{AE} e \drawProportionalLine{EB,BG},\\
and $\left\{
\begin{aligned}
\drawProportionalLine{FD,DG} & \parallel \drawProportionalLine{CE} \\
\drawProportionalLine{EF} & \parallel \drawProportionalLine{CB,BD}\\
\end{aligned}
\right\}$ \byref{prop:I.XXXI}\\
draw \drawProportionalLine{AG} also.

$\drawAngle{A} = \drawAngle{CEA}$ \byref{prop:I.V} $=$ half a right angle \byref{prop:I.XXXII}\\
$\drawAngle{CBE} = \drawAngle{BEC}$ \byref{prop:I.V} $=$ half a right angle \byref{prop:I.XXXII}\\
$\therefore \drawAngle{CEA,BEC} = $ a right angle.

$\drawAngle{DBG} = \drawAngle{CBE} = \drawAngle{BEC} = \drawAngle{FEB} = \drawAngle{G} =$\\
meio ângulo reto \byref{prop:I.V,prop:I.XXXII,prop:I.XXIX,prop:I.XXXIV},\\% melhoria: no DBG original é girado de forma semelhante a A; provavelmente só precisamos de prop:I.XXIX e prop:I.XV aqui
and $\drawProportionalLine{BD} = \drawProportionalLine{DG}$, $\drawProportionalLine{CB,BD} = \drawProportionalLine{EF} = \drawProportionalLine{FD,DG}$, \byref{prop:I.VI,prop:I.XXXIV}.

Hence by \byref{prop:I.XLVII}

$\drawProportionalLine{AG}^2 = \left\{
\begin{aligned}
& \drawProportionalLine{AC,CB,BD}^2 + \drawProportionalLine{DG}^2 \mbox{ ou } \drawProportionalLine{BD}^2 \\
& \left\{
	\begin{aligned}
		& + \drawProportionalLine{AE}^2 = 2 \cdot \drawProportionalLine{AC}^2 \\
		& + \drawProportionalLine{EB,BG}^2 = 2\cdot \drawProportionalLine{EF}^2 \\
	\end{aligned}
	\right. \\
\end{aligned}
\right.$\\
$\therefore \drawProportionalLine{AC,CB,BD}^2 + \drawProportionalLine{BD}^2 = 2 \cdot \drawProportionalLine{CB}^2 + 2 \cdot \drawProportionalLine{CB,BD}^2$.
\end{center}

\qed

\startproblem{Prop. XI. prob.}\label{prop:II.XI}
\defineNewPicture{
pair A, B, C, D, E, F, G, H, K;
numeric w;
w := 7/2u;
A := (0, 0);
B := (w, 0);
C := (0, w);
D := (w, w);
E := 1/2[A, C];
F := E shifted (0, -abs(E-B));
G := F shifted (abs(F-A), 0);
H = whatever[A, B] = (xpart(G), whatever);
K = whatever[G, H] = whatever[C, D];
draw byPolygon(A,H,K,C)(byyellow);
draw byPolygon(H,B,D,K)(byblue);
draw byPolygon(A,H,G,F)(byblue);
draw byLine(K, G, byblack, DASHED_LINE, REGULAR_WIDTH);
byLineDefine(A, H, byred, SOLID_LINE, REGULAR_WIDTH);
byLineDefine(H, B, byred, DASHED_LINE, REGULAR_WIDTH);
byLineDefine(B, E, byblack, SOLID_LINE, REGULAR_WIDTH);
byLineDefine(C, E, byblue, DASHED_LINE, REGULAR_WIDTH);
byLineDefine(E, A, byblue, SOLID_LINE, REGULAR_WIDTH);
byLineDefine(A, F, byyellow, SOLID_LINE, REGULAR_WIDTH);
draw byNamedLineSeq(-1)(BE,HB,AH);
draw byNamedLineSeq(0)(CE,EA,AF);
draw byLabelsOnPolygon(F, A, E, C, K, D, B, H, G)(ALL_LABELS, 0);
}
\drawCurrentPictureInMargin
\problem{P}{ara}{dividir uma determinada linha reta \drawProportionalLine{AH,HB}
$\drawProportionalLine{AH,HB} \cdot \drawProportionalLine{HB} = \drawProportionalLine{AH}^2$
}

\begin{center}
Describe \drawPolygon[bottom]{AHKC,HBDK} \byref{prop:I.XLVI},\\
make $\drawProportionalLine{EA} = \drawProportionalLine{CE}$ \byref{prop:I.X},\\
draw \drawProportionalLine{BE},\\
take $\drawProportionalLine{EA,AF} = \drawProportionalLine{BE}$ \byref{prop:I.III},\\
on \drawProportionalLine{AF} describe \drawPolygon[bottom]{AHGF} \byref{prop:I.XLVI}.

Produce \drawProportionalLine{KG} \byref{post:I.II}.

Then, \byref{prop:II.VI} $\drawProportionalLine{CE,EA,AF} \cdot \drawProportionalLine{AF} + \drawProportionalLine{EA}^2 = \drawProportionalLine{EA,AF}^2 = \drawProportionalLine{BE}^2 = \drawProportionalLine{AH,HB}^2 + \drawProportionalLine{EA}^2 \therefore \drawProportionalLine{CE,EA,AF} \cdot \drawProportionalLine{AF} = \drawProportionalLine{AH,HB}^2$, or,\\
$\drawPolygon[bottom]{AHKC,AHGF} = \drawPolygon[bottom]{AHKC,HBDK} \therefore \drawPolygon[bottom]{AHGF} = \drawPolygon[bottom]{HBDK}$

$\therefore \drawProportionalLine{AH,HB} \cdot \drawProportionalLine{HB} = \drawProportionalLine{AH}^2$.

\end{center}

\qed

\starttheorem{Prop. XII. Teor.}\label{prop:II.XII}
\defineNewPicture{
pair A, B, C, D;
numeric w, h;
w := 7/2u;
h := 3u;
A := (3/5w, 0);
B := (w, h);
C := (0, 0);
D := (w, 0);
draw byLine(B, A, byred, SOLID_LINE, REGULAR_WIDTH);
byLineDefine(C, A, byblack, SOLID_LINE, REGULAR_WIDTH);
byLineDefine(A, D, byblack, DASHED_LINE, REGULAR_WIDTH);
byLineDefine(D, B, byyellow, SOLID_LINE, REGULAR_WIDTH);
byLineDefine(B, C, byblue, SOLID_LINE, REGULAR_WIDTH);
draw byNamedLineSeq(0)(DB,AD,CA,BC);
draw byLabelsOnPolygon(C, B, D, A)(ALL_LABELS, 0);
}
\drawCurrentPictureInMargin
\problem{I}{n}{qualquer triângulo de ângulo obtuso, o quadrado do lado que subtende o ângulo obtuso excede a soma dos quadrados dos lados que contêm o ângulo obtuso, em duas vezes o retângulo compreendido por de qualquer um desses lados e as partes produzidas do mesmo do ângulo obtuso à perpendicular caem sobre ele a partir do ângulo agudo oposto.\\
$\drawProportionalLine{BC}^2 > \drawProportionalLine{CA}^2 + \drawProportionalLine{BA}^2$ by $2 \cdot \drawProportionalLine{CA} \cdot \drawProportionalLine{AD}$
}

\begin{center}
By \byref{prop:II.IV}\\
$\drawProportionalLine{CA,AD}^2 = \drawProportionalLine{CA}^2 + \drawProportionalLine{AD}^2 + 2 \cdot \drawProportionalLine{CA} \cdot \drawProportionalLine{AD}$:

add $\drawProportionalLine{DB}^2$ to both\\
$\drawProportionalLine{CA,AD}^2 + \drawProportionalLine{DB}^2 = \drawProportionalLine{BC}^2$ \byref{prop:I.XLVII}\\
$= 2 \cdot \drawProportionalLine{CA} \cdot \drawProportionalLine{AD} + \drawProportionalLine{CA}^2 + \left\{
	\begin{aligned}
		&\drawProportionalLine{AD}^2 \\
		&\drawProportionalLine{DB}^2\\
	\end{aligned}
\right\}$ or\\
$+ \drawProportionalLine{BA}^2$ \byref{prop:I.XLVII}. % a apresentação original não é muito clara. isso deve ser AD ^ 2 + DB ^ 2, que é, por I.XLVII = BA ^ 2

$\therefore \drawProportionalLine{BC}^2 = 2 \cdot \drawProportionalLine{CA} \cdot \drawProportionalLine{AD} + \drawProportionalLine{CA}^2 + \drawProportionalLine{BA}^2$: hence $ \drawProportionalLine{BC}^2 > \drawProportionalLine{CA}^2 + \drawProportionalLine{BA}^2$ by $2 \cdot \drawProportionalLine{CA} \cdot \drawProportionalLine{AD}$.
\end{center}

\qed

\starttheorem{Prop. XIII. Teor.}\label{prop:II.XIII}
\defineNewPicture[1/4]{
pair A,B,C,D,E,F,G,H, d;
numeric w, h;
w := 3u;
h := 3u;
A := (2/5w, h);
B:= (0, 0);
C := (w, 0);
D = whatever[B, C] = (xpart(A), whatever);
d := (0, -h -4/3u);
E := (w, h) shifted d;
F := (0, 0) shifted d;
G := (2/5w, 0) shifted d;
H = whatever[F, G] = (xpart(E), whatever);
draw byLine(A, D, byyellow, SOLID_LINE, REGULAR_WIDTH);
byLineDefine(A, B, byred, SOLID_LINE, REGULAR_WIDTH);
byLineDefine(B, D, byblack, SOLID_LINE, REGULAR_WIDTH);
byLineDefine(D, C, byblack, DASHED_LINE, REGULAR_WIDTH);
byLineDefine(C, A, byblue, SOLID_LINE, REGULAR_WIDTH);
draw byNamedLineSeq(0)(AB,BD,DC,CA);
draw byLine(E, G, byblue, SOLID_LINE, REGULAR_WIDTH);
byLineDefine(E, F, byred, SOLID_LINE, REGULAR_WIDTH);
byLineDefine(F, G, byblack, SOLID_LINE, REGULAR_WIDTH);
byLineDefine(G, H, byblack, DASHED_LINE, REGULAR_WIDTH);
byLineDefine(H, E, byyellow, SOLID_LINE, REGULAR_WIDTH);
draw byNamedLineSeq(0)(EF,FG,GH,HE);
label.top(btex First etex, (xpart(1/2[B, C]), ypart(A) + 1/4u));
label.top(btex Second etex, (xpart(1/2[F, H]), ypart(E)));
draw byLabelsOnPolygon(B, A, C, D)(ALL_LABELS, 0);
draw byLabelsOnPolygon(F, E, H, G)(ALL_LABELS, 0);
}
\drawCurrentPictureInMargin
\problem{I}{n}{qualquer triângulo, o quadrado do lado que subtende um ângulo agudo, é menor que a soma dos quadrados dos lados que contêm esse ângulo, por duas vezes o retângulo compreendido por qualquer um desses lados, e a parte dele interceptada entre o pé da perpendicular deixada cair sobre ele do ângulo oposto, e o ponto angular do ângulo agudo.\\
First.\\
$\drawProportionalLine{CA}^2 < \drawProportionalLine{BD,DC}^2 + \drawProportionalLine{AB}^2$ by $2 \cdot \drawProportionalLine{BD,DC} \cdot \drawProportionalLine{BD}$.\\
Second.\\
$\drawProportionalLine{EG}^2 < \drawProportionalLine{EF}^2 + \drawProportionalLine{FG}^2$ by $2 \cdot \drawProportionalLine{FG} \cdot \drawProportionalLine{FG,GH}$.
}

\begin{center}
Primeiro, suponha que a perpendicular caia dentro do triângulo, então \byref{prop:II.VII}\\
$\drawProportionalLine{BD,DC}^2 + \drawProportionalLine{BD}^2 = 2 \cdot \drawProportionalLine{BD,DC} \cdot \drawProportionalLine{BD} + \drawProportionalLine{DC}^2$,\\
add to each $\drawProportionalLine{AD}^2$ then,\\
$\drawProportionalLine{BD,DC}^2 + \drawProportionalLine{BD}^2 + \drawProportionalLine{AD}^2 = 2 \cdot \drawProportionalLine{BD,DC} \cdot \drawProportionalLine{BD} + \drawProportionalLine{DC}^2 + \drawProportionalLine{AD}^2$\\
$\therefore$ \byref{prop:I.XLVII}\\
$\drawProportionalLine{BD,DC}^2 + \drawProportionalLine{AB}^2 = 2 \cdot \drawProportionalLine{BD,DC} \cdot \drawProportionalLine{BD} + \drawProportionalLine{CA}^2$,\\
and $\therefore \drawProportionalLine{CA}^2 < \drawProportionalLine{BD,DC}^2 + \drawProportionalLine{AB}^2$ by $2 \cdot \drawProportionalLine{BD,DC} \cdot \drawProportionalLine{BD}$

Em seguida, suponha que a perpendicular caia sem o triângulo, então \byref{prop:II.VII}\\
$\drawProportionalLine{FG,GH}^2 + \drawProportionalLine{FG}^2 = 2 \cdot \drawProportionalLine{FG,GH} \cdot \drawProportionalLine{FG} + \drawProportionalLine{GH}^2$,\\
add to each $\drawProportionalLine{HE}^2$ then\\
$\drawProportionalLine{FG,GH}^2 + \drawProportionalLine{FG}^2 + \drawProportionalLine{HE}^2= 2 \cdot \drawProportionalLine{FG,GH} \cdot \drawProportionalLine{FG} + \drawProportionalLine{GH}^2 + \drawProportionalLine{HE}^2$\\
$\therefore$ \byref{prop:I.XLVII}\\
$\drawProportionalLine{EF} + \drawProportionalLine{FG}^2 = 2 \cdot \drawProportionalLine{FG,GH} \cdot \drawProportionalLine{FG}^2 + \drawProportionalLine{EG}^2$,\\
$\therefore \drawProportionalLine{EG}^2 < \drawProportionalLine{EF}^2 + \drawProportionalLine{FG}^2$ by $2 \cdot \drawProportionalLine{FG,GH} \cdot \drawProportionalLine{FG}$.
\end{center}

\qed

\startproblem{Prop. XIV. prob.}\label{prop:II.XIV}
\defineNewPicture[1/3]{
path A;
pair a, b, c, d, e, f;
pair B, C, D, E, F, G, H;
numeric w, h, s, r;
a := (0, 0);
b := (u, 1/4u);
c := (2u, 1/5u);
d := (11/5u, -u);
e := (7/8u, -2u);
f := (1/8u, -3/4u);
A := a--b--c--d--e--f--cycle;
s := 0;
for i := 1 step 1 until length(A) - 1:
	s := s + 1/2(
		abs((point 0 of A) - (point i of A))
		*distanceToLine((point i + 1 of A), (point 0 of A)--(point i of A))
		);
endfor;
w := 5/2u;
h := s/w;
B := (w, 0);
C := (w, -h);
D := (0, -h);
E := (0, 0);
F := (-h, 0);
G := 1/2[B, F];
r := abs(B - G);
H := (D--(E shifted ((E-D)*10))) intersectionpoint ((fullcircle scaled 2r) shifted G);
byPointLabelRemove(a,b,c,d,e,f);
forsuffixes i=a,b,c,d,e,f:
	i := i shifted (xpart(G)-xpart(1/2[urcorner(A),ulcorner(A)]), r - ypart(lrcorner(A)) + 1/2u);
endfor;
draw byPolygon(a,b,c,d,e,f)(byyellow);
draw byPolygon(B,C,D,E)(byred);
draw byLine(H, G, byred, SOLID_LINE, REGULAR_WIDTH);
draw byLine(H, E, byblue, SOLID_LINE, REGULAR_WIDTH);
draw byLine(E, D, byyellow, SOLID_LINE, REGULAR_WIDTH);
byLineDefine(H, F, byblack, SOLID_LINE, THIN_WIDTH);
byLineDefine(H, B, byblack, SOLID_LINE, THIN_WIDTH);
byLineDefine(F, E, byblack, DASHED_LINE, REGULAR_WIDTH);
byLineDefine(E, G, byblue, DASHED_LINE, REGULAR_WIDTH);
byLineDefine(G, B, byblack, SOLID_LINE, REGULAR_WIDTH);
draw byNamedLineSeq(1)(FE,EG,GB,HB,HF);
draw byArc.G(G, B, F, r, byred, 0, 0, 0, 0);
byLineDefine(B, F, byblack, SOLID_LINE, REGULAR_WIDTH);
draw byLabelsOnPolygon(B, C, D, E, F, H)(ALL_LABELS, 0);
draw byLabelsOnPolygon(H, G, B)(OMIT_FIRST_LABEL+OMIT_LAST_LABEL, 0);
}
\drawCurrentPictureInMargin
\problem[4]{T}{o}{desenhe uma linha reta cujo quadrado seja igual a uma determinada figura retilínea.\\
Para traçar \drawSizedLine{HE} tal que, $\drawSizedLine{HE}^2 = \drawPolygon{abcdef}$
}

\begin{center}
Make $\drawPolygon{BCDE} = \drawPolygon{abcdef}$ \byref{prop:I.XLV},
produce \drawSizedLine{EG,GB} until $\drawSizedLine{FE} = \drawSizedLine{ED}$;\\
take $\drawSizedLine{FE,EG} = \drawSizedLine{GB}$ \byref{prop:I.X}.

Describe
\drawFromCurrentPicture[bottom]{
startTempScale(1/3);
draw byNamedLineFull(B, F, 0, 0,  0, 1)(BF);
startAutoLabeling;
draw byNamedArc(G);
stopAutoLabeling;
stopTempScale;
}
\byref{post:I.III},\\
e produza \drawSizedLine{ED} para atendê-lo: trace \drawSizedLine{HG}.

$\drawSizedLine{BG}^2 \mbox{ ou } \drawSizedLine{HG}^2 = \drawSizedLine{FE} \cdot \drawSizedLine{EG,GB} + \drawSizedLine{EG}^2$ \byref{prop:II.V},\\
but $\drawSizedLine{HG}^2 = \drawSizedLine{HE} ^2 + \drawSizedLine{EG}^2$ \byref{prop:I.XLVII};

$\therefore \drawSizedLine{HE}^2 + \drawSizedLine{EG}^2 = \drawSizedLine{FE} \cdot \drawSizedLine{EG,GB} + \drawSizedLine{EG}^2$

$\therefore \drawSizedLine{HE}^2 = \drawSizedLine{FE} \cdot \drawSizedLine{EG,GB}$, and

$\therefore \drawSizedLine{HE}^2 = \drawPolygon{BCDE} = \drawPolygon{abcdef}$.
\end{center}

\qed

\part{Livro III}

\chapter*{Definições}

\startdefinition{}\label{def:III.I}
Círculos iguais são aqueles cujos diâmetros são iguais.

\startdefinition{}\label{def:III.II}
\defineNewPicture{
	pair O, A, B;
	numeric r;
	r := 3/4u;
	O := (0, 0);
	A := (-r, -r);
	B := (r, -r);
	draw byCircleR(O, r, byblack, 0, 0, -1);
	draw byLine(A, B, byblack, SOLID_LINE, REGULAR_WIDTH);
}\drawCurrentPictureInMargin
Diz-se que uma linha reta toca um círculo quando encontra o círculo e, ao ser produzida, não o corta.

\startdefinition{}\label{def:III.III}
\defineNewPicture{
	pair A, B, C;
	numeric r[];
	r1 := 2/3u;
	r2 := 1/2r1;
	r3 := 2/5r1;
	A := (0, 0);
	B := A shifted (dir(110) scaled (r1-r2));
	C := A shifted (dir(-130) scaled (r1+r3));
	fill (fullcircle scaled 2r1) shifted A withcolor byyellow;
	draw byCircleR(A, r1, byyellow, 0, 0, 0);
	fill (fullcircle scaled 2r2) shifted B withcolor byred;
	draw byCircleR(B, r2, byblack, 0, 0, -1/2);
	fill (fullcircle scaled 2r3) shifted C withcolor byblue;
	draw byCircleR(C, r3, byblue, 0, 0, -1/2);
}\drawCurrentPictureInMargin
Diz-se que os círculos se tocam, mas não se cortam.

\startdefinition{}\label{def:III.IV}
\defineNewPicture{
	pair O, A, B, C, D, E, F;
	numeric r;
	r := 3/4u;
	O := (0, 0);
	A := dir(170) scaled r;
	B := dir(-110) scaled r;
	C := A xscaled -1;
	D := B xscaled -1;
	E := 1/2[A, B];
	F := 1/2[C, D];
	draw byLine(A, B, byblack, SOLID_LINE, REGULAR_WIDTH);
	draw byLine(C, D, byblack, SOLID_LINE, REGULAR_WIDTH);
	byLineDefine(O, E, byred, SOLID_LINE, REGULAR_WIDTH);
	byLineDefine(O, F, byblue, SOLID_LINE, REGULAR_WIDTH);
	draw byNamedLineSeq(0)(OE,OF);
	draw byCircleR(O, r, byblack, 0, 0, 0);
}\drawCurrentPictureInMargin
Diz-se que as linhas retas estão igualmente distantes do centro de um círculo quando as perpendiculares traçadas a elas a partir do centro são iguais.

\startdefinition{}\label{def:III.V}
E diz-se que a linha reta sobre a qual cai a perpendicular maior está mais distante do centro.

\startdefinition{}\label{def:III.VI}
\defineNewPicture{
	pair A, B;
	numeric r;
	r := 3/4u;
	A := (0, 0);
	B := (0, -1/8u);
	draw byFilledCircleSegment(A, r, 1/2, 4 - 1/2, byred);
	draw byFilledCircleSegment(B, r, 4 - 1/2, 8 + 1/2, byblue);
}\drawCurrentPictureInMargin
Um segmento de círculo é a figura contida por uma reta e pela parte da circunferência que ela corta.

\startdefinition{}\label{def:III.VII}
Um ângulo de um segmento é compreendido por uma linha reta e uma circunferência de um círculo.

% improvement: originally, seventh definition was missing

\startdefinition{}\label{def:III.VIII}
\defineNewPicture{
	pair O, A, B, C, D;
	numeric r, b, e;
	r := u;
	b := -1;
	e := 5;
	O := (0, 0);
	A := dir(100) scaled r;
	B := dir(30) scaled r;
	C := point b of fullcircle scaled 2r;
	D := point e of fullcircle scaled 2r;
	byAngleDefine(C, A, D, byblue, SOLID_SECTOR);
	byAngleDefine(C, B, D, byyellow, SOLID_SECTOR);
	draw byNamedAngleResized();
	draw byArcBE(O, b, e, r, byblack, 0, 0, 0, 0);
	draw byLine(C, A, byblack, SOLID_LINE, REGULAR_WIDTH);
	draw byLine(C, B, byblack, SOLID_LINE, REGULAR_WIDTH);
	draw byLine(D, A, byblack, SOLID_LINE, REGULAR_WIDTH);
	draw byLine(D, B, byblack, SOLID_LINE, REGULAR_WIDTH);
	draw byLine(C, D, byblack, SOLID_LINE, REGULAR_WIDTH);
}\drawCurrentPictureInMargin
Um ângulo em um segmento é o ângulo compreendido por duas linhas retas traçadas de qualquer ponto da circunferência do segmento até as extremidades da linha reta que é a base do segmento.

\startdefinition{}\label{def:III.IX}
\defineNewPicture{
	pair O, A, B, C;
	numeric r, b, e;
	r := 3/4u;
	b := -3/2;
	e := 9/2;
	O := (0, 0);
	A := dir(80) scaled r;
	B := point b of fullcircle scaled 2r;
	C := point e of fullcircle scaled 2r;
	byAngleDefine(C, A, B, byblue, SOLID_SECTOR);
	draw byNamedAngleResized();
	draw byArcBE.Op(O, b, e, r, byblack, 1, 0, 0, 0);
	draw byArcBE.Om(O, e, b + 8, r, byblack, 0, 0, 0, 0);
	draw byLine(C, A, byblack, SOLID_LINE, REGULAR_WIDTH);
	draw byLine(B, A, byblack, SOLID_LINE, REGULAR_WIDTH);
}\drawCurrentPictureInMargin
Diz-se que um ângulo está na parte da circunferência, ou do arco, interceptada entre as linhas retas que contêm o ângulo.

\startdefinition{}\label{def:III.X}
\defineNewPicture{
	pair O;
	numeric r, b, e;
	r := 2/3u;
	b := 1;
	e := 3;
	O := (0, 0);
	draw byFilledCircleSector(O, r, b, e, byyellow);
	draw byArcBE(O, e, b + 8, r, byblack, 0, 0, -1, 0);
}\drawCurrentPictureInMargin
Um setor de círculo é a figura contida por dois raios e pelo arco entre eles.

\startdefinition{}\label{def:III.XI}
\defineNewPicture{
	pair M, N, A, B, C, D, E, F;
	numeric r[], b, e;
	r1 := 3/2u;
	r2 := u;
	b := 1;
	e := 3;
	M := (0, 0);
	N := (0, -1/3u);
	A := (point b of fullcircle scaled 2r1) shifted M;
	B := (point 1/3[b,e] of fullcircle scaled 2r1) shifted M;
	C := (point e of fullcircle scaled 2r1) shifted M;
	D := (point b of fullcircle scaled 2r2) shifted N;
	E := (point 1/3[b,e] of fullcircle scaled 2r2) shifted N;
	F := (point e of fullcircle scaled 2r2) shifted N;
	draw byPolygon(A,B,C)(byred);
	draw byPolygon(D,E,F)(byred);
	draw byArcBE(M, b, e, r1, byblack, 0, 0, 0, 1);
	draw byArcBE(N, b, e, r2, byblack, 0, 0, 0, 1);
}\drawCurrentPictureInMargin
Segmentos semelhantes de círculos são aqueles que contêm ângulos iguais.

\startdefinition{}\label{def:III.XII}
\defineNewPicture{
	pair O;
	numeric r[];
	r1 := 1/3u;
	r2 := 1/2u;
	r3 := 3/4u;
	O := (0, 0);
	draw byFilledCircleSegment(O, r3, 0, 8, byred);
	draw byFilledCircleSegment(O, r2, 0, 8, white);
	draw byCircleR.OII(O, r2, byblack, 0, 0, 0);
	draw byFilledCircleSegment(O, r1, 0, 8, byblue);
}\drawCurrentPictureInMargin
Os círculos que têm o mesmo centro são chamados de \emph{círculos concêntricos}.

% improvement: this last definition was not marked as definition in Byrne's book, and is not present in conventional editions of the Elements, but since it sort of is, let it be

\vfill\pagebreak

\startproblem{Prop. I. prob.}\label{prop:III.I}
\defineNewPicture{
pair A, B, C, D, E, F, G;
numeric r, a;
r := 9/4u;
F := (0, 0);
A := F shifted (dir(170)*r);
B := F shifted (dir(-95)*r);
D := 1/2[A, B];
C := F shifted (dir(angle(A-B) - 90)*r);
E := F shifted (dir(angle(A-B) +90)*r);
G := F shifted (dir(-45)*1/2r);
a := -angle(G-D);
forsuffixes i=A, B, D, C, E, F, G:
	i := i rotated a;
endfor;
byAngleDefine(A, D, F, byblue, SOLID_SECTOR);
byAngleDefine(F, D, G, byyellow, SOLID_SECTOR);
byAngleDefine(G, D, B, byblack, SOLID_SECTOR);
draw byNamedAngleResized();
draw byLine(D, G)(byblue, DASHED_LINE, REGULAR_WIDTH);
draw byLine(A, D)(byred, SOLID_LINE, REGULAR_WIDTH);
draw byLine(D, B)(byred, DASHED_LINE, REGULAR_WIDTH);
draw byLine(E, C)(byblack, SOLID_LINE, REGULAR_WIDTH);
draw byMarkLine(1/2, byblack)(EC);
byLineDefine(A, G, byblue, SOLID_LINE, REGULAR_WIDTH);
byLineDefine(B, G, byblack, DASHED_LINE, REGULAR_WIDTH);
draw byNamedLineSeq(0)(AG,BG);
draw byCircleR(F, r, byblue, 0, 0, 0);
draw byLabelsOnPolygon(A, G, B)(ALL_LABELS, 0);
draw byLabelsOnPolygon(E, D, A)(OMIT_FIRST_LABEL+OMIT_LAST_LABEL, 0);
draw byLabelsOnPolygon(E, F, C)(OMIT_FIRST_LABEL+OMIT_LAST_LABEL, 4);
draw byLabelLineEnd(E, C, 0);
draw byLabelLineEnd(C, E, 0);
}
\drawCurrentPictureInMargin
\problem{P}{ara}{encontrar o centro de um determinado círculo \drawCircle[middle][1/4]{F}.}

Trace dentro do círculo qualquer linha reta \drawUnitLine{AD,DB}, faça $\drawUnitLine{AD} = \drawUnitLine{DB}$, trace $\drawUnitLine{EC} \perp \drawUnitLine{AD,DB}$; bissecta \drawUnitLine{EC}, e o ponto de bissecção é o centro.

Pois, se for possível, deixe qualquer outro ponto como ponto de encontro de \drawUnitLine{AG}, \drawUnitLine{DG} e \drawUnitLine{BG} ser o centro.

Porque em \drawLine[bottom][triangleADG]{AG,DG,AD} e \drawLine[middle][triangleDGB]{DB,DG,BG} $\drawUnitLine{AG} = \drawUnitLine{BG}$ \byref{\hypref,def:I.XV}, $\drawUnitLine{AD} = \drawUnitLine{DB}$ \byref{\constref} e \drawUnitLine{DG} comuns, $\drawAngle{ADF,FDG} = \drawAngle{GDB}$ \byref{prop:I.VIII} e, portanto, são ângulos retos; mas $\drawAngle{FDG,GDB} = \drawRightAngle$ \byref{\constref} $\drawAngle{GDB} = \drawAngle{FDG,GDB}$ \byref{ax:I.XI} o que é um absurdo; portanto o ponto assumido não é o centro do círculo; e da mesma maneira pode ser provado que nenhum outro ponto que não esteja em \drawUnitLine{EC} é o centro, portanto o centro está em \drawUnitLine{EC} e, portanto, o ponto onde \drawUnitLine{EC} é dividido ao meio é o centro.

\qed

\starttheorem{Prop. II. Teor.}\label{prop:III.II}
\defineNewPicture{
pair A, B, D, E, F;
numeric r;
r := 9/4u;
D := (0, 0);
A := (dir(185) scaled r) shifted D;
B := (dir(-70) scaled r) shifted D;
E := 3/5[A, B];
F := (dir(angle(E-D)) scaled r) shifted D;
byAngleDefine(D, A, B, byblue, SOLID_SECTOR);
byAngleDefine(A, E, D, byyellow, SOLID_SECTOR);
byAngleDefine(A, B, D, byblack, SOLID_SECTOR);
draw byNamedAngleResized();
draw byLine(D, E, byblack, SOLID_LINE, REGULAR_WIDTH);
draw byLine(E, F, byblack, DASHED_LINE, REGULAR_WIDTH);
draw byLine(A, B, byred, SOLID_LINE, REGULAR_WIDTH);
byLineDefine(A, D, byyellow, SOLID_LINE, REGULAR_WIDTH);
byLineDefine(B, D, byblue, SOLID_LINE, REGULAR_WIDTH);
draw byNamedLineSeq(0)(AD,BD);
draw byCircleR(D, r, byred, 0, 0, 0);
draw byLabelsOnPolygon(B, F, A, D)(ALL_LABELS, 0);
draw byLabelsOnPolygon(A, E, F)(OMIT_FIRST_LABEL+OMIT_LAST_LABEL, 0);
}
\drawCurrentPictureInMargin
\problem[2]{U}{ma}{linha reta (\drawSizedLine{AB}) que une dois pontos na circunferência de um círculo \drawCircle[middle][1/4]{D} está totalmente dentro do círculo.}

\begin{center}
Encontre o centro de \circleD\ \byref{prop:III.I};\\
do centro trace \drawSizedLine{DE} para qualquer ponto em \drawSizedLine{AB}, encontrando a circunferência a partir do centro;\\
trace \drawSizedLine{AD} e \drawSizedLine{BD}.

Then $\drawAngle{A} = \drawAngle{B}$ \byref{prop:I.V}\\
but $\drawAngle{E} > \drawAngle{A} \mbox{ ou } > \drawAngle{B}$ \byref{prop:I.XVI}

$\therefore \drawSizedLine{AD} > \drawSizedLine{DE}$ \byref{prop:I.XIX}\\
but $\drawSizedLine{AD} = \drawSizedLine{DE,EF}$,

$\therefore \drawSizedLine{DE,EF} > \drawSizedLine{DE}$;

$\therefore \drawSizedLine{DE} < \drawSizedLine{DE,EF}$;

$\therefore$ every point in \drawSizedLine{AB} lies within the circle.
\end{center}

\qed

\starttheorem{Prop. III. Teor.}\label{prop:III.III}
\defineNewPicture[1/2]{
pair A, B, C, D, E, F;
numeric r;
r := 9/4u;
E := (0, 0);
A := (dir(-90 - 60) scaled r) shifted E;
B := (dir(-90 + 60) scaled r) shifted E;
C := (dir(90) scaled r) shifted E;
D := (dir(-90) scaled r) shifted E;
F = whatever[A, B] = whatever[C, D];
byAngleDefine(E, A, F, byblue, SOLID_SECTOR);
byAngleDefine(A, F, E, byblack, SOLID_SECTOR);
byAngleDefine(E, F, B, byyellow, SOLID_SECTOR);
byAngleDefine(F, B, E, byred, SOLID_SECTOR);
draw byNamedAngleResized();
draw byLine(C, E, byblack, DASHED_LINE, REGULAR_WIDTH);
draw byLine(E, F, byblack, SOLID_LINE, REGULAR_WIDTH);
draw byLine(F, D, byblack, DASHED_LINE, REGULAR_WIDTH);
byLineDefine(A, F, byred, SOLID_LINE, REGULAR_WIDTH);
byLineDefine(F, B, byred, DASHED_LINE, REGULAR_WIDTH);
draw byNamedLineSeq(0)(AF,FB);
byLineDefine(A, E, byyellow, SOLID_LINE, REGULAR_WIDTH);
byLineDefine(E, B, byblue, SOLID_LINE, REGULAR_WIDTH);
draw byNamedLineSeq(0)(AE,EB);
draw byCircleR(E, r, byblue, 0, 0, 0);
draw byLabelsOnCircle(A, B)(E);
draw byLabelsOnPolygon(D, F, A)(OMIT_FIRST_LABEL+OMIT_LAST_LABEL, 2);
draw byLabelsOnPolygon(A, E, C)(OMIT_FIRST_LABEL+OMIT_LAST_LABEL, 0);
}
\drawCurrentPictureInMargin
\problem[3]{S}{e}{uma linha reta (\drawUnitLine{EF}) traçada através do centro de um círculo \drawCircle[middle][1/6]{E} corta ao meio uma corda (\drawUnitLine{AF,FB}) que não passa pelo centro, ela é perpendicular a ela; ou, se for perpendicular a ele, ele o divide ao meio.}

\begin{center}
Trace \drawUnitLine{AE} e \drawUnitLine{EB} no centro do círculo.

In \drawLine[bottom][triangleAEF]{AE,EF,AF} and
\drawLine[bottom][triangleFEB]{EF,EB,FB}\\
$\drawUnitLine{AE} = \drawUnitLine{EB}$, \drawUnitLine{EF} common,\\
and $\drawUnitLine{AF} = \drawUnitLine{FB} \therefore \drawAngle{AFE} = \drawAngle{EFB}$ \byref{prop:I.VIII}\\
and $\therefore \drawUnitLine{EF} \perp \drawUnitLine{AF,FB}$ \byref{def:I.X} % improvement: definition 7 in the original is irrelevant

Again let $\drawUnitLine{EF} \perp \drawUnitLine{AF,FB}$\\
Then in \triangleAEF\ and \triangleFEB\\
$\drawAngle{A} = \drawAngle{B}$ \byref{prop:I.V}\\
$\drawAngle{AFE} = \drawAngle{EFB}$ \byref{\hypref}\\
and $\drawUnitLine{AE} = \drawUnitLine{EB}$

$\therefore \drawUnitLine{AF} = \drawUnitLine{FB}$ \byref{prop:I.XXVI}

and $\therefore \drawUnitLine{EF}$ bisects \drawUnitLine{AF,FB}.
\end{center}

\qed

\starttheorem{Prop. IV. Teor.}\label{prop:III.IV}
\defineNewPicture{
pair A, B, C, D, E, F;
numeric r;
r := 9/4u;
F := (0, 0);
A := (dir(-175)*r) shifted F;
B := (dir(-140)*r) shifted F;
C := (dir(-50)*r) shifted F;
D := (dir(-10)*r) shifted F;
E = whatever[A, C] = whatever[B, D];
byAngleDefine(F, E, D, byblue, SOLID_SECTOR);
byAngleDefine(D, E, C, byyellow, SOLID_SECTOR);
draw byNamedAngleResized();
draw byLine(E, F, byblack, DASHED_LINE, REGULAR_WIDTH);
draw byLine(B, D, byred, SOLID_LINE, REGULAR_WIDTH);
draw byLine(A, C, byblack, SOLID_LINE, REGULAR_WIDTH);
draw byCircleR(F, r, byblue, 0, 0, 0);
draw byLabelsOnCircle(A, B, C, D)(F);
draw byLabelLineEnd(F, E, 0);
draw byLabelsOnPolygon(C, E, B)(OMIT_FIRST_LABEL+OMIT_LAST_LABEL, 0);
}
\drawCurrentPictureInMargin
\problem{S}{e}{em um círculo duas linhas retas se cortam, mas não passam ambas pelo centro, elas não se cortam ao meio.}

Se uma das linhas passa pelo centro, é evidente que não pode ser cortada ao meio pela outra, que não passa pelo centro.

But if neither of the lines \drawUnitLine{AC} or \drawUnitLine{BD} pass through the centre, draw \drawUnitLine{EF} from the centre to their intersection.

\begin{center}
Se \drawUnitLine{AC} for dividido ao meio, \drawUnitLine{EF} $\perp$ para ele \byref{prop:III.III}\\
$\therefore \drawAngle{FED,DEC} = \drawRightAngle$\\
and if \drawUnitLine{BD} be bisected, $\drawUnitLine{EF} \perp \drawUnitLine{BD}$ \byref{prop:III.III}\\
$\therefore \drawAngle{FED} = \drawRightAngle$;

and $\therefore \drawAngle{FED} = \drawAngle{FED,DEC}$;\\
uma parte igual ao todo, o que é um absurdo:

$\therefore$ \drawUnitLine{AC} and \drawUnitLine{BD} do not bisect one another.
\end{center}

\qed

\starttheorem{Prop. V. Teor.}\label{prop:III.V}
\defineNewPicture{
pair M, N, E, F, G, C;
numeric r[], s;
path c[];
r1 := 2u;
r2 := 2u;
s := u;
M := (1/2s, 0);
N := (-1/2s, 0);
c1 := (fullcircle scaled 2r1) shifted M;
c2 := (fullcircle scaled 2r2) shifted N;
E := 1/2[M, N];
C := (subpath(0, 4) of c1) intersectionpoint (subpath(0, 4) of c2);
G := point 7/2 of c2;
F := c1 intersectionpoint (E--G);
byLineDefine(C, E, byyellow, SOLID_LINE, REGULAR_WIDTH);
byLineDefine(E, F, byblack, SOLID_LINE, REGULAR_WIDTH);
byLineDefine(F, G, byblack, DASHED_LINE, REGULAR_WIDTH);
draw byNamedLineSeq(0)(CE,EF,FG);
draw byArcBE.Ma(M, 4, 0, r1, byred, 0, 0, 0, 0);
draw byArcBE.Na(N, 4, 0, r2, byblue, 0, 0, 0, 0);
draw byArcBE.Nb(N, 4, 8, r2, byblue, 0, 0, 0, 0);
draw byArcBE.Mb(M, 4, 8, r1, byred, 0, 0, 0, 0);
draw byLabelLineEnd(C, E, 0);
draw byLabelLineEnd(G, E, 0);
draw byLabelsOnPolygon(C, E, F)(OMIT_FIRST_LABEL+OMIT_LAST_LABEL, 0);
draw byLabelsOnPolygon(C, F, E)(OMIT_FIRST_LABEL+OMIT_LAST_LABEL, 1);
}
\drawCurrentPictureInMargin
\problem{S}{e}{dois círculos interceptam \drawArc{Ma,Na,Nb,Mb}, eles não têm o mesmo centro.}

Suponhamos que seja possível que dois círculos que se cruzam tenham um centro comum; a partir desse suposto centro trace \drawUnitLine{CE} até o ponto de intersecção, e \drawUnitLine{EF,FG} encontrando as circunferências dos círculos.

\begin{center}
Then $\drawUnitLine{CE} = \drawUnitLine{EF}$ \byref{def:I.XV}\\
and $\drawUnitLine{CE} = \drawUnitLine{EF,FG}$ \byref{def:I.XV}

$\therefore \drawUnitLine{EF} = \drawUnitLine{EF,FG}$\\
uma parte igual ao todo, o que é um absurdo:

Os círculos $\therefore$ que deveriam se cruzar em qualquer ponto não podem ter o mesmo centro.
\end{center}

\qed

\starttheorem{Prop. VI. Teor.}\label{prop:III.VI}
\defineNewPicture[1/4]{
pair M, N, B, C, E, F;
numeric r[], a;
path c[];
a := 80;
r1 := 9/4u;
r2 := 7/4u;
M := (0, 0);
N := M shifted (dir(a)*(r1-r2));
c1 := (fullcircle scaled 2r1) shifted M;
c2 := (fullcircle scaled 2r2) shifted N;
F := 1/2[M, N];
C :=c1 intersectionpoint (M--(M shifted (dir(a)*2r1)));
B := point -3/2 of c1;
E := c2 intersectionpoint (F--B);
byLineDefine(C, F, byyellow, SOLID_LINE, REGULAR_WIDTH);
byLineDefine(F, E, byblue, DASHED_LINE, REGULAR_WIDTH);
byLineDefine(E, B, byblue, SOLID_LINE, REGULAR_WIDTH);
draw byNamedLineSeq(0)(CF,FE,EB);
draw byCircle.M(M, C, byred, 0, 0, 0);
draw byCircle.N(N, C, byblack, 0, 0, -1);
draw byLabelsOnCircle(B, C)(M);
draw byLabelsOnPolygon(E, F, C)(OMIT_FIRST_LABEL+OMIT_LAST_LABEL, 0);
draw byLabelPoint(E, angle(B-F)+45, 2);
byPointLabelRemove(M, N);
}
\drawCurrentPictureInMargin
\problem{S}{e}{dois círculos \drawCircle{N,M} se tocam internamente, eles não têm o mesmo centro.}

Pois, se for possível, que ambos os círculos tenham o mesmo centro; de tal suposto centro, desenhe \drawUnitLine{FE,EB} e \drawUnitLine{CF} até o ponto de contato.

\begin{center}
Then $\drawUnitLine{CF} = \drawUnitLine{FE}$ \byref{def:I.XV}\\
and $\drawUnitLine{CF} = \drawUnitLine{FE,EB}$ \byref{def:I.XV}

$\therefore \drawUnitLine{FE} = \drawUnitLine{FE,EB}$;
\end{center}
\noindent uma parte igual ao todo, o que é um absurdo; portanto, o ponto assumido não é o centro de ambos os círculos, e da mesma maneira pode ser demonstrado que nenhum outro ponto o é.

\qed

\starttheorem{Prop. VII. Teor.}\label{prop:III.VII}
\defineNewPicture[1/2]{
pair A, B, C, D, E, F, G, H, K, M, N, d;
numeric r;
r := 2u;
E := (0, 0);
A := E shifted (dir(90)*r);
D := E shifted (dir(-90)*r);
F := 2/3[E, D];
B := E shifted (dir(20)*r);
C := E shifted (dir(-5)*r);
G := E shifted (dir(15)*r);
H := E shifted (dir(-170)*r);
K := E shifted (dir(175)*r);
M = whatever[E, H] = whatever[F, K];
N := D;
byAngleDefine(B, E, C, byblack, SOLID_SECTOR);
byAngleDefine(C, E, F, byyellow, SOLID_SECTOR);
draw byNamedAngleResized(BEC, CEF);
draw byLine(F, B, byred, SOLID_LINE, REGULAR_WIDTH);
draw byLine(E, C, byblue, DASHED_LINE, REGULAR_WIDTH);
draw byLine(F, C, byblue, SOLID_LINE, REGULAR_WIDTH);
draw byLine(E, B, byred, DASHED_LINE, REGULAR_WIDTH);
byLineDefine(A, E, byblack, DASHED_LINE, REGULAR_WIDTH);
byLineDefine(E, F, byblack, SOLID_LINE, REGULAR_WIDTH);
byLineDefine(F, D, byyellow, SOLID_LINE, REGULAR_WIDTH);
draw byNamedLineSeq(0)(AE,EF,FD);
draw byCircleR(E, r, byblue, 0, 0, 0);
draw byLabelsOnCircle(A, B, C, D)(E);
draw byLabelsOnPolygon(D, F, E, A)(OMIT_FIRST_LABEL+OMIT_LAST_LABEL, 0);
label.top(btex Fig. 1 etex, (0, r + 1/4cm));
d := (0, -2r -2u);
draw image(
byAngleDefine(G, F, E, byyellow, SOLID_SECTOR);
byAngleDefine(K, F, E, byred, SOLID_SECTOR);
draw byNamedAngleResized(GFE,KFE);
draw byLine(F, G, byred, SOLID_LINE, REGULAR_WIDTH);
draw byLine(E, G, byred, DASHED_LINE, REGULAR_WIDTH);
draw byLine(E, M, byyellow, SOLID_LINE, REGULAR_WIDTH);
draw byLine(M, H, byyellow, DASHED_LINE, REGULAR_WIDTH);
byLineDefine(F, M, byblue, SOLID_LINE, REGULAR_WIDTH);
byLineDefine(M, K, byblue, DASHED_LINE, REGULAR_WIDTH);
draw byNamedLineSeq(0)(FM,MK);
byLineDefine(A, E, byblack, DASHED_LINE, REGULAR_WIDTH);
byLineDefine(E, F, byblack, SOLID_LINE, REGULAR_WIDTH);
byLineDefine(F, N, byblack, SOLID_LINE, REGULAR_WIDTH);
draw byNamedLineSeq(0)(AE,EF,FN);
draw byCircleR(E, r, byblue, 0, 0, 0);
draw byLabelsOnCircle(G, K, H)(E);
draw byLabelsOnPolygon(D, F, M, H)(OMIT_FIRST_LABEL+OMIT_LAST_LABEL, 0);
draw byLabelsOnPolygon(M, E, A)(OMIT_FIRST_LABEL+OMIT_LAST_LABEL, 0);
label.top(btex Fig. 2 etex, (0, r+1/4cm));
) shifted d;
}
\drawCurrentPictureInMargin
\problem[3]{I}{f}{from any point within a circle \drawFromCurrentPicture{
startTempScale(1/3);
draw byNamedCircle(E);
draw byLabelPoint(F, 0, 0);
stopTempScale;
} que não é o centro, linhas
$\left\{\vcenter{
\nointerlineskip\hbox{\drawUnitLine{EF,AE}}
\nointerlineskip\hbox{\drawUnitLine{FB}}
\nointerlineskip\hbox{\drawUnitLine{FC}}}\right.$
são atraídos pela circunferência; a maior dessas linhas é aquela (\drawUnitLine{EF,AE}) que passa pelo centro, e a menor é a parte restante (\drawUnitLine{FD}) do diâmetro.\\
Dos outros, aquele (\drawUnitLine{FB}) que está mais próximo da linha que passa pelo centro, é maior que aquele (\drawUnitLine{FC}) que está mais remoto.\\
Figura 2. As duas retas (\drawUnitLine{FM,MK} e \drawUnitLine{FG}) que fazem ângulos iguais com aquela que passa pelo centro, em lados opostos deste, são iguais entre si; e não pode ser traçada uma terceira linha igual a eles, do mesmo ponto à circunferência.}

\startsubproposition{Figure I.}
No centro do círculo trace \drawUnitLine{EB} e \drawUnitLine{EC}; então $\drawUnitLine{AE} = \drawUnitLine{EB}$ \byref{def:I.XV} $\drawUnitLine{EF,AE} = \drawUnitLine{EF} + \drawUnitLine{EB} > \drawUnitLine{FB}$ \byref{prop:I.XX} da mesma maneira \drawUnitLine{EF,AE} pode ser mostrado como sendo maior que \drawUnitLine{FC}, ou qualquer outra linha traçada do mesmo ponto até a circunferência. Novamente, por \byref{prop:I.XX} $\drawUnitLine{EF} + \drawUnitLine{FC} > \drawUnitLine{EC} = \drawUnitLine{FD} + \drawUnitLine{EF}$, pegue \drawUnitLine{EF} de ambos; $\therefore \drawUnitLine{FC} > \drawUnitLine{FD}$ \byref{ax:I.III}, e da mesma maneira pode ser mostrado que \drawUnitLine{FD} é menor que qualquer outra linha traçada do mesmo ponto até a circunferência. Novamente, em \drawLine[middle][triangleEFB]{FB,EF,EB} e \drawLine[middle][triangleEFC]{FC,EF,EC}, \drawUnitLine{EF} comum, $\drawAngle{BEC,CEF} > \drawAngle{CEF}$ e $\drawUnitLine{EB} = \drawUnitLine{EC} \therefore \drawUnitLine{FB} > \drawUnitLine{FC}$ \byref{prop:I.XXIV} e \drawUnitLine{FB} podem, da mesma maneira, ser provados maiores do que qualquer outra linha traçada do mesmo ponto até a circunferência mais distante de \drawUnitLine{EF,AE}.

\startsubproposition{Figure II.}
\begin{center}
If $\drawAngle{KFE} = \drawAngle{GFE}$ then $\drawUnitLine{FM,MK} = \drawUnitLine{FG}$,\\
se não, pegue $\drawUnitLine{FM} = \drawUnitLine{FG}$, desenhe \drawUnitLine{EM,MH}\\
then in \drawLine[middle][triangleEFM]{EM,EF,FM} and \drawLine[middle][triangleEFG]{FG,EF,EG}, \drawUnitLine{EF} common,\\
$\drawAngle{KFE} = \drawAngle{GFE}$ and $\drawUnitLine{FG} = \drawUnitLine{FM}$

$\therefore \drawUnitLine{EG} = \drawUnitLine{EM}$ \byref{prop:I.IV}

$\therefore \drawUnitLine{EG} = \drawUnitLine{EM,MH} = \drawUnitLine{EM}$\\
uma parte igual ao todo, o que é um absurdo:
\end{center}

\noindent $\therefore \drawUnitLine{FG} = \drawUnitLine{FM,MK}$; e nenhuma outra linha é igual a \drawUnitLine{FG} traçada do mesmo ponto até a circunferência; pois se estivesse mais próximo daquele que passa pelo centro seria maior, e se estivesse mais distante seria menor.

\qed

\starttheorem{Prop. VIII. Teor.}\label{prop:III.VIII}
\problem{O}{}{texto original desta proposição é aqui dividido em três partes.}

\startsubproposition{I.}
\defineNewPicture[1/4]{
pair M, D, A, E, F;
numeric r;
r := 7/4u;
M := (0, 0);
D := M shifted (dir(90)*3/2r);
A := (dir(-90)*r) shifted M;
E := (dir(-140)*r) shifted M;
F := (dir(-170)*r) shifted M;
byAngleDefine(D, M, F, byyellow, SOLID_SECTOR);
byAngleDefine(F, M, E, byblack, SOLID_SECTOR);
draw byNamedAngleResized();
draw byLine(D, E, byred, SOLID_LINE, REGULAR_WIDTH);
draw byLine(M, E, byred, DASHED_LINE, REGULAR_WIDTH);
draw byLine(M, F, byblue, DASHED_LINE, REGULAR_WIDTH);
byLineDefine(M, A, byblack, DASHED_LINE, REGULAR_WIDTH);
byLineDefine(D, M, byblack, SOLID_LINE, REGULAR_WIDTH);
byLineDefine(D, F, byblue, SOLID_LINE, REGULAR_WIDTH);
draw byNamedLineSeq(0)(MA, DM,DF);
draw byCircle.M(M, E, byblack, 0, 0, 0);
draw byLabelsOnCircle(F, E, A)(M);
draw byLabelsOnPolygon(F, D, M, A)(OMIT_FIRST_LABEL+OMIT_LAST_LABEL, 0);
}
\drawCurrentPictureInMargin
Se de um ponto sem círculo, linhas retas
$\left\{\vcenter{
\nointerlineskip\hbox{\drawProportionalLine{DM,MA}}
\nointerlineskip\hbox{\drawProportionalLine{DE}}
\nointerlineskip\hbox{\drawProportionalLine{DF}}}\right\}$
são atraídos pela circunferência; daquelas que caem sobre a circunferência côncava, a maior é aquela (\drawUnitLine{DM,MA}) que passa pelo centro, e a linha (\drawUnitLine{DE}) que está mais próxima da maior é maior que aquela (\drawUnitLine{DF}) que está mais remota.

\begin{center}
Trace \drawUnitLine{MF} e \drawUnitLine{ME} no centro.\\
Então, \drawUnitLine[0.75cm]{DM,MA} que passa pelo centro é maior;\\
for since $\drawUnitLine[0.75cm]{MA} = \drawUnitLine[0.75cm]{ME}$, if \drawUnitLine[0.75cm]{DM} be added to both $\drawUnitLine{DM,MA} = \drawUnitLine[0.75cm]{DM} + \drawUnitLine[0.75cm]{ME}$; but $> \drawUnitLine[0.75cm]{DE}$ \byref{prop:I.XX}\\
$\therefore$ \drawUnitLine{DM,MA} é maior que qualquer outra linha traçada do mesmo ponto até a circunferência côncava.\\
Again in \drawLine[middle][triangleDFM]{DM,MF,DF} and \drawLine[middle][triangleDEM]{DM,ME,DE}, $\drawUnitLine[0.5cm]{MF} = \drawUnitLine[0.5cm]{ME}$, and \drawUnitLine[0.5cm]{DM} common,\\
but $\drawAngle{DMF,FME} > \drawAngle{DMF}$, $\therefore \drawUnitLine{DE} > \drawUnitLine{DF}$ \byref{prop:I.XXIV};\\
e da mesma maneira \drawUnitLine{DE} pode ser mostrado como $>$ do que qualquer outra linha mais remota de \drawUnitLine{DM,MA}.
\end{center}

\defineNewPicture{
pair M, D, G, H, K;
numeric r;
r := 7/4u;
M := (0, 0);
D := M shifted (dir(90)*2r);
G := (dir(90)*r) shifted M;
H := (dir(130)*r) shifted M;
K := (dir(110)*r) shifted M;
draw byLine(G, M, byblack, SOLID_LINE, REGULAR_WIDTH);
draw byLine(H, M, byblue, SOLID_LINE, REGULAR_WIDTH);
draw byLine(D, K, byred, DASHED_LINE, REGULAR_WIDTH);
draw byLine(K, M, byred, SOLID_LINE, REGULAR_WIDTH);
byLineDefine(D, G, byblack, DASHED_LINE, REGULAR_WIDTH);
byLineDefine(D, H, byblue, DASHED_LINE, REGULAR_WIDTH);
draw byNamedLineSeq(0)(DG,DH);
draw byCircle.M(M, G, byblack, 0, 0, 0);
draw byLabelsOnPolygon(K, H, M)(OMIT_FIRST_LABEL+OMIT_LAST_LABEL, 2);
draw byLabelsOnPolygon(G, K, M)(OMIT_FIRST_LABEL+OMIT_LAST_LABEL, 2);
draw byLabelsOnPolygon(D, G, K)(OMIT_FIRST_LABEL+OMIT_LAST_LABEL, 2);
draw byLabelsOnPolygon(H, D, G)(OMIT_FIRST_LABEL+OMIT_LAST_LABEL, 0);
draw byLabelsOnPolygon(G, M, H)(OMIT_FIRST_LABEL+OMIT_LAST_LABEL, 0);
}
\drawCurrentPictureInMargin
\startsubproposition{II.}
Das linhas que caem na circunferência convexa, a menor é aquela (\drawUnitLine{DG}) que está sendo produzida passaria pelo centro, e a linha que está mais próxima do menor é menor que aquela que está mais remota.

\begin{center}
For, since $\drawUnitLine{KM} + \drawUnitLine{DK} > \drawUnitLine{DG,GM}$ \byref{prop:I.XX}\\
and $\drawUnitLine{KM} = \drawUnitLine{GM}$, $\therefore \drawUnitLine{DK} > \drawUnitLine{DG}$ \byref{ax:I.V}

And again, since $\drawUnitLine[0.75cm]{HM} + \drawUnitLine[0.75cm]{DH} > \drawUnitLine[0.75cm]{KM} + \drawUnitLine[0.75cm]{DK}$ \byref{prop:I.XXI}, and $\drawUnitLine[0.75cm]{HM} = \drawUnitLine[0.75cm]{KM}$,

$\therefore \drawUnitLine{DK} < \drawUnitLine{DH}$. E o mesmo acontece com outros.
\end{center}

\startsubproposition{III.}
\defineNewPicture[1/4]{
pair M, D, B, G, H, N, O;
numeric r;
r := 7/4u;
M := (0, 0);
D := M shifted (dir(90)*2r);
G := (dir(90)*r) shifted M;
B := (dir(90 - 20)*r) shifted M;
H := (dir(90 + 45)*r) shifted M;
N := (dir(90 - 45)*r) shifted M;
O = whatever[D, N] = whatever[M, B];
byAngleDefine(H, D, M, byyellow, SOLID_SECTOR);
byAngleDefine(M, D, N, byblue, SOLID_SECTOR);
draw byNamedAngleResized();
draw byLine(B, O, byred, SOLID_LINE, REGULAR_WIDTH);
draw byLine(O, N, byyellow, SOLID_LINE, REGULAR_WIDTH);
draw byLine(D, M, byblack, SOLID_LINE, REGULAR_WIDTH);
draw byLine(H, M, byblue, SOLID_LINE, REGULAR_WIDTH);
draw byLine(B, M, byblack, DASHED_LINE, REGULAR_WIDTH);
byLineDefine(N, M, byyellow, DASHED_LINE, REGULAR_WIDTH);
byLineDefine(D, H, byblue, DASHED_LINE, REGULAR_WIDTH);
byLineDefine(D, O, byred, DASHED_LINE, REGULAR_WIDTH);
draw byNamedLineSeq(0)(DH,DO,NM);
draw byCircle.M(M, H, byblack, 0, 0, 0);
draw byLabelsOnPolygon(H, D, O, N)(OMIT_FIRST_LABEL+OMIT_LAST_LABEL, 0);
draw byLabelsOnPolygon(N, M, H)(OMIT_FIRST_LABEL+OMIT_LAST_LABEL, 0);
draw byLabelsOnPolygon(G, H, M)(OMIT_FIRST_LABEL+OMIT_LAST_LABEL, 2);
draw byLabelsOnPolygon(M, B, G)(OMIT_FIRST_LABEL+OMIT_LAST_LABEL, 2);
draw byLabelsOnPolygon(M, N, B)(OMIT_FIRST_LABEL+OMIT_LAST_LABEL, 2);
}
\drawCurrentPictureInMargin
Também são iguais as linhas que fazem ângulos iguais com aquela que passa pelo centro, quer caiam na circunferência côncava ou convexa; e nenhuma terceira linha igual a eles pode ser traçada do mesmo ponto até a circunferência.

\begin{center}
For if $\drawUnitLine{DO,ON} > \drawUnitLine{DH}$, but making $\drawAngle{HDM} = \drawAngle{MDN}$;\\
faça $\drawUnitLine{DO} = \drawUnitLine{DH}$ e trace \drawUnitLine{BM,BO}.

Then in \drawLine[middle][triangleDMO]{DM,DO,BO,BM} and \drawLine[middle][triangleDMH]{HM,DH,DM} we have $\drawUnitLine{DO} = \drawUnitLine{DH}$\\
and \drawUnitLine{DM} common, and also $\drawAngle{MDN} = \drawAngle{HDM}$,

$\therefore \drawUnitLine{BM,BO} = \drawUnitLine{HM}$ \byref{prop:I.IV};\\
but $\drawUnitLine{HM} = \drawUnitLine{BM}$;\\
$\therefore \drawUnitLine{BM} = \drawUnitLine{BM,BO}$, which is absurd.

$\therefore \drawUnitLine{DH} \mbox{ is } \neq \drawUnitLine{DO}$, nem a qualquer parte do \drawUnitLine{DO,ON}, $\therefore \drawUnitLine{DO,ON} \mbox{ is } \ngtr \drawUnitLine{DH}$. % \mbox{ não é } =, \mbox{ não é } >

Nem o $\drawUnitLine{DH} > \drawUnitLine{DO,ON}$, eles são $\therefore =$ um para o outro.
\end{center}
E qualquer outra linha traçada do mesmo ponto até a circunferência deve estar no mesmo lado de uma dessas linhas e estar mais ou menos distante da linha que passa pelo centro e não pode, portanto, ser igual a ela.

\qed

\starttheorem{Prop. IX. Teor.}\label{prop:III.IX}
\defineNewPicture[1/4]{
pair D, A, B, C, F, L, H;
numeric r;
r := 7/4u;
D := (0, 0);
A := (dir(170)*r) shifted D;
B := (dir(-90)*r) shifted D;
C := (dir(-45)*r) shifted D;
L := (dir(45)*r) shifted D;
H := (dir(45 + 180)*r) shifted D;
F := 2/4[D, L];
draw byLine(D, A, byyellow, DASHED_LINE, REGULAR_WIDTH);
draw byLine(D, B, byyellow, SOLID_LINE, REGULAR_WIDTH);
draw byLine(D, C, byblue, SOLID_LINE, REGULAR_WIDTH);
byLineDefine(D, F, byblack, SOLID_LINE, REGULAR_WIDTH);
byLineDefine(F, L, byred, DASHED_LINE, REGULAR_WIDTH);
byLineDefine(D, H, byred, SOLID_LINE, REGULAR_WIDTH);
draw byNamedLineSeq(0)(DH, DF, FL);
draw byCircle.D(D, A, byblue, 0, 0, 0);
draw byLabelsOnCircle(A, B, C, H, L)(D);
draw byLabelsOnPolygon(A, D, F, L)(OMIT_FIRST_LABEL+OMIT_LAST_LABEL, 0);
}
\drawCurrentPictureInMargin
\problem[5]{S}{e}{for tomado um ponto dentro de um círculo \drawCircle[middle]{D}, a partir do qual mais de duas linhas retas iguais (\drawUnitLine{DA}, \drawUnitLine{DB}, \drawUnitLine{DC}) podem ser traçadas até a circunferência, esse ponto deve ser o centro do círculo.}

Pois se for suposto que o ponto \drawPointL[middle][DH,DA,DF]{D} no qual mais de duas retas iguais se encontram não é o centro, algum outro ponto \drawFromCurrentPicture[middle][pontoF]{
startGlobalRotation(-lineAngle.DF);
draw byNamedPointLines(F)("");
stopGlobalRotation;
} deve ser; junte esses dois pontos por \drawUnitLine{DF} e produza-os em ambos os sentidos da circunferência.

Então, como mais de duas linhas retas iguais são traçadas de um ponto que não é o centro até a circunferência, pelo menos duas delas devem estar no mesmo lado do diâmetro \drawUnitLine{DH,DF,FL}; e como a partir de um ponto \drawPointL[middle][DA,DF]{D}, que não é o centro, traçam-se retas até a circunferência; o maior é \drawUnitLine{DF,FL}, que passa pelo centro: e \drawUnitLine{DC} que está mais próximo de \drawUnitLine{DF,FL}, $> \drawUnitLine{DB}$ que é mais remoto \byref{prop:III.VIII}; mas $\drawUnitLine{DC} = \drawUnitLine{DB}$ \byref{\hypref} o que é um absurdo.

O mesmo pode ser demonstrado para qualquer outro ponto, diferente de \drawPointL[middle][DA,DF]{D}, que deve ser o centro do círculo.

\qed

\starttheorem{Prop. X. Teor.}\label{prop:III.X}
\defineNewPicture[1/5]{
pair P, G, H, B, d, dd;
pair Pd, Gd, Hd, Bd, Pdd;
numeric r, t[];
path cr[], crd[];
r := 7/4u;
d := (0, -9/4r);
P := (0, 0);
cr1 := ((fullcircle scaled 7/3r xscaled 4/5) rotated 45) shifted P;
cr2 := ((fullcircle scaled 7/3r xscaled 4/5) rotated -45) shifted P;
H := (subpath (0, 2) of cr1) intersectionpoint cr2;
B := (subpath (0, -2) of cr1) intersectionpoint cr2;
G := (subpath (-2, -4) of cr1) intersectionpoint cr2;
Pd := P shifted d;
crd1 := (fullcircle scaled 2r) shifted Pd;
dd := (0, -3/2r);
crd2 := crd1 shifted dd;
Pdd := Pd shifted dd;
t1 := xpart(crd2 intersectiontimes (subpath (-2, -4) of crd1));
t2 := xpart(crd2 intersectiontimes (subpath (0, -2) of crd1));
Bd := point t1 of crd2;
Hd := point t2 of crd2;
Gd := point -2 of crd1;
crd2 := subpath (t1, t2 + 8) of crd2;
crd2 := crd2 .. (point -2 of crd1) .. cycle;
draw byLine(P, B, byyellow, SOLID_LINE, REGULAR_WIDTH);
draw byLine(P, G, byblack, SOLID_LINE, REGULAR_WIDTH);
draw byLine(P, H, byblue, SOLID_LINE, REGULAR_WIDTH);
draw byArbitraryFigure.fI(cr1, byred, 0, 0);
draw byArbitraryFigure.fII(cr2, byblue, 0, 0);
draw byLine(Pdd, Gd, byblack, SOLID_LINE, REGULAR_WIDTH);
byLineDefine(Pdd, Bd, byyellow, SOLID_LINE, REGULAR_WIDTH);
byLineDefine(Pdd, Hd, byblue, SOLID_LINE, REGULAR_WIDTH);
draw byNamedLineSeq(0)(PddBd,PddHd);
draw byArbitraryFigure.fdI(crd1, byred, 0, 0);
draw byArbitraryFigure.fdII(crd2, byblue, 0, 0);
byCircleDefineR.PI(P, r, byred, 0, 0, 0);
byCircleDefineR.PII(P, r, byblue, 0, 0, 0);
draw byLabelsOnCircle(G, B, H)(PI);
draw byLabelsOnPolygon(G, P, H)(OMIT_FIRST_LABEL+OMIT_LAST_LABEL, 0);
byPointLabelDefine(Gd, "G");
byPointLabelDefine(Hd, "H");
byPointLabelDefine(Bd, "B");
byPointLabelDefine(Dd, "D");
byPointLabelDefine(Pdd, "P");
draw byLabelLineEnd(Gd, Pdd, 0);
draw byLabelPoint(Bd, 180, 2);
draw byLabelPoint(Hd, 0, 2);
draw byLabelsOnPolygon(Hd, Pdd, Bd)(OMIT_FIRST_LABEL+OMIT_LAST_LABEL, 0);
}
\drawCurrentPictureInMargin
\problem{O}{ne}{circle \drawFromCurrentPicture{draw byNamedCircle(PII);} cannot intersect another \drawFromCurrentPicture[middle][circlePI]{draw byNamedCircle(PI);} in more points than two.}

Pois, se for possível, que se cruze em três pontos; do centro de \drawCircle{PII} desenhe \drawUnitLine{PG}, \drawUnitLine{PB} e \drawUnitLine{PH} até os pontos de intersecção;

$\therefore \drawUnitLine{PG} = \drawUnitLine{PB} = \drawUnitLine{PH}$ \byref{def:I.XV}, but as the circles intersect, they have not the same centre \byref{prop:III.V}:

$\therefore$ o ponto assumido não é o centro de \circlePI e $\therefore$ como \drawUnitLine{PG}, \drawUnitLine{PB} e \drawUnitLine{PH} são traçados a partir de um ponto que não é o centro, eles não são iguais a \byref{prop:III.VII,prop:III.VIII}; mas antes foi demonstrado que eram iguais, o que é um absurdo; os círculos, portanto, não se cruzam em três pontos.

\qed

\starttheorem{Prop. XI. Teor.}\label{prop:III.XI}
\defineNewPicture[1/2]{
pair M, N, A, D, F, G, H, K;
numeric r[];
path cr[];
r1 := 9/4u;
r2 := 2/3r1;
M := (0, 0);
N := M shifted (0, +r1-r2);
cr1 := (fullcircle scaled 2r1) shifted M;
cr2 := (fullcircle scaled 2r2) shifted N;
A := (0, r1) shifted M;
F := 1/2[M,N] shifted (dir(-20)*1/3r2);
G := 1/2[M,N] shifted (dir(-20 + 180)*1/3r2);
D := cr2 intersectionpoint (F--10[F, G]);
H := cr1 intersectionpoint (F--10[F, G]);
K := cr1 intersectionpoint (F--10[G, F]);
draw byPolygon(A,F,G)(byyellow);
draw byLine(A, G, byred, SOLID_LINE, REGULAR_WIDTH);
draw byLine(A, F, byblue, DASHED_LINE, REGULAR_WIDTH);
byLineDefine(H, D, byyellow, SOLID_LINE, REGULAR_WIDTH);
byLineDefine(D, G, byyellow, DASHED_LINE, REGULAR_WIDTH);
byLineDefine(G, F, byblack, SOLID_LINE, REGULAR_WIDTH);
byLineDefine(F, K, byblue, SOLID_LINE, REGULAR_WIDTH);
draw byNamedLineSeq(0)(HD, DG, GF, FK);
byPointLabelDefine(M, "F");
byPointLabelDefine(N, "G");
draw byCircle.M(M, A, byblack, 0, 0, 0.5);
draw byCircle.N(N, A, byblue, 0, 0, -0.5);
draw byLabelsOnPolygon(K, F, G, D)(OMIT_FIRST_LABEL+OMIT_LAST_LABEL, 0);
draw byLabelsOnPolygon(A, D, N)(OMIT_FIRST_LABEL+OMIT_LAST_LABEL, 0);
draw byLabelsOnCircle(A, H)(M);
}
\drawCurrentPictureInMargin
\problem[4]{S}{e}{duas circunferências \drawCircle[middle][1/4]{N} e \drawCircle[middle][1/5]{M} se tocam internamente, a linha direita que une seus centros, ao ser produzida, deverá passar por um ponto de contato.}

Pois, se for possível, deixe o \drawSizedLine{GF} unir seus centros e produzi-lo nos dois sentidos; de um ponto de contato desenhe \drawSizedLine{AG} para o centro de \circleN, e do mesmo ponto de contato desenhe \drawSizedLine{AF} para o centro de \circleM.

\begin{center}
Because in
\drawFromCurrentPicture{
startAutoLabeling;
draw byNamedPolygon(AFG);
stopAutoLabeling;
draw byNamedLineFull(A, A, 1, 1,  0, 0)(GF);
};
$\drawSizedLine{GF} + \drawSizedLine{AG} > \drawSizedLine{AF}$ \byref{prop:I.XX},\\
e $\drawSizedLine{AF} = \drawSizedLine{HD,DG,GF}$ pois são raios de \circleM,\\
but $\drawSizedLine{GF} + \drawSizedLine{AG} > \drawSizedLine{HD,DG,GF}$;\\
take away \drawSizedLine{GF} which is common,\\
and $\drawSizedLine{AG} > \drawSizedLine{HD,DG}$;\\
but $\drawSizedLine{AG} = \drawSizedLine{DG}$,\\
porque eles são raios de \circleN,\\
e $\therefore \drawSizedLine{DG} > \drawSizedLine{HD,DG}$ uma parte maior que o todo, o que é um absurdo.
\end{center}

Os centros não estão, portanto, posicionados de modo que uma linha que os una possa passar por qualquer ponto, exceto um ponto de contato.

\qed

\starttheorem{Prop. XII. Teor.}\label{prop:III.XII}
\defineNewPicture[1/4]{
pair M, N, A, C, D, F, G;
numeric r[];
path cr[];
r1 := 3/2u;
r2 := 2u;
M := (0, 0);
N := (0, -r1-r2);
cr1 := (fullcircle scaled 2r1) shifted M;
cr2 := (fullcircle scaled 2r2) shifted N;
A := M shifted (0, -r1);
F := M shifted (dir(185)*1/2r1);
G := N shifted (dir(175)*1/2r2);
C := cr1 intersectionpoint (F--G);
D := cr2 intersectionpoint (F--G);
byLineDefine(F, C, byred, SOLID_LINE, REGULAR_WIDTH);
byLineDefine(C, D, byblack, SOLID_LINE, REGULAR_WIDTH);
byLineDefine(D, G, byblue, SOLID_LINE, REGULAR_WIDTH);
byLineDefine(A, F, byyellow, DASHED_LINE, REGULAR_WIDTH);
byLineDefine(A, G, byyellow, SOLID_LINE, REGULAR_WIDTH);
draw byNamedLineSeq(0)(FC,CD,DG,AG,AF);
draw byCircle.M(M, A, byblue, 0, 0, -1/2);
draw byCircle.N(N, A, byred, 0, 0, -1/2);
byPointLabelDefine(M, "F");
byPointLabelDefine(N, "G");
draw byLabelsOnPolygon(C, F, A)(OMIT_FIRST_LABEL+OMIT_LAST_LABEL, 0);
draw byLabelsOnPolygon(A, G, D)(OMIT_FIRST_LABEL+OMIT_LAST_LABEL, 0);
draw byLabelsOnPolygon(A, C, F)(OMIT_FIRST_LABEL+OMIT_LAST_LABEL, 0);
draw byLabelsOnPolygon(G, D, D)(OMIT_FIRST_LABEL+OMIT_LAST_LABEL, 0);
draw byLabelPoint(A, angle(F-A)-45, 2);
}
\drawCurrentPictureInMargin
\problem[4]{S}{e}{dois círculos \drawCircle{M} e \drawCircle{N} se tocam externamente, a linha reta \drawUnitLine{FC,CD,DG} que une seus centros passa pelo ponto de contato.}

Se for possível, seja o \drawUnitLine{FC,CD,DG} ingressar nos centros, e não passar por um ponto de contato; em seguida, a partir de um ponto de contato, trace \drawUnitLine{AF} e \drawUnitLine{AG} para os centros.

\begin{center}
$\because \drawUnitLine{AF} + \drawUnitLine{AG} > \drawUnitLine{FC,CD,DG}$ \byref{prop:I.XX}\\ %because
and $\drawUnitLine{FC} = \drawUnitLine{AF}$ \byref{def:I.XV},\\
and $\drawUnitLine{DG} = \drawUnitLine{AG}$ \byref{def:I.XV},

$\therefore \drawUnitLine{FC} + \drawUnitLine{DG} > \drawUnitLine{FC,CD,DG}$,\\
uma parte maior que o todo, o que é um absurdo.
\end{center}

Os centros não estão, portanto, posicionados de modo que a linha que os une possa passar por qualquer ponto, exceto o ponto de contato.

\qed

\starttheorem{Prop. XIII. Teor.}\label{prop:III.XIII}
\problem{U}{m}{círculo não pode tocar outro, seja externamente ou internamente, em mais pontos do que um}

\defineNewPicture{
pair M, N, F, G, H, B, D;
numeric r[];
path cr[];
r1 := 7/4u;
r2 := 3/4r1;
M := (0, 0);
N := (dir(120)*(r1-r2)) shifted M;
cr1 := (fullcircle scaled 2r1) shifted M;
cr2 := (fullcircle scaled 2r2) shifted N;
t1 := xpart(cr1 intersectiontimes (M--10[M, N]));
t2 := xpart(cr2 intersectiontimes (M--10[M, N]));
cr2 := (subpath (t2 + 2/3, t2 - 2/3 + 8) of cr2) .. tension 3/2 .. cycle;
B := point (t1 - 1/2) of cr1;
D := point (t1 + 1/2) of cr1;
G := 3/4[B, 1/2[M, N]];
H := 5/4[B, 1/2[M, N]];
draw byLine(D, G, byblack, SOLID_LINE, REGULAR_WIDTH);
byLineDefine(D, H, byred, SOLID_LINE, REGULAR_WIDTH);
byLineDefine(B, G, byblue, DASHED_LINE, REGULAR_WIDTH);
byLineDefine(G, H, byblue, SOLID_LINE, REGULAR_WIDTH);
draw byNamedLineSeq(0)(BG,GH,DH);
draw byCircleR(M, r1, byyellow, 0, 0, 0.5);
draw byArbitraryFigure.fI(cr2, byblue, 0, 0);
byCircleDefineR(M, r1, byyellow, 0, 0, 0);
byCircleDefineR(N, r2, byblue, 0, 0, 0);
byPointLabelDefine(M, "H");
byPointLabelDefine(N, "G");
draw byLabelsOnCircle(D, B)(M);
draw byLabelsOnPolygon(B, G, H, D)(OMIT_FIRST_LABEL+OMIT_LAST_LABEL, 0);
}

\startsubproposition{Figure I.}\drawCurrentPictureInMargin
Pois, se for possível, deixe \drawCircle{M} e \drawCircle{N} se tocarem internamente em dois pontos; desenhar \drawUnitLine{GH} unindo seus centros, e produzi-lo até passar por um dos pontos de contato \byref{prop:III.XI};

\begin{center}
trace \drawUnitLine{DH} e \drawUnitLine{DG}.

But $\drawUnitLine{BG} = \drawUnitLine{DG}$ \byref{def:I.XV},\\
$\therefore$ se \drawUnitLine{GH} for adicionado a ambos,\\
$\drawUnitLine{GH,BG} = \drawUnitLine{GH} + \drawUnitLine{DG}$;

but $\drawUnitLine{GH,BG} = \drawUnitLine{DH}$ \byref{def:I.XV},\\
and $\therefore \drawUnitLine{GH} + \drawUnitLine{DG} = \drawUnitLine{DH}$;

but $\therefore \drawUnitLine{GH} + \drawUnitLine{DG} > \drawUnitLine{DH}$ \byref{prop:I.XX},\\
which is absurd.
\end{center}

\defineNewPicture{
pair P, Q, A, C;
numeric r[], t[];
path cr[];
r1 := u;
P := (0, 0);
Q := P shifted (1/4r1, 0);
cr3 := (fullcircle scaled 2r1) shifted P;
cr4 := (fullcircle scaled 2r1) shifted Q;
t3 := xpart(cr3 intersectiontimes (subpath (0, 4) of cr4));
t4 := xpart(cr3 intersectiontimes (subpath (4, 8) of cr4));
t5 := xpart(cr4 intersectiontimes (subpath (0, 4) of cr3));
t6 := xpart(cr4 intersectiontimes (subpath (4, 8) of cr3));
A := point t3 of cr3;
C := point t4 of cr3;
draw byLine(A, C, byred, SOLID_LINE, REGULAR_WIDTH);
draw byMarkLine(3/7, byblack)(AC);
draw byMarkLine(4/7, byblack)(AC);
draw byArcBE.PI(P, t3, t4, r1, byblue, 0, 0, 0, 0);
draw byArcBE.PII(P, t4, t3 + 8, r1, byblack, 0, 0, 0, 0);
draw byArcBE.QI(Q, t5, t6 + 8, r1, byblue, 0, 0, 0, 0);
draw byArcBE.QII(Q, t5, t6, r1, byblack, 0, 0, 0, 0);
}
\startsubproposition{Figure II.}\drawCurrentPictureInMargin
Mas se os pontos de contato forem as extremidades da linha direita que une os centros, esta linha reta deve ser dividida ao meio em dois pontos diferentes para os dois centros; porque é o diâmetro dos dois círculos, o que é um absurdo.

\vfill\pagebreak

\defineNewPicture{
pair G, H, A, C;
numeric r[];
path cr[];
r1 := 7/4u;
r2 := 3/4r1;
H := (0, 0);
G := H shifted (0, r1 + r2);
cr5 := (fullcircle scaled 2r1) shifted H;
cr6 := (fullcircle scaled 2r2) shifted G;
A := 1/2[point 2 of cr5, point 6 of cr6];
C := A shifted (1/3r2, 0);
cr5 := (subpath (2 + 2/4, 2 - 2/4 + 8) of cr5) .. tension 2 .. cycle;
cr6 := (subpath (6 + 2/4, 6 - 2/4 + 8) of cr6) .. tension 2 .. cycle;
byLineDefine(H, A, byred, SOLID_LINE, REGULAR_WIDTH);
byLineDefine(A, G, byred, DASHED_LINE, REGULAR_WIDTH);
byLineDefine(C, H, byblue, DASHED_LINE, REGULAR_WIDTH);
byLineDefine(C, G, byblack, SOLID_LINE, REGULAR_WIDTH);
draw byNamedLineSeq(-1/3)(HA,AG,CG,CH);
draw byArbitraryFigure.fII(cr5, byyellow, 0, 0);
draw byArbitraryFigure.fIII(cr6, byblue, 0, 0);
byCircleDefineR(H, r1, byyellow, 0, 0, 0);
byCircleDefineR(G, r2, byblue, 0, 0, 0);
draw byLabelsOnPolygon(H, G, C)(ALL_LABELS, 0);
draw byLabelPoint(A, angle(G-A)+45, 3);
}
\startsubproposition{Figure III.}\drawCurrentPictureInMargin
A seguir, se for possível, deixe \drawCircle{H} e \drawCircle{G} se tocarem externamente em dois pontos; desenhe \drawUnitLine{HA,AG} unindo os centros dos círculos, e passando por um dos pontos de contato, e desenhe \drawUnitLine{CH} e \drawUnitLine{CG}.

\begin{center}
$\drawUnitLine{CH} = \drawUnitLine{HA}$ \byref{def:I.XV};\\
and $\drawUnitLine{AG} = \drawUnitLine{CG}$ \byref{def:I.XV};

$\therefore \drawUnitLine{CG} + \drawUnitLine{CH} = \drawUnitLine{HA,AG}$;

but $\drawUnitLine{CG} + \drawUnitLine{CH} > \drawUnitLine{HA,AG}$ \byref{prop:I.XX}, which is absurd.
\end{center}

There is therefore no case in which two circles can touch one another in two points.

\qed

\starttheorem{Prop. XIV. Teor.}\label{prop:III.XIV}
\defineNewPicture{
pair A, B, C, D, E, F, G;
numeric r;
r := 9/4u;
E := (0, 0);
A := (dir(90-20)*r) shifted E;
B := (dir(90-130)*r) shifted E;
C := (dir(90+20)*r) shifted E;
D := (dir(90+130)*r) shifted E;
F = whatever[A, B] = whatever[E, E shifted ((A-B) rotated 90)];
G = whatever[C, D] = whatever[E, E shifted ((C-D) rotated 90)];
byAngleDefine(E, F, A, byyellow, SOLID_SECTOR);
byAngleDefine(E, G, C, byblack, ARC_SECTOR);
draw byNamedAngleResized();
draw byLine(E, A, byblack, SOLID_LINE, REGULAR_WIDTH);
draw byLine(E, C, byblue, SOLID_LINE, REGULAR_WIDTH);
byLineDefine(E, F, byblack, DASHED_LINE, REGULAR_WIDTH);
byLineDefine(E, G, byblue, DASHED_LINE, REGULAR_WIDTH);
draw byNamedLineSeq(0)(EF,EG);
byLineDefine(A, F, byred, SOLID_LINE, REGULAR_WIDTH);
byLineDefine(F, B, byred, DASHED_LINE, REGULAR_WIDTH);
draw byNamedLineSeq(0)(AF, FB);
byLineDefine(C, G, byyellow, SOLID_LINE, REGULAR_WIDTH);
byLineDefine(G, D, byyellow, DASHED_LINE, REGULAR_WIDTH);
draw byNamedLineSeq(0)(CG, GD);
draw byCircleR(E, r, byblue, 0, 0, 0);
draw byLabelsOnCircle(A, B, C, D)(E);
draw byLabelsOnPolygon(F, E, G)(OMIT_FIRST_LABEL+OMIT_LAST_LABEL, 0);
draw byLabelLineEnd(G, E, 0);
draw byLabelLineEnd(F, E, 0);
}
\drawCurrentPictureInMargin
\problem{E}{qual}{straight lines
$\left(\vcenter{\nointerlineskip\hbox{\drawProportionalLine{AF,FB}}\nointerlineskip\hbox{\drawProportionalLine{CG,GD}}}\right)$
inscrito em um círculo estão igualmente distantes do centro; e também, linhas retas igualmente distantes do centro são iguais.}

\begin{center}
Do centro de \drawCircle[middle][1/5]{E} trace \drawProportionalLine{EF} $\perp$ para \drawProportionalLine{AF,FB} e $\drawProportionalLine{EG} \perp \drawProportionalLine{CG,GD}$, junte \drawProportionalLine{EA} e \drawProportionalLine{EC}.

Then $\drawProportionalLine{CG} = \mbox{ half } \drawProportionalLine{CG,GD}$ \byref{prop:III.III}\\
and $\drawProportionalLine{AF} = \frac{1}{2} \drawProportionalLine{AF,FB}$ \byref{prop:III.III},\\
since $\drawProportionalLine{CG,GD} = \drawProportionalLine{AF,FB}$ \byref{\hypref}\\
$\therefore \drawProportionalLine{CG} = \drawProportionalLine{AF}$,\\
and $\drawProportionalLine{EA} = \drawProportionalLine{EC}$ \byref{def:I.XV}\\
$\therefore \drawProportionalLine{EA}^2 = \drawProportionalLine{EC}^2$;\\
but since \drawAngle{F} is a right angle\\
$\drawProportionalLine{EA}^2 = \drawProportionalLine{EF}^2 + \drawProportionalLine{AF}^2$ \byref{prop:I.XLVII}\\
and $\drawProportionalLine{EC}^2 = \drawProportionalLine{EG}^2 + \drawProportionalLine{CG}^2$ for the same reason,\\
$\therefore \drawProportionalLine{EF}^2 + \drawProportionalLine{AF}^2 = \drawProportionalLine{EG}^2 + \drawProportionalLine{CG}^2$\\
$\therefore \drawProportionalLine{EF}^2 = \drawProportionalLine{EG}^2$\\
$\therefore \drawProportionalLine{EF} = \drawProportionalLine{EG}$
\end{center}

Além disso, se as linhas \drawProportionalLine{AF,FB} e \drawProportionalLine{CG,GD} estiverem igualmente distantes do centro; isto é, se as perpendiculares \drawProportionalLine{EF} e \drawProportionalLine{EG} forem iguais, então $\drawProportionalLine{AF,FB} = \drawProportionalLine{CG,GD}$.

\begin{center}
For, as in the preceding case,\\
$\drawProportionalLine{EG}^2 + \drawProportionalLine{CG}^2 = \drawProportionalLine{AF}^2 + \drawProportionalLine{EF}^2$\\
but $\drawProportionalLine{EG}^2 = \drawProportionalLine{EF}^2$

$\therefore \drawProportionalLine{CG} = \drawProportionalLine{AF}$,% de melhoria: no original havia quadrados desses aqui
e os duplos destes \drawProportionalLine{AF,FB} e \drawProportionalLine{CG,GD} também são iguais.
\end{center}

\qed

\starttheorem{Prop. XV. Teor.}\label{prop:III.XV}
\problem{O}{}{diâmetro é a maior linha reta num círculo: e, de todas as outras, aquela que está mais próxima do centro é maior que a mais remota.}

\startsubproposition{Figure I.}
\defineNewPicture[2/3]{
pair A, D, E, F, G, M, N;
numeric r;
r := 9/4u;
E := (0, 0);
A := (r, 0) shifted E;
D := (-r, 0) shifted E;
F := (dir(90-30)*r) shifted E;
G := (dir(90+30)*r) shifted E;
M := (dir(90-60)*r) shifted E;
N := (dir(90+60)*r) shifted E;
byAngleDefine(N, E, G, byred, SOLID_SECTOR);
byAngleDefine(G, E, F, byyellow, SOLID_SECTOR);
byAngleDefine(F, E, M, byred, SOLID_SECTOR);
draw byNamedAngleResized();
draw byLine(E, M, byyellow, DASHED_LINE, REGULAR_WIDTH);
draw byLine(E, N, byyellow, SOLID_LINE, REGULAR_WIDTH);
draw byLine(E, F, byblue, DASHED_LINE, REGULAR_WIDTH);
draw byLine(E, G, byblack, DASHED_LINE, REGULAR_WIDTH);
draw byLine(D, E, byred, SOLID_LINE, REGULAR_WIDTH);
draw byLine(E, A, byblack, SOLID_LINE, REGULAR_WIDTH);
draw byLine(F, G, byred, DASHED_LINE, REGULAR_WIDTH);
draw byLine(M, N, byblue, SOLID_LINE, REGULAR_WIDTH);
draw byCircleR(E, r, byblack, 0, 0, 0);
draw byLabelsOnCircle(A, D, M, N, F, G)(E);
draw byLabelsOnPolygon(A, E, D)(OMIT_FIRST_LABEL+OMIT_LAST_LABEL, 0);
}
\drawCurrentPictureInMargin
\begin{center}
O diâmetro \drawUnitLine{DE,EA} é $>$ qualquer linha \drawUnitLine{MN}.

Para traçar \drawUnitLine{EN} e \drawUnitLine{EM}.

Then $\drawUnitLine{EM} = \drawUnitLine{EA}$\\
and $\drawUnitLine{EN} = \drawUnitLine{DE}$,\\
$\therefore \drawUnitLine{EN} + \drawUnitLine{EM} = \drawUnitLine{DE,EA}$\\
but $\drawUnitLine{EN} + \drawUnitLine{EM} > \drawUnitLine{MN}$ \byref{prop:I.XX}

$\therefore \drawUnitLine{DE,EA} > \drawUnitLine{MN}$.
\end{center}

Novamente, a linha que está mais próxima do centro é maior que a que está mais distante.

Primeiro, sejam as linhas fornecidas \drawUnitLine{MN} e \drawUnitLine{FG}, que estão no mesmo lado do centro e não se cruzam;

\begin{center}
draw $\left\{
\vcenter{
\nointerlineskip\hbox{\drawUnitLine{EN},}
\nointerlineskip\hbox{\drawUnitLine{EM},}
\nointerlineskip\hbox{\drawUnitLine{EG},}
\nointerlineskip\hbox{\drawUnitLine{EF}.}
}
\right.$

In
\drawFromCurrentPicture{
draw byNamedAngle(NEG,GEF,FEM);
startAutoLabeling;
draw byNamedLineSeq(0)(MN,EM,EN);
stopAutoLabeling;
}
and
\drawFromCurrentPicture{
draw byNamedAngle(GEF);
startAutoLabeling;
draw byNamedLineSeq(0)(FG,EF,EG);
stopAutoLabeling;
},\\
$\drawUnitLine{EN} \mbox{ e } \drawUnitLine{EM} = \drawUnitLine{EG} \mbox{ e } \drawUnitLine{EF}$;\\
but $\drawAngle{NEG,GEF,FEM} > \drawAngle{GEF}$,

$\therefore \drawUnitLine{MN} > \drawUnitLine{EF}$ \byref{prop:I.XXIV}
\end{center}

\vfill\pagebreak

\startsubproposition{Figure II.}
\defineNewPicture{
pair A, B, C, D, E, F, G, M, N, K, L, H;
numeric r;
r := 9/4u;
E := (0, 0);
A := (r, 0) shifted E;
D := (-r, 0) shifted E;
F := (dir(90-30)*r) shifted E;
G := (dir(90+30)*r) shifted E;
B := (dir(-90-60)*r) shifted E;
C := (dir(-90+60)*r) shifted E;
M := (dir(90-60)*r) shifted E;
N := (dir(90+60)*r) shifted E;
H := 1/2[B, C];
K := 1/2[F, G];
L := 1/2[M, N];
draw byLine(E, L, byyellow, DASHED_LINE, REGULAR_WIDTH);
draw byLine(L, K, byred, DASHED_LINE, REGULAR_WIDTH);
draw byLine(E, H, byblue, DASHED_LINE, REGULAR_WIDTH);
draw byLine(F, G, byyellow, SOLID_LINE, REGULAR_WIDTH);
draw byLine(M, N, byblue, SOLID_LINE, REGULAR_WIDTH);
draw byLine(B, C, byred, SOLID_LINE, REGULAR_WIDTH);
draw byLine(D, A, byblack, SOLID_LINE, REGULAR_WIDTH);
draw byCircleR(E, r, byblack, 0, 0, 0);
draw byLabelsOnCircle(A, D, B, C, F, G, M, N)(E);
draw byLabelLineEnd(K, H, 0);
draw byLabelLineEnd(H, K, 0);
draw byLabelPoint(E, angle(E-H)-45, 2);
draw byLabelPoint(L, angle(E-H)-45, 2);
}
\drawCurrentPictureInMargin
Sejam as linhas fornecidas \drawUnitLine{BC} e \drawUnitLine{FG} que estão em lados diferentes do centro ou se cruzam; do centro trace \drawUnitLine{EL,LK} e $\drawUnitLine{EH} \perp \drawUnitLine{FG} \mbox{ e } \drawUnitLine{BC}$,

\begin{center}
make $\drawUnitLine{EH} = \drawUnitLine{EL}$,\\
e trace $\drawUnitLine{MN} \perp \drawUnitLine{EL,LK}$.

Como \drawUnitLine{BC} e \drawUnitLine{MN} estão igualmente distantes do centro, $\drawUnitLine{BC} = \drawUnitLine{MN}$ \byref{prop:III.XIV};\\
but $\drawUnitLine{MN} > \drawUnitLine{FG}$ \byref{prop:III.XV},

$\therefore \drawUnitLine{BC} > \drawUnitLine{FG}$.
\end{center}

\qed

\starttheorem{Prop. XVI. Teor.}\label{prop:III.XVI}
\defineNewPicture[1/5]{
pair A, B, C, D, E, F, G, H, K;
numeric r;
r :=2u;
D := (0, 0);
A := (0, -r);
B := (0, r);
C := (dir(190)*r) shifted D;
E := (4/3r, -r);
F := (4/3r, -1/3r);
G := 11/12[A, F];
H := (dir(angle(G-D))*r) shifted D;
K := (-r, -r);
byAngleDefine(C, A, D, byyellow, SOLID_SECTOR);
byAngleDefine(D, A, G, byblue, SOLID_SECTOR);
byAngleDefine(G, A, E, byred, SOLID_SECTOR);
byAngleDefine(D, C, A, byblack, SOLID_SECTOR);
byAngleDefine(A, G, D, byblack, ARC_SECTOR);
draw byNamedAngleResized();
draw byLine(A, C, byred, SOLID_LINE, REGULAR_WIDTH);
draw byLine(D, C, byblue, SOLID_LINE, REGULAR_WIDTH);
draw byLine(D, H, byblue, DASHED_LINE, REGULAR_WIDTH);
draw byLine(H, G, byblack, DASHED_LINE, REGULAR_WIDTH);
draw byLine(A, G, byred, DASHED_LINE, REGULAR_WIDTH);
draw byLine(G, F, byblack, DASHED_LINE, REGULAR_WIDTH);
draw byLine(B, D, byyellow, DASHED_LINE, REGULAR_WIDTH);
draw byLine(D, A, byblack, SOLID_LINE, REGULAR_WIDTH);
draw byLineFull(E, K, byyellow, 0, 0)(E, K, 0, 0, -1);
draw byCircle.D(D, A, byblue, 0, 0, 0);
draw byLabelsOnCircle(B, C)(D);
draw byLabelsOnPolygon(E, A, K, noPoint)(ALL_LABELS, 1);
draw byLabelsOnPolygon(F, G, A)(OMIT_FIRST_LABEL+OMIT_LAST_LABEL, 0);
draw byLabelsOnPolygon(C, D, B)(OMIT_FIRST_LABEL+OMIT_LAST_LABEL, 0);
draw byLabelPoint(H, angle(G-H)+45, 2);
}
\drawCurrentPictureInMargin
\problem{A}{}{reta \drawUnitLine{EK} traçada a partir da extremidade do diâmetro \drawUnitLine{BD,DA}
de um círculo perpendicular a ele cai sem o círculo\\
E se qualquer linha reta \drawUnitLine{AG} for traçada a partir de um ponto perpendicular ao ponto de contato, ela corta o círculo.}

\startsubproposition{Parte I.}
Se for possível, seja \drawUnitLine{AC}, que encontra o círculo novamente, ser $\perp \drawUnitLine{DA}$ e trace \drawUnitLine{DC}.

\begin{center}
Then, $\because \drawUnitLine{DA} = \drawUnitLine{DC}$,  $\drawAngle{CAD} = \drawAngle{C}$ \byref{prop:I.V},\\
e $\therefore$ cada um desses ângulos é agudo \byref{prop:I.XVII}\\
but $\drawAngle{CAD} = \drawRightAngle$ \byref{\hypref}, which is absurd,

portanto, \drawUnitLine{AC} traçado $\perp \drawUnitLine{DA}$ não encontra o círculo novamente.
\end{center}

\startsubproposition{Parte II.}
Seja \drawUnitLine{EK} $\perp \drawUnitLine{DA}$ e seja \drawUnitLine{AG} desenhado a partir de um ponto \drawFromCurrentPicture{draw byNamedPointLines(G,"GF");} entre \drawUnitLine{EK} e o círculo, que, se for possível, não corta o círculo.

\begin{center}
$\because \drawAngle{DAG,GAE} = \drawRightAngle$,\\ %Because
$\therefore \drawAngle{DAG}$ is an acute angle;\\
suponha que $\drawUnitLine{DH,HG} \perp \drawUnitLine{AG}$, traçado a partir do centro do círculo, deve cair ao lado de \drawAngle{DAG} o ângulo agudo.

$\therefore$ \drawAngle{G} que deveria ser um ângulo reto, é $> \drawAngle{DAG}$,

$\therefore \drawUnitLine{DA} > \drawUnitLine{DH,HG}$;

but $\drawUnitLine{DH} = \drawUnitLine{DA}$,
\end{center}

\noindent e $\therefore \drawUnitLine{DH} > \drawUnitLine{DH,HG}$, uma parte maior que o todo, o que é um absurdo. Portanto, o ponto não fica fora do círculo e, portanto, a linha reta \drawUnitLine{AG} corta o círculo.

\qed

\startproblem{Prop. XVII. prob.}\label{prop:III.XVII}
\defineNewPicture[1/2]{
pair A, B, D, E, F;
numeric r[], a;
path cr[];
r1 := 6/4u;
r2 := 9/4u;
E := (0, 0);
cr1 := (fullcircle scaled 2r1) shifted E;
cr2 := (fullcircle scaled 2r2) shifted E;
A := (dir(50)*r2) shifted E;
D := (dir(50)*r1) shifted E;
F := cr2 intersectionpoint (D--D shifted (dir(angle(A-E) - 90)*r2));
B := cr1 intersectionpoint (E--F);
a := angle(B-E);
forsuffixes i=A, B, D, F:
i := ((i shifted -E) rotated -a) shifted E;
endfor;
byAngleDefine(A, B, E, byyellow, SOLID_SECTOR);
byAngleDefine(F, D, E, byyellow, SOLID_SECTOR);
byAngleDefine(F, E, A, byblue, SOLID_SECTOR);
draw byNamedAngleResized();
draw byLine(A, B, byblue, SOLID_LINE, REGULAR_WIDTH);
draw byLine(F, D, byblue, DASHED_LINE, REGULAR_WIDTH);
byLineDefine(A, D, byred, SOLID_LINE, REGULAR_WIDTH);
byLineDefine(D, E, byred, DASHED_LINE, REGULAR_WIDTH);
byLineDefine(F, B, byblack, SOLID_LINE, REGULAR_WIDTH);
byLineDefine(B, E, byblack, DASHED_LINE, REGULAR_WIDTH);
draw byNamedLineSeq(0)(AD,DE,BE,FB);
draw byCircleR.EI(E, r1, byred, 0, 0, -1);
draw byCircleR.EII(E, r2, byyellow, 0, 0, 0);
draw byLabelsOnCircle(A, F)(EII);
draw byLabelsOnPolygon(F, E, A)(OMIT_FIRST_LABEL+OMIT_LAST_LABEL, 0);
draw byLabelPoint(D, angle(A-E)+45, 2);
draw byLabelPoint(B, angle(F-E)-45, 2);
}
\drawCurrentPictureInMargin
\problem{P}{ara}{traçar uma tangente a um determinado círculo \drawCircle[middle][1/5]{EI} a partir de um determinado ponto, dentro ou fora de sua circunferência.}

Se um determinado ponto estiver na circunferência, como em \drawPointL[middle][FB]{B}, é claro que a reta $\drawUnitLine{AB} \perp \drawUnitLine{BE}$, o raio, será a tangente necessária \byref{prop:III.XVI}.

Mas se o ponto fornecido \drawPointL[middle][FB]{A} estiver fora da circunferência,

\begin{center}
trace \drawUnitLine{DE,AD} dele para o centro, cortando \circleEI;\\
e trace $\drawUnitLine{FD} \perp \drawUnitLine{DE}$,\\
describe \drawCircle[middle][1/6]{EII} concentric with \circleEI\ radius $= \drawUnitLine{DE,AD}$,\\ % improvement: These two lines were not originally there, but without them this proof makes no sense
trace \drawUnitLine{BE,FB} para o centro a partir do ponto onde \drawUnitLine{FD} cai na circunferência \circleEII\,\\
trace $\drawUnitLine{AB} \perp \drawUnitLine{BE,FB}$ a partir do ponto onde corta \circleEI,\\
%
Then \drawUnitLine{AB} will be the tangent required.\\
For in \drawLine[bottom]{FD,FB,BE,DE} and \drawLine[bottom]{BE,DE,AD,AB} $\drawUnitLine{AD,DE} = \drawUnitLine{FB,BE}$, \drawAngle{E} common,\\
and $\drawUnitLine{DE} = \drawUnitLine{BE}$,\\
$\therefore \mbox{ \byref{prop:I.IV} } \drawAngle{B} = \drawAngle{D} = \drawRightAngle$,\\
$\therefore \drawUnitLine{FD}$ é uma tangente a \circleEI.
\end{center}

\qed

\starttheorem{Prop. XVIII. Teor.}\label{prop:III.XVIII}
\defineNewPicture{
pair B, C, D, F, G;
numeric r;
r := 7/4u;
F := (0, 0);
C := (r, 0);
G := (r, 6/5r);
D := 6/5[C, G];
B := (dir(angle(G-F))*r) shifted F;
draw byCircle.F(F, C, byyellow, 0, 0, 0);
byAngleDefine(F, C, G, byred, SOLID_SECTOR);
byAngleDefine(C, G, F, byyellow, SOLID_SECTOR);
draw byNamedAngleResized();
byLineDefine(F, B, byred, SOLID_LINE, REGULAR_WIDTH);
byLineDefine(B, G, byred, DASHED_LINE, REGULAR_WIDTH);
byLineDefine(C, D, byblue, DASHED_LINE, REGULAR_WIDTH);
byLineDefine(F, C, byblue, SOLID_LINE, REGULAR_WIDTH);
draw byNamedLineSeq(0)(BG,FB,FC,CD);
draw byLabelsOnCircle(C)(F);
draw byLabelsOnPolygon(C, F, G)(OMIT_FIRST_LABEL+OMIT_LAST_LABEL, 0);
draw byLabelPoint(B, angle(G-F)+45, 2);
draw byLabelPoint(G, angle(D-C)-90, 1);
draw byLabelPoint(D, angle(D-C)-90, 1);
}
\drawCurrentPictureInMargin
\problem[4]{S}{e}{uma reta \drawUnitLine{CD} for tangente a um círculo, a reta \drawUnitLine{FC} traçada do centro ao ponto de contato é perpendicular a ela.}

\begin{center}
Pois, se for possível, seja \drawUnitLine{FB,BG} $\perp \drawUnitLine{CD}$,\\
then $\because \drawAngle{G} = \drawRightAngle$, \drawAngle{C} is acute \byref{prop:I.XVII} %because

$\therefore \drawUnitLine{FC} > \drawUnitLine{FB,BG}$ \byref{prop:I.XIX};

but $\drawUnitLine{FC} = \drawUnitLine{FB}$,\\
e $\therefore \drawUnitLine{FB} > \drawUnitLine{FB,BG}$, uma parte maior que o todo, o que é um absurdo.

$\therefore \drawUnitLine{FB,BG}$ is not $\perp \drawUnitLine{CD}$;
\end{center}

\noindent e da mesma maneira pode ser demonstrado que nenhuma outra linha exceto \drawUnitLine{FC} é perpendicular a \drawUnitLine{CD}.

\qed

\starttheorem{Prop. XIX. Teor.}\label{prop:III.XIX}
\defineNewPicture{
pair A, C, E, F, G;
numeric r;
r := 7/4u;
G := (0, 0);
A := (-r, 0) shifted G;
C := (r, 0) shifted G;
E := (r, 6/5r) shifted G;
F := (-1/7r, 1/2r) shifted G;
byAngleDefine(A, C, F, byblue, SOLID_SECTOR);
byAngleDefine(F, C, E, byyellow, SOLID_SECTOR);
draw byNamedAngleResized();
draw byLine(A, C, byyellow, SOLID_LINE, REGULAR_WIDTH);
draw byLine(C, E, byblue, SOLID_LINE, REGULAR_WIDTH);
draw byLine(C, F, byred, DASHED_LINE, REGULAR_WIDTH);
draw byCircleR(G, r, byred, 0, 0, 0);
draw byLabelsOnCircle(A, C)(G);
draw byLabelLineEnd(F, C, 0);
draw byLabelLineEnd(E, C, 0);
}
\drawCurrentPictureInMargin
\problem{S}{e}{uma linha reta \drawUnitLine{CE} for tangente a um círculo, a linha reta \drawUnitLine{AC}, traçada perpendicularmente a ela a partir do ponto de contato, passa pelo centro do círculo.}

Pois, se for possível, deixe o centro sem \drawUnitLine{AC}, e desenhe \drawUnitLine{CF} do suposto centro até o ponto de contato.

\begin{center}
$\because \drawUnitLine{CF} \perp \drawUnitLine{CE}$ \byref{prop:III.XVIII}\\ % Because
$\therefore \drawAngle{FCE} = \drawRightAngle$, a right angle;\\
but $\drawAngle{ACF,FCE} = \drawRightAngle$ \byref{\hypref},

e $\therefore \drawAngle{FCE} = \drawAngle{ACF,FCE}$, uma parte igual ao todo, o que é um absurdo.
\end{center}

Portanto o ponto assumido não é o centro; e da mesma maneira pode ser demonstrado que nenhum outro ponto sem \drawUnitLine{AC} é o centro.

\qed

\starttheorem{Prop. XX. Teor.}\label{prop:III.XX}
\problem[4]{O}{}{ângulo no centro de um círculo é o dobro do ângulo na circunferência, quando eles têm a mesma parte da circunferência como base.}

\startsubproposition{Figure I.}
\defineNewPicture{
pair A, E, F, C;
r := 7/4u;
E := (0, 0);
A := (dir(80)*r) shifted E;
F := (dir(80 + 180)*r) shifted E;
C := (dir(-30)*r) shifted E;
byAngleDefine(C, A, F, byyellow, SOLID_SECTOR);
byAngleDefine(C, E, F, byblue, SOLID_SECTOR);
byAngleDefine(E, C, A, byred, SOLID_SECTOR);
draw byNamedAngleResized();
draw byLine(A, C, byblack, SOLID_LINE, THIN_WIDTH);
draw byLine(E, C, byblack, SOLID_LINE, REGULAR_WIDTH);
byLineDefine(E, F, byred, DASHED_LINE, REGULAR_WIDTH);
byLineDefine(E, A, byred, SOLID_LINE, REGULAR_WIDTH);
draw byNamedLineSeq(0)(EF,EA);
draw byCircleR(E, r, byblue, 0, 0, 0);
draw byLabelsOnCircle(F, A, C)(E);
draw byLabelsOnPolygon(F, E, A)(OMIT_FIRST_LABEL+OMIT_LAST_LABEL, 0);
}
\drawCurrentPictureInMargin
\begin{center}
Seja o centro do círculo estar em \drawUnitLine{EF,EA}\\
a side of \drawAngle{CAF}.

$\because \drawUnitLine{EC} = \drawUnitLine{EA}$,\\ %Because
$\drawAngle{CAF} = \drawAngle{C}$ \byref{prop:I.V}.

But $\drawAngle{CEF} = \drawAngle{CAF} + \drawAngle{C}$,\\
or $\drawAngle{CEF} = \mbox{ twice } \drawAngle{CAF}$ \byref{prop:I.XXXII}.
\end{center}

\startsubproposition{Figure II.}
\defineNewPicture{
pair A, E, F, C, B;
r := 7/4u;
E := (0, 0);
A := (dir(80)*r) shifted E;
F := (dir(80 + 180)*r) shifted E;
B := (dir(185)*r) shifted E;
C := (dir(-30)*r) shifted E;
byAngleDefine(C, A, F, byyellow, SOLID_SECTOR);
byAngleDefine(C, E, F, byblue, SOLID_SECTOR);
byAngleDefine(E, C, A, byyellow, SOLID_SECTOR);
byAngleDefine(B, A, F, byred, SOLID_SECTOR);
byAngleDefine(B, E, F, byblack, SOLID_SECTOR);
byAngleDefine(E, B, A, byred, SOLID_SECTOR);
draw byNamedAngleResized();
draw byLine(A, C, byblack, SOLID_LINE, THIN_WIDTH);
draw byLine(E, C, byblack, SOLID_LINE, THIN_WIDTH);
draw byLine(A, B, byblack, SOLID_LINE, THIN_WIDTH);
draw byLine(E, B, byblack, SOLID_LINE, THIN_WIDTH);
draw byLine(A, F, byblack, SOLID_LINE, REGULAR_WIDTH);
draw byCircleR(E, r, byblue, 0, 0, 0);
draw byLabelsOnCircle(F, A, B, C)(E);
draw byLabelsOnPolygon(B, E, A)(OMIT_FIRST_LABEL+OMIT_LAST_LABEL, 0);
}
\drawCurrentPictureInMargin
\begin{center}
Seja o centro estar dentro de \drawAngle{BAF,CAF}, o ângulo na circunferência;

trace \drawUnitLine{AF} do ponto angular até o centro do círculo;\\
depois $\drawAngle{B} = \drawAngle{BAF}$ e $\drawAngle{C} = \drawAngle{CAF}$, devido à igualdade dos lados \byref{prop:I.V}.

Hence $\drawAngle{BAF} + \drawAngle{B} + \drawAngle{CAF} + \drawAngle{C} = \mbox{ twice } \drawAngle{BAF,CAF}$.

But $\drawAngle{BEF} = \drawAngle{BAF} + \drawAngle{B}$,\\
and $\drawAngle{CEF} = \drawAngle{CAF} + \drawAngle{C}$,\\
$\therefore \drawAngle{BEF,CEF} = \mbox{ twice } \drawAngle{BAF,CAF}$.
\end{center}

\vfill\pagebreak

\startsubproposition{Figure III.}
\defineNewPicture{
pair E, C, F, G, D;
r := 7/4u;
E := (0, 0);
F := (dir(80 + 180)*r) shifted E;
C := (dir(-30)*r) shifted E;
D := (dir(30)*r) shifted E;
G := (dir(30 + 180)*r) shifted E;
byAngleDefine(G, E, F, byblue, SOLID_SECTOR);
byAngleDefine(F, E, C, byyellow, SOLID_SECTOR);
byAngleDefine(G, D, F, byblack, SOLID_SECTOR);
byAngleDefine(F, D, C, byred, SOLID_SECTOR);
draw byNamedAngleResized();
draw byLine(D, C, byblack, SOLID_LINE, THIN_WIDTH);
draw byLine(D, F, byblack, SOLID_LINE, THIN_WIDTH);
draw byLine(E, C, byblack, SOLID_LINE, THIN_WIDTH);
draw byLine(E, F, byblack, SOLID_LINE, THIN_WIDTH);
draw byLine(D, G, byred, SOLID_LINE, REGULAR_WIDTH);
draw byCircleR(E, r, byblue, 0, 0, 0);
draw byLabelsOnCircle(D, G, C, F)(E);
draw byLabelsOnPolygon(G, E, D)(OMIT_FIRST_LABEL+OMIT_LAST_LABEL, 0);
}
\drawCurrentPictureInMargin\begin{center}
Seja o centro ficar sem \drawAngle{FDC}\\
e trace \drawUnitLine{DG}, o diâmetro.

$\because \drawAngle{GEF,FEC} = \mbox{ twice } \drawAngle{GDF,FDC}$;\\ %Because
and $\drawAngle{GEF} = \mbox{ twice } \drawAngle{GDF}$ (case I.);

$\therefore \drawAngle{FEC} = \mbox{ twice } \drawAngle{FDC}$.
\end{center}

\qed

\starttheorem{Prop. XXI. Teor.}\label{prop:III.XXI}
\problem{O}{s}{ângulos no mesmo segmento de um círculo são iguais.}

\defineNewPicture{
pair A, B, D, E, F;
numeric r;
r := 7/4u;
F := (0, 0);
B := (dir(-90 - 50)*r) shifted F;
D := (dir(-90 + 50)*r) shifted F;
A := (dir(90 + 25)*r) shifted F;
E := (dir(90 - 35)*r) shifted F;
byAngleDefine(B, A, D, byred, SOLID_SECTOR);
byAngleDefine(B, E, D, byblue, SOLID_SECTOR);
byAngleDefine(B, F, D, byyellow, SOLID_SECTOR);
draw byNamedAngleResized();
byLineDefine(B, F, byblue, SOLID_LINE, REGULAR_WIDTH);
byLineDefine(D, F, byred, SOLID_LINE, REGULAR_WIDTH);
draw byNamedLineSeq(0)(BF,DF);
draw byLine(B, D, byblack, DASHED_LINE, REGULAR_WIDTH);
draw byLine(B, A, byblack, SOLID_LINE, THIN_WIDTH);
draw byLine(B, E, byblack, SOLID_LINE, THIN_WIDTH);
draw byLine(D, A, byblack, SOLID_LINE, THIN_WIDTH);
draw byLine(D, E, byblack, SOLID_LINE, THIN_WIDTH);
draw byCircleR(F, r, byblue, 0, 0, 0);
draw byLabelsOnCircle(B, D, E, A)(F);
draw byLabelsOnPolygon(B, F, D)(OMIT_FIRST_LABEL+OMIT_LAST_LABEL, 0);
}
\startsubproposition{Figure I.}
\drawCurrentPictureInMargin
Seja o segmento ser maior que um semicírculo e trace \drawUnitLine{DF} e \drawUnitLine{BF} no centro.

\begin{center}
$\drawAngle{F} = \mbox{ twice } \drawAngle{A} \mbox{ or twice } = \drawAngle{E}$ \byref{prop:III.XX};

$\therefore \drawAngle{A} = \drawAngle{E}$
\end{center}

\defineNewPicture{
pair A, B, D, E, F, G;
numeric r;
path cr;
r := 7/4u;
F := (0, 0);
cr := (fullcircle scaled 2r) shifted F;
B := (dir(90 + 85)*r) shifted F;
D := (dir(90 - 85)*r) shifted F;
A := (dir(90 + 25)*r) shifted F;
E := (dir(90 - 35)*r) shifted F;
G := (dir(-90 + 20)*r) shifted F;
byAngleDefine(B, A, G, byyellow, SOLID_SECTOR);
byAngleDefine(G, A, D, byred, SOLID_SECTOR);
byAngleDefine(B, E, G, byblue, SOLID_SECTOR);
byAngleDefine(G, E, D, byblack, SOLID_SECTOR);
draw byNamedAngleResized();
draw byFilledCircleSegment.BG(F, r, xpart(cr intersectiontimes (F -- 2[F, B])), xpart(cr intersectiontimes (F -- 2[F, G])), byblue);
draw byFilledCircleSegment.GD(F, r, xpart(cr intersectiontimes (F -- 2[F, G])), 8 + xpart(cr intersectiontimes (F -- 2[F, D])), byyellow);
draw byLine(G, A, byblue, SOLID_LINE, REGULAR_WIDTH);
draw byLine(G, E, byred, SOLID_LINE, REGULAR_WIDTH);
draw byLine(B, D, byblack, DASHED_LINE, REGULAR_WIDTH);
draw byLine(B, A, byblack, SOLID_LINE, THIN_WIDTH);
draw byLine(B, E, byblack, SOLID_LINE, THIN_WIDTH);
draw byLine(D, A, byblack, SOLID_LINE, THIN_WIDTH);
draw byLine(D, E, byblack, SOLID_LINE, THIN_WIDTH);
draw byCircleR(F, r, byblue, 0, 0, 0);
draw byLabelsOnCircle(B, D, E, A, G)(F);
}
\startsubproposition{Figure II.}
\drawCurrentPictureInMargin
Seja o segmento um semicírculo, ou menos que um semicírculo, trace \drawUnitLine{GA} o diâmetro, trace também \drawUnitLine{GE}.

\begin{center}
$\drawAngle{BAG} = \drawAngle{BEG}$ and $\drawAngle{GAD} = \drawAngle{GED}$ (case I.)

$\therefore \drawAngle{BAG,GAD} = \drawAngle{BEG,GED}$.
\end{center}

\qed

\starttheorem{Prop. XXII. Teor.}\label{prop:III.XXII}
\defineNewPicture{
pair A, B, C, D, E;
numeric r;
r := 7/4u;
E := (0, 0);
A := (dir(80)*r) shifted E;
B := (dir(10)*r) shifted E;
C := (dir(-100)*r) shifted E;
D := (dir(150)*r) shifted E;
byAngleDefine(D, A, C, byred, SOLID_SECTOR);
byAngleDefine(C, A, B, byblue, SOLID_SECTOR);
byAngleDefine(A, B, D, byyellow, SOLID_SECTOR);
byAngleDefine(D, B, C, byred, SOLID_SECTOR);
byAngleDefine(B, C, A, byblack, SOLID_SECTOR);
byAngleDefine(A, C, D, byyellow, SOLID_SECTOR);
byAngleDefine(C, D, B, byblue, SOLID_SECTOR);
byAngleDefine(B, D, A, byblack, SOLID_SECTOR);
draw byNamedAngleResized();
draw byLine(A, B, byblack, SOLID_LINE, THIN_WIDTH);
draw byLine(B, C, byblack, SOLID_LINE, THIN_WIDTH);
draw byLine(C, D, byblack, SOLID_LINE, THIN_WIDTH);
draw byLine(D, A, byblack, SOLID_LINE, THIN_WIDTH);
draw byLine(A, C, byred, SOLID_LINE, REGULAR_WIDTH);
draw byLine(B, D, byblack, SOLID_LINE, REGULAR_WIDTH);
draw byCircleR(E, r, byred, 0, 0, 1/2);
draw byLabelsOnCircle(A, B, C, D)(E);
}
\drawCurrentPictureInMargin
\problem{O}{s}{ângulos opostos \drawAngle{DAC,CAB} e \drawAngle{BCA,ACD}, \drawAngle{CDB,BDA} e \drawAngle{ABD,DBC} de qualquer figura quadrilátera inscrito em um círculo, são juntos iguais a dois ângulos retos.}

\begin{center}
Trace \drawUnitLine{AC} e \drawUnitLine{BD} as diagonais;\\
e porque os ângulos no mesmo segmento são iguais\\
$\drawAngle{CDB} = \drawAngle{CAB}$,\\
and $\drawAngle{DAC} = \drawAngle{DBC}$;

add \drawAngle{BCA,ACD} to both.\\
$\drawAngle{DAC,CAB} + \drawAngle{BCA,ACD} = \drawAngle{BCA,ACD} + \drawAngle{CDB} + \drawAngle{DBC} = \mbox{ two right angles }$ \byref{prop:I.XXXII}.

In like manner it may be shown that,\\
$\drawAngle{CDB,BDA} + \drawAngle{ABD,DBC} = \drawTwoRightAngles$.
\end{center}

\qed

\starttheorem{Prop. XXIII. Teor.}\label{prop:III.XXIII}
\defineNewPicture{
pair A, B, C, D, M, N;
numeric r[], t[];
path cr[];
r1 := 14/6u;
r2 := 15/6u;
M := (0, 0);
N := (0, 1/3r1);
cr1 := (fullcircle scaled 2r1) shifted M;
cr2 := (fullcircle scaled 2r2) shifted N;
t1 := xpart(cr1 intersectiontimes (subpath (-2, 2) of cr2));
t2 := xpart(cr1 intersectiontimes (subpath (2, 6) of cr2));
t3 := xpart(cr2 intersectiontimes (subpath (-2, 2) of cr1));
t4 := xpart(cr2 intersectiontimes (subpath (2, 6) of cr1));
A := point t1 of cr1;
B := point t2 of cr1;
C := point 3/2 of cr1;
D := cr2 intersectionpoint (1/2[A,C]--2[A,C]);
byAngleDefine(A, C, B, byyellow, SOLID_SECTOR);
byAngleDefine(A, D, B, byblue, SOLID_SECTOR);
draw byNamedAngleResized();
draw byLineFull(A, D, byred, 0, 0)(B, D, 1, 0, 0);
draw byLineFull(B, C, byblue, 0, 0)(A, C, 1, 0, 0);
draw byLineFull(B, D, byyellow, 0, 0)(A, D, 1, 0, 0);
draw byLineFull(A, B, byblack, 0, 0)(A, B, 0, 0, 1);
draw byArc.M(M, A, B, r1, byred, 0, 0, 0, 1);
draw byArc.N(N, A, B, r2, byblue, 0, 0, 0, 1);
draw byLabelsOnPolygon(A, B, noPoint)(ALL_LABELS, 0);
draw byLabelsOnPolygon(B, D, A)(OMIT_FIRST_LABEL+OMIT_LAST_LABEL, 0);
draw byLabelPoint(C, (angle(C-M)-90+angle(D-A))/2, 1);
}
\drawCurrentPictureInMargin
\problem{S}{obre}{a mesma linha reta e do mesmo lado dela, não podem ser construídos dois segmentos semelhantes de círculos que não coincidam.}

\begin{center}
Pois se for possível, deixemos dois segmentos semelhantes\\
\drawFromCurrentPicture[bottom]{
draw byNamedLine(AB);
draw byNamedArcExact(N);
draw byLabelsOnPolygon(A, B, noPoint)(ALL_LABELS, 0);
}
and
\drawFromCurrentPicture[bottom]{
draw byNamedLine(AB);
draw byNamedArcExact(M);
draw byLabelsOnPolygon(A, B, noPoint)(ALL_LABELS, 0);
}
be constructed;

trace qualquer linha direita \drawUnitLine{AD} cortando ambos os segmentos,

trace \drawUnitLine{BC} e \drawUnitLine{BD}.

Como os segmentos são semelhantes,\\
$\drawAngle{C} = \drawAngle{D}$ \byref{def:III.XI}, % improvement: definition 10 in the original seems to be irrelevant

but $\drawAngle{C} > \drawAngle{D}$ \byref{prop:I.XVI}\\
which is absurd;

portanto, nenhum ponto em nenhum dos segmentos cai sem o outro e, portanto, os segmentos coincidem.
\end{center}

\qed

\starttheorem{Prop. XXIV. Teor.}\label{prop:III.XXIV}
\defineNewPicture{
pair A, B, C, D, M, N, d;
numeric r, t[];
path cr[];
r := 5/2u;
t1 := 2-1;
t2 := 2+1;
d := (0, -3/2u);
M := (0, 0);
N := M shifted d;
cr1 := ((subpath (t1, t2) of fullcircle) scaled 2r) shifted M;
cr2 := ((subpath (t1, t2) of fullcircle) scaled 2r) shifted N;
A := point length(cr1) of cr1;
B := point 0 of cr1;
C := point length(cr2) of cr2;
D := point 0 of cr2;
draw byFilledCircleSegment(M, r, t1, t2, byred);
draw byLineFull(A, B, byblack, 0, 0)(A, B, 0, 0, -1);
draw byArc.M(M, B, A, r, byblue, 0, 0, 1/2, 1);
draw byFilledCircleSegment(N, r, t1, t2, byyellow);
draw byLineFull(C, D, byblue, 0, 0)(C, D, 0, 0, -1);
draw byArc.N(N, D, C, r, byred, 0, 0, 1/2, 1);
draw byLabelsOnPolygon(B, A, noPoint)(ALL_LABELS, 0);
draw byLabelsOnPolygon(D, C, noPoint)(ALL_LABELS, 0);
}
\drawCurrentPictureInMargin
\problem{S}{imilar}{segments
\drawFromCurrentPicture[bottom]{
draw byNamedFilledCircleSegment(M);
draw byNamedLine(AB);
draw byNamedArcExact(M);
draw byLabelsOnPolygon(B, A, noPoint)(ALL_LABELS, 0);
}
and
\drawFromCurrentPicture[bottom]{
draw byNamedFilledCircleSegment(N);
draw byNamedLine(CD);
draw byNamedArcExact(N);
draw byLabelsOnPolygon(D, C, noPoint)(ALL_LABELS, 0);
}, de círculos sobre linhas retas iguais (\drawUnitLine{AB} e \drawUnitLine{CD}) são iguais entre si.}

\begin{center}
For, if
\drawFromCurrentPicture[bottom]{
draw byNamedFilledCircleSegment(N);
draw byLabelsOnPolygon(D, C, noPoint)(ALL_LABELS, 0);
}
be so applied to
\drawFromCurrentPicture[bottom]{
draw byNamedFilledCircleSegment(M);
draw byLabelsOnPolygon(B, A, noPoint)(ALL_LABELS, 0);
},\\
that \drawUnitLine{CD} may fall on \drawUnitLine{AB},\\
as extremidades do \drawUnitLine{CD} podem estar nas extremidades do \drawUnitLine{AB}\\
e \drawArc{M} no mesmo lado que \drawArc{N};\\
$\because \drawUnitLine{CD} = \drawUnitLine{AB}$,\\ %because
\drawUnitLine{CD} must wholly coincide with \drawUnitLine{AB};
\end{center}

\noindent e os segmentos semelhantes estando então na mesma reta e no mesmo lado dela, também devem coincidir com \byref{prop:III.XXIII} e, portanto, são iguais.

\qed

\startproblem{Prop. XXV. prob.}\label{prop:III.XXV}
\defineNewPicture{
pair A, B, C, D, E, F, O;
numeric r;
r := 7/4u;
O := (0, 0);
A := (dir(-20)*r) shifted O;
B := (dir(85)*r) shifted O;
C := (dir(180)*r) shifted O;
D := 1/2[A, B];
E := 1/2[B, C];
F = whatever[D, D shifted ((A-B) rotated 90)] = whatever[E, E shifted ((B-C) rotated 90)];
byLineDefine(D, F, byred, SOLID_LINE, REGULAR_WIDTH);
byLineDefine(E, F, byyellow, SOLID_LINE, REGULAR_WIDTH);
draw byNamedLineSeq(0)(DF,EF);
draw byLine(A, B, byblack, SOLID_LINE, REGULAR_WIDTH);
draw byLine(B, C, byblue, SOLID_LINE, REGULAR_WIDTH);
draw byArcBE(O, -6/5, 5, r, byblue, 0, 0, 1/2, 0);
draw byLabelsOnCircle(A, B, C)(O);
draw byLabelLineEnd(D, F, 0);
draw byLabelLineEnd(E, F, 0);
draw byLabelsOnPolygon(D, F, E)(OMIT_FIRST_LABEL+OMIT_LAST_LABEL, 0);
}
\drawCurrentPictureInMargin
\problem{U}{m}{segmento de um círculo sendo dado, para descrever o círculo do qual ele é o segmento.}

\begin{center}
De qualquer ponto do segmento trace \drawUnitLine{BC} e \drawUnitLine{AB},\\
divida-os ao meio e dos pontos de bissecção\\
draw $\drawUnitLine{EF} \perp \drawUnitLine{BC}$\\
and $\drawUnitLine{DF} \perp \drawUnitLine{AB}$\\
onde eles se encontram é o centro do círculo.
\end{center}

Como \drawUnitLine{BC} terminado no círculo é dividido ao meio perpendicularmente por \drawUnitLine{EF}, ele passa pelo centro \byref{prop:III.I}, da mesma forma que \drawUnitLine{DF} passa pelo centro, portanto o centro está na interseção dessas perpendiculares.

\qed

\starttheorem{Prop. XXVI. Teor.}\label{prop:III.XXVI}
\defineNewPicture{
pair A, B, C, D, E, F, G, H, d;
numeric r, t[];
path cr[];
r := 7/4u;
t1 := 5;
t2 := 7;
G := (0, 0);
cr1 := (fullcircle scaled 2r) shifted G;
A := point 3/2 of cr1;
B := point t1 of cr1;
C := point t2 of cr1;
d := (0, -5/2r);
H := G shifted d;
cr2 := (fullcircle scaled 2r) shifted H;
D := point 5/2 of cr2;
E := point t1 of cr2;
F := point t2 of cr2;
byAngleDefine(B, A, C, byred, SOLID_SECTOR);
byAngleDefine(B, G, C, byyellow, SOLID_SECTOR);
draw byNamedAngleResized(BAC, BGC);
draw byFilledCircleSegment(G, r, t1, t2, byyellow);
draw byLine(B, C, byblack, SOLID_LINE, REGULAR_WIDTH);
draw byLine(A, B, byblack, SOLID_LINE, THIN_WIDTH);
draw byLine(A, C, byblack, SOLID_LINE, THIN_WIDTH);
byLineDefine(B, G, byblue, SOLID_LINE, REGULAR_WIDTH);
byLineDefine(C, G, byred, SOLID_LINE, REGULAR_WIDTH);
draw byNamedLineSeq(0)(BG,CG);
draw byArc.G(G, C, B, r, byblue, 0, 0, 0, 0);
draw byArc.Gb(G, B, C, r, byblack, 0, 0, 0, 0);
byCircleDefineR(G, r, byblue, 0, 0, 0);
byAngleDefine(E, D, F, byblue, SOLID_SECTOR);
byAngleDefine(E, H, F, byblack, SOLID_SECTOR);
draw byNamedAngleResized(EDF, EHF);
draw byFilledCircleSegment(H, r, t1, t2, byyellow);
draw byLine(E, F, byblack, DASHED_LINE, REGULAR_WIDTH);
draw byLine(D, E, byblack, SOLID_LINE, THIN_WIDTH);
draw byLine(D, F, byblack, SOLID_LINE, THIN_WIDTH);
byLineDefine(E, H, byblue, DASHED_LINE, REGULAR_WIDTH);
byLineDefine(F, H, byred, DASHED_LINE, REGULAR_WIDTH);
draw byNamedLineSeq(0)(EH,FH);
draw byArc.H(H, F, E, r, byred, 0, 0, 0, 0);
draw byArc.Hb(H, E, F, r, byblack, 0, 0, 0, 0);
byCircleDefineR(H, r, byred, 0, 0, 0);
draw byLabelsOnCircle(A, B, C)(G);
draw byLabelsOnCircle(D, E, F)(H);
draw byLabelsOnPolygon(B, G, C)(OMIT_FIRST_LABEL+OMIT_LAST_LABEL, 0);
draw byLabelsOnPolygon(E, H, F)(OMIT_FIRST_LABEL+OMIT_LAST_LABEL, 0);
}
\drawCurrentPictureInMargin
\problem[4]{C}{írculos}{desiguais \drawCircle[middle][1/5]{G} e \drawCircle[middle][1/5]{H}, os arcos \drawArc{Gb}, \drawArc{Hb} nos quais estão ângulos iguais, seja no centro da circunferência, são iguais.}

\begin{center}
Pois, seja $\drawAngle{G} = \drawAngle{H}$ no centro,\\
trace \drawUnitLine{BC} e \drawUnitLine{EF}.\\
Then since $\circleG = \circleH$,\\
\drawLine[bottom]{BG,CG,BC}
and
\drawLine[bottom]{EH,FH,EF}
have\\
$\drawUnitLine{BG} = \drawUnitLine{CG} = \drawUnitLine{EH} = \drawUnitLine{FH}$,\\
and $\drawAngle{G} = \drawAngle{H}$,\\
$\therefore \drawUnitLine{BC} = \drawUnitLine{EF}$ \byref{prop:I.IV}.\\
But $\drawAngle{A} = \drawAngle{D}$ \byref{prop:III.XX};\\
$\therefore$
\drawFromCurrentPicture{
startTempScale(1/5);
draw byNamedArc(G);
draw byNamedLine(BC);
draw byLabelsOnCircle(B, C)(G);
stopTempScale;
}
and
\drawFromCurrentPicture{
startTempScale(1/5);
draw byNamedArc(H);
draw byNamedLine(EF);
draw byLabelsOnCircle(E, F)(H);
stopTempScale;
}
são semelhantes \byref{def:III.XI};\\ % melhoria: a definição 10 no original parece ser irrelevante
eles também são iguais \byref{prop:III.XXIV}
\end{center}

Se, portanto, os segmentos iguais forem retirados dos círculos iguais, os segmentos restantes serão iguais;

\begin{center}
hence $
\drawFromCurrentPicture{
draw byNamedFilledCircleSegment(G);
draw byNamedArcLabel(Gb);
}
=
\drawFromCurrentPicture{
draw byNamedFilledCircleSegment(H);
draw byNamedArcLabel(Hb);
}
$ \byref{ax:I.III};\\
and $\therefore \drawArc{Gb} = \drawArc{Hb}$.
\end{center}

Mas se os ângulos iguais dados estiverem na circunferência, é evidente que os ângulos no centro, sendo o dobro daqueles na circunferência, também são iguais e, portanto, os arcos sobre os quais eles se apoiam são iguais.

\qed

\starttheorem{Prop. XXVII. Teor.}\label{prop:III.XXVII}
\defineNewPicture[1/2]{
pair A, B, C, D, E, F, G, H, K, d;
numeric r, t[];
path cr[];
r := 7/4u;
t1 := 5;
t2 := 7;
t3 := 13/2;
G := (0, 0);
cr1 := (fullcircle scaled 2r) shifted G;
A := point 3/2 of cr1;
B := point t1 of cr1;
C := point t2 of cr1;
K := point t3 of cr1;
d := (0, -5/2r);
H := G shifted d;
cr2 := (fullcircle scaled 2r) shifted H;
D := point 3/2 of cr2;
E := point t1 of cr2;
F := point t3 of cr2;
byAngleDefine(B, A, K, byyellow, SOLID_SECTOR);
byAngleDefine(B, G, K, byyellow, SOLID_SECTOR);
byAngleDefine(K, A, C, byblue, SOLID_SECTOR);
byAngleDefine(K, G, C, byblue, SOLID_SECTOR);
draw byNamedAngleResized(BAK, BGK, KAC, KGC);
draw byLine(A, K, byblack, SOLID_LINE, THIN_WIDTH);
draw byLine(A, B, byblack, SOLID_LINE, THIN_WIDTH);
draw byLine(A, C, byblack, SOLID_LINE, THIN_WIDTH);
byLineDefine(G, B, byblack, SOLID_LINE, THIN_WIDTH);
byLineDefine(G, C, byblack, SOLID_LINE, THIN_WIDTH);
draw byNamedLineSeq(0)(GB,GC);
draw byLine(G, K, byblack, SOLID_LINE, THIN_WIDTH);
byArcDefine.G(G, C, B, r, byred, 0, 0, 0, 0);
byArcDefine.GbI(G, B, K, r, byblack, 0, 0, 0, 0);
byArcDefine.GbII(G, K, C, r, byred, 1, 0, 0, 0);
draw byNamedArcSeq(0)(G, GbI, GbII);
byCircleDefineR(G, r, byred, 0, 0, 0);
byAngleDefine(E, D, F, byred, SOLID_SECTOR);
byAngleDefine(E, H, F, byred, SOLID_SECTOR);
draw byNamedAngleResized(EDF, EHF);
draw byLine(D, F, byblack, SOLID_LINE, THIN_WIDTH);
draw byLine(D, E, byblack, SOLID_LINE, THIN_WIDTH);
byLineDefine(H, E, byblack, SOLID_LINE, THIN_WIDTH);
byLineDefine(H, F, byblack, SOLID_LINE, THIN_WIDTH);
draw byNamedLineSeq(0)(HE,HF);
byArcDefine.H(H, F, E, r, byblue, 0, 0, 0, 0);
byArcDefine.Hb(H, E, F, r, byblack, 1, 0, 0, 0);
draw byNamedArcSeq(0)(H, Hb);
byCircleDefineR(H, r, byblue, 0, 0, 0);
draw byLabelsOnCircle(A, B, C, K)(G);
draw byLabelsOnCircle(D, E, F)(H);
draw byLabelsOnPolygon(B, G, C)(OMIT_FIRST_LABEL+OMIT_LAST_LABEL, 0);
draw byLabelsOnPolygon(E, H, F)(OMIT_FIRST_LABEL+OMIT_LAST_LABEL, 0);
}
\drawCurrentPictureInMargin
\problem[3]{C}{írculos}{desiguais \drawCircle[middle][1/3]{G} e \drawCircle[middle][1/3]{H} os ângulos \drawAngle{BAK} e \drawAngle{D} que estão sobre arcos iguais são iguais, estejam eles nos centros ou nas circunferências.}

\begin{center}
Pois se for possível, deixe um deles\\
\drawAngle{D} be greater than the other \drawAngle{BAK}\\
and make $\drawAngle{BAK,KAC} = \drawAngle{D}$

$\therefore \drawArc{GbI,GbII} = \drawArc{Hb}$ \byref{prop:III.XXVI}

but $\drawArc{GbI} = \drawArc{Hb}$ \byref{\hypref}\\
$\therefore \drawArc{GbI} = \drawArc{GbI,GbII}$ uma parte igual ao todo, o que é um absurdo;

$\therefore$ nenhum ângulo é maior que o outro,

e $\therefore$ eles são iguais.
\end{center}

\qed

\starttheorem{Prop. XXVIII. Teor.}\label{prop:III.XXVIII}
\defineNewPicture[1/2]{
pair A, B, D, E, K, L, d;
numeric r, t[];
path cr[];
r := 7/4u;
t1 := 5;
t2 := 7;
K := (0, 0);
cr1 := (fullcircle scaled 2r) shifted K;
A := point t1 of cr1;
B := point t2 of cr1;
d := (0, -5/2r);
L := K shifted d;
cr2 := (fullcircle scaled 2r) shifted L;
D := point t1 of cr2;
E := point t2 of cr2;
byAngleDefine(A, K, B, byred, SOLID_SECTOR);
draw byNamedAngleResized(AKB);
draw byLine(A, B, byred, SOLID_LINE, REGULAR_WIDTH);
byLineDefine(A, K, byblack, SOLID_LINE, REGULAR_WIDTH);
byLineDefine(B, K, byblue, SOLID_LINE, REGULAR_WIDTH);
draw byNamedLineSeq(0)(AK,BK);
draw byArc.K(K, B, A, r, byyellow, 0, 0, 0, 0);
draw byArc.Kb(K, A, B, r, byblue, 0, 0, 0, 0);
byCircleDefineR(K, r, byyellow, 0, 0, 0);
byAngleDefine(D, L, E, byyellow, SOLID_SECTOR);
draw byNamedAngleResized(DLE);
draw byLine(D, E, byred, DASHED_LINE, REGULAR_WIDTH);
byLineDefine(D, L, byblack, DASHED_LINE, REGULAR_WIDTH);
byLineDefine(E, L, byblue, DASHED_LINE, REGULAR_WIDTH);
draw byNamedLineSeq(0)(DL,EL);
draw byArc.L(L, E, D, r, byblack, 0, 0, 0, 0);
draw byArc.Lb(L, D, E, r, byred, 0, 0, 0, 0);
byCircleDefineR(L, r, byblack, 0, 0, 0);
draw byLabelsOnCircle(A, B)(K);
draw byLabelsOnCircle(D, E)(L);
draw byLabelsOnPolygon(A, K, B)(OMIT_FIRST_LABEL+OMIT_LAST_LABEL, 0);
draw byLabelsOnPolygon(D, L, E)(OMIT_FIRST_LABEL+OMIT_LAST_LABEL, 0);
}
\drawCurrentPictureInMargin
\problem{C}{írculos}{desiguais \drawCircle[middle][1/3]{K} e \drawCircle[middle][1/3]{L}, cordas iguais \drawUnitLine{AB}, \drawUnitLine{DE} cortam arcos iguais.}

\begin{center}
Dos centros dos círculos iguais,\\
trace \drawUnitLine{AK}, \drawUnitLine{BK} e \drawUnitLine{DL}, \drawUnitLine{EL};

and $\because \circleK = \circleL$\\ % because
$\drawUnitLine{AK}, \drawUnitLine{BK} = \drawUnitLine{DL}, \drawUnitLine{EL}$\\
also $\drawUnitLine{AB} = \drawUnitLine{DE}$ \byref{\hypref}

$\therefore \drawAngle{K} = \drawAngle{L}$

$\therefore \drawArc{Kb} = \drawArc{Lb}$ \byref{prop:III.XXVI}

and $\therefore \drawArc[bottom][1/3]{K} = \drawArc[bottom][1/3]{L}$ \byref{ax:I.III}.
\end{center}

\qed

\starttheorem{Prop. XXIX. Teor.}\label{prop:III.XXIX}
\defineNewPicture[1/2]{
pair A, B, D, E, K, L, d;
numeric r, t[];
path cr[];
r := 7/4u;
t1 := 5;
t2 := 7;
K := (0, 0);
cr1 := (fullcircle scaled 2r) shifted K;
A := point t1 of cr1;
B := point t2 of cr1;
d := (0, -5/2r);
L := K shifted d;
cr2 := (fullcircle scaled 2r) shifted L;
D := point t1 of cr2;
E := point t2 of cr2;
byAngleDefine(A, K, B, byred, SOLID_SECTOR);
draw byNamedAngleResized(AKB);
draw byLine(A, B, byred, SOLID_LINE, REGULAR_WIDTH);
byLineDefine(A, K, byblack, SOLID_LINE, REGULAR_WIDTH);
byLineDefine(B, K, byblue, SOLID_LINE, REGULAR_WIDTH);
draw byNamedLineSeq(0)(AK,BK);
draw byArc.K(K, B, A, r, byyellow, 0, 0, 0, 0);
draw byArc.Kb(K, A, B, r, byblue, 0, 0, 0, 0);
byCircleDefineR(K, r, byyellow, 0, 0, 0);
byAngleDefine(D, L, E, byyellow, SOLID_SECTOR);
draw byNamedAngleResized(DLE);
draw byLine(D, E, byred, DASHED_LINE, REGULAR_WIDTH);
byLineDefine(D, L, byblack, DASHED_LINE, REGULAR_WIDTH);
byLineDefine(E, L, byblue, DASHED_LINE, REGULAR_WIDTH);
draw byNamedLineSeq(0)(DL,EL);
draw byArc.L(L, E, D, r, byblack, 0, 0, 0, 0);
draw byArc.Lb(L, D, E, r, byred, 0, 0, 0, 0);
byCircleDefineR(L, r, byblack, 0, 0, 0);
draw byLabelsOnCircle(A, B)(K);
draw byLabelsOnCircle(D, E)(L);
draw byLabelsOnPolygon(A, K, B)(OMIT_FIRST_LABEL+OMIT_LAST_LABEL, 0);
draw byLabelsOnPolygon(D, L, E)(OMIT_FIRST_LABEL+OMIT_LAST_LABEL, 0);
}
\drawCurrentPictureInMargin
\problem{C}{írculos}{desiguais \drawCircle[middle][1/3]{K} e \drawCircle[middle][1/3]{L} as cordas \drawUnitLine{AB} e \drawUnitLine{DE} que subtendem arcos iguais são iguais.}

\begin{center}
Se os arcos iguais forem semicírculos, a proposição é evidente. Mas se não,

seja \drawUnitLine{AK}, \drawUnitLine{BK} e \drawUnitLine{DL}, \drawUnitLine{EL} serem atraídos para os centros;

$\because \drawArc{Kb} = \drawArc{Lb}$ \byref{\hypref}\\ % because
and $\drawAngle{K} = \drawAngle{L}$ \byref{prop:III.XXVII};

but $\drawUnitLine{AK} \mbox{ e } \drawUnitLine{BK} = \drawUnitLine{DL} \mbox{ e } \drawUnitLine{EL}$

$\therefore \drawUnitLine{AB} = \drawUnitLine{DE}$ \byref{prop:I.IV};

mas estes são os acordes que subentendem os arcos iguais.
\end{center}

\qed

\startproblem{Prop. XXX. prob.}\label{prop:III.XXX}
\defineNewPicture{
pair A, B, C, D, E;
numeric r, t[];
path cr;
r := 9/4u;
t1 := -1/2;
t2 := 4 + 1/2;
t3 := 2;
E := (0, 0);
cr := (fullcircle scaled 2r) shifted E;
A := point t2 of cr;
B := point t1 of cr;
D := point t3 of cr;
C := 1/2[A, B];
byAngleDefine(A, C, D, byblue, SOLID_SECTOR);
byAngleDefine(D, C, B, byred, SOLID_SECTOR);
draw byNamedAngleResized();
draw byLine(D, C, byyellow, SOLID_LINE, REGULAR_WIDTH);
byLineDefine(D, A, byblue, SOLID_LINE, REGULAR_WIDTH);
byLineDefine(D, B, byblue, DASHED_LINE, REGULAR_WIDTH);
byLineDefine(A, C, byblack, SOLID_LINE, REGULAR_WIDTH);
byLineDefine(C, B, byblack, DASHED_LINE, REGULAR_WIDTH);
draw byNamedLineSeq(0)(AC, CB, DB, DA);
byArcDefine.Er(E, B, D, r, byred, 1, 0, 0, 0);
byArcDefine.El(E, D, A, r, byred, 0, 0, 0, 0);
draw byNamedArcSeq(0.5)(Er,El);
draw byLabelsOnCircle(A, B, D)(Er);
draw byLabelLineEnd(C, D, 0);
}
\drawCurrentPictureInMargin
\problem{P}{ara}{bisseccionar um determinado arco \drawArc[middle][1/5]{El,Er}.}

\begin{center}
Draw \drawUnitLine{AC,CB};\\
make $\drawUnitLine{AC} = \drawUnitLine{CB}$,\\
trace $\drawUnitLine{DC} \perp \drawUnitLine{AC,CB}$ e ele divide o arco ao meio.

Trace \drawUnitLine{DA} e \drawUnitLine{DB}.

$\drawUnitLine{AC} = \drawUnitLine{CB}$ \byref{\constref}\\
\drawUnitLine{DC} é comum,\\ % comum a quê? provavelmente triângulos ACD e CBD
and $\drawAngle{ACD} = \drawAngle{DCB}$ \byref{\constref}

$\therefore \drawUnitLine{DA} = \drawUnitLine{DB}$ \byref{prop:I.IV}\\
$
\drawFromCurrentPicture{
startGlobalRotation(-arcAngle.El);
startAutoLabeling;
draw byNamedArc(El);
stopAutoLabeling;
stopGlobalRotation;
}
=
\drawFromCurrentPicture{
startGlobalRotation(-arcAngle.Er);
startAutoLabeling;
draw byNamedArc(Er);
stopAutoLabeling;
stopGlobalRotation;
}
$ \byref{prop:III.XXVIII},\\
e, portanto, o arco dado é dividido ao meio.
\end{center}

\qed

\starttheorem{Prop. XXXI. Teor.}\label{prop:III.XXXI}
\problem{E}{m}{um círculo, o ângulo de um semicírculo é reto, o ângulo de um segmento maior que um semicírculo é agudo e o ângulo de um segmento menor que um semicírculo é obtuso.}

\defineNewPicture{
pair A, B, C, E;
numeric r;
r := 3/2u;
E := (0, 0);
A := (dir(110)*r) shifted E;
B := (dir(180)*r) shifted E;
C := (dir(0)*r) shifted E;
byAngleDefine(A, B, C, byred, SOLID_SECTOR);
byAngleDefine(B, C, A, byblue, SOLID_SECTOR);
byAngleDefine(E, A, B, byyellow, SOLID_SECTOR);
byAngleDefine(C, A, E, byblack, SOLID_SECTOR);
draw byNamedAngleResized();
draw byLine(A, E, byred, SOLID_LINE, REGULAR_WIDTH);
draw byLine(A, B, byblack, SOLID_LINE, THIN_WIDTH);
draw byLine(A, C, byblack, SOLID_LINE, THIN_WIDTH);
draw byLine(B, E, byblue, SOLID_LINE, REGULAR_WIDTH);
draw byLine(E, C, byblack, SOLID_LINE, REGULAR_WIDTH);
draw byCircleR(E, r, byblack, 0, 0, 0);
draw byLabelsOnCircle(A, B, C)(E);
draw byLabelsOnPolygon(C, E, B)(OMIT_FIRST_LABEL+OMIT_LAST_LABEL, 0);
}
\startsubproposition{Figure I.}
\drawCurrentPictureInMargin
\begin{center}
O ângulo \drawAngle{EAB,CAE} em um semicírculo é um ângulo reto.

Trace \drawUnitLine{AE} e \drawUnitLine{BE,EC}\\
$\drawAngle{B}=\drawAngle{EAB}$ and $\drawAngle{C} = \drawAngle{CAE}$ \byref{prop:I.V}\\
$\drawAngle{C} + \drawAngle{B} = \drawAngle{EAB,CAE} = \mbox{ a metade de } \drawTwoRightAngles= \drawRightAngle$ \byref{prop:I.XXXII}
\end{center}

\defineNewPicture{
pair A, B, C, E, D;
numeric r;
r := 3/2u;
E := (0, 0);
A := (dir(110)*r) shifted E;
B := (dir(180 + 30)*r) shifted E;
C := (dir(0 + 30)*r) shifted E;
D := (dir(0 - 30)*r) shifted E;
byAngleDefine(B, A, D, byblue, SOLID_SECTOR);
byAngleDefine(D, A, C, byred, SOLID_SECTOR);
draw byNamedAngleResized();
draw byLine(A, D, byblack, SOLID_LINE, THIN_WIDTH);
draw byLine(A, B, byblack, SOLID_LINE, THIN_WIDTH);
draw byLine(B, D, byblack, SOLID_LINE, THIN_WIDTH);
draw byLine(A, C, byblue, SOLID_LINE, REGULAR_WIDTH);
draw byLine(B, C, byred, SOLID_LINE, REGULAR_WIDTH);
draw byCircleR(E, r, byblue, 0, 0, 0);
draw byLabelsOnCircle(A, B, C, D)(E);
}
\startsubproposition{Figure II.}
\drawCurrentPictureInMargin
\begin{center}
O ângulo \drawAngle{BAD} em um segmento maior que um semicírculo é agudo

Trace \drawUnitLine{BC} o diâmetro e \drawUnitLine{AC}\\
$\therefore \drawAngle{BAD,DAC} = \mbox{ a right angle}$, $\therefore$ \drawAngle{BAD} is acute.
\end{center}

\defineNewPicture{
pair A, B, C, E, D;
numeric r;
r := 3/2u;
E := (0, 0);
A := (dir(140)*r) shifted E;
B := (dir(-110)*r) shifted E;
C := (dir(5)*r) shifted E;
D := (dir(80)*r) shifted E;
byAngleDefine(A, D, C, byred, SOLID_SECTOR);
byAngleDefine(C, B, A, byyellow, SOLID_SECTOR);
draw byNamedAngleResized();
draw byLine(A, B, byred, SOLID_LINE, REGULAR_WIDTH);
draw byLine(B, C, byblue, SOLID_LINE, REGULAR_WIDTH);
draw byLine(C, D, byblack, SOLID_LINE, THIN_WIDTH);
draw byLine(D, A, byblack, SOLID_LINE, THIN_WIDTH);
draw byCircleR(E, r, byblue, 0, 0, 0);
draw byLabelsOnCircle(A, B, C, D)(E);
}
\startsubproposition{Figure III.}
\drawCurrentPictureInMargin
\begin{center}
O ângulo \drawAngle{D} em um segmento menor que um semicírculo é obtuso.

Pegue qualquer ponto na circunferência oposta, para o qual trace \drawUnitLine{BC} e \drawUnitLine{AB}.\\
$\because \drawAngle{B} + \drawAngle{D} = \drawTwoRightAngles$ \byref{prop:III.XXII}\\ % Because
but $\drawAngle{B} < \drawRightAngle$ (figure II.), $\therefore$ \drawAngle{D} is obtuse.
\end{center}

\qed

\starttheorem{Prop. XXXII. Teor.}\label{prop:III.XXXII}
\defineNewPicture{
pair A, B, C, D, E, F, O;
numeric r;
r := 9/4u;
O := (0, 0);
A := (0, r) shifted O;
B := (0, -r) shifted O;
C := (dir(-20)*r) shifted O;
D := (dir(30)*r) shifted O;
E := (-r, -r) shifted O;
F := (r, -r) shifted O;
byAngleDefine(A, B, E, byblack, ARC_SECTOR);
byAngleDefine(D, B, A, byblue, SOLID_SECTOR);
byAngleDefine(F, B, D, byyellow, SOLID_SECTOR);
byAngleDefine(B, A, D, byyellow, SOLID_SECTOR);
byAngleDefine(A, D, B, byblack, SOLID_SECTOR);
byAngleDefine(C, D, B, byblack, ARC_SECTOR);
byAngleDefine(D, C, B, byred, SOLID_SECTOR);
draw byNamedAngleResized();
draw byLine(A, D, byblack, SOLID_LINE, THIN_WIDTH);
draw byLine(D, C, byblack, SOLID_LINE, THIN_WIDTH);
draw byLine(C, B, byblack, SOLID_LINE, THIN_WIDTH);
draw byLine(D, B, byred, SOLID_LINE, REGULAR_WIDTH);
draw byLineFull(E, F, byblue, 0, 0)(E, F, 0, 0, 1);
draw byLine(A, B, byblack, SOLID_LINE, REGULAR_WIDTH);
draw byCircleR(O, r, byred, 0, 0, 0);
draw byLabelsOnCircle(A, B, C, D)(O);
draw byLabelsOnPolygon(E, F, noPoint)(ALL_LABELS, 0);
}
\drawCurrentPictureInMargin
\problem{S}{e}{uma linha reta \drawUnitLine{EF} for uma tangente a um círculo, e do ponto de contato for traçada uma linha reta \drawUnitLine{DB} cortando o círculo, o ângulo \drawAngle{FBD} feito por esta linha com a tangente é igual ao ângulo \drawAngle{A} no segmento alternativo do círculo.}

Se a corda passar pelo centro, é evidente que os ângulos são iguais, pois cada um deles é um ângulo reto. \byref{prop:III.XVI,prop:III.XXXI}

Mas se não, desenhe $\drawUnitLine{AB} \perp \drawUnitLine{EF}$ a partir do ponto de contato, ele deve passar pelo centro do círculo, \byref{prop:III.XIX}

\begin{center}
$\therefore \drawAngle{ADB} = \drawAngle{ABE} $ \byref{prop:III.XXXI}

$\drawAngle{A} + \drawAngle{DBA} = \drawAngle{ABE} = \drawAngle{FBA}$ \byref{prop:I.XXXII}

$\therefore \drawAngle{A} = \drawAngle{FBD}$ \byref{ax:I.III}. % de melhoria: parece ser o axioma III aqui e ali embaixo

Again $\drawAngle{B} = \drawTwoRightAngles = \drawAngle{A} + \drawAngle{C}$ \byref{prop:III.XXII}

$\therefore \drawAngle{DBE} = \drawAngle{C}$, \byref{ax:I.III}, que é o ângulo no segmento alternativo.
\end{center}

\qed

\startproblem{Prop. XXXIII. prob.}\label{prop:III.XXXIII}
\defineNewPicture{
pair A, B, D, E, G, K;
numeric r;
r := 7/4u;
G := (0, 0);
A := (0, -r) shifted G;
B := (dir(10)*r) shifted G;
E := (0, r) shifted G;
D := (r, -r) shifted G;
K := (-r, -r) shifted G;
byAngleDefine(E, A, K, byblack, ARC_SECTOR);
byAngleDefine(B, A, E, byred, SOLID_SECTOR);
byAngleDefine(D, A, B, byyellow, SOLID_SECTOR);
byAngleDefine(G, B, A, byred, SOLID_SECTOR);
draw byNamedAngleResized(EAK,BAE,DAB,GBA);
draw byLine(A, B, byblack, SOLID_LINE, REGULAR_WIDTH);
draw byLine(G, B, byyellow, SOLID_LINE, REGULAR_WIDTH);
byLineDefine(A, G, byblue, SOLID_LINE, REGULAR_WIDTH);
byLineDefine(G, E, byblue, DASHED_LINE, REGULAR_WIDTH);
draw byNamedLineSeq(0)(AG,GE);
draw byLineFull(K, D, byred, 0, 0)(K, D, 0, 0, 1);
draw byCircle.G(G, B, byblue, 0, 0, 0);
byAngleDefine.givenRight(E, A, K, byblack, ARC_SECTOR);
byAngleDefine.givenObtuse(B, A, K, byblack, ARC_SECTOR);
byAngleDefine.givenAcute(D, A, B, byblack, ARC_SECTOR);
setAttribute("angle","Standalone","givenRight",2);
setAttribute("angle","Standalone","givenObtuse",2);
setAttribute("angle","Standalone","givenAcute",2);
draw byLabelsOnCircle(A, B)(G);
draw byLabelsOnPolygon(A, G, E)(OMIT_FIRST_LABEL+OMIT_LAST_LABEL, 0);
draw byLabelsOnPolygon(K, D, noPoint)(ALL_LABELS, 0);
}
\drawCurrentPictureInMargin
\problem{U}{ma}{dada linha reta \drawUnitLine{AB} para descrever um segmento de círculo que deve conter um ângulo igual a um determinado ângulo \drawAngle{givenRight}, \drawAngle{givenObtuse}, \drawAngle{givenAcute}.}

Se um determinado ângulo for reto, dividir a reta dada ao meio e descrever um semicírculo sobre ela, evidentemente conterá um ângulo reto. \byref{prop:III.XXXI}

Se o ângulo dado for agudo ou obtuso, faça com a linha dada, em sua extremidade,

\begin{center}
$\drawAngle{DAB} = \drawAngle{givenAcute}$,

draw $\drawUnitLine{AG} \perp \drawUnitLine{KD}$\\
and make $\drawAngle{B} = \drawAngle{BAE}$,

descreva \drawCircle[middle][1/5]{G} com \drawUnitLine{AG} ou \drawUnitLine{GB} como raio, pois eles são iguais.

\drawUnitLine{KD} é tangente a \circleG\ \byref{prop:III.XVI}

$\therefore$ \drawUnitLine{AB} divide o círculo em dois segmentos capazes de conter ângulos iguais a \drawAngle{EAK,BAE} e \drawAngle{DAB} que foram tornados respectivamente iguais a \drawAngle{givenObtuse} e \drawAngle{givenAcute} \byref{prop:III.XXXII}.
\end{center}

\qed

\startproblem{Prop. XXXIV. prob.}\label{prop:III.XXXIV}
\defineNewPicture{
pair A, B, F, E, O;
numeric r;
path cr;
r := 7/4u;
O := (0, 0);
cr := (fullcircle scaled 2r) shifted O;
A := point 1 of cr;
B := point -2 of cr;
E := (r, -r) shifted O;
F := (-r, -r) shifted O;
draw byFilledCircleSegment(O, r, -2, 1, byyellow);
byAngleDefine.B(F, B, A, byblue, SOLID_SECTOR);
byAngleDefine.givenAngle(F, B, A, byred, SOLID_SECTOR);
setAttribute("angle","Standalone","givenAngle",2);
draw byNamedAngleResized(B);
draw byLine(A, B, byblack, SOLID_LINE, REGULAR_WIDTH);
draw byLineFull(E,F, byred, 0, 0)(E, F, 0, 0, -1);
draw byCircleR(O, r, byblue, 0, 0, 0);
draw byLabelsOnCircle(A)(O);
draw byLabelsOnPolygon(E, B, F, noPoint)(ALL_LABELS, 2);
}
\drawCurrentPictureInMargin
\problem{C}{ortar}{de um determinado círculo \drawCircle{O} um segmento que deve conter um ângulo igual a um determinado ângulo \drawAngle{givenAngle}.}

\begin{center}
Trace \drawUnitLine{EF} \byref{prop:III.XVII}, uma tangente ao círculo em qualquer ponto;

no ponto de contato faça $\drawAngle{B} = \drawAngle{givenAngle}$ o ângulo dado;

and \drawFromCurrentPicture[middle][segmentO]{
draw byNamedFilledCircleSegment(O);
draw byLabelsOnCircle(A, B)(O);
} contém um ângulo $=$ do ângulo fornecido.

Because \drawUnitLine{EF} is a tangent,\\
and \drawUnitLine{AB} cuts it,\\
o ângulo em $\segmentO\ = \drawAngle{B}$ \byref{prop:III.XXXII},

but $\drawAngle{B} = \drawAngle{givenAngle}$ \byref{\constref}.
\end{center}

\qed

\starttheorem{Prop. XXXV. Teor.}\label{prop:III.XXXV}
\defineNewPicture{
pair A, B, C, D, E;
numeric r;
r := 5/4u;
E := (0, 0);
A := (dir(50)*r) shifted E;
B := (dir(100)*r) shifted E;
C := (dir(50 + 180)*r) shifted E;
D := (dir(100 + 180)*r) shifted E;
draw byLine(A, E, byblack, SOLID_LINE, REGULAR_WIDTH);
draw byLine(E, C, byblack, DASHED_LINE, REGULAR_WIDTH);
byLineDefine(D, E, byblue, SOLID_LINE, REGULAR_WIDTH);
byLineDefine(E, B, byblue, DASHED_LINE, REGULAR_WIDTH);
draw byNamedLineSeq(0)(DE, EB);
draw byCircleR(E, r, byyellow, 0, 0, 0);
draw byLabelsOnCircle(A, B, C, D)(E);
draw byLabelsOnPolygon(C, E, B)(OMIT_FIRST_LABEL+OMIT_LAST_LABEL, 0);
}
\problem{I}{f}{two chords
$\left\{\vcenter{
\nointerlineskip\hbox{\drawSizedLine{AE,EC}}
\nointerlineskip\hbox{\drawSizedLine{DE,EB}}}\right\}$
em um círculo se cruzam, o retângulo compreendido pelos segmentos de um é igual ao retângulo compreendido pelos segmentos do outro.}
\startsubproposition{Figure I.}
\drawCurrentPictureInMargin
Se as linhas retas dadas passam pelo centro, elas são cortadas ao meio no ponto de intersecção, portanto os retângulos sob seus segmentos são quadrados de suas metades e, portanto, iguais.

\defineNewPicture{
pair A, B, C, D, E, H;
numeric r;
r := 5/4u;
E := (0, 0);
A := (dir(30)*r) shifted E;
B := (dir(170)*r) shifted E;
C := (dir(30 + 180)*r) shifted E;
D := (dir(-50)*r) shifted E;
H = whatever[A, C] = whatever[B, D];
draw byLine(D, E, byred, SOLID_LINE, REGULAR_WIDTH);
draw byLine(E, B, byyellow, SOLID_LINE, REGULAR_WIDTH);
draw byLine(B, H, byblue, SOLID_LINE, REGULAR_WIDTH);
draw byLine(H, D, byblue, DASHED_LINE, REGULAR_WIDTH);
byLineDefine(A, E, byblack, SOLID_LINE, REGULAR_WIDTH);
byLineDefine(E, H, byred, DASHED_LINE, REGULAR_WIDTH);
byLineDefine(H, C, byblack, DASHED_LINE, REGULAR_WIDTH);
draw byNamedLineSeq(0)(AE,EH,HC);
draw byCircleR(E, r, byred, 0, 0, 0);
draw byLabelsOnCircle(B, D, C, A)(E);
draw byLabelsOnPolygon(D, H, C)(OMIT_FIRST_LABEL+OMIT_LAST_LABEL, 0);
draw byLabelsOnPolygon(B, E, A)(OMIT_FIRST_LABEL+OMIT_LAST_LABEL, 0);
}
\startsubproposition{Figure II.}
\drawCurrentPictureInMargin
\begin{center}
Seja \drawSizedLine{HC,EH,AE} passar pelo centro e \drawSizedLine{BH,HD} não; trace \drawSizedLine{EB} e \drawSizedLine{DE}.

Then $\drawSizedLine{BH} \times \drawSizedLine{HD} = \drawSizedLine{EB}^2 - \drawSizedLine{EH}^2$ \byref{prop:II.VI}, or $\drawSizedLine{BH} \times \drawSizedLine{HD} = \drawSizedLine{HC,EH}^2 - \drawSizedLine{EH}^2$.

$\therefore \drawSizedLine{BH} \times \drawSizedLine{HD} = \drawSizedLine{EH,AE} \times \drawSizedLine{AE}$ \byref{prop:II.V}.
\end{center}

\defineNewPicture{
pair A, B, C, D, E, F, K, L;
numeric r;
path cr;
r := 5/4u;
F := (0, 0);
cr := (fullcircle scaled 2r) shifted F;
A := (dir(00)*r) shifted F;
B := (dir(170)*r) shifted F;
C := (dir(-150)*r) shifted F;
D := (dir(-50)*r) shifted F;
E = whatever[A, C] = whatever[B, D];
K := cr intersectionpoint (F -- 10[E, F]);
L := cr intersectionpoint (F -- 10[F, E]);
draw byLine(K, E, byred, DASHED_LINE, REGULAR_WIDTH);
draw byLine(E, L, byred, SOLID_LINE, REGULAR_WIDTH);
draw byLine(B, E, byblue, SOLID_LINE, REGULAR_WIDTH);
draw byLine(E, D, byblue, DASHED_LINE, REGULAR_WIDTH);
byLineDefine(A, E, byblack, SOLID_LINE, REGULAR_WIDTH);
byLineDefine(E, C, byblack, DASHED_LINE, REGULAR_WIDTH);
draw byNamedLineSeq(0)(AE,EC);
draw byCircleR(F, r, byblue, 0, 0, 0);
draw byLabelsOnCircle(B, D, K, L, A, C)(F);
draw byLabelsOnPolygon(B, E, K)(OMIT_FIRST_LABEL+OMIT_LAST_LABEL, 0);
}
\startsubproposition{Figure III.}
\drawCurrentPictureInMargin
Seja nenhuma das linhas fornecidas passar pelo centro, trace através de sua interseção um diâmetro \drawSizedLine{KE,EL},

\begin{center}
and $\drawSizedLine{KE} \times \drawSizedLine{EL} = \drawSizedLine{BE} \times \drawSizedLine{ED}$ (part 2.),\\
also $\drawSizedLine{KE} \times \drawSizedLine{EL} = \drawSizedLine{AE} \times \drawSizedLine{EC}$ (part 2.);

$\therefore \drawSizedLine{BE} \times \drawSizedLine{ED} = \drawSizedLine{AE} \times \drawSizedLine{EC}$.
\end{center}

\qed

\starttheorem{Prop. XXXVI. Teor.}\label{prop:III.XXXVI}
\defineNewPicture{
pair A, B, C, D, F;
numeric r;
path cr[];
r := 7/4u;
F := (0, 0);
cr1 := (fullcircle scaled 2r) shifted F;
D := (r, 3/2r) shifted F;
C := (dir(angle(D-F)))*r;
A := (dir(angle(D-F) + 180))*r;
cr2 := ((fullcircle scaled abs(D-F)) rotated angle(D-F)) shifted 1/2[D, F];
B := cr1 intersectionpoint (subpath (0, 4) of cr2);
draw byLine(B, F, byyellow, SOLID_LINE, REGULAR_WIDTH);
byLineDefine(C, F, byred, DASHED_LINE, REGULAR_WIDTH);
byLineDefine(F, A, byblack, SOLID_LINE, REGULAR_WIDTH);
byLineDefine(D, B, byblue, SOLID_LINE, REGULAR_WIDTH);
byLineDefine(D, C, byred, SOLID_LINE, REGULAR_WIDTH);
draw byNamedLineSeq(0)(DB, DC, CF, FA);
draw byCircleR(F, r, byblack, 0, 0, 0);
draw byLabelsOnCircle(A, B)(F);
draw byLabelsOnPolygon(B, D, F, A)(OMIT_FIRST_LABEL+OMIT_LAST_LABEL, 0);
draw byLabelPoint(C, angle(C-F) - 45, 2);
}
\problem{S}{e}{a partir de um ponto sem círculo forem traçadas duas linhas retas até ele, uma das quais \drawSizedLine{DB} é tangente ao círculo e a outra \drawSizedLine{FA,CF,DC} o corta; o retângulo sob toda a linha de corte \drawSizedLine{FA,CF,DC} e o segmento externo \drawSizedLine{DC} é igual ao quadrado da tangente \drawSizedLine{DB}.}
\drawCurrentPictureInMargin
\startsubproposition{Figure I.}
\begin{center}
Seja \drawSizedLine{FA,CF,DC} passar pelo centro;

trace \drawSizedLine{BF} do centro até o ponto de contato;

$\drawSizedLine{DB}^2 = \drawSizedLine{CF,DC}^2 - \drawSizedLine{BF}^2$ \byref{prop:I.XLVII}, %\mbox{ minus }

or $\drawSizedLine{DB}^2 = \drawSizedLine{CF,DC}^2 - \drawSizedLine{CF}^2$, %\mbox{ minus }

$\therefore \drawSizedLine{DB}^2 = \drawSizedLine{FA,CF,DC} \times \drawSizedLine{DC}$ \byref{prop:II.VI}.
\end{center}

\defineNewPicture{
pair A, B, C, D, E, F;
numeric r;
path cr[];
r := 7/4u;
E := (0, 0);
cr1 := (fullcircle scaled 2r) shifted E;
D := (r, 3/2r) shifted E;
A := (dir(angle(D-E) + 120))*r;
C := cr1 intersectionpoint (A--D);
cr2 := ((fullcircle scaled abs(D-E)) rotated angle(D-E)) shifted 1/2[D, E];
B := cr1 intersectionpoint (subpath (0, 4) of cr2);
draw byLine(E, B, byyellow, SOLID_LINE, REGULAR_WIDTH);
draw byLine(E, C, byblue, DASHED_LINE, REGULAR_WIDTH);
draw byLine(D, C, byred, SOLID_LINE, REGULAR_WIDTH);
draw byLine(C, A, byred, DASHED_LINE, REGULAR_WIDTH);
draw byLine(E, B, byyellow, SOLID_LINE, REGULAR_WIDTH);
byLineDefine(E, A, byyellow, DASHED_LINE, REGULAR_WIDTH);
byLineDefine(D, B, byblue, SOLID_LINE, REGULAR_WIDTH);
byLineDefine(D, E, byblack, SOLID_LINE, REGULAR_WIDTH);
draw byNamedLineSeq(0)(DB,DE,EA);
draw byCircleR(E, r, byblack, 0, 0, 0);
draw byLabelsOnCircle(A, B)(E);
draw byLabelsOnPolygon(B, D, E, A)(OMIT_FIRST_LABEL+OMIT_LAST_LABEL, 0);
draw byLabelPoint(C, angle(C-E) + 45, 2);
}\drawCurrentPictureInMargin
\startsubproposition{Figure II.}
\begin{center}
Se \drawSizedLine{CA,DC} não passar pelo centro,

trace \drawSizedLine{EA} e \drawSizedLine{EC}.

Then $\drawSizedLine{CA,DC} \times \drawSizedLine{CA} = \drawSizedLine{DE}^2 - \drawSizedLine{EC}^2$ \byref{prop:II.VI}, %\mbox{ minus }

that is, $\drawSizedLine{CA,DC} \times \drawSizedLine{CA} = \drawSizedLine{DE}^2 - \drawSizedLine{EB}^2$, %\mbox{ minus }

$\therefore \drawSizedLine{CA,DC} \times \drawSizedLine{CA} = \drawSizedLine{DB}^2$ \byref{prop:III.XVIII}.
\end{center}

\qed

\starttheorem{Prop. XXXVII. Teor.}\label{prop:III.XXXVII}
\defineNewPicture{
pair A, B, C, D, E, F;
numeric r;
path cr[];
r := 2u;
F := (0, 0);
cr1 := (fullcircle scaled 2r) shifted F;
D := (4/3r, 4/3r) shifted F;
cr2 := ((fullcircle scaled abs(D-F)) rotated angle(D-F)) shifted 1/2[D, F];
B := cr1 intersectionpoint (subpath (0, 4) of cr2);
E := cr1 intersectionpoint (subpath (4, 8) of cr2);
A := (dir(angle(D-F) + 140))*r;
C := cr1 intersectionpoint (D--A);
byAngleDefine(D, B, F, byblue, SOLID_SECTOR);
byAngleDefine(F, E, D, byred, SOLID_SECTOR);
draw byNamedAngleResized();
draw byLine(D, C, byblack, DASHED_LINE, REGULAR_WIDTH);
draw byLine(C, A, byblack, SOLID_LINE, REGULAR_WIDTH);
draw byLine(D, F, byyellow, SOLID_LINE, REGULAR_WIDTH);
byLineDefine(B, F, byred, DASHED_LINE, REGULAR_WIDTH);
byLineDefine(E, F, byblue, DASHED_LINE, REGULAR_WIDTH);
byLineDefine(D, B, byred, SOLID_LINE, REGULAR_WIDTH);
byLineDefine(D, E, byblue, SOLID_LINE, REGULAR_WIDTH);
draw byNamedLineSeq(0)(BF,DB,DE,EF);
draw byCircleR(F, r, byblack, 0, 0, 0);
draw byLabelsOnCircle(A, B, E)(F);
draw byLabelsOnPolygon(E, F, B)(OMIT_FIRST_LABEL+OMIT_LAST_LABEL, 0);
draw byLabelsOnPolygon(B, D, E)(OMIT_FIRST_LABEL+OMIT_LAST_LABEL, 0);
draw byLabelPoint(C, angle(C-F) + 45, 2);
}
\drawCurrentPictureInMargin
\problem{S}{e}{a partir de um ponto fora de um círculo forem traçadas duas linhas retas, uma \drawSizedLine{CA,DC} cortando o círculo, a outra \drawSizedLine{DB} encontrando-o, e se o retângulo compreendido por toda a linha de corte \drawSizedLine{CA,DC} e seu segmento externo \drawSizedLine{DC} forem iguais ao quadrado da linha que encontra o círculo, o último \drawSizedLine{DB} é uma tangente ao círculo.}

\begin{center}
Trace a partir do ponto dado \drawSizedLine{DE}, uma tangente ao círculo,\\
e trace a partir do centro \drawSizedLine{DF}, \drawSizedLine{BF} e \drawSizedLine{EF},\\
$\drawSizedLine{DE}^2 = \drawSizedLine{CA,DC} \times \drawSizedLine{DC}$ \byref{prop:III.XXXVI}\\
but $\drawSizedLine{DB}^2 = \drawSizedLine{CA,DC} \times \drawSizedLine{DC}$ \byref{\hypref},\\
$\therefore \drawSizedLine{DB}^2 = \drawSizedLine{DE}^2$,\\
and $\therefore \drawSizedLine{DB} = \drawSizedLine{DE}$.

Then in \drawLine{DB,DF,BF} and \drawLine{EF,DF,DE}\\
$\drawSizedLine{BF} \mbox{ e } \drawSizedLine{DB} = \drawSizedLine{EF} \mbox{ e } \drawSizedLine{DE}$,\\
e \drawSizedLine{DF} é comum,\\
$\therefore \drawAngle{B} = \drawAngle{E}$ \byref{prop:I.VIII};\\
but $\drawAngle{E} = \drawRightAngle$ a right angle \byref{prop:III.XVIII},

$\therefore \drawAngle{B} = \drawRightAngle$ a right angle,\\
e $\therefore$ \drawSizedLine{DB} é uma tangente ao círculo \byref{prop:III.XVI}.
\end{center}

\qed

\part{Livro IV}

\chapter*{Definições}

\startdefinition{}\label{def:IV.I}
\defineNewPicture{
pair A, B, C, D, E, F, G, H;
A := (0, 0);
B := (1/2u, u);
C := (2u, 3/2u);
D := (4/3u, -1/2u);
E := 1/2[A, B];
F := 1/2[B, C];
G := 1/2[C, D];
H := 1/2[D, A];
draw byPolygon(E,F,G,H)(byred);
byLineDefine(A, B, byblack, SOLID_LINE, REGULAR_WIDTH);
byLineDefine(B, C, byblack, SOLID_LINE, REGULAR_WIDTH);
byLineDefine(C, D, byblack, SOLID_LINE, REGULAR_WIDTH);
byLineDefine(D, A, byblack, SOLID_LINE, REGULAR_WIDTH);
draw byNamedLineSeq(1)(AB,BC,CD,DA);
}\drawCurrentPictureInMargin[inside]
Uma figura retilínea é dita \emph{inscrito em} outra, quando todos os pontos angulares da figura inscrita estão nos lados da figura na qual ela está inscrita.

\startdefinition{}\label{def:IV.II}
Diz-se que uma figura está descrita sobre outra figura, quando todos os lados da figura circunscrita passam pelos pontos angulares da outra figura.

\startdefinition{}\label{def:IV.III}
\defineNewPicture{
angleScale := 3/4;
pair A, B, C, D, O;
numeric r;
r := 4/5u;
A := (r, 0);
B := (0, r);
C := (-r, 0);
D := (0, -r);
O := (0, 0);
byAngleDefine(A, B, C, byred, SOLID_SECTOR);
byAngleDefine(B, C, D, byred, SOLID_SECTOR);
byAngleDefine(C, D, A, byred, SOLID_SECTOR);
byAngleDefine(D, A, B, byred, SOLID_SECTOR);
draw byNamedAngleResized();
byLineDefine(A, B, byred, SOLID_LINE, REGULAR_WIDTH);
byLineDefine(B, C, byred, SOLID_LINE, REGULAR_WIDTH);
byLineDefine(C, D, byred, SOLID_LINE, REGULAR_WIDTH);
byLineDefine(D, A, byred, SOLID_LINE, REGULAR_WIDTH);
draw byNamedLineSeq(-1)(DA,CD,BC,AB);
draw byCircleR(O, r, byblack, 0, 0, 1);
}
\drawCurrentPictureInMargin[inside]
Uma figura retilínea é dita \emph{inscrito em} um círculo, quando o vértice de cada ângulo da figura está na circunferência do círculo.

\startdefinition{}\label{def:IV.IV}
\defineNewPicture{
pair A, B, C, D, O;
numeric r;
r := 2/3u;
A := (r, r);
B := (-r, r);
C := (-r, -r);
D := (r, -r);
O := (0, 0);
draw byPolygon(A,B,C,D)(byred);
draw byArbitraryFigure.ABCDo(A--B--C--D--cycle, byred, 0, 0);
draw byFilledCircleSector(O, r, 0, 8, white);
}\drawCurrentPictureInMargin[inside]
Diz-se que uma figura retilínea é \emph{circunscrito a} um círculo, quando cada um dos seus lados é uma tangente ao círculo.

\startdefinition{}\label{def:IV.V}
\defineNewPicture{
pair A, B, C, D, E, F, O;
numeric r[];
r1 := 3/4u;
A := dir(0)*r1;
B := dir(60)*r1;
C := dir(120)*r1;
D := dir(180)*r1;
E := dir(240)*r1;
F := dir(300)*r1;
r2 := abs(1/2[A, B]);
O := (0, 0);
draw byPolygon(A,B,C,D,E,F)(byblue);
draw byArbitraryFigure.ABCDEFo(A--B--C--D--E--F--cycle, byblue, 0, 0);
draw byFilledCircleSector(O, r2, 0, 8, white);
}\drawCurrentPictureInMargin[inside]
Diz-se que um círculo é \emph{inscrito em} uma figura retilínea, quando cada lado da figura é uma tangente ao círculo.

\vfill\pagebreak

\startdefinition{}\label{def:IV.VI}
\defineNewPicture[1/2]{
pair A, B, C, D, E, F, G, H, K, I;
numeric r;
A := (0, 0);
B := (3u, 7/3u);
C := (7/2u, 0);
D = whatever[A, B] = whatever[C, C shifted dir(angle(C-A)) shifted dir(angle(C-B))];
E = whatever[B, C] = whatever[A, A shifted dir(angle(A-B)) shifted dir(angle(A-C))];
F = whatever[C, A] = whatever[B, B shifted dir(angle(B-C)) shifted dir(angle(B-A))];
I = whatever[C, D] = whatever[A, E];
G = whatever[A, B] = whatever[I, I shifted ((A-B) rotated 90)];
H = whatever[B, C] = whatever[I, I shifted ((B-C) rotated 90)];
K = whatever[C, A] = whatever[I, I shifted ((C-A) rotated 90)];
r := abs(D-I);
draw byPolygon(G,H,K)(byblue);
byLineDefine(A, B, byblack, SOLID_LINE, REGULAR_WIDTH);
byLineDefine(B, C, byblack, SOLID_LINE, REGULAR_WIDTH);
byLineDefine(C, A, byblack, SOLID_LINE, REGULAR_WIDTH);
draw byNamedLineSeq(1)(AB,BC,CA);
draw byCircleR(I, r, byblack, 0, 0, -1);
}\drawCurrentPictureInMargin[inside]
Diz-se que um círculo é \emph{circunscrito a} uma figura retilínea, quando a circunferência passa pelo vértice de cada ângulo da figura.

\drawPolygon{GHK} is circumscribed.

\startdefinition{}\label{def:IV.VII}
\defineNewPicture{
pair A, B, O;
numeric r;
r := 3/4u;
O := (0, 0);
A := dir(0)*r;
B := dir(100)*r;
draw byLine(A, B, byblack, SOLID_LINE, REGULAR_WIDTH);
draw byCircleR(O, r, byblue, 0, 0, 0);
}\drawCurrentPictureInMargin[inside]
Diz-se que uma linha reta é \emph{inscrito em} um círculo, quando suas extremidades estão na circunferência.

\vskip 1cm

O Quarto Livro dos Elementos é dedicado à solução de problemas, principalmente relativos à inscrição e circunscrição de polígonos e círculos regulares.

Um polígono regular é aquele cujos ângulos e lados são iguais.

\vfill\pagebreak

\startproblem{Prop. I. prob.}\label{prop:IV.I}
\defineNewPicture{
pair A, B, C, E, G, H, O;
numeric r[];
path cr[];
r1 := 4/3u;
r2 := 7/4u;
O := (0, 0);
B := (r1, 0) shifted O;
C := (-r1, 0) shifted O;
cr1 := (fullcircle scaled 2r1) shifted O;
cr2 := (fullcircle scaled 2r2) shifted C;
A := cr1 intersectionpoint (subpath (0, 2) of cr2);
E := (B--C) intersectionpoint cr2;
G := (xpart(lrcorner(cr1)), ypart(lrcorner(cr2)));
H := G shifted (-r2, 0);
draw byLine(B, E, byred, SOLID_LINE, REGULAR_WIDTH);
draw byLine(E, C, byred, DASHED_LINE, REGULAR_WIDTH);
draw byLine(A, C, byyellow, SOLID_LINE, REGULAR_WIDTH);
draw byLine(G, H, byblue, SOLID_LINE, REGULAR_WIDTH);
draw byCircleR(O, r1, byyellow, 0, 0, 0);
draw byCircle.C(C, E, byblue, 0, 0, 0);
draw byLabelsOnCircle(C, B)(O);
draw byLabelsOnPolygon(C, A, E)(OMIT_FIRST_LABEL+OMIT_LAST_LABEL, 0);
draw byLabelsOnPolygon(C, E, A)(OMIT_FIRST_LABEL+OMIT_LAST_LABEL, 1);
draw byLabelsOnPolygon(G, H, noPoint)(ALL_LABELS, 0);
}
\drawCurrentPictureInMargin
\problem{N}{um}{determinado círculo \drawCircle{O} colocar uma reta, igual a uma dada reta (\drawUnitLine{GH}), não maior que o diâmetro do círculo.}

\begin{center}
Trace \drawUnitLine{EC,BE}, o diâmetro de \circleO;\\
e se for $\drawUnitLine{EC,BE} = \drawUnitLine{GH}$, o problema está resolvido.

Mas se \drawUnitLine{EC,BE} não for igual a \drawUnitLine{GH},\\
$\drawUnitLine{EC,BE} > \drawUnitLine{GH}$ \byref{\hypref};

make $\drawUnitLine{EC} = \drawUnitLine{GH}$ \byref{prop:I.III}\\
with \drawUnitLine{EC} as radius,\\
describe \drawCircle{C}, cutting \circleO, \\
e trace \drawUnitLine{AC}, que é a linha necessária.

For $\drawUnitLine{AC} = \drawUnitLine{EC} = \drawUnitLine{GH}$ \byref{prop:I.XV,\constref}.
\end{center}

\qed

\startproblem{Prop. II. prob.}\label{prop:IV.II}
\defineNewPicture{
pair A, B, C, D, E, F, G, H, O, d;
numeric r;
path cr;
r := 7/4u;
O := (0, 0);
cr := (fullcircle scaled 2r) shifted O;
A := (0, -r) shifted O;
G := (-r, -r) shifted O;
H := (r, -r) shifted O;
D := (0, 0);
E := (-1/2r, 4/3r);
F := (3/4r, r);
d := -1/2[ulcorner(D--E--F), lrcorner(D--E--F)] shifted (0, -2r);
D := D shifted d;
E := E shifted d;
F := F shifted d;
C := (subpath (-1, 5) of cr) intersectionpoint (A -- A shifted (dir(angleValue(F, E, D))*2r));
B := (subpath (-1, 5) of cr) intersectionpoint (A -- A shifted (dir(180-angleValue(D, F, E))*2r));
byAngleDefine(C, B, A, byblack, SOLID_SECTOR);
byAngleDefine(A, C, B, byblack, ARC_SECTOR);
byAngleDefine(G, A, B, byyellow, SOLID_SECTOR);
byAngleDefine(B, A, C, byred, SOLID_SECTOR);
byAngleDefine(C, A, H, byblue, SOLID_SECTOR);
byAngleDefine(F, E, D, byblue, SOLID_SECTOR);
byAngleDefine(E, D, F, byred, SOLID_SECTOR);
byAngleDefine(D, F, E, byyellow, SOLID_SECTOR);
draw byNamedAngleResized();
draw byArbitraryFigure.BAC(B--A--C, byblack, 0, 1);
draw byArbitraryFigure.BAC(D--E--F--cycle, byblack, 0, 1);
draw byLine(B, C, byyellow, SOLID_LINE, REGULAR_WIDTH);
draw byLine(G, H, byred, SOLID_LINE, REGULAR_WIDTH);
draw byCircleR(O, r, byblack, 0, 0, 0);
draw byLabelsOnCircle(A, B, C)(O);
draw byLabelsOnPolygon(D, E, F)(ALL_LABELS, 0);
draw byLabelsOnPolygon(H, G, noPoint)(ALL_LABELS, 0);
}
\drawCurrentPictureInMargin
\problem{E}{m}{um determinado círculo \drawCircle{O} para inscrever um triângulo equiângulo a um determinado triângulo.}

\begin{center}
Para qualquer ponto do círculo dado trace \drawUnitLine{GH}, uma tangente \byref{prop:III.XVII};

e no ponto de contato faça $\drawAngle{CAH} = \drawAngle{E}$ \byref{prop:I.XXIII}

and in like manner $\drawAngle{GAB} = \drawAngle{F}$,

e trace \drawUnitLine{BC}.

$\because \drawAngle{CAH} = \drawAngle{E}$ \byref{\constref}\\ %Because
and $\drawAngle{CAH} = \drawAngle{B}$ \byref{prop:III.XXXII}\\
$\therefore \drawAngle{B} = \drawAngle{E}$;\\
also $\drawAngle{C} = \drawAngle{F}$ for the same reason.

$\therefore \drawAngle{BAC} = \drawAngle{D}$ \byref{prop:I.XXXII},
\end{center}

\noindent e portanto o triângulo inscrito em o círculo é equiângulo ao dado.

\qed

\startproblem{Prop. III. prob.}\label{prop:IV.III}
\defineNewPicture[3/5]{
byAngleMacroName[3] := "byAngleMSectors";
vardef byAngleMSectors (expr angleArc, angleColor)(suffix angleOptionalColors) =
    save p;
    path p[];
    p1 := (subpath(0, arctime (1/3arclength(angleArc)) of angleArc) of angleArc) scaled (angleSize*angleScale);
    p2 := (subpath(arctime (1/3arclength(angleArc)) of angleArc, arctime(2/3arclength(angleArc)) of angleArc) of angleArc) scaled (angleSize*angleScale);
    p3 := (subpath(arctime (2/3arclength(angleArc)) of angleArc, length(angleArc)) of angleArc) scaled (angleSize*angleScale);
    image(
        fill ((0, 0)--p1--cycle) withcolor angleOptionalColors[0];
        fill ((0, 0)--p2--cycle) withcolor angleColor;
        fill ((0, 0)--p3--cycle) withcolor angleOptionalColors[1];
    )
enddef;
pair A, B, C, D, E, F, G, H, K, L, M, N, d;
numeric r;
r := 5/4u;
d := (-4/3r, 11/5r);
D := (0, 0) shifted d;
E := (1/3r, -r) shifted d;
F := (-5/6r, -r) shifted d;
G := 4/3[F, E];
H := 4/3[E, F];
K := (0, 0);
C := (0, -r) shifted K;
A := (dir(-90+angleValue(D, F, H))*r) shifted K;
B := (dir(-90-angleValue(G, E, D))*r) shifted K;
L = whatever[C, C shifted (dir(angle(C-K) + 90))] = whatever[A, A shifted (dir(angle(A-K) + 90))];
M = whatever[B, B shifted (dir(angle(B-K) + 90))] = whatever[A, L];
N = whatever[B, M] = whatever[C, L];
byAngleDefine(D, F, H, byyellow, SOLID_SECTOR);
byAngleDefine(C, K, A, byyellow, SOLID_SECTOR);
byAngleDefine(G, E, D, byblue, SOLID_SECTOR);
byAngleDefine(B, K, C, byblue, SOLID_SECTOR);
byAngleDefine(M, L, N, byred, SOLID_SECTOR);
byAngleDefine(D, F, E, byred, SOLID_SECTOR);
byAngleDefine(L, N, M, byblack, ARC_SECTOR);
byAngleDefine(F, E, D, byblack, ARC_SECTOR);
byAngleDefineExtended(N, M, L, byblack, 3)(byblue, byred);
byAngleDefineExtended(E, D, F, byblack, 3)(byblue, byred);
byAngleDefine(K, A, L, byblack, SOLID_SECTOR);
byAngleDefine(L, C, K, byblack, SOLID_SECTOR);
draw byNamedAngleResized();
draw byArbitraryFigure.FDE(F--D--E, byblack, 0, 1);
byLineDefine(N, L, byblue, SOLID_LINE, REGULAR_WIDTH);
byLineDefine(L, M, byyellow, SOLID_LINE, REGULAR_WIDTH);
byLineDefine(M, N, byred, DASHED_LINE, REGULAR_WIDTH);
draw byLine(G, H, byblack, SOLID_LINE, REGULAR_WIDTH);
draw byLine(K, C, byred, SOLID_LINE, REGULAR_WIDTH);
draw byArbitraryFigure.AKB(A--K--B, byblack, 0, 1);
draw byCircleR(K, r, byred, 0, 0, -1);
draw byNamedLineSeq(0)(NL,LM,MN);
byArbitraryFigureDefine.LAKC(L--A--K--C--cycle, byblack, 0, 1);
draw byLabelsOnPolygon(L, A, M, B, N, C)(ALL_LABELS, 0);
draw byLabelsOnPolygon(G, E, F, H, noPoint)(ALL_LABELS, 0);
draw byLabelsOnPolygon(F, D, E)(OMIT_FIRST_LABEL+OMIT_LAST_LABEL, 0);
draw byLabelsOnPolygon(A, K, B)(OMIT_FIRST_LABEL+OMIT_LAST_LABEL, 0);
}
\drawCurrentPictureInMargin
\problem{S}{obre}{um determinado círculo \drawCircle[middle][1/3]{K} para circunscrever um triângulo equiângulo a um determinado triângulo.}

\begin{center}
Produza qualquer lado \drawUnitLine{GH}, do triângulo dado em ambos os sentidos;

do centro do círculo fornecido, trace \drawUnitLine{KC}, qualquer raio.

Make $\drawAngle{CKA} = \drawAngle{DFH}$ \byref{prop:I.XXIII}\\
and $\drawAngle{BKC} = \drawAngle{GED}$.

Nas extremidades dos raios, trace \drawUnitLine{NL}, \drawUnitLine{LM} e \drawUnitLine{MN}, tangentes ao círculo dado. \byref{prop:III.XVII}

Os quatro ângulos de
\drawFromCurrentPicture[bottom]{
startTempAngleScale(2/3);
draw byNamedAngleResized(A, C, L, CKA);
draw byNamedArbitraryFigure(LAKC);
draw byLabelsOnPolygon(L, A, K, C)(ALL_LABELS, 0);
stopTempAngleScale;
},
juntos são iguais a quatro ângulos retos. \byref{prop:I.XXXII}

mas \drawAngle{C} e \drawAngle{A} são ângulos retos \byref{\constref}

$\therefore \drawAngle{L} + \drawAngle{CKA} = \drawTwoRightAngles$, two right angles

but $\drawAngle{DFH,DFE} = \drawTwoRightAngles$ \byref{prop:I.XIII}\\
and $\drawAngle{CKA} = \drawAngle{DFH}$ \byref{\constref} and $\therefore \drawAngle{L} = \drawAngle{DFE}$.

In the same manner it can be demonstrated that\\
$\drawAngle{N} = \drawAngle{FED}$;

$\therefore \drawAngle{M} = \drawAngle{D}$ \byref{prop:I.XXXII}\\
e, portanto, o triângulo circunscrito a do círculo dado é equiângulo ao triângulo dado.
\end{center}

\qed

\startproblem{Prop. IV. prob.}\label{prop:IV.IV}
\defineNewPicture{
pair A, B, C, D, E, F, G;
numeric r;
A := (0, 0);
B := (-3/2u, -3u);
C := (3u, -3u);
D = whatever[A, A shifted (unitvector(B-A) + unitvector(C-A))] = whatever[B, B shifted (unitvector(A-B) + unitvector(C-B))];
E = whatever[D, D shifted ((A-B) rotated 90)] = whatever[A, B];
F = whatever[D, D shifted ((B-C) rotated 90)] = whatever[B, C];
G = whatever[D, D shifted ((C-A) rotated 90)] = whatever[A, C];
r := abs(E-D);
byAngleDefine(B, E, D, byred, SOLID_SECTOR);
byAngleDefine(D, F, B, byred, SOLID_SECTOR);
byAngleDefine(E, B, D, byblue, SOLID_SECTOR);
byAngleDefine(D, B, F, byyellow, SOLID_SECTOR);
byAngleDefine(G, C, D, byblack, SOLID_SECTOR);
byAngleDefine(D, C, F, byblack, ARC_SECTOR);
byAngleDefine(E, D, F, byblack, ARC_SECTOR);
draw byNamedAngleResized();
draw byLine(E, D, byyellow, DASHED_LINE, REGULAR_WIDTH);
draw byLine(G, D, byred, DASHED_LINE, REGULAR_WIDTH);
draw byLine(F, D, byblack, DASHED_LINE, REGULAR_WIDTH);
byLineDefine(C, D, byblue, SOLID_LINE, REGULAR_WIDTH);
byLineDefine(B, D, byblue, DASHED_LINE, REGULAR_WIDTH);
draw byNamedLineSeq(0)(CD,BD);
byLineDefine(A, B, byyellow, SOLID_LINE, REGULAR_WIDTH);
byLineDefine(B, C, byblack, SOLID_LINE, REGULAR_WIDTH);
byLineDefine(C, A, byred, SOLID_LINE, REGULAR_WIDTH);
draw byNamedLineSeq(0)(AB,BC,CA);
draw byCircle.D(D, G, byblack, 0, 0, -1);
byLineDefine(B, F, byblack, SOLID_LINE, REGULAR_WIDTH);
byLineDefine(B, E, byyellow, SOLID_LINE, REGULAR_WIDTH);
draw byLabelsOnPolygon(B, E, A, G, C, F)(ALL_LABELS, 0);
draw byLabelsOnPolygon(E, D, G)(OMIT_FIRST_LABEL+OMIT_LAST_LABEL, 0);
}
\drawCurrentPictureInMargin
\problem{I}{n}{um dado triângulo \drawFromCurrentPicture[bottom]{
startTempScale(1/4);
startAutoLabeling;
draw byNamedLineSeq(0)(CA,BC,AB);
stopAutoLabeling;
stopTempScale;
} to inscribe a circle.}

Divida \drawAngle{EBD,DBF} e \drawAngle{GCD,DCF} \byref{prop:I.IX} por \drawUnitLine{BD} e \drawUnitLine{CD}; a partir do ponto onde essas linhas se encontram, trace \drawUnitLine{FD}, \drawUnitLine{ED} e \drawUnitLine{GD} respectivamente perpendiculares a \drawUnitLine{BC}, \drawUnitLine{AB} e \drawUnitLine{CA}.

\begin{center}
In
\drawFromCurrentPicture{
startTempScale(3/4);
draw byNamedAngle(F, DBF);
startAutoLabeling;
draw byNamedLineSeq(0)(BD,FD,BF);
stopAutoLabeling;
stopTempScale;
}
and
\drawFromCurrentPicture{
startTempScale(3/4);
draw byNamedAngle(E, EBD);
startAutoLabeling;
draw byNamedLineSeq(0)(BE,ED,BD);
stopAutoLabeling;
stopTempScale;
}\\
$\drawAngle{DBF} = \drawAngle{EBD}$, $\drawAngle{F} = \drawAngle{E}$ and \drawUnitLine{BD} common,\\
$\therefore \drawUnitLine{ED} = \drawUnitLine{FD}$ \byref{prop:I.IV,prop:I.XXVI}.

In like manner, it may be shown also that $\drawUnitLine{GD} = \drawUnitLine{FD}$,

$\therefore \drawUnitLine{FD} = \drawUnitLine{ED} = \drawUnitLine{GD}$;
\end{center}

portanto, com qualquer uma dessas linhas como raio, descreva \drawCircle[middle][1/2]{D} e ela passará pelas extremidades das outras duas; e os lados do triângulo dado, sendo perpendiculares aos três raios em suas extremidades, tocam o círculo \byref{prop:III.XVI}, que é portanto inscrito em o triângulo dado.

\qed

\startproblem{Prop. V. prob.}\label{prop:IV.V}
\defineNewPicture[1]{
def proptmp (expr a, b, c, s) =
pair A, B, C, D, E, F;
numeric r;
r := 7/4u;
F := s;
A := (dir(a)*r) shifted F;
B := (dir(b)*r) shifted F;
C := (dir(c)*r) shifted F;
D := 1/2[A, B];
E := 1/2[A, C];
byAngleDefine(A, D, F, byred, SOLID_SECTOR);
byAngleDefine(F, D, B, byblack, SOLID_SECTOR);
draw byNamedAngleResized();
draw byLine(D, F, byyellow, SOLID_LINE, REGULAR_WIDTH);
draw byLine(E, F, byyellow, DASHED_LINE, REGULAR_WIDTH);
draw byLine(A, F, byblack, DASHED_LINE, REGULAR_WIDTH);
byLineDefine(A, D, byblue, SOLID_LINE, REGULAR_WIDTH);
byLineDefine(D, B, byblue, DASHED_LINE, REGULAR_WIDTH);
draw byNamedLineSeq(0)(AD,DB);
if (sind(b)<>-sind(c)):
byLineDefine(B, F, byblack, SOLID_LINE, REGULAR_WIDTH);
byLineDefine(C, F, byblack, SOLID_LINE, THIN_WIDTH);
draw byNamedLineSeq(0)(BF,CF);
draw byLine(B, C, byred, SOLID_LINE, REGULAR_WIDTH);
else:
draw byLine(B, C, byblack, SOLID_LINE, REGULAR_WIDTH);
fi;
byLineDefine(C, E, byred, SOLID_LINE, REGULAR_WIDTH);
byLineDefine(E, A, byred, DASHED_LINE, REGULAR_WIDTH);
draw byNamedLineSeq(0)(CE,EA);
draw byCircleR(F, r, byblack, 0, 0, 0);
draw byLabelsOnCircle(A, B, C)(F);
draw byLabelsOnPolygon(B, D, A)(OMIT_FIRST_LABEL+OMIT_LAST_LABEL, 0);
draw byLabelsOnPolygon(A, E, C)(OMIT_FIRST_LABEL+OMIT_LAST_LABEL, 0);
draw byLabelsOnPolygon(C, F, B)(OMIT_FIRST_LABEL+OMIT_LAST_LABEL, 0);
enddef;
proptmp(100, 170, 10, (0, -8u));
proptmp(80, 160, -20, (0, -4u));
proptmp(100, 200, -30, (0, 0));
}
\drawCurrentPictureInMargin
\problem{P}{ara}{descrever um círculo em torno de um determinado triângulo.}

Make $\drawUnitLine{AD} = \drawUnitLine{DB}$ and $\drawUnitLine{CE} = \drawUnitLine{EA}$ \byref{prop:I.X}

Dos pontos de bissecção trace \drawUnitLine{DF} e \drawUnitLine{EF} $\perp$ \drawUnitLine{AD} e \drawUnitLine{CE} respectivamente \byref{prop:I.XI}, e de seu ponto de encontro trace \drawUnitLine{BF}, \drawUnitLine{AF} e \drawUnitLine{CF} e descreva um círculo com qualquer um deles, e será o círculo necessário.

\begin{center}
In
\drawFromCurrentPicture{
draw byNamedAngle(FDB);
startAutoLabeling;
draw byNamedLineSeq(0)(DF,BF,DB);
stopAutoLabeling;
}
and
\drawFromCurrentPicture{
draw byNamedAngle(ADF);
startAutoLabeling;
draw byNamedLineSeq(0)(AF,DF,AD);
stopAutoLabeling;
}\\
$\drawUnitLine{DB} = \drawUnitLine{AD}$ \byref{\constref},\\
\drawUnitLine{DF} common,\\
$\drawAngle{FDB} = \drawAngle{ADF}$ \byref{\constref},

$\therefore \drawUnitLine{AF} = \drawUnitLine{BF}$ \byref{prop:I.IV}.

In like manner it may be shown that $\drawUnitLine{CF} = \drawUnitLine{AF}$.
\end{center}

$\therefore \drawUnitLine{AF} = \drawUnitLine{BF} = \drawUnitLine{CF}$; e, portanto, um círculo descrito a partir do encontro dessas três linhas com qualquer uma delas como um raio circunscreverá o triângulo dado.

\qed

\startproblem{Prop. VI. prob.}\label{prop:IV.VI}
\defineNewPicture{
pair A, B, C, D, E;
numeric r;
r := 7/4u;
E := (0, 0);
A := (0, r);
B := (-r, 0);
C := (0, -r);
D := (r, 0);
byAngleDefine(D, A, B, byyellow, SOLID_SECTOR);
byAngleDefine(C, D, A, byblack, SOLID_SECTOR);
byAngleDefine(A, E, B, byblue, SOLID_SECTOR);
byAngleDefine(D, E, A, byred, SOLID_SECTOR);
draw byNamedAngleResized();
draw byLine(A, E, byblack, DASHED_LINE, REGULAR_WIDTH);
draw byLine(C, E, byblack, SOLID_LINE, THIN_WIDTH);
byLineDefine(B, E, byred, DASHED_LINE, REGULAR_WIDTH);
byLineDefine(D, E, byblue, DASHED_LINE, REGULAR_WIDTH);
draw byNamedLineSeq(0)(BE, DE);
byLineDefine(A, B, byblack, SOLID_LINE, REGULAR_WIDTH);
byLineDefine(B, C, byyellow, SOLID_LINE, REGULAR_WIDTH);
byLineDefine(C, D, byblue, SOLID_LINE, REGULAR_WIDTH);
byLineDefine(D, A, byred, SOLID_LINE, REGULAR_WIDTH);
draw byNamedLineSeq(0)(AB,BC,CD,DA);
draw byCircleR(E, r, byred, 0, 0, 1);
draw byLabelsOnCircle(A, B, C, D)(E);
draw byLabelsOnPolygon(C, E, B)(OMIT_FIRST_LABEL+OMIT_LAST_LABEL, 0);
}
\drawCurrentPictureInMargin
\problem{N}{um}{determinado círculo \drawCircle{E} para inscrever um quadrado.}

Trace os dois diâmetros do círculo $\perp$ entre si e trace \drawUnitLine{BC}, \drawUnitLine{AB}, \drawUnitLine{DA} e \drawUnitLine{CD}.

\begin{center}
\drawLine[middle][squareABCD]{DA,CD,BC,AB} is a square.

Pois, como \drawAngle{A} e \drawAngle{D} estão, cada um deles, em semicírculo, são ângulos retos \byref{prop:III.XXXI},

$\therefore \drawUnitLine{CD} \parallel \drawUnitLine{AB}$ \byref{prop:I.XXVIII}:

and in like manner $\drawUnitLine{DA} \parallel \drawUnitLine{BC}$.

And $\because \drawAngle{AEB} = \drawAngle{DEA}$ \byref{\constref}, %because

and $\drawUnitLine{BE} = \drawUnitLine{AE} = \drawUnitLine{DE}$ \byref{def:I.XV}.

$\therefore \drawUnitLine{AB} = \drawUnitLine{DA}$ \byref{prop:I.IV};
\end{center}

e como os lados e ângulos adjacentes do paralelogramo \squareABCD\ são iguais, eles são todos iguais \byref{prop:I.XXXIV}; e $\therefore$ \squareABCD, inscrito em o círculo dado, é um quadrado.

\qed

\startproblem{Prop. VII. prob.}\label{prop:IV.VII}
\defineNewPicture[1/4]{
pair A, B, C, D, E, F, G, H, K;
numeric r;
r := 7/4u;
E := (0, 0);
A := (0, r);
B := (-r, 0);
C := (0, -r);
D := (r, 0);
F := (r, r);
G := (-r, r);
H := (-r, -r);
K := (r, -r);
byAngleDefine(F, G, H, byred, SOLID_SECTOR);
byAngleDefine(G, H, K, byred, SOLID_SECTOR);
byAngleDefine(H, K, F, byred, SOLID_SECTOR);
byAngleDefine(K, F, G, byred, SOLID_SECTOR);
byAngleDefine(B, E, C, byblack, SOLID_SECTOR);
byAngleDefine(E, C, H, byyellow, SOLID_SECTOR);
draw byNamedAngleResized();
draw byLine(A, C, byred, DASHED_LINE, REGULAR_WIDTH);
draw byLine(B, D, byblue, DASHED_LINE, REGULAR_WIDTH);
byLineDefine(F, G, byblack, SOLID_LINE, REGULAR_WIDTH);
byLineDefine(G, H, byred, SOLID_LINE, REGULAR_WIDTH);
byLineDefine(H, K, byblue, SOLID_LINE, REGULAR_WIDTH);
byLineDefine(K, F, byyellow, SOLID_LINE, REGULAR_WIDTH);
draw byNamedLineSeq(0)(KF,HK,GH,FG);
draw byCircleR(E, r, byblue, 0, 0, -1);
draw byLabelsOnPolygon(H, B, G, A, F, D, K, C)(ALL_LABELS, 0);
draw byLabelsOnPolygon(A, E, D)(OMIT_FIRST_LABEL+OMIT_LAST_LABEL, 0);
}
\drawCurrentPictureInMargin
\problem{S}{obre}{um determinado círculo \drawCircle{E} para circunscrever um quadrado.}

Trace dois diâmetros do círculo dado perpendiculares entre si e, através de suas extremidades, trace as tangentes \drawUnitLine{HK}, \drawUnitLine{GH}, \drawUnitLine{FG} e \drawUnitLine{KF} ao círculo;

\begin{center}
e \drawLine[bottom][squareFGHK]{KF,HK,GH,FG} é um quadrado.

$\drawAngle{C} = \drawRightAngle$ a right angle, \byref{prop:III.XVIII}

also $\drawAngle{E} = \drawRightAngle$ \byref{\constref},

$\therefore \drawUnitLine{HK} \parallel \drawUnitLine{BD}$; da mesma forma pode-se demonstrar que $\drawUnitLine{FG} \parallel \drawUnitLine{BD}$, e também que $\drawUnitLine{GH} \mbox{ e } \drawUnitLine{KF} \parallel \drawUnitLine{AC}$;

$\therefore$ \squareFGHK\ is a parallelogram,

e $\because \drawAngle{C} = \drawAngle{G} = \drawAngle{F} = \drawAngle{K} = \drawAngle{H}$ são todos ângulos retos \byref{prop:I.XXXIV}; % porque

também é evidente que \drawUnitLine{HK}, \drawUnitLine{GH}, \drawUnitLine{FG} e \drawUnitLine{KF} são iguais.

$\therefore$ \squareFGHK\ is a square.
\end{center}

\qed

\startproblem{Prop. VIII. prob.}\label{prop:IV.VIII}
\defineNewPicture[1/2]{
pair A, B, C, D, E, F, G, H, K;
numeric r;
r := 7/4u;
G := (0, 0);
A := (-r, r);
B := (-r,-r);
C := (r, -r);
D := (r, r);
E := (0, r);
F := (-r, 0);
H := (0, -r);
K := (r, 0);
draw byPolygon(E,D,K,G)(byblack);
draw byPolygon(F,G,H,B)(byred);
draw byPolygon(G,K,C,H)(byblue);
draw byLine(F, G, byyellow, SOLID_LINE, REGULAR_WIDTH);
draw byLine(G, K, byyellow, DASHED_LINE, REGULAR_WIDTH);
draw byLine(G, H, byblack, SOLID_LINE, REGULAR_WIDTH);
draw byLine(E, G, byblack, DASHED_LINE, REGULAR_WIDTH);
byLineDefine(A, B, byblack, SOLID_LINE, REGULAR_WIDTH);
byLineDefine(A, E, byblue, DASHED_LINE, REGULAR_WIDTH);
byLineDefine(E, D, byblue, SOLID_LINE, REGULAR_WIDTH);
byLineDefine(D, K, byred, DASHED_LINE, REGULAR_WIDTH);
byLineDefine(K, C, byred, SOLID_LINE, REGULAR_WIDTH);
draw byNamedLineSeq(0)(AB,AE,ED,DK,KC);
draw byCircleR(G, r, byyellow, 0, 0, -1);
draw byLabelsOnPolygon(B, F, A, E, D, K, C, H)(ALL_LABELS, 0);
draw byLabelsOnPolygon(F, G, E)(OMIT_FIRST_LABEL+OMIT_LAST_LABEL, 0);
}
\drawCurrentPictureInMargin
\problem{I}{nscrever}{um círculo em um determinado quadrado.}

\begin{center}
Make $\drawUnitLine{ED} = \drawUnitLine{AE}$,\\
and $\drawUnitLine{KC} = \drawUnitLine{DK}$,\\
draw $\drawUnitLine{FG,GK} \parallel \drawUnitLine{AE,ED}$,\\
and $\drawUnitLine{EG,GH} \parallel \drawUnitLine{DK,KC}$\\
\byref{prop:I.XXXI}

$\therefore$ \drawPolygon[bottom]{EDKG} is a parallelogram;

and since $\drawUnitLine{AE,ED} = \drawUnitLine{DK,KC}$ \byref{\hypref}\\
$\drawUnitLine{ED} = \drawUnitLine{DK}$

$\therefore$ \polygonEDKG\ is equilateral \byref{prop:I.XXXIV}

In like manner it can be shown that $ \drawPolygon[bottom]{GKCH} = \drawPolygon[bottom]{FGHB}$ are equilateral parallelograms;

$\therefore \drawUnitLine{AE} = \drawUnitLine{GK} = \drawUnitLine{GH} = \drawUnitLine{FG}$.
\end{center}

\noindent e portanto se um círculo for descrito a partir do encontro destas linhas com qualquer uma delas como raio, será inscrito em o quadrado dado \byref{prop:III.XVI}.

\qed

\startproblem{Prop. IX. prob.}\label{prop:IV.IX}
\defineNewPicture{
pair A, B, C, D, E;
numeric r;
r := 7/4u;
E := (0, 0);
A := (dir(45+90)*r) shifted E;
B := (dir(45+180)*r) shifted E;
C := (dir(45+270)*r) shifted E;
D := (dir(45+360)*r) shifted E;
draw byPolygon(A,C,D)(byred);
draw byPolygon(A,B,C)(byyellow);
byAngleDefine(D, A, E, byyellow, SOLID_SECTOR);
byAngleDefine(E, A, B, byblack, SOLID_SECTOR);
byAngleDefine(A, B, E, byred, SOLID_SECTOR);
byAngleDefine(E, B, C, byblue, SOLID_SECTOR);
draw byNamedAngleResized();
draw byLine(A, E, byblue, SOLID_LINE, REGULAR_WIDTH);
draw byLine(E, C, byblue, DASHED_LINE, REGULAR_WIDTH);
byLineDefine(B, E, byblack, SOLID_LINE, REGULAR_WIDTH);
byLineDefine(E, D, byblack, DASHED_LINE, REGULAR_WIDTH);
draw byNamedLineSeq(0)(BE, ED);
draw byCircleR(E, r, byblack, 0, 0, 0);
draw byLabelsOnCircle(A, B, C, D)(E);
draw byLabelsOnPolygon(B, E, A)(OMIT_FIRST_LABEL+OMIT_LAST_LABEL, 1);
}
\drawCurrentPictureInMargin
\problem{P}{ara}{descrever um círculo em torno de um determinado quadrado \drawPolygon{ACD,ABC}.}

Trace as diagonais \drawUnitLine{AE,EC} e \drawUnitLine{BE,ED} que se cruzam;

então, como \drawPolygon{ACD} e \drawPolygon{ABC} têm seus lados iguais, e a base \drawUnitLine{AE,EC} comum a ambos,

\begin{center}
$\drawAngle{DAE} = \drawAngle{EAB}$ \byref{prop:I.VIII},\\
or \drawAngle{DAE,EAB} is bisected: \\
in like manner it can be shown that \drawAngle{ABE,EBC} is bisected;

but $\drawAngle{DAE,EAB} = \drawAngle{ABE,EBC}$,

hence $\drawAngle{EAB} = \drawAngle{ABE}$ their halves,

$\therefore \drawUnitLine{BE} = \drawUnitLine{AE}$ \byref{prop:I.VI};

and in like manner it can be proved that \\
$\drawUnitLine{AE} = \drawUnitLine{BE} = \drawUnitLine{EC} = \drawUnitLine{ED}$.
\end{center}

Se a partir da confluência dessas linhas com qualquer uma delas como raio, for descrito um círculo, ele circunscreverá o quadrado dado.

\qed

\startproblem{Prop. X. prob.}\label{prop:IV.X}
\defineNewPicture[1/2]{
pair A, B, C, D, E, F;
numeric r[];
path cr[];
r1 := 11/4u;
A := (0, 0);
cr1 := (fullcircle scaled 2r1) shifted A;
B := (r1, 0);
r2 := r1*sqrt(5)/2 - 1/2r1;
cr2 := (fullcircle scaled 2r2) shifted B;
C := (r2, 0);
D := cr1 intersectionpoint (subpath (0, 4) of cr2);
F = whatever[1/2[A, C], 1/2[A, C] shifted ((A-C) rotated 90)]
 = whatever[1/2[A, D], 1/2[A, D] shifted ((A-D) rotated 90)];
r3 := abs(F-A);
byAngleDefine(D, A, B, byblack, ARC_SECTOR);
byAngleDefine(C, D, A, byblack, SOLID_SECTOR);
byAngleDefine(B, D, C, byyellow, SOLID_SECTOR);
byAngleDefine(D, C, B, byblue, SOLID_SECTOR);
byAngleDefine(C, B, D, byred, SOLID_SECTOR);
draw byNamedAngleResized();
draw byLine(A, D, byyellow, SOLID_LINE, REGULAR_WIDTH);
draw byLine(B, D, byblue, SOLID_LINE, REGULAR_WIDTH);
draw byLine(C, D, byred, SOLID_LINE, REGULAR_WIDTH);
draw byLine(A, C, byblack, SOLID_LINE, REGULAR_WIDTH);
draw byLine(C, B, byblack, DASHED_LINE, REGULAR_WIDTH);
draw byCircleABC.F(A, C, D, byblue, 0, 0, 0);
draw byCircle.A(A, B, byred, 0, 0, 0);
draw byLabelsOnCircle(B)(A);
draw byLabelsOnCircle(A)(F);
draw byLabelPoint(D, angle(D-A)-15, 2);
draw byLabelsOnPolygon(A, C, D)(OMIT_FIRST_LABEL+OMIT_LAST_LABEL, 0);
}
\drawCurrentPictureInMargin
\problem{C}{onstruir}{um triângulo isósceles, em que cada um dos ângulos da base seja o dobro do ângulo vertical.}

\begin{center}
Pegue qualquer linha reta \drawProportionalLine{AC,CB} \\
and divide it so that $\drawProportionalLine{AC,CB} \times \drawProportionalLine{CB} = \drawProportionalLine{AC}^2$ \byref{prop:II.XI}

Tendo \drawProportionalLine{AC,CB} como raio, descreva \drawCircle[middle][1/5]{A} e coloque nele a partir da extremidade do raio, $\drawProportionalLine{BD} = \drawProportionalLine{AC}$ \byref{prop:IV.I};\\
draw \drawProportionalLine{AD}.

Então \drawLine{AD,BD,CB,AC} é o triângulo necessário.

Pára, trace \drawProportionalLine{CD} e descreva \drawCircle[middle][1/3]{F} sobre \drawLine{AD,CD,AC} \byref{prop:IV.V}

Since $\drawProportionalLine{AC,CB} \times \drawProportionalLine{CB} = \drawProportionalLine{AC}^2 = \drawProportionalLine{BD}^2$,\\
$\therefore$ \drawProportionalLine{BD} é tangente a \circleF\ \byref{prop:III.XXXVII}\\
$\therefore \drawAngle{BDC} = \drawAngle{A}$ \byref{prop:III.XXXII},\\
add \drawAngle{CDA} to each, $\therefore \drawAngle{BDC} + \drawAngle{CDA} = \drawAngle{A} + \drawAngle{CDA}$;

but $\drawAngle{BDC} + \drawAngle{CDA} \mbox{ ou } \drawAngle{BDC,CDA} = \drawAngle{B}$ \byref{prop:I.V}:\\
since $\drawProportionalLine{AD} = \drawProportionalLine{AC,CB}$ \byref{prop:I.V}\\
consequently $\drawAngle{B} = \drawAngle{A} + \drawAngle{CDA} = \drawAngle{C}$ \byref{prop:I.XXXII}\\
$\therefore \drawProportionalLine{CD} = \drawProportionalLine{BD}$ \byref{prop:I.VI}\\
$\therefore \drawProportionalLine{BD} = \drawProportionalLine{AC} = \drawProportionalLine{CD}$ \byref{\constref}\\
$\therefore \drawAngle{A} = \drawAngle{CDA}$ \byref{prop:I.V}

$\therefore \drawAngle{B} = \drawAngle{BDC,CDA} = \drawAngle{C} = \drawAngle{A} + \drawAngle{CDA} = \mbox{ twice } \drawAngle{A}$;

e conseqüentemente cada ângulo na base é o dobro do ângulo vertical.
\end{center}

\qed

\startproblem{Prop. XI. prob.}\label{prop:IV.XI}
\defineNewPicture[1/4]{
pair A, B, C, D, E, O;
numeric r;
r := 9/4u;
O := (0, 0);
A := (dir(90+0/5(360))*r) shifted O;
B := (dir(90+1/5(360))*r) shifted O;
C := (dir(90+2/5(360))*r) shifted O;
D := (dir(90+3/5(360))*r) shifted O;
E := (dir(90+4/5(360))*r) shifted O;
draw byArbitraryFigure.pg(C--E--B--D, byblack, 0, 1);
draw byPolygon(A,C,D)(byyellow);
byAngleDefineExtended(A, B, C, byblack, -1)(bytransparent);
byAngleDefineExtended(D, E, A, byblack, -1)(bytransparent);
byAngleDefine.Cl(B, C, A, byblack, -ARC_SECTOR);
byAngleDefine.Dr(A, D, E, byblack, -ARC_SECTOR);
byAngleDefine.Al(B, A, C, byblack, -ARC_SECTOR);
byAngleDefine.Ar(D, A, E, byblack, -ARC_SECTOR);
byAngleDefine(C, A, D, byblack, SOLID_SECTOR);
byAngleDefineExtended(E, C, A, byblack, 1)(byyellow);
byAngleDefine(D, C, E, byblue, SOLID_SECTOR);
byAngleDefine(A, D, B, byblack, ARC_SECTOR);
byAngleDefine(B, D, C, byred, SOLID_SECTOR);
draw byNamedAngleResized();
draw byNamedAngleDummySides(ADB,BDC);
byLineDefine(A, B, byblue, SOLID_LINE, REGULAR_WIDTH);
byLineDefine(B, C, byred, SOLID_LINE, REGULAR_WIDTH);
byLineDefine(C, D, byblack, SOLID_LINE, REGULAR_WIDTH);
byLineDefine(D, E, byred, DASHED_LINE, REGULAR_WIDTH);
byLineDefine(E, A, byyellow, SOLID_LINE, REGULAR_WIDTH);
draw byNamedLineSeq(0)(AB,BC,CD,DE,EA);
draw byCircleR(O, r, byblue, 0, 0, 1);
draw byLabelsOnCircle(A, B, C, D, E)(O);
}
\drawCurrentPictureInMargin
\problem[2]{E}{m}{um determinado círculo \drawCircle[middle][1/4]{O} para inscrever um pentágono equilátero e equiângulo.}

Construa um triângulo isósceles, em que cada um dos ângulos da base seja o dobro do ângulo do vértice, e inscreva no círculo dado um triângulo \drawPolygon[bottom]{ACD} equiângulo a ele \byref{prop:IV.II};

\begin{center}
Bisect \drawAngle{ECA,DCE} and \drawAngle{ADB,BDC} \byref{prop:I.IX},\\
trace \drawUnitLine{BC}, \drawUnitLine{AB}, \drawUnitLine{EA} e \drawUnitLine{DE}.
\end{center}

Como cada um dos ângulos \drawAngle{ECA}, \drawAngle{DCE}, \drawAngle{CAD}, \drawAngle{BDC} e \drawAngle{ADB} são iguais, os arcos sobre os quais eles se apoiam são iguais \byref{prop:III.XXVI}; e $\therefore$ \drawUnitLine{CD}, \drawUnitLine{BC}, \drawUnitLine{AB}, \drawUnitLine{EA} e \drawUnitLine{DE} que subentendem que esses arcos são iguais \byref{prop:III.XXIX} e $\therefore$ o pentágono é equilátero, também é equiângulo, pois cada um de seus ângulos está sobre arcos iguais \byref{prop:III.XXVII}.

\qed

\startproblem{Prop. XII. prob.}\label{prop:IV.XII}
\defineNewPicture[1/2]{
pair B, C, D, G, H, K, L, M, F;
numeric r[];
r1 := 5/2u;
F := (0, 0);
G := (dir(90+0/5(360))*r1) shifted F;
H := (dir(90+1/5(360))*r1) shifted F;
K := (dir(90+2/5(360))*r1) shifted F;
L := (dir(90+3/5(360))*r1) shifted F;
M := (dir(90+4/5(360))*r1) shifted F;
B := 1/2[H, K];
C := 1/2[K, L];
D := 1/2[L, M];
r2 := abs(F-B);
byAngleDefine(B, F, K, byred, SOLID_SECTOR);
byAngleDefine(K, F, C, byyellow, SOLID_SECTOR);
byAngleDefine(C, F, L, byblue, SOLID_SECTOR);
byAngleDefine(L, F, D, byred, SOLID_SECTOR);
byAngleDefine(F, K, B, byyellow, ARC_SECTOR);
byAngleDefine(C, K, F, byblack, SOLID_SECTOR);
byAngleDefine(F, C, K, byblue, ARC_SECTOR);
byAngleDefine(L, C, F, byblue, ARC_SECTOR);
byAngleDefine(F, L, C, byblack, SOLID_SECTOR);
byAngleDefine(D, L, F, byblack, ARC_SECTOR);
draw byNamedAngleResized();
draw byLine(F, K, byblue, SOLID_LINE, REGULAR_WIDTH);
draw byLine(F, C, byblack, DASHED_LINE, REGULAR_WIDTH);
draw byLine(F, L, byblack, SOLID_LINE, REGULAR_WIDTH);
byLineDefine(F, B, byred, DASHED_LINE, REGULAR_WIDTH);
byLineDefine(F, D, byyellow, DASHED_LINE, REGULAR_WIDTH);
draw byNamedLineSeq(0)(FB,FD);
byLineDefine(G, H, byblack, SOLID_LINE, REGULAR_WIDTH);
byLineDefine(H, B, byblue, DASHED_LINE, REGULAR_WIDTH);
byLineDefine(B, K, byblack, SOLID_LINE, REGULAR_WIDTH);
byLineDefine(K, C, byred, SOLID_LINE, REGULAR_WIDTH);
byLineDefine(C, L, byyellow, SOLID_LINE, REGULAR_WIDTH);
byLineDefine(L, M, byblack, SOLID_LINE, REGULAR_WIDTH);
byLineDefine(M, G, byblack, SOLID_LINE, REGULAR_WIDTH);
draw byNamedLineSeq(0)(GH,HB,BK,KC,CL,LM,MG);
draw byCircleR(F, r2, byred, 0, 0, -1);
draw byLabelsOnPolygon(M, D, L, C, K, B, H, G)(OMIT_FIRST_LABEL+OMIT_LAST_LABEL, 0);
draw byLabelsOnPolygon(B, F, D)(OMIT_FIRST_LABEL+OMIT_LAST_LABEL, 0);
}
\drawCurrentPictureInMargin
\problem{P}{ara}{descrever um pentágono equilátero e equiângulo em torno de um determinado círculo \drawCircle[middle][1/4]{F}.}

Desenhe cinco tangentes através dos vértices dos ângulos de qualquer pentágono regular inscrito em do círculo dado \circleF\ \byref{prop:III.XVII}.

These five tangents will form the required pentagon.

\begin{center}
Draw
$\left\{\vcenter{
\nointerlineskip\hbox{\drawProportionalLine{FB}}
\nointerlineskip\hbox{\drawProportionalLine{FK}}
\nointerlineskip\hbox{\drawProportionalLine{FC}}
\nointerlineskip\hbox{\drawProportionalLine{FL}}
\nointerlineskip\hbox{\drawProportionalLine{FD}}
}\right.$.

In \drawLine{FB,FK,BK} and \drawLine{FK,FC,KC}\\
$\drawProportionalLine{BK} = \drawProportionalLine{KC}$ \byref{prop:I.XLVII} \\
$\drawProportionalLine{FC} = \drawProportionalLine{FB}$, and \drawProportionalLine{FK} common;

$\therefore \drawAngle{FKB} = \drawAngle{CKF}$ and $\therefore \drawAngle{BFK} = \drawAngle{KFC}$ \byref{prop:I.VIII}

$\therefore \drawAngle{FKB,CKF} = \mbox{ twice } \drawAngle{CKF}$, and $\drawAngle{BFK,KFC} = \mbox{ twice } \drawAngle{KFC}$.

In the same manner it can be demonstrated that\\
$\drawAngle{FLC,DLF} = \mbox{ twice } \drawAngle{FLC}$, and $\drawAngle{CFL,LFD} = \mbox{ twice } \drawAngle{CFL}$;

but $\drawAngle{BFK,KFC} = \drawAngle{CFL,LFD}$ \byref{prop:III.XXVII},

$\therefore$ their halves $\drawAngle{KFC} = \drawAngle{CFL}$, also $\drawAngle{FCK} = \drawAngle{LCF}$,\\
and \drawProportionalLine{FC} common;

$\therefore \drawAngle{CKF} = \drawAngle{FLC}$ and $\drawProportionalLine{KC} = \drawProportionalLine{CL}$,

$\therefore \drawProportionalLine{KC,CL} = \mbox{ twice } \drawProportionalLine{KC}$.

In the same manner it can be demonstrated that $\drawProportionalLine{BK,HB} = \mbox{ twice } \drawProportionalLine{BK}$,\\
but $\drawProportionalLine{BK} = \drawProportionalLine{KC}$

$\therefore \drawProportionalLine{BK,HB} = \drawProportionalLine{KC,CL}$;
\end{center}

Da mesma forma pode-se demonstrar que os outros lados são iguais e, portanto, o pentágono é equilátero, também é equiângulo, por

\begin{center}
$\drawAngle{FLC,DLF} = \mbox{ twice } \drawAngle{FLC}$ and $\drawAngle{FKB,CKF} = \mbox{ twice } \drawAngle{CKF}$,

and therefore $ \drawAngle{CKF} = \drawAngle{FLC}$,

$\therefore \drawAngle{FLC,DLF} = \drawAngle{FKB,CKF}$;

da mesma maneira pode-se demonstrar que os demais ângulos do pentágono descrito são iguais.
\end{center}

\qed

\startproblem{Prop. XIII. prob.}\label{prop:IV.XIII}
\defineNewPicture{
pair A, B, C, D, E, F, G, H, K, L, M;
numeric r[];
r1 := 5/2u;
F := (0, 0);
A := (dir(90+0/5(360))*r1) shifted F;
B := (dir(90+1/5(360))*r1) shifted F;
C := (dir(90+2/5(360))*r1) shifted F;
D := (dir(90+3/5(360))*r1) shifted F;
E := (dir(90+4/5(360))*r1) shifted F;
G := 1/2[A, B];
H := 1/2[B, C];
K := 1/2[C, D];
L := 1/2[D, E];
M := 1/2[E, A];
r2 := abs(F-G);
byAngleDefine(G, B, F, byyellow, SOLID_SECTOR);
byAngleDefine(F, B, H, byyellow, SOLID_SECTOR);
byAngleDefine(F, H, C, byblack, SOLID_SECTOR);
byAngleDefine(H, C, F, byblue, SOLID_SECTOR);
byAngleDefine(F, C, K, byblue, SOLID_SECTOR);
byAngleDefine(C, K, F, byblack, SOLID_SECTOR);
byAngleDefine(K, D, F, byred, SOLID_SECTOR);
byAngleDefine(F, D, E, byred, SOLID_SECTOR);
byAngleDefine(D, E, F, byblack, ARC_SECTOR);
byAngleDefine(F, E, M, byblack, ARC_SECTOR);
draw byNamedAngleResized();
draw byLine(F, G, byblack, SOLID_LINE, THIN_WIDTH);
draw byLine(F, M, byblack, SOLID_LINE, THIN_WIDTH);
draw byLine(F, A, byblack, SOLID_LINE, THIN_WIDTH);
draw byLine(F, H, byblue, DASHED_LINE, REGULAR_WIDTH);
draw byLine(F, K, byblue, SOLID_LINE, REGULAR_WIDTH);
draw byLine(F, C, byblack, SOLID_LINE, REGULAR_WIDTH);
draw byLine(F, D, byyellow, DASHED_LINE, REGULAR_WIDTH);
byLineDefine(F, B, byred, SOLID_LINE, REGULAR_WIDTH);
byLineDefine(F, E, byred, DASHED_LINE, REGULAR_WIDTH);
draw byNamedLineSeq(0)(FB,FE);
byLineDefine(A, B, byblack, SOLID_LINE, THIN_WIDTH);
byLineDefine(B, H, byblack, DASHED_LINE, REGULAR_WIDTH);
byLineDefine(H, C, byyellow, SOLID_LINE, REGULAR_WIDTH);
byLineDefine(C, K, byyellow, SOLID_LINE, REGULAR_WIDTH);
byLineDefine(K, D, byblack, DASHED_LINE, REGULAR_WIDTH);
byLineDefine(D, E, byblack, SOLID_LINE, THIN_WIDTH);
byLineDefine(E, A, byblack, SOLID_LINE, THIN_WIDTH);
draw byNamedLineSeq(0)(AB,BH,HC,CK,KD,DE,EA);
draw byCircle.F(F, H, byyellow, 0, 0, -1);
draw byLabelsOnPolygon(C, H, B, A, E, D, K)(ALL_LABELS, 0);
draw byLabelsOnPolygon(E, F, D)(OMIT_FIRST_LABEL+OMIT_LAST_LABEL, 2);
}
\drawCurrentPictureInMargin
\problem{I}{nscrever}{um círculo em um determinado pentágono equiângulo e equilátero.}

Let
\drawLine[bottom]{EA,DE,KD,CK,HC,BH,AB}
ser um dado pentágono equilátero e equilátero; é necessário inscrever um círculo nele.

\begin{center}
Make $\drawAngle{HCF} = \drawAngle{FCK}$, and $\drawAngle{FDE} = \drawAngle{KDF}$ \byref{prop:I.IX}

Draw \drawUnitLine{FD}, \drawUnitLine{FC}, \drawUnitLine{FB}, \drawUnitLine{FE}, \&c.

$\because \drawUnitLine{KD,CK} = \drawUnitLine{HC,BH}$, $\drawAngle{HCF} = \drawAngle{FCK}$ e \drawUnitLine{FC} comuns aos dois triângulos %Porque

\drawFromCurrentPicture{
startTempScale(1/2);
draw byNamedAngle(FBH,HCF);
startAutoLabeling;
draw byNamedLineSeq(0)(FB,FC,HC,BH);
stopAutoLabeling;
stopTempScale;
}
and
\drawFromCurrentPicture{
startTempScale(1/2);
draw byNamedAngle(FCK,KDF);
startAutoLabeling;
draw byNamedLineSeq(0)(FC,FD,KD,CK);
stopAutoLabeling;
stopTempScale;
};

$\therefore \drawUnitLine{FB} = \drawUnitLine{FD}$ and $\drawAngle{FBH} = \drawAngle{KDF}$ \byref{prop:I.IV}

And $\because \drawAngle{B} = \drawAngle{D} = \mbox{ twice } \drawAngle{KDF}$ %because

$\therefore \mbox{ twice } \drawAngle{FBH}$, hence \drawAngle{B} is bisected by \drawUnitLine{FB}.

Da mesma forma, pode ser demonstrado que \drawAngle{DEF,FEM} é dividido ao meio por \drawUnitLine{FE} e que o ângulo restante do polígono é dividido ao meio de maneira semelhante.

Trace \drawUnitLine{FK}, \drawUnitLine{FH}, \&c.\ perpendiculares aos lados do pentágono.

Then in the two triangles
\drawFromCurrentPicture{
startTempScale(2/3);
draw byNamedAngle(H,HCF);
startAutoLabeling;
draw byNamedLineSeq(0)(FH,FC,HC);
stopAutoLabeling;
stopTempScale;
}
and
\drawFromCurrentPicture{
startTempScale(2/3);
draw byNamedAngle(FCK,K);
startAutoLabeling;
draw byNamedLineSeq(0)(CK,FC,FK);
stopAutoLabeling;
stopTempScale;
}\\
we have $\drawAngle{HCF} = \drawAngle{FCK}$ \byref{\constref}, \drawUnitLine{FC} common,\\
e $\drawAngle{H} = \drawAngle{K} = \mbox{ a right angle }$;

$\therefore \drawUnitLine{FK} = \drawUnitLine{FH}$ \byref{prop:I.XXVI}
\end{center}

Da mesma forma, pode-se mostrar que as cinco perpendiculares aos lados do pentágono são iguais entre si.

Descreva \drawCircle[middle][1/5]{F} com qualquer uma das perpendiculares como raio, e será o círculo inscrito necessário. Pois se não tocar os lados do pentágono, mas cortá-los, então uma linha traçada a partir da extremidade perpendicularmente ao diâmetro de um círculo cairá dentro do círculo, o que se mostrou absurdo \byref{prop:III.XVI}.

\qed

\startproblem{Prop. XIV. prob.}\label{prop:IV.XIV}
\defineNewPicture[1/2]{
pair A, B, C, D, E, F;
numeric r;
r := 9/4u;
F := (0, 0);
A := (dir(90+0/5(360))*r) shifted F;
B := (dir(90+1/5(360))*r) shifted F;
C := (dir(90+2/5(360))*r) shifted F;
D := (dir(90+3/5(360))*r) shifted F;
E := (dir(90+4/5(360))*r) shifted F;
byAngleDefine(C, F, B, byblue, SOLID_SECTOR);
byAngleDefine(D, F, C, byblack, ARC_SECTOR);
byAngleDefine(B, C, F, byblack, SOLID_SECTOR);
byAngleDefine(F, C, D, byyellow, SOLID_SECTOR);
byAngleDefine(C, D, F, byyellow, SOLID_SECTOR);
byAngleDefine(F, D, E, byred, SOLID_SECTOR);
draw byNamedAngleResized();
draw byLine(F, A, byyellow, DASHED_LINE, REGULAR_WIDTH);
draw byLine(F, C, byred, DASHED_LINE, REGULAR_WIDTH);
draw byLine(F, D, byblue, DASHED_LINE, REGULAR_WIDTH);
byLineDefine(F, B, byblack, SOLID_LINE, REGULAR_WIDTH);
byLineDefine(F, E, byyellow, SOLID_LINE, REGULAR_WIDTH);
draw byNamedLineSeq(0)(FB,FE);
byLineDefine(A, B, byblack, SOLID_LINE, THIN_WIDTH);
byLineDefine(B, C, byblue, SOLID_LINE, REGULAR_WIDTH);
byLineDefine(C, D, byred, SOLID_LINE, REGULAR_WIDTH);
byLineDefine(D, E, byblack, DASHED_LINE, REGULAR_WIDTH);
byLineDefine(E, A, byblack, SOLID_LINE, THIN_WIDTH);
draw byNamedLineSeq(0)(AB,BC,CD,DE,EA);
draw byCircleR(F, r, byred, 0, 0, 1);
draw byLabelsOnCircle(A, B, C, D, E)(F);
draw byLabelsOnPolygon(A, F, E)(OMIT_FIRST_LABEL+OMIT_LAST_LABEL, 1);
}
\drawCurrentPictureInMargin
\problem{P}{ara}{descrever um círculo em torno de um determinado pentágono equilátero e equiângulo.}

\begin{center}
Divida \drawAngle{BCF,FCD} e \drawAngle{CDF,FDE} por \drawUnitLine{FC} e \drawUnitLine{FD} e, a partir do ponto de corte, trace \drawUnitLine{FE}, \drawUnitLine{FA} e \drawUnitLine{FB}.

$\drawAngle{BCF,FCD} = \drawAngle{CDF,FDE}$,\\
$\drawAngle{FCD} = \drawAngle{CDF}$,

$\therefore \drawUnitLine{FD} = \drawUnitLine{FC}$ \byref{prop:I.VI};

and since in
\drawFromCurrentPicture{
draw byNamedAngle(BCF);
startAutoLabeling;
draw byNamedLineSeq(0)(FB,FC,BC);
stopAutoLabeling;
}
and
\drawFromCurrentPicture{
draw byNamedAngle(FCD);
startAutoLabeling;
draw byNamedLineSeq(0)(FC,FD,CD);
stopAutoLabeling;
},\\
$\drawUnitLine{BC} = \drawUnitLine{CD}$, and \drawUnitLine{FC} common,\\
also $\drawAngle{BCF} = \drawAngle{FCD}$;

$\therefore \drawUnitLine{FB} = \drawUnitLine{FD}$ \byref{prop:I.IV}.

In like manner it may be proved that\\
$\drawUnitLine{FA} = \drawUnitLine{FE} = \drawUnitLine{FB}$.

and therefore $\drawUnitLine{FA} = \drawUnitLine{FB} = \drawUnitLine{FC} = \drawUnitLine{FD} = \drawUnitLine{FE}$.
\end{center}

Portanto, se um círculo for descrito a partir do ponto onde estas cinco linhas se encontram, com qualquer uma delas como raio, ele circunscreverá o pentágono dado.

\qed

\startproblem{Prop. XV. prob.}\label{prop:IV.XV}
\defineNewPicture[1/2]{
pair A, B, C, D, E, F, G, H;
numeric r;
r := 9/4u;
G := (0, 0);
A := (dir(90)*r) shifted G;
B := (dir(150)*r) shifted G;
C := (dir(210)*r) shifted G;
D := (dir(270)*r) shifted G;
E := (dir(330)*r) shifted G;
F := (dir(30)*r) shifted G;
H := (dir(270)*r) shifted D;
byAngleDefine(A, G, F, byblack, -ARC_SECTOR);
byAngleDefine(B, G, A, byblack, -ARC_SECTOR);
byAngleDefine(C, G, B, byblack, -ARC_SECTOR);
byAngleDefine(D, G, C, byred, SOLID_SECTOR);
byAngleDefine(E, G, D, byblue, SOLID_SECTOR);
byAngleDefine(F, G, E, byblack, SOLID_SECTOR);
draw byNamedAngleResized();
draw byLine(D, H, byred, SOLID_LINE, REGULAR_WIDTH);
draw byLine(A, D, byblack, SOLID_LINE, REGULAR_WIDTH);
draw byLine(B, E, byyellow, SOLID_LINE, REGULAR_WIDTH);
draw byLine(C, F, byblue, SOLID_LINE, REGULAR_WIDTH);
draw byLine(A, B, byblack, SOLID_LINE, REGULAR_WIDTH);
draw byLine(B, C, byblack, SOLID_LINE, REGULAR_WIDTH);
draw byLine(C, D, byblack, DASHED_LINE, REGULAR_WIDTH);
draw byLine(D, E, byred, DASHED_LINE, REGULAR_WIDTH);
draw byLine(E, F, byblue, DASHED_LINE, REGULAR_WIDTH);
draw byLine(F, A, byblack, SOLID_LINE, REGULAR_WIDTH);
draw byCircle.D(D, G, byred, 0, 0, 0);
draw byCircleR(G, r, byyellow, 0, 0, 0);
byLineDefine(G, C, lineColor.CF, SOLID_LINE, REGULAR_WIDTH);
byLineDefine(G, D, lineColor.AD, SOLID_LINE, REGULAR_WIDTH);
byLineDefine(G, E, lineColor.BE, SOLID_LINE, REGULAR_WIDTH);
draw byLabelsOnCircle(A, B, F)(G);
draw byLabelPoint(G, angle(A + F - 2G), 3);
draw byLabelPoint(D, angle(D-G) + 45, 2);
draw byLabelsOnPolygon(D, C, G)(OMIT_FIRST_LABEL+OMIT_LAST_LABEL, 0);
draw byLabelsOnPolygon(G, E, D)(OMIT_FIRST_LABEL+OMIT_LAST_LABEL, 0);
}
\drawCurrentPictureInMargin
\problem[2]{P}{ara}{inscrever um hexágono equilátero e equiângulo em um determinado círculo \drawCircle[middle][1/4]{G}.}

A partir de qualquer ponto da circunferência do círculo dado, descreva \drawCircle[middle][1/4]{D} passando pelo seu centro, e trace os diâmetros \drawUnitLine{AD}, \drawUnitLine{CF} e \drawUnitLine{BE}; trace \drawUnitLine{CD}, \drawUnitLine{DE}, \drawUnitLine{EF}, \&c.\ e o hexágono necessário é inscrito em o círculo fornecido.

Como o \drawUnitLine{GD} passa pelos centros dos círculos,
\drawFromCurrentPicture{
draw byNamedAngle(DGC);
startAutoLabeling;
draw byNamedLineSeq(0)(GC,GD,CD);
stopAutoLabeling;
}
and
\drawFromCurrentPicture{
draw byNamedAngle(EGD);
startAutoLabeling;
draw byNamedLineSeq(0)(GD,GE,DE);
stopAutoLabeling;
}
são triângulos equiláteros, portanto $\drawAngle{DGC} = \drawAngle{EGD} = \mbox{ one-third of } \drawTwoRightAngles$ \byref{prop:I.XXXII} mas $\drawAngle{DGC,EGD,FGE} = \drawTwoRightAngles$ \byref{prop:I.XIII};

$\therefore \drawAngle{DGC} = \drawAngle{EGD} = \drawAngle{FGE} = \mbox{ one-third of } \drawTwoRightAngles$ \byref{prop:I.XXXII}, e os ângulos verticalmente opostos a estes são todos iguais entre si \byref{prop:I.XV}, e estão em arcos iguais \byref{prop:III.XXVI}, que são subtendidos por cordas iguais \byref{prop:III.XXIX}; e como cada um dos ângulos do hexágono é o dobro do ângulo de um triângulo equilátero, ele também é equiângulo.

\qed

\startproblem{Prop. XVI. prob.}\label{prop:IV.XVI}
\defineNewPicture{
pair A, B, C, D, E, F, G, H, O;
numeric r;
r := 9/4u;
O := (0, 0);
A := (dir(90 + 0*120)*r) shifted O;
C := (dir(90 + 1*120)*r) shifted O;
D := (dir(90 + 2*120)*r) shifted O;
B := (dir(90 + 1*72)*r) shifted O;
F := (dir(90 + 2*72)*r) shifted O;
G := (dir(90 + 3*72)*r) shifted O;
H := (dir(90 + 4*72)*r) shifted O;
draw byLine(C, F, byblack, DASHED_LINE, REGULAR_WIDTH);
byLineDefine(A, B, byred, SOLID_LINE, REGULAR_WIDTH);
byLineDefine(B, F, byblue, SOLID_LINE, REGULAR_WIDTH);
byLineDefine(F, G, byblack, SOLID_LINE, THIN_WIDTH);
byLineDefine(G, H, byblack, SOLID_LINE, THIN_WIDTH);
byLineDefine(H, A, byblack, SOLID_LINE, THIN_WIDTH);
byLineDefine(A, C, byyellow, SOLID_LINE, REGULAR_WIDTH);
byLineDefine(C, D, byblack, SOLID_LINE, THIN_WIDTH);
byLineDefine(D, A, byblack, SOLID_LINE, THIN_WIDTH);
draw byNamedLineSeq(0)(AB,BF,FG,GH,HA);
draw byNamedLineSeq(0)(AC,CD,DA);
draw byCircleR(O, r, byblue, 0, 0, 1);
draw byLabelsOnCircle(A, B, C, F)(O);
}
\drawCurrentPictureInMargin
\problem{I}{nscrever}{um círculo em um quindecágono equilátero e equiângulo em um determinado círculo.}

Sejam \drawUnitLine{AB} e \drawUnitLine{BF} os lados de um pentágono equilátero inscrito em o círculo dado, e \drawUnitLine{AC} o lado de um triângulo equilátero inscrito.

\begin{center}
$$
\left.\vcenter{\hbox{The arc subtended} \hbox{por \drawUnitLine{AB} e \drawUnitLine{BF}}}\right\} = \dfrac{2}{5} = \dfrac{6}{15} \left\{\vcenter{\hbox{of the whole} \hbox{circumference}}\right.
$$

$$
\left.\vcenter{\hbox{The arc subtended} \hbox{by \drawUnitLine{AC}}}\right\} = \dfrac{1}{3} = \dfrac{5}{15} \left\{\vcenter{\hbox{of the whole} \hbox{circumference}}\right.
$$

Their difference $ = \dfrac{1}{15}$
\end{center}

$\therefore$ o arco subtendido pela diferença $\drawUnitLine{CF} = \dfrac{1}{15}$ de toda a circunferência.

Portanto, se linhas retas iguais a \drawUnitLine{CF} forem colocadas no círculo \byref{prop:IV.I}, um quindecágono equilátero e equiangular será, portanto, inscrito em o círculo.

\qed

\part{Livro V}

\chapter*{Definições}

\startdefinition[1]{}\label{def:V.I}
Diz-se que uma magnitude menor é uma parte alíquota ou submúltiplo de uma magnitude maior, quando a menor mede a maior; isto é, quando o menos está contido um certo número de vezes exatamente no maior.

\startdefinition[2]{}\label{def:V.II}
Diz-se que uma grandeza maior é múltiplo de uma menor, quando o maior é medido pela menor; isto é, quando o maior contém o menor um certo número de vezes exatamente.

\startdefinition[3]{}\label{def:V.III}
Razão é a relação que uma quantidade mantém com outra da mesma espécie, no que diz respeito à magnitude.

\startdefinition[4]{}\label{def:V.IV}
Diz-se que as magnitudes têm uma razão entre si, quando são da mesma espécie; e aquele que não for o maior pode ser multiplicado de modo a exceder o outro.

\vskip \baselineskip

\begin{center}
\emph{As outras definições serão dadas ao longo do livro onde seu auxílio for exigido pela primeira vez.}
\end{center}

\chapter*{Axiomas}

\startaxiom{}\label{ax:V.I}
Equimúltiplos ou equissubmúltiplos iguais ou de magnitudes iguais são iguais.

\begin{center}
$
\begin{aligned}
	\mbox{Se } A &= B \mbox{, then}\\
	\mbox{duas vezes } A &= \mbox{duas vezes } B \mbox{,}\\
	\mbox{isto é, }2A &= 2B \mbox{;}\\
	3A &= 3B \mbox{;}\\
	4A &= 4B \mbox{;}\\
	\mbox{\&c.} & \mbox{ \&c.}\\
	\mbox{e } \frac{1}{2} \mbox{ de } A &= \frac{1}{2} \mbox{ de } B \mbox{;}\\
	\frac{1}{3} \mbox{ de } A &= \frac{1}{3} \mbox{ de } B \mbox{;}\\
	\frac{1}{4} \mbox{ de } A &= \frac{1}{4} \mbox{ de } B \mbox{;}\\
	\mbox{\&c. } & \mbox{ \&c.}\\
\end{aligned}
$
\end{center}

\startaxiom{}\label{ax:V.II}
Um múltiplo de maior magnitude é maior que o mesmo múltiplo de menor.

\begin{center}
$\begin{aligned}
	\mbox{Let } A &> B \mbox{,} \\
	\mbox{ then } 2A &> 2B \mbox{;}\\
	3A&> 3B \mbox{;} \\
	4A &> 4B \mbox{;}\\
	\mbox{\&c. } &\mbox{\&c.} \\
\end{aligned}
$
\end{center}

\startaxiom{}\label{ax:V.III}
Aquela grandeza, da qual um múltiplo é maior que o mesmo múltiplo de outro, é maior que o outro.

\begin{center}
$\begin{aligned}
\mbox{Let } 2A &> 2B \mbox{,}\\
\mbox{then } A &> B \mbox{;}\\
\mbox{or, let } 3A &> 3B \mbox{,}\\
\mbox{then } A &> B \mbox{;}\\
\mbox{or, let } mA &> mB  \mbox{,}\\
\mbox{then } A &> B \mbox{.}\\
\end{aligned}$
\end{center}

\vfill\pagebreak

\starttheorem{Prop. I. Theor.}\label{prop:V.I}
\defineNewPicture{
byMagnitudeSymbolDefine.ab("semicircleUp", byred, 1);
byMagnitudeSymbolDefine.cd("wedgeDown", byyellow, 0);
byMagnitudeSymbolDefine.ef("sectorDown", byblue, 1);
byMagnitudeDefine.A(0, false)(5)(ab);
byMagnitudeDefine.B(0, false)(1)(ab);
byMagnitudeDefine.C(0, false)(5)(cd);
byMagnitudeDefine.D(0, false)(1)(cd);
byMagnitudeDefine.E(0, false)(5)(ef);
byMagnitudeDefine.F(0, false)(1)(ef);
}
\problem{S}{e}{qualquer número de grandezas for equimúltiplo de tantas outras, cada uma de cada: qualquer que seja o múltiplo da primeira de sua parte, o mesmo múltiplo das primeiras grandezas tomadas em conjunto será o de todas as outras tomadas em conjunto.}

\begin{center}
Seja \drawMagnitude{A} o mesmo múltiplo de \drawMagnitude{B},\\
que \drawMagnitude{C} é de \drawMagnitude{D},\\
que \drawMagnitude{E} é de \drawMagnitude{F}.\\

Then is evident that\\
$\left.\vcenter{
\nointerlineskip\hbox{\drawMagnitude{A}}
\nointerlineskip\hbox{\drawMagnitude{C}}
\nointerlineskip\hbox{\drawMagnitude{E}}
}\right\}$
é o mesmo múltiplo de
$\left\{\vcenter{
\nointerlineskip\hbox{\drawMagnitude{B}}
\nointerlineskip\hbox{\drawMagnitude{D}}
\nointerlineskip\hbox{\drawMagnitude{F}}
}\right.$\\
qual \drawMagnitude{A} é de \drawMagnitude{B};

because there are as many magnitudes\\
in
$\left\{\vcenter{
\nointerlineskip\hbox{\drawMagnitude{A}}
\nointerlineskip\hbox{\drawMagnitude{C}}
\nointerlineskip\hbox{\drawMagnitude{E}}
}\right\}
=
\left\{\vcenter{
\nointerlineskip\hbox{\drawMagnitude{B}}
\nointerlineskip\hbox{\drawMagnitude{D}}
\nointerlineskip\hbox{\drawMagnitude{F}}
}\right.$\\
as there are in $\drawMagnitude{A} = \drawMagnitude{B}$.
\end{center}

A mesma demonstração é válida para qualquer número de magnitudes, que aqui foi aplicada a três.

$\therefore$ Se houver qualquer número de magnitudes, \&c.

\vfill\pagebreak

\starttheorem{Prop. II. Theor.}\label{prop:V.II}
\defineNewPicture{
byMagnitudeSymbolDefine.ab("circle", byyellow, 0);
byMagnitudeSymbolDefine.cd("sectorUp", byred, 1);
byMagnitudeSymbolDefine.e("circle", byblue, 0);
byMagnitudeSymbolDefine.f("sectorUp", byblack, 1);
byMagnitudeDefine.A(0, false)(3)(ab);
byMagnitudeDefine.B(0, false)(1)(ab);
byMagnitudeDefine.C(0, false)(3)(cd);
byMagnitudeDefine.D(0, false)(1)(cd);
byMagnitudeDefine.E(0, false)(4)(e);
byMagnitudeDefine.F(0, false)(4)(f);
}
\problem{S}{e}{a primeira grandeza for o mesmo múltiplo da segunda que a terceira da quarta, e a quinta for o mesmo múltiplo da segunda que a sexta é da quarta, então a primeira, juntamente com a quinta, será o mesmo múltiplo da segunda que a terceira, juntamente com a sexta, é da quarta.}

Seja \drawMagnitude{A}, o primeiro, o mesmo múltiplo de \drawMagnitude{B}, o segundo, que \drawMagnitude{C}, o terceiro, é de \drawMagnitude{D}, o quarto; e seja \drawMagnitude{E}, o quinto, o mesmo múltiplo de \drawMagnitude{B}, o segundo, que \drawMagnitude{F}, o sexto, é de \drawMagnitude{D}, o quarto.

Then it is evident, that
$\left\{\vcenter{
\nointerlineskip\hbox{\drawMagnitude{A}}
\nointerlineskip\hbox{\drawMagnitude{E}}
}\right\}$,
o primeiro e o quinto juntos, é o mesmo múltiplo de \drawMagnitude{B}, o segundo, que
$\left\{\vcenter{
\nointerlineskip\hbox{\drawMagnitude{C}}
\nointerlineskip\hbox{\drawMagnitude{F}}
}\right\}$,
o terceiro e o sexto juntos são do mesmo múltiplo de \drawMagnitude{D}, o quarto; porque existem tantas magnitudes em
$\left\{\vcenter{
\nointerlineskip\hbox{\drawMagnitude{A}}
\nointerlineskip\hbox{\drawMagnitude{E}}
}\right\} = \drawMagnitude{B}$
as there are in
$\left\{\vcenter{
\nointerlineskip\hbox{\drawMagnitude{C}}
\nointerlineskip\hbox{\drawMagnitude{F}}
}\right\} =\drawMagnitude{D}$.

$\therefore$ If the first magnitude, \&c.

\vfill\pagebreak

\starttheorem{Prop. III. Theor.}\label{prop:V.III}
\defineNewPicture{
byMagnitudeSymbolDefine.a("square", byyellow, 0);
byMagnitudeSymbolDefine.be("square", byred, 0);
byMagnitudeSymbolDefine.c("rhombus", byblack, 0);
byMagnitudeSymbolDefine.df("rhombus", byblue, 0);
byMagnitudeDefine.A(0, false)(1, 2, 1)(a);
byMagnitudeDefine.B(0, false)(1)(be);
byMagnitudeDefine.C(0, false)(2, 2)(c);
byMagnitudeDefine.D(0, false)(1)(df);
byMagnitudeDefine.E(0, false)(4, 4, 4, 4)(be);
byMagnitudeDefine.F(0, false)(4, 4, 4, 4)(df);
}
\problem{S}{e}{a primeira das quatro grandezas for o mesmo múltiplo da segunda que a terceira é para a quarta, e se quaisquer equimúltiplos da primeira e da terceira forem tomados, esses serão equimúltiplos; um do segundo e outro do quarto.}

\begin{center}
Seja $\left\{\drawMagnitude{A}\right\}$ o mesmo múltiplo de \drawMagnitude{B}\\
qual $\left\{\drawMagnitude{C}\right\}$ é de \drawMagnitude{D};\\
pegue $\left\{\drawMagnitude{E}\right\}$ o mesmo múltiplo de $\left\{\drawMagnitude{A}\right.$,\\
qual $\left\{\drawMagnitude{F}\right\}$ é de $\left\{\drawMagnitude{C}\right.$.

Then it is evident,\\
que $\left\{\drawMagnitude{E}\right\}$ é o mesmo múltiplo de \drawMagnitude{B}\\
which $\left\{\drawMagnitude{F}\right\}$ if of \drawMagnitude{D};\\
$\because \left\{\drawMagnitude{E}\right\}$ contains $\left\{\drawMagnitude{A}\right\}$ contains \drawMagnitude{B} \\
as many times as \\
$\left.\drawMagnitude{F}\right\}$ contains $\left\{\drawMagnitude{C}\right\}$ contains \drawMagnitude{D}.\\

O mesmo raciocínio é aplicável em todos os casos.

$\therefore$ Se for o primeiro de quatro, \&c.
\end{center}

\vfill\pagebreak

\startdefinition{Definição V.}\label{def:V.V} %ownnumber=5
\defineNewPicture{
byMagnitudeSymbolDefine.a("circle", byred, 0);
byMagnitudeSymbolDefine.b("square", byyellow, 0);
byMagnitudeSymbolDefine.c("rhombus", byblue, 0);
byMagnitudeSymbolDefine.d("wedgeDown", byblack, 0);
byMagnitudeDefine.A(0, false)(1)(a);
byMagnitudeDefine.B(0, false)(1)(b);
byMagnitudeDefine.C(0, false)(1)(c);
byMagnitudeDefine.D(0, false)(1)(d);
byMagnitudeDefine.Am(1, false)(2, 3, 4, 5, 6)(a);
byMagnitudeDefine.Bm(1, false)(2, 3, 4, 5, 6)(b);
byMagnitudeDefine.Cm(1, false)(2, 3, 4, 5, 6)(c);
byMagnitudeDefine.Dm(1, false)(2, 3, 4, 5, 6)(d);
}

Quatro magnitudes, \drawMagnitude{A}, \drawMagnitude{B}, \drawMagnitude{C}, \drawMagnitude{D}, são consideradas proporcionais quando cada equimúltiplo do primeiro e do terceiro for tomado, e cada equimúltiplo do segundo e do quarto, como,

\sepSpace

\vbox{\hfill\hbox{
\vbox{\hbox{of the first}\hbox{\drawMagnitude{Am}}\hbox{\&c.}}
\vbox{\hbox{of the third}\hbox{\drawMagnitude{Cm}}\hbox{\&c.}}
}\hfill\ }

\sepSpace

\vbox{\hfill\hbox{
\vbox{\hbox{of the second}\hbox{\drawMagnitude{Bm}}\hbox{\&c.}}
\vbox{\hbox{of the fourth}\hbox{\drawMagnitude{Dm}}\hbox{\&c.}}
}\hfill\ }

\sepSpace

Então, tomando cada par de equimúltiplos do primeiro e do terceiro, e cada par de equimúltiplos do segundo e do quarto,

\begin{center}
If
$\left\{\vcenter{
\nointerlineskip\hbox{$\drawMagnitude[middle][1]{Am} >, = \mbox{or} < \drawMagnitude[middle][1]{Bm}$}
\nointerlineskip\hbox{$\drawMagnitude[middle][1]{Am} >, = \mbox{or} < \drawMagnitude[middle][2]{Bm}$}
\nointerlineskip\hbox{$\drawMagnitude[middle][1]{Am} >, = \mbox{or} < \drawMagnitude[middle][3]{Bm}$}
\nointerlineskip\hbox{$\drawMagnitude[middle][1]{Am} >, = \mbox{or} < \drawMagnitude[middle][4]{Bm}$}
\nointerlineskip\hbox{$\drawMagnitude[middle][1]{Am} >, = \mbox{or} < \drawMagnitude[middle][5]{Bm}$}
}\right.$

$\vcenter{\hbox{then}\hbox{will}}\left\{\vcenter{
\nointerlineskip\hbox{$\drawMagnitude[middle][1]{Cm} >, = \mbox{or} < \drawMagnitude[middle][1]{Dm}$}
\nointerlineskip\hbox{$\drawMagnitude[middle][1]{Cm} >, = \mbox{or} < \drawMagnitude[middle][2]{Dm}$}
\nointerlineskip\hbox{$\drawMagnitude[middle][1]{Cm} >, = \mbox{or} < \drawMagnitude[middle][3]{Dm}$}
\nointerlineskip\hbox{$\drawMagnitude[middle][1]{Cm} >, = \mbox{or} < \drawMagnitude[middle][4]{Dm}$}
\nointerlineskip\hbox{$\drawMagnitude[middle][1]{Cm} >, = \mbox{or} < \drawMagnitude[middle][5]{Dm}$}
}\right.$
\end{center}

Ou seja, se o dobro do primeiro for maior, igual ou menor que o dobro do segundo, o dobro do terceiro será maior, igual ou menor que o dobro do quarto; ou, se duas vezes o primeiro for maior, igual ou menor que três vezes o segundo, duas vezes o terceiro será maior, igual ou menor que três vezes o quarto, e assim por diante, conforme expresso acima.

\begin{center}
If
$\left\{\vcenter{
\nointerlineskip\hbox{$\drawMagnitude[middle][2]{Am} >, = \mbox{or} < \drawMagnitude[middle][1]{Bm}$}
\nointerlineskip\hbox{$\drawMagnitude[middle][2]{Am} >, = \mbox{or} < \drawMagnitude[middle][2]{Bm}$}
\nointerlineskip\hbox{$\drawMagnitude[middle][2]{Am} >, = \mbox{or} < \drawMagnitude[middle][3]{Bm}$}
\nointerlineskip\hbox{$\drawMagnitude[middle][2]{Am} >, = \mbox{or} < \drawMagnitude[middle][4]{Bm}$}
\nointerlineskip\hbox{$\drawMagnitude[middle][2]{Am} >, = \mbox{or} < \drawMagnitude[middle][5]{Bm}$}
}\right.$

$\vcenter{\hbox{then}\hbox{will}}\left\{\vcenter{
\nointerlineskip\hbox{$\drawMagnitude[middle][2]{Cm} >, = \mbox{or} < \drawMagnitude[middle][1]{Dm}$}
\nointerlineskip\hbox{$\drawMagnitude[middle][2]{Cm} >, = \mbox{or} < \drawMagnitude[middle][2]{Dm}$}
\nointerlineskip\hbox{$\drawMagnitude[middle][2]{Cm} >, = \mbox{or} < \drawMagnitude[middle][3]{Dm}$}
\nointerlineskip\hbox{$\drawMagnitude[middle][2]{Cm} >, = \mbox{or} < \drawMagnitude[middle][4]{Dm}$}
\nointerlineskip\hbox{$\drawMagnitude[middle][2]{Cm} >, = \mbox{or} < \drawMagnitude[middle][5]{Dm}$}
}\right.$
\end{center}

Em outros termos, se três vezes o primeiro for maior, igual ou menor que duas vezes o segundo, três vezes o terceiro será maior, igual ou menor que duas vezes o quarto; ou, se três vezes o primeiro for maior, igual ou menor que três vezes o segundo, então três vezes o terceiro será maior, igual ou menor que três vezes o quarto; ou se três vezes o primeiro for maior, igual ou menor que quatro vezes o segundo, então três vezes o terceiro será maior, igual ou menor que quatro vezes o quarto, e assim por diante. De novo,

\begin{center}
If
$\left\{\vcenter{
\nointerlineskip\hbox{$\drawMagnitude[middle][3]{Am} >, = \mbox{or} < \drawMagnitude[middle][1]{Bm}$}
\nointerlineskip\hbox{$\drawMagnitude[middle][3]{Am} >, = \mbox{or} < \drawMagnitude[middle][2]{Bm}$}
\nointerlineskip\hbox{$\drawMagnitude[middle][3]{Am} >, = \mbox{or} < \drawMagnitude[middle][3]{Bm}$}
\nointerlineskip\hbox{$\drawMagnitude[middle][3]{Am} >, = \mbox{or} < \drawMagnitude[middle][4]{Bm}$}
\nointerlineskip\hbox{$\drawMagnitude[middle][3]{Am} >, = \mbox{or} < \drawMagnitude[middle][5]{Bm}$}
}\right.$

$\vcenter{\hbox{then}\hbox{will}}\left\{\vcenter{
\nointerlineskip\hbox{$\drawMagnitude[middle][3]{Cm} >, = \mbox{or} < \drawMagnitude[middle][1]{Dm}$}
\nointerlineskip\hbox{$\drawMagnitude[middle][3]{Cm} >, = \mbox{or} < \drawMagnitude[middle][2]{Dm}$}
\nointerlineskip\hbox{$\drawMagnitude[middle][3]{Cm} >, = \mbox{or} < \drawMagnitude[middle][3]{Dm}$}
\nointerlineskip\hbox{$\drawMagnitude[middle][3]{Cm} >, = \mbox{or} < \drawMagnitude[middle][4]{Dm}$}
\nointerlineskip\hbox{$\drawMagnitude[middle][3]{Cm} >, = \mbox{or} < \drawMagnitude[middle][5]{Dm}$}
}\right.$
\end{center}

E assim por diante, com quaisquer outros equimúltiplos das quatro grandezas, tomadas da mesma maneira.

\vskip \baselineskip

Euclid expresses this definition as follows :—

Diz-se que a primeira das quatro grandezas tem a mesma razão para a segunda, que a terceira tem para a quarta, quando são tomados quaisquer equimúltiplos da primeira e da terceira, e quaisquer equimúltiplos da segunda e da quarta; se o múltiplo do primeiro for menor que o do segundo e o múltiplo do terceiro for também menor que o do quarto; ou, se o múltiplo do primeiro for igual ao do segundo, o múltiplo do terceiro será também igual ao do quarto; ou, se o múltiplo do primeiro for maior que o do segundo, o múltiplo do terceiro também será maior que o do quarto.

In future we shall express this definition generally, thus:

\begin{center}
$\begin{aligned}
	\mbox{Se }M \drawMagnitude{A} &>, = \mbox{ ou } < m \drawMagnitude{B} \mbox{,}\\
	\mbox{when }M \drawMagnitude{C} &>, = \mbox{ ou } < m \drawMagnitude{D} \mbox{,}\\
\end{aligned}$
\end{center}

Então inferimos que \drawMagnitude{A}, o primeiro, tem a mesma proporção do \drawMagnitude{B}, o segundo, que \drawMagnitude{C}, o terceiro, tem para \drawMagnitude{D} o quarto; expresso nas demonstrações seguintes assim:

\begin{center}
$\begin{aligned}
\drawMagnitude{A} : \drawMagnitude{B} & :: \drawMagnitude{C} : \drawMagnitude{D} \mbox{;} \\
\mbox{or thus, } \drawMagnitude{A} : \drawMagnitude{B} &= \drawMagnitude{C} : \drawMagnitude{D} \mbox{;} \\
\mbox{or thus, } \dfrac{\drawMagnitude{A}}{\drawMagnitude{B}} &= \dfrac{\drawMagnitude{C}}{\drawMagnitude{D}} \mbox{:} \\
\end{aligned}$

e é lido, \\
\enquote{assim como \drawMagnitude{A} está para \drawMagnitude{B}, \drawMagnitude{C} está para \drawMagnitude{D}.}

And if $\drawMagnitude{A} : \drawMagnitude{B} :: \drawMagnitude{C} : \drawMagnitude{D}$ we shall infer if \\
$M \drawMagnitude{A} >, = \mbox{ ou } < m \drawMagnitude{B}$, then will \\
$M \drawMagnitude{C} >, = \mbox{ ou } < m \drawMagnitude{D}$.
\end{center}

Isto é, se o primeiro está para o segundo, assim como o terceiro está para o quarto; então, se $M$ vezes o primeiro for maior, igual ou menor que $m$ vezes o segundo, então $M$ vezes o terceiro será maior, igual ou menor que $m$ vezes o quarto, em que $M$ e $m$ não devem ser considerados múltiplos particulares, mas qualquer par de múltiplos; nem marcas como \drawMagnitude{A}, \drawMagnitude{D}, \drawMagnitude{B}, \&c.\ devem ser consideradas mais do que representantes de magnitudes geométricas.

The student should throughly understand this definition before proceeding further.

\vfill\pagebreak

\starttheorem{Prop. IV. Theor.}\label{prop:V.IV}
\defineNewPicture{
byMagnitudeSymbolDefine.I("circle", byyellow, 0);
byMagnitudeSymbolDefine.II("square", byblack, 0);
byMagnitudeSymbolDefine.III("rhombus", byred, 0);
byMagnitudeSymbolDefine.IV("wedgeDown", byblue, 0);
}
\problem{S}{e}{a primeira das quatro grandezas tiver a mesma razão para a segunda, que a terceira tem para a quarta, então quaisquer equimúltiplos da primeira e da terceira terão a mesma razão para quaisquer equimúltiplos da segunda e da quarta; a saber, o equimúltiplo do primeiro terá a mesma razão para o do segundo, que o equimúltiplo do terceiro tem para o do quarto.}

Seja $\drawMagnitude{I} : \drawMagnitude{II} :: \drawMagnitude{III} : \drawMagnitude{IV}$, então $3\drawMagnitude{I} : 2\drawMagnitude{II} :: 3\drawMagnitude{III} : 2\drawMagnitude{IV}$, cada equimúltiplo de $3\drawMagnitude{I}$ e $3\drawMagnitude{III}$ são equimúltiplos de \drawMagnitude{I} e \drawMagnitude{III}, e cada equimúltiplo de $2\drawMagnitude{II}$ e $2\drawMagnitude{IV}$, são equimúltiplos de \drawMagnitude{II} e \drawMagnitude{IV} \byref{prop:V.III}

Ou seja, $M$ vezes $3\drawMagnitude{I}$ e $M$ vezes $3\drawMagnitude{III}$ são equimúltiplos de \drawMagnitude{I} e \drawMagnitude{III}, e $m$ vezes $2\drawMagnitude{II}$ e $m 2\drawMagnitude{IV}$ são equimúltiplos de $2\drawMagnitude{II}$ e $2\drawMagnitude{IV}$; mas $\drawMagnitude{I} : \drawMagnitude{II} :: \drawMagnitude{III} : \drawMagnitude{IV}$ \byref{\hypref}; $\therefore$ se $M 3\drawMagnitude{I} <, =, \mbox{ or } > m 2 \drawMagnitude{II}$, então $M 3\drawMagnitude{III} <, =, \mbox{ or } > m 2 \drawMagnitude{IV}$ \byref{def:V.V} e portanto $3\drawMagnitude{I} : 2\drawMagnitude{II} :: 3\drawMagnitude{III} : 2\drawMagnitude{IV}$ \byref{def:V.V}

O mesmo raciocínio é válido se qualquer outro equimúltiplo do primeiro e do terceiro for tomado, qualquer outro equimúltiplo do segundo e do quarto.

$\therefore$ If the first four magnitudes, \&c.

\vfill\pagebreak

\starttheorem{Prop. V. Theor.}\label{prop:V.V}
\defineNewPicture{
byMagnitudeSymbolDefine.a("sectorDown", byblue, 1);
byMagnitudeSymbolDefine.b("semicircleDown", byyellow, 1);
byMagnitudeSymbolDefine.c("miniTriangleUp", byblack, 0);
byMagnitudeSymbolDefine.d("miniSquare", byred, 0);
byMagnitudeDefine.I(0, false)(1, 2, 1)(a, a, b);
byMagnitudeDefine.II(0, false)(1, 1)(c, d);
}
\problem{S}{e}{uma grandeza for o mesmo múltiplo de outra, e uma grandeza tirada da primeira for uma grandeza tirada da outra, o resto será o mesmo múltiplo do resto, de modo que o todo é do todo.}

\begin{center}
Let $\drawMagnitude{I} = M' \drawMagnitude{II}$\\
and $\drawMagnitude[middle][3]{I} = M' \drawMagnitude[middle][2]{II}$,

$\therefore \drawMagnitude{I} - \drawMagnitude[middle][3]{I} = M' \drawMagnitude{II} - M' \drawMagnitude[middle][2]{II}$, % two times \mbox{ minus }

$\therefore \drawMagnitude[middle][-3]{I} = M' (\drawMagnitude{II} - \drawMagnitude[middle][2]{II})$,\\ %\mbox{ minus }
and $\therefore \drawMagnitude[middle][-3]{I} = M' \drawMagnitude[middle][-2]{II}$.

$\therefore$ If one magnitude, \&c.
\end{center}

\vfill\pagebreak

\starttheorem{Prop. VI. Theor.}\label{prop:V.VI}
\defineNewPicture{
byMagnitudeSymbolDefine.a("sectorDown", byyellow, 1);
byMagnitudeSymbolDefine.b("miniSquare", byred, 0);
byMagnitudeSymbolDefine.c("semicircleDown", byblack, 1);
byMagnitudeSymbolDefine.d("miniTriangleUp", byblue, 0);
byMagnitudeDefine.I(0, false)(1, 2, 1)(a);
byMagnitudeDefine.II(0, false)(1)(b);
byMagnitudeDefine.III(0, false)(2)(c);
byMagnitudeDefine.IV(0, false)(1)(d);
}
\problem{S}{e}{duas grandezas são equimúltiplos de outras duas, e se equimúltiplos destas forem tomados das duas primeiras, os restos ou serão iguais a estas outras, ou equimúltiplos delas.}

\begin{center}
Seja $\drawMagnitude{I} = M'\drawMagnitude{II}$; e $\drawMagnitude{III} = M'\drawMagnitude{IV}$

then $\drawMagnitude{I} - m'\drawMagnitude{II} = M'\drawMagnitude{II} -m'\drawMagnitude{II} = (M' - m') \drawMagnitude{II}$\\ % three times \mbox{ minus }
and $\drawMagnitude{III} - m'\drawMagnitude{IV} = M'\drawMagnitude{IV} - m'\drawMagnitude{IV} = (M' - m') \drawMagnitude{IV}$. %three times \mbox{ minus }

Portanto, $(M' - m') \drawMagnitude{II}$ e $(M' - m') \drawMagnitude{IV}$ são equimúltiplos de \drawMagnitude{II} e \drawMagnitude{IV}, e iguais a \drawMagnitude{II} e \drawMagnitude{IV}, quando $M' - m' = 1$. %três vezes \mbox{ menos }

$\therefore$ Se duas magnitudes forem equimúltiplas, \&c.
\end{center}

\vfill\pagebreak

\starttheoremAZ{Prop. A. Theor.}\label{prop:V.A} % Proposition with letters
\defineNewPicture{
byMagnitudeSymbolDefine.a("circle", byred, 0);
byMagnitudeSymbolDefine.b("square", byblack, 0);
byMagnitudeSymbolDefine.c("wedgeDown", byblue, 0);
byMagnitudeSymbolDefine.d("rhombus", byyellow, 0);
byMagnitudeDefine.I(0, false)(1)(a);
byMagnitudeDefine.II(0, false)(1)(b);
byMagnitudeDefine.III(0, false)(1)(c);
byMagnitudeDefine.IV(0, false)(1)(d);
byMagnitudeDefine.dI(0, false)(2)(a);
byMagnitudeDefine.dII(0, false)(2)(b);
byMagnitudeDefine.dIII(0, false)(2)(c);
byMagnitudeDefine.dIV(0, false)(2)(d);
}
\problem{S}{e}{a primeira das quatro grandezas tiver com a segunda a mesma razão que a terceira tem com a quarta, então, se a primeira for maior que a segunda, a terceira também será maior que a quarta; e se for igual, igual; se menos, menos.}

\begin{center}
Deixe $\drawMagnitude{I} : \drawMagnitude{II} :: \drawMagnitude{III} : \drawMagnitude{IV}$; portanto, pela quinta definição, se $\drawMagnitude{dI} > \drawMagnitude{dII}$, então será $\drawMagnitude{dIII} > \drawMagnitude{dIV}$;

mas se $\drawMagnitude{I} > \drawMagnitude{II}$, então $\drawMagnitude{dI} > \drawMagnitude{dII}$ e $\drawMagnitude{dIII} > \drawMagnitude{dIV}$,\\
and $\therefore \drawMagnitude{III} > \drawMagnitude{IV}$.

Similarly, if $\drawMagnitude{I} =, \mbox{ ou } < \drawMagnitude{II}$, then will $\drawMagnitude{III} =, \mbox{ ou } < \drawMagnitude{IV}$.

$\therefore$ Se for o primeiro de quatro, \&c.
\end{center}

\startdefinition[14]{Definição XIV}\label{def:V.XIV}
\def\varA{\textcolor{byred}{A}}
\def\varB{\textcolor{byblack}{B}}
\def\varC{\textcolor{byblue}{C}}
\def\varD{\textcolor{byyellow}{D}}
Os geômetras fazem uso do termo técnico \enquote{Invertendo,} por inversão, quando há quatro proporcionais, e infere-se que a segunda está para a primeira como a quarta está para a terceira.

Seja $\varA : \varB :: \varC : \varD$, então, por \enquote{invertendo} infere-se $\varB : \varA :: \varD : \varC$

\vfill\pagebreak

\starttheoremAZ{Prop. B. Theor.}\label{prop:V.B} % Proposition with letters
\defineNewPicture{
byMagnitudeSymbolDefine.I("wedgeDown", byblue, 0);
byMagnitudeSymbolDefine.II("semicircleDown", byblack, 1);
byMagnitudeSymbolDefine.III("square", byred, 0);
byMagnitudeSymbolDefine.IV("rhombus", byyellow, 0);
}
\problem{S}{e}{quatro grandezas são proporcionais, elas também o são quando tomadas inversamente.}

\begin{center}
Let $\drawMagnitude{I} : \drawMagnitude{II} :: \drawMagnitude{III} : \drawMagnitude{IV}$,\\
then, inversely, $\magnitudeII : \magnitudeI :: \magnitudeIV : \magnitudeIII$.

If $M \magnitudeI < m \magnitudeII$, then $M \magnitudeIII < m \magnitudeIV$\\
by the fifth definition.

Let $M \magnitudeI < m \magnitudeII$, that is, $m \magnitudeII > M \magnitudeI$,\\
$\therefore M \magnitudeIII < m \magnitudeIV$, or, $m \magnitudeIV > M \magnitudeIII$;

$\therefore$ if $m \magnitudeII > M \magnitudeI$, then will $m \magnitudeIV > M \magnitudeIII$.

In the same manner it may be shown,\\
that if $m \magnitudeII = \mbox{ ou } < M \magnitudeI$,\\
then will $m \magnitudeIV =, \mbox{ ou } < M \magnitudeIII$;

and therefore, by the fifth definition, we infer\\
that $\magnitudeII : \magnitudeI :: \magnitudeIV : \magnitudeIII$.

$\therefore$ If four magnitudes, \&c.
\end{center}

\vfill\pagebreak

\starttheoremAZ{Prop. C. Theor.},\label{prop:V.C} % Proposition with letters
\defineNewPicture{
byMagnitudeSymbolDefine.i("square", byblue, 0);
byMagnitudeDefine.I(0, false)(2, 2)(i);
byMagnitudeSymbolDefine.II("circle", byblack, 0);
byMagnitudeSymbolDefine.iii("rhombus", byyellow, 0);
byMagnitudeDefine.III(0, false)(2, 2)(iii);
byMagnitudeSymbolDefine.IV("wedgeUp", byred, 0);
}
\problem{S}{e}{o primeiro for o mesmo múltiplo do segundo, ou a mesma parte dele, o terceiro é do quarto; o primeiro está para o segundo, assim como o terceiro está para o quarto.}

\begin{center}
Seja \drawMagnitude{I}, o primeiro, o mesmo múltiplo de \drawMagnitude{II}, o segundo, \\
aquele \drawMagnitude{III}, o terceiro, é do \drawMagnitude{IV}, o quarto.

Then $\magnitudeI : \magnitudeII :: \magnitudeIII : \magnitudeIV$

take $M \magnitudeI$, $m \magnitudeII$, $M \magnitudeIII$, $m \magnitudeIV$;

porque \magnitudeI\ é o mesmo múltiplo de \magnitudeII \\
que \magnitudeIII é de \magnitudeIV\ (de acordo com a hipótese);

e $M \magnitudeI$ se for tomado o mesmo múltiplo de \magnitudeI \\
que $M \magnitudeIII$ é de \magnitudeIII,

$\therefore$ (de acordo com a terceira proposição),\\
$M \magnitudeI$ é o mesmo múltiplo de \magnitudeII \\
que $M \magnitudeIII$ é de \magnitudeIV.

Portanto, se $M \magnitudeI$ for de \magnitudeII\ um múltiplo maior que $m \magnitudeII$ é, \\
então $M \magnitudeIII$ é um múltiplo maior de \magnitudeIV\ do que $m \magnitudeIV$;

that is, if $M \magnitudeI$ be greater than $m \magnitudeII$, then $M \magnitudeIII$ will be greater than $m \magnitudeIV$;

da mesma maneira que pode ser mostrado, se $M \magnitudeI$ for igual a $m \magnitudeII$, então $M \magnitudeIII$ será igual a $m \magnitudeIV$.

And, generally, if $M \magnitudeI >, = \mbox{ ou } < m \magnitudeII$\\
than $M \magnitudeIII \mbox{ will be } >, = \mbox{ ou } < m \magnitudeIV$;

$\therefore$ by the fifth definition,\\
$\magnitudeI : \magnitudeII :: \magnitudeIII : \magnitudeIV$

A seguir, seja \magnitudeII a mesma parte de \magnitudeI \\
que \magnitudeIV\ é de \magnitudeIII.

In this case also $\magnitudeII : \magnitudeI :: \magnitudeIV : \magnitudeIII$.

Pois, porque \magnitudeII\ é a mesma parte de \magnitudeI\ que \magnitudeIV\ é de \magnitudeIII,

portanto \magnitudeI\ é o mesmo múltiplo de \magnitudeII \\
que \magnitudeIII\ é de \magnitudeIV.

Therefore, by the preceding case, \\
$\magnitudeI : \magnitudeII :: \magnitudeIII : \magnitudeIV$;

and $\therefore \magnitudeII : \magnitudeI :: \magnitudeIV : \magnitudeIII$, \\
by proposition B.

$\therefore$ Se o primeiro for o mesmo múltiplo, \&c.
\end{center}

\vfill\pagebreak

\starttheoremAZ{Prop. D. Theor.}\label{prop:V.D}
\defineNewPicture{
byMagnitudeSymbolDefine.i("circle", byyellow, 0);
byMagnitudeSymbolDefine.ii("square", byblack, 0);
byMagnitudeSymbolDefine.iii("rhombus", byred, 0);
byMagnitudeSymbolDefine.iv("wedgeDown", byblue, 0);
byMagnitudeSymbolDefine.va("semicircleDown", byred, 1);
byMagnitudeSymbolDefine.vb("semicircleUp", byred, 1);
byMagnitudeSymbolDefine.via("sectorDown", byblack, 1);
byMagnitudeSymbolDefine.vib("sectorUp", byblack, 1);
byMagnitudeDefine.I(0, false)(1, 2)(i);
byMagnitudeDefine.II(0, false)(1)(ii);
byMagnitudeDefine.III(0, false)(2, 2)(iii);
byMagnitudeDefine.IV(0, false)(1)(iv);
byMagnitudeDefine.V(0, false)(1, 2)(va, vb);
byMagnitudeDefine.VI(0, false)(2, 2)(via, vib);
}
\problem{S}{e}{o primeiro estiver para o segundo como o terceiro para o quarto, e se o primeiro for um múltiplo ou parte do segundo; o terceiro é o mesmo múltiplo ou a mesma parte do quarto.}

\begin{center}
Let $\drawMagnitude{I} : \drawMagnitude{II} :: \drawMagnitude{III} : \drawMagnitude{IV}$;\\
e primeiro, seja \magnitudeI\ um múltiplo \magnitudeII;\\
\magnitudeIII\ será o mesmo múltiplo de \magnitudeIV.

%\unprotect
%\ \hfill\vbox{\halign{\hfil # \hfil & \hfil # \hfil & \hfil # \hfil & \hfil # \hfil \\
\begin{tabular}{c c c c}
{\footnotesize First} & {\footnotesize Second} & {\footnotesize Third} & {\footnotesize Fourth} \\
\magnitudeI\ & \magnitudeII\ & \magnitudeIII\ & \magnitudeIV \\
&\drawMagnitude{V}\ &\drawMagnitude{VI}\ & \\
\end{tabular}
%}}\hfill\
%\protect

Take $\magnitudeV = \magnitudeI$.

Qualquer miltiple \magnitudeI\ é de \magnitudeII\\
tome \magnitudeVI\ o mesmo múltiplo de \magnitudeIV, \\
then, $\because \magnitudeI : \magnitudeII :: \magnitudeIII : \magnitudeIV$\\ %because
e do segundo e do quarto, tomamos equimúltiplos,\\
\magnitudeI\ and \magnitudeVI\ therefore \byref{prop:V.IV}\\
 $\magnitudeI : \magnitudeV :: \magnitudeIII : \magnitudeVI$, but \byref{\constref},\\
 $\magnitudeI = \magnitudeV \therefore$ \byref{prop:V.A} $\magnitudeIII = \magnitudeVI$\\
 e \magnitudeVI\ é o mesmo múltiplo de \magnitudeIV\\
 que \magnitudeI\ é de \magnitudeII.

 Next, let $\magnitudeII : \magnitudeI :: \magnitudeIV : \magnitudeIII$,\\
 e também \magnitudeII\ uma parte de \magnitudeI;\\
 então \magnitudeIV\ será a mesma parte de \magnitudeIII.

Inversamente \byref{prop:V.B}, $\magnitudeI : \magnitudeII :: \magnitudeIII : \magnitudeIV$,\\ % melhoria: no original havia uma referência a "B. 5." o que não faz muito sentido
mas \magnitudeII\ faz parte de \magnitudeI;\\
isto é, \magnitudeI\ é um múltiplo de \magnitudeII;\\
$\therefore$ pelo caso anterior, \magnitudeIII\ é o mesmo múltiplo de \magnitudeIV\\
isto é, \magnitudeIV\ é a mesma parte de \magnitudeIII\\
que \magnitudeII\ é de \magnitudeI.

$\therefore$ se o primeiro for igual ao segundo, \&c.
\end{center}

\vfill\pagebreak

\starttheorem{Prop. VII. Theor.}\label{prop:V.VII}
\defineNewPicture{
byMagnitudeSymbolDefine.I("circle", byred, 0);
byMagnitudeSymbolDefine.II("rhombus", byblue, 0);
byMagnitudeSymbolDefine.III("square", byyellow, 0);
}
\problem{M}{agnitudes}{iguais têm a mesma razão para a mesma magnitude, e as mesmas têm a mesma razão para magnitudes iguais.}

\begin{center}
Seja $\drawMagnitude{I} = \drawMagnitude{II}$ e \drawMagnitude{III} qualquer outra magnitude;\\
then $\magnitudeI : \magnitudeIII = \magnitudeII : \magnitudeIII$ and $\magnitudeIII : \magnitudeI = \magnitudeIII : \magnitudeII$.

$\because \magnitudeI = \magnitudeII$,\\ %Because
$\therefore M \magnitudeI = M \magnitudeII$;

$\therefore$ if $M \magnitudeI >, = \mbox{ ou } < m \magnitudeIII$, then\\
$M \magnitudeII >, = \mbox{ ou } < m \magnitudeIII$,\\
and $\therefore \magnitudeI : \magnitudeIII = \magnitudeII : \magnitudeIII$ \byref{def:V.V}.

Do raciocínio acima fica evidente que,\\
if $m \magnitudeIII >, = \mbox{ ou } < M \magnitudeI$, then\\
$m \magnitudeIII >, = \mbox{ ou } < M \magnitudeII$\\
$\therefore \magnitudeIII : \magnitudeI = \magnitudeIII : \magnitudeII$ \byref{def:V.V}.

$\therefore$ Equal magnitudes, \&c.
\end{center}

\vfill\pagebreak

\startdefinition{Definição VII.}\label{def:V.VII} %ownnumber=7
\defineNewPicture{
magnitudeScale := 4/5;
byMagnitudeSymbolDefine.i("circle", byred, 0);
byMagnitudeSymbolDefine.ii("square", byyellow, 0);
byMagnitudeSymbolDefine.iii("rhombus", byblue, 0);
byMagnitudeSymbolDefine.iv("wedgeDown", byblack, 0);
byMagnitudeDefine.I(0, false)(5)(i);
byMagnitudeDefine.II(0, false)(4)(ii);
byMagnitudeDefine.III(0, false)(5)(iii);
byMagnitudeDefine.IV(0, false)(4)(iv);
byMagnitudeSymbolDefine.Ia("wedgeDown", byred, 0);
byMagnitudeSymbolDefine.IIa("semicircleDown", byblack, 1);
byMagnitudeSymbolDefine.IIIa("square", byblue, 0);
byMagnitudeSymbolDefine.IVa("rhombus", byyellow, 0);
}

Quando dos equimúltiplos de quatro grandezas (tomadas como na quinta definição), o múltiplo da primeira é maior que o da segunda, mas o múltiplo da terceira não é maior que o múltiplo da quarta; então diz-se que a primeira tem para a segunda uma razão maior do que a terceira grandeza tem para a quarta: e, pelo contrário, diz-se que a terceira tem para a quarta uma razão menor do que a primeira tem para a segunda.

Se, entre os equimúltiplos de quatro grandezas, comparados na quinta definição, encontrarmos $\drawMagnitude{I} > \drawMagnitude{II}$, mas $\drawMagnitude{III} = \mbox{ or } > \drawMagnitude{IV}$, ou se encontrarmos qualquer múltiplo particular $M'$ da primeira e terceira, e um múltiplo particular $m'$ da segunda e quarta, tal, que $M'$ vezes o primeiro é $> m'$ vezes o segundo, mas $M'$ vezes o terceiro não é $> m'$ vezes o quarto, ou seja.\ e.\ $= \mbox{ or } < m'$ vezes o quarto; então diz-se que o primeiro tem para o segundo uma proporção maior do que o terceiro para o quarto; ou a terceira tem para a quarta, sob tais circunstâncias, uma razão menor do que a primeira tem para a segunda: embora vários outros equimúltiplos possam tender a mostrar que as quatro grandezas são proporcionais.

This definition will in future be expressed thus :—

\begin{center}
If $M' \drawMagnitude{Ia} > m' \drawMagnitude{IIa}$, but $M' \drawMagnitude{IIIa} = \mbox{ ou } < m' \drawMagnitude{IVa}$,\\
then $\magnitudeIa : \magnitudeIIa > \magnitudeIIIa : \magnitudeIVa$.
\end{center}

Na expressão geral acima, $M'$ e $m'$ devem ser considerados múltiplos particulares, não como os múltiplos $M$ e $m$ introduzidos na quinta definição, que nessa definição são considerados todos os pares de múltiplos que podem ser tomados. Também deve ser observado aqui que \magnitudeIa, \magnitudeIIa, \magnitudeIIIa e símbolos semelhantes devem ser considerados meramente representantes de magnitudes geométricas.

In a partial arithmetical way, this may be set forth as follows:

Tomemos os quatro números $\textcolor{byred}{8}$, $\textcolor{byblack}{7}$, $\textcolor{byblue}{10}$ e $\textcolor{byyellow}{9}$.

\begin{center}
%\unprotect
%\ \hfill\vbox{
%\offinterlineskip
%\tabskip=0pt
%\halign{\vrule height1.9ex depth0.9ex
%	\hfil \textcolor{byred} { # } \hfil &
%	\hfil \textcolor{byblack} { # } \hfil &
%	\hfil \textcolor{byblue} { # } \hfil &
%	\hfil \textcolor{byyellow} { # } \hfil \vrule \\
%
\begin{tabular}{>{\color{byred}}c >{\color{byblack}}c >{\color{byblue}}c >{\color{byyellow}}c}
	\noalign{\hrule}
	\textcolor{byblack}{First}&
	\textcolor{byblack}{Second}&
	\textcolor{byblack}{Third}&
	\textcolor{byblack}{Fourth}\\
	8 	& 7 		& 10 	& 9 \\
	\noalign{\hrule}
	16 	& 14 	& 20 	& 18 \\
	24 	& 21 	& 30 	& 27 \\
	32 	& 28 	& 40 	& 36 \\
	40 	& 35 	& 50 	& 45 \\
	48 	& 42 	& 60 	& 54 \\
	56 	& 49 	& 70 	& 63 \\
	64 	& 56 	& 80 	& 72 \\
	72 	& 63 	& 90 	& 81 \\
	80 	& 70 	& 100 	& 90 \\
	88 	& 77 	& 110 	& 99 \\
	96 	& 84 	& 120 	& 108 \\
	104 	& 91 	& 130 	& 117 \\
	112 	& 98 	& 140 	& 126 \\
	\textcolor{byblack}{ \&c.} &
	\textcolor{byblack}{ \&c.} &
	\textcolor{byblack}{ \&c.} &
	\textcolor{byblack}{ \&c.} \\
	\noalign{\hrule}
\end{tabular}
%}}\hfill\
%\protect
\end{center}

Entre os múltiplos acima encontramos $\textcolor{byred}{16} > \textcolor{byblack}{14}$ e $\textcolor{byblue}{20} > \textcolor{byyellow}{18}$; isto é, duas vezes o primeiro é maior que duas vezes o segundo, e duas vezes o terceiro é maior que duas vezes o quarto; e $\textcolor{byred}{16} < \textcolor{byblack}{21}$ e $\textcolor{byblue}{20} < \textcolor{byyellow}{27}$; isto é, duas vezes o primeiro é menor que três vezes o segundo, e duas vezes o terceiro é menor que três vezes o quarto; e entre os mesmos múltiplos podemos encontrar $\textcolor{byred}{72} > \textcolor{byblack}{56}$ e $\textcolor{byblue}{90} > \textcolor{byyellow}{72}$; isto é, 9 vezes o primeiro é maior que 8 vezes o segundo e 9 vezes o terceiro é maior que 8 vezes o quarto. Muitos outros equimúltiplos poderiam ser selecionados, o que tenderia a mostrar que os números $\textcolor{byred}{8}$, $\textcolor{byblack}{7}$, $\textcolor{byblue}{10}$ e $\textcolor{byyellow}{9}$ eram proporcionais, mas não são, pois podemos encontrar múltiplos do primeiro $>$ um múltiplo do segundo, mas o mesmo múltiplo do terceiro que foi tirado do primeiro, não $>$ o mesmo múltiplo do quarto que foi levado para o segundo; por exemplo, 9 vezes o primeiro é $>$ 10 vezes o segundo, mas 9 vezes o terceiro não é $>$ 10 vezes o quarto, ou seja, $\textcolor{byred}{72} > \textcolor{byblack}{70}$, mas $\textcolor{byblue}{90} \ngtr \textcolor{byyellow}{90}$, ou 8 vezes o primeiro encontramos $>$ 9 vezes o segundo, mas 8 vezes o terceiro não é maior que 9 vezes o quarto, ou seja, $\textcolor{byred}{64} > \textcolor{byblack}{63}$, mas $\textcolor{byblue}{80} \mbox{ not } > \textcolor{byyellow}{81}$. Quando quaisquer múltiplos como esses podem ser encontrados, diz-se que o primeiro $\textcolor{byred}{(8)}$ tem o segundo $\textcolor{byblack}{(7)}$ uma proporção maior do que o terceiro $\textcolor{byblue}{(10)}$ tem para o quarto $\textcolor{byyellow}{(9)}$ e, pelo contrário, diz-se que o terceiro $\textcolor{byblue}{(10)}$ tem para o quarto $\textcolor{byyellow}{(9)}$ uma proporção menor do que o primeiro $\textcolor{byred}{(8)}$ tem para o segundo $\textcolor{byblack}{(7)}$.

\vfill\pagebreak

\starttheorem{Prop. VIII. Theor.}\label{prop:V.VIII}
\defineNewPicture{
byMagnitudeSymbolDefine.Ia("miniTriangleUp", byblack, 0);
byMagnitudeSymbolDefine.Ib("square", byred, 0);
byMagnitudeDefine.I(0, false)(1, 1)(Ia, Ib);
byMagnitudeSymbolDefine.II("square", byyellow, 0);
byMagnitudeSymbolDefine.III("circle", byblue, 0);
}
\problem{D}{e}{grandezas desiguais, a maior tem uma razão maior para o mesmo do que a menor: e a mesma grandeza tem uma razão maior para o menor do que para o maior.}

\begin{center}
Sejam \drawMagnitude[bottom]{I} e \drawMagnitude{II} duas magnitudes desiguais,\\
and \drawMagnitude{III} any other.
\end{center}

Provaremos primeiro que a \magnitudeI\ que é a maior das duas grandezas desiguais, tem uma razão maior para a \magnitudeIII\ do que a \magnitudeII, a menor, tem para a \magnitudeIII; isto é, $\magnitudeI : \magnitudeIII > \magnitudeII : \magnitudeIII$;

\begin{center}
take $M' \magnitudeI$, $m' \magnitudeIII$, $M' \magnitudeII$, $m' \magnitudeIII$;\\
such, that $M' \drawMagnitude{Ia}$ and $M' \drawMagnitude{Ib}$ shall be each $>$ \magnitudeIII;\\
considere também $m' \magnitudeIII$ o menor múltiplo de \magnitudeIII,\\
which will make $m' \magnitudeIII > M' \magnitudeII = M' \magnitudeIb$;

$\therefore M' \magnitudeII \mbox{ is } \ngtr m' \magnitudeIII$,\\ %\mbox{is not} >
but $M' \magnitudeI \mbox{ is } > m' \magnitudeIII$, for,\\
como $m' \magnitudeIII$ é o primeiro múltiplo que primeiro se torna $> M' \magnitudeIb$, então $(m' - 1)\magnitudeIII$ ou $m' \magnitudeIII - \magnitudeIII$ é $\ngtr M' \magnitudeIa$, e $\magnitudeIII \ngtr M' \magnitudeIa$ %duas vezes\mbox{ menos } e \mbox{não é} >

$\therefore m' \magnitudeIII - \magnitudeIII + \magnitudeIII$ must be $< M' \magnitudeIb + M' \magnitudeIa$;\\ %\mbox{ minus }
that is, $m' \magnitudeIII$ must be $< M' \magnitudeI$;

$\therefore M' \magnitudeI \mbox{ is } > m' \magnitudeIII$; mas foi mostrado acima que $M' \magnitudeII \mbox{ is } \ngtr m' \magnitudeIII$, portanto, pela sétima definição, \magnitudeI\ tem para \magnitudeIII\ uma proporção maior que $\magnitudeII : \magnitudeIII$. % \mbox{ não é } >

A seguir provaremos que \magnitudeIII\ tem uma razão maior para \magnitudeII, quanto menor, do que para \magnitudeI, maior;\\
or, $\magnitudeIII : \magnitudeII > \magnitudeIII : \magnitudeI$.

Take $m' \magnitudeIII$, $M' \magnitudeII$, $m' \magnitudeIII$ and $M' \magnitudeI$,\\
the same as in the first case, such that\\
$M' \magnitudeIa$ e $M' \magnitudeIb$ serão cada um $> \magnitudeIII$, e $m' \magnitudeIII$ o menor múltiplo de \magnitudeIII, que primeiro se torna maior que $M' \magnitudeIb = M' \magnitudeII$.

$\therefore m' \magnitudeIII - \magnitudeIII \mbox{ is } \ngtr M' \magnitudeIb$,\\ % \mbox{ minus },  \mbox{ is not } >
e \magnitudeIII é $\ngtr M' \magnitudeIa$; consequentemente\\ % não $>
$m' \magnitudeIII - \magnitudeIII + \magnitudeIII \mbox{ is } < M' \magnitudeIb + M' \magnitudeIa$; % \mbox{ minus }

$\therefore m' \magnitudeIII \mbox{ is } < M' \magnitudeI$ e $\therefore$ pela sétima definição, \magnitudeIII\ tem para \magnitudeII\ uma proporção maior do que \magnitudeIII\ tem para \magnitudeI.

$\therefore$ Of unequal magnitudes, \&c.
\end{center}

O artifício empregado nesta proposição para encontrar entre os múltiplos tomados, como na quinta definição, um múltiplo do primeiro maior que o múltiplo da segunda, mas o mesmo múltiplo do terceiro que foi tirado do primeiro, não maior que o mesmo múltiplo do quarto que foi tirado do segundo, pode ser ilustrado numericamente da seguinte forma: -

O número $\textcolor{byblack}{9}$ tem uma proporção maior para $\textcolor{byblue}{7}$ do que $\textcolor{byyellow}{8}$ tem para $\textcolor{byblue}{7}$: ou seja, $\textcolor{byblack}{9} : \textcolor{byblue}{7} > \textcolor{byyellow}{8} : \textcolor{byblue}{7}$; ou $\textcolor{byred}{8} + \textcolor{byblack}{1} : \textcolor{byblue}{7} > \textcolor{byyellow}{8} : \textcolor{byblue}{7}$.

O múltiplo de $\textcolor{byblack}{1}$, que primeiro se torna maior que $\textcolor{byblue}{7}$, é $8$ vezes, portanto, podemos multiplicar o primeiro e o terceiro por $8$, $9$, $10$, ou qualquer outro número maior; neste caso, multipliquemos o primeiro e o terceiro por $8$, e teremos $\textcolor{byred}{64} + \textcolor{byblack}{8}$ e $\textcolor{byyellow}{64}$: novamente, o primeiro múltiplo de $\textcolor{byblue}{7}$ que se torna maior que $64$ é $10$ vezes; então, multiplicando o segundo e o quarto por $10$, teremos $\textcolor{byblue}{70}$ e $\textcolor{byblue}{70}$; então, organizando esses múltiplos, temos—

\begin{center}
%\unprotect
%\ \hfill\vbox{
%\offinterlineskip\lineskip3pt
%\halign{\hfil # \hfil & \hfil # \hfil & \hfil # \hfil & \hfil # \hfil\\
{\renewcommand{\arraystretch}{0.75}
\begin{tabular}{c c c c}
{\footnotesize 8 times} & {\footnotesize 10 times} & {\footnotesize 8 times} & {\footnotesize 10 times} \\
{\footnotesize the first} & {\footnotesize the second} & {\footnotesize the third} & {\footnotesize the fourth} \\
$\textcolor{byred}{64} + \textcolor{byblack}{8}$ & $\textcolor{byblue}{70}$ & $\textcolor{byyellow}{64}$ & $\textcolor{byblue}{70}$ \\
\end{tabular}
}
%}}\hfill\
%\protect
\end{center}

Consequentemente, $\textcolor{byred}{64} + \textcolor{byblack}{8}$, ou $72$, é maior que $\textcolor{byblue}{70}$, mas $\textcolor{byyellow}{64}$ não é maior que $\textcolor{byblue}{70}$, $\therefore$ pela sétima definição $\textcolor{byblack}{9}$ tem uma proporção maior para $\textcolor{byblue}{7}$ do que $\textcolor{byyellow}{8}$ tem para $\textcolor{byblue}{7}$.

O que foi dito acima é meramente ilustrativo da demonstração anterior, pois esta propriedade poderia ser mostrada facilmente destes ou de outros números da seguinte maneira; porque, se um antecedente contém seu consequente um número maior de vezes do que outro antecedente contém seu consequente, ou quando uma fração é formada por um antecedente para o numerador, e seu consequente para o denominador é maior do que outra fração que é formada por outro antecedente para o numerador e seu consequente para o denominador, a razão do primeiro antecedente para o seu consequente é maior do que a razão do último antecedente para o seu consequente.

Assim, o número $9$ tem uma proporção maior para $7$, do que $8$ tem para $7$, pois $\frac{9}{7}$ é maior que $\frac{8}{7}$.

Novamente, $17 : 19$ é uma proporção maior que $13 : 15$, porque $\frac{17}{19} = \frac{17 \times 15}{19 \times 15} = \frac{255}{185}$ e $\frac{13}{15} = \frac{13 \times 19}{15 \times 19} = \frac{247}{185}$, portanto, é evidente que $\frac{255}{185}$ é maior que $\frac{247}{185}$, $\therefore \frac{17}{19}$ é maior que $\frac{13}{15}$ e, de acordo com o que foi mostrado acima, $17$ tem para $19$ uma proporção maior que $13$ tem que $15$.

De modo que os termos gerais sobre os quais existe uma proporção maior, igual ou menor são os seguintes: -

Se $\frac{A}{B}$ for maior que $\frac{C}{D}$, diz-se que $A$ tem para $B$ uma proporção maior do que $C$ tem para $D$; se $\frac{A}{B}$ for igual a $\frac{C}{D}$, então $A$ terá para $B$ a mesma proporção que $C$ tem para $D$; e se $\frac{A}{B}$ for menor que $\frac{C}{D}$, diz-se que $A$ tem para $B$ uma proporção menor do que $C$ tem para $D$.

O estudante deve compreender perfeitamente tudo até esta proposição antes de prosseguir, a fim de compreender plenamente as seguintes proposições deste livro. Portanto, recomendamos fortemente que o aluno comece de novo, e leia até isso lentamente, e raciocine cuidadosamente a cada passo, à medida que avança, protegendo-se especialmente contra o sistema malicioso de depender inteiramente da memória. Seguindo estas instruções, ele descobrirá que as partes que geralmente apresentam dificuldades consideráveis ​​não apresentarão qualquer dificuldade no prosseguimento do estudo deste importante livro.

\vfill\pagebreak

\starttheorem{Prop. IX. Theor.}\label{prop:V.IX}
\defineNewPicture{
byMagnitudeSymbolDefine.I("rhombus", byblue, 0);
byMagnitudeSymbolDefine.II("circle", byred, 0);
byMagnitudeSymbolDefine.III("square", byyellow, 0);
}
\problem{M}{agnitudes}{que têm a mesma proporção e a mesma magnitude são iguais entre si; e aqueles para os quais a mesma magnitude tem a mesma proporção são iguais entre si.}

\begin{center}
Let $\drawMagnitude{I} : \drawMagnitude{III} :: \drawMagnitude{II} : \drawMagnitude{III}$, then $\magnitudeI = \magnitudeII$.

Pois, se não, deixe $\magnitudeI > \magnitudeII$, então irá\\
$\magnitudeI : \magnitudeIII > \magnitudeII : \magnitudeIII$ \byref{prop:V.VIII},
o que é absurdo de acordo com a hipótese.

$\therefore \magnitudeI \mbox{ is } \ngtr \magnitudeII$. % \mbox{ is not } >

In the same manner it may be shown, that\\
$\magnitudeII \mbox{ is } \ngtr \magnitudeI$, % \mbox{ is not } >

$\therefore \magnitudeI = \magnitudeII$.

Again, let $\magnitudeIII : \magnitudeI :: \magnitudeIII : \magnitudeII$,\\
then will $\magnitudeI = \magnitudeII$.

For (invert.) $\magnitudeI : \magnitudeIII :: \magnitudeII : \magnitudeIII$,\\
therefore, by the first case, $\magnitudeI = \magnitudeII$.

$\therefore$ Magnitudes que possuem a mesma proporção, \&c.
\end{center}

This may be shown otherwise, as follows :—

Seja $\textcolor{byyellow}{A} : \textcolor{byblue}{B} = \textcolor{byyellow}{A} : \textcolor{byred}{C}$, depois $\textcolor{byblue}{B} = \textcolor{byred}{C}$, pois, como a fração $\frac{\textcolor{byyellow}{A}}{\textcolor{byblue}{B}} = \mbox{ the fraction } \frac{\textcolor{byyellow}{A}}{\textcolor{byred}{C}}$, e o numerador de uma é igual ao numerador da outra, portanto, os denominadores dessas frações são iguais, ou seja, $\textcolor{byblue}{B} = \textcolor{byred}{C}$.

Again, if $\textcolor{byblue}{B} : \textcolor{byyellow}{A} = \textcolor{byred}{C} : \textcolor{byyellow}{A}$, $\textcolor{byblue}{B} = \textcolor{byred}{C}$. For, as $\frac{\textcolor{byblue}{B}}{\textcolor{byyellow}{A}} = \frac{\textcolor{byred}{C}}{\textcolor{byyellow}{A}}$, $\textcolor{byblue}{B} \mbox{ must } = \textcolor{byred}{C}$.

\vfill\pagebreak

\starttheorem{Prop. X. Theor.}\label{prop:V.X}
\defineNewPicture{
byMagnitudeSymbolDefine.I("wedgeDown", byblue, 0);
byMagnitudeSymbolDefine.II("circle", byred, 0);
byMagnitudeSymbolDefine.III("square", byyellow, 0);
}
\problem{A}{quela}{magnitude que tem uma proporção maior do que outra tem a mesma magnitude, é a maior das duas; e aquela grandeza com a qual a mesma tem uma razão maior do que tem com outra grandeza, é a menor das duas.}

\begin{center}
Let $\drawMagnitude{I} : \drawMagnitude{III} > \drawMagnitude{II} : \drawMagnitude{III}$, then $\magnitudeI > \magnitudeII$.

Pois se não, deixe $\magnitudeI = \mbox{ or } < \magnitudeII$;\\
then, $\magnitudeI : \magnitudeIII = \magnitudeII : \magnitudeIII$ \byref{prop:V.VII} or\\
$\magnitudeI : \magnitudeIII < \magnitudeII : \magnitudeIII$ \byref{prop:V.VIII} e (invertido.), o que é um absurdo segundo a hipótese.

$\therefore \magnitudeI \mbox{ is } \neq \mbox{ ou } < \magnitudeII$, and\\ % \mbox{ is not } =
$\therefore \magnitudeI \mbox{ must be } >\magnitudeII$.

Again, let $\magnitudeIII : \magnitudeII > \magnitudeIII : \magnitudeI$,\\
then, $\magnitudeII < \magnitudeI$.

For if not, $\magnitudeII \mbox{ must be } > \mbox{ ou } = \magnitudeI$,\\
then $\magnitudeIII : \magnitudeII < \magnitudeIII : \magnitudeI$ \byref{prop:V.VIII} and (invert.);\\
or $\magnitudeIII : \magnitudeII = \magnitudeIII : \magnitudeI$ \byref{prop:V.VII}, which is absurd \byref{\hypref};

$\therefore \magnitudeII \mbox{ is } \ngtr \mbox{ ou } = \magnitudeI$,\\ %  \mbox{ is not } >
and $\therefore \magnitudeII \mbox{ must be } < \magnitudeI$.

$\therefore$ That magnitude which has, \&c.
\end{center}

\vfill\pagebreak

\starttheorem{Prop. XI. Theor.}\label{prop:V.XI}
\defineNewPicture{
byMagnitudeSymbolDefine.I("rhombus", byblue, 0);
byMagnitudeSymbolDefine.II("square", byblue, 0);
byMagnitudeSymbolDefine.III("miniTriangleUp", byblack, 0);
byMagnitudeSymbolDefine.IV("miniCircle", byblack, 0);
byMagnitudeSymbolDefine.V("circle", byred, 0);
byMagnitudeSymbolDefine.VI("wedgeDown", byyellow, 0);
}
\problem{R}{azões}{que são iguais para a mesma proporção são iguais entre si.}

\begin{center}
Seja $\drawMagnitude{I} : \drawMagnitude{II} = \drawMagnitude{V} : \drawMagnitude{VI}$ e $\magnitudeV : \magnitudeVI = \drawMagnitude{III} : \drawMagnitude{IV}$,\\
then will $\magnitudeI : \magnitudeII = \magnitudeIII : \magnitudeIV$.

For if $M \magnitudeI >, = \mbox{ ou } < m \magnitudeII$,\\
then $M \magnitudeV >, = \mbox{ ou } < m \magnitudeVI$,\\
and if $M \magnitudeV >, = \mbox{ ou } < m \magnitudeVI$,\\
then $M \magnitudeIII >, = \mbox{ ou } < m \magnitudeIV$ \byref{def:V.V};

$\therefore$ if $M \magnitudeI >, = \mbox{ ou } < m \magnitudeII$, $M \magnitudeIII >, = \mbox{ ou } < m \magnitudeIV$\\
and $\therefore$ \byref{def:V.V} $\magnitudeI : \magnitudeII = \magnitudeIII : \magnitudeIV$.

$\therefore$ Razões iguais \&c.
\end{center}

\vfill\pagebreak

\starttheorem{Prop. XII. Theor.}\label{prop:V.XII}
\defineNewPicture{
byMagnitudeSymbolDefine.I("square", byred, 0);
byMagnitudeSymbolDefine.II("circle", byred, 0);
byMagnitudeSymbolDefine.III("semicircleDown", byblack, 1);
byMagnitudeSymbolDefine.IV("sectorUp", byblack, 1);
byMagnitudeSymbolDefine.V("rhombus", byyellow, 0);
byMagnitudeSymbolDefine.VI("wedgeDown", byyellow, 0);
byMagnitudeSymbolDefine.VII("miniCircle", byblue, 0);
byMagnitudeSymbolDefine.VIII("miniTriangleDown", byblue, 0);
byMagnitudeSymbolDefine.IX("miniTriangleUp", byblack, 0);
byMagnitudeSymbolDefine.X("miniCircle", byblack, 0);
}
\problem{S}{e}{qualquer número de grandezas for proporcional, como um dos antecedentes é para o seu conseqüente, assim também todos os antecedentes tomados em conjunto serão para todos os consequentes.}

\begin{center}
Let $\drawMagnitude{I} : \drawMagnitude{II} =
\drawMagnitude{III} : \drawMagnitude{IV} =
\drawMagnitude{V} : \drawMagnitude{VI} =
\drawMagnitude{VII} : \drawMagnitude{VIII} =
\drawMagnitude{IX} : \drawMagnitude{X}$;\\
then will $\magnitudeI : \magnitudeII =
\magnitudeI + \magnitudeIII + \magnitudeV + \magnitudeVII + \magnitudeIX :
\magnitudeII + \magnitudeIV + \magnitudeVI + \magnitudeVIII + \magnitudeX$.

For if $M \magnitudeI > m \magnitudeII$, then $M \magnitudeIII > m \magnitudeIV$,\\
and $M \magnitudeV > m \magnitudeVI$, $M \magnitudeVII > m \magnitudeVIII$,\\
also $M \magnitudeIX > m \magnitudeX$ \byref{def:V.V}

Therefore, if $M \magnitudeI > m \magnitudeII$, then will\\
$M \magnitudeI + M \magnitudeIII + M \magnitudeV + M \magnitudeVII + M \magnitudeIX$, or $M (\magnitudeI + \magnitudeIII + \magnitudeV + \magnitudeVII + \magnitudeIX)$ be greater than $m \magnitudeII + m \magnitudeIV + m \magnitudeVI + m \magnitudeVIII + m \magnitudeX$, or $m (\magnitudeII + \magnitudeIV + \magnitudeVI + \magnitudeVIII + \magnitudeX)$.
\end{center}

Da mesma forma pode ser mostrado, se $M$ vezes um dos antecedentes for igual ou menor que $m$ vezes um dos consequentes, $M$ vezes todos os antecedentes tomados em conjunto, será igual ou menor que $m$ vezes todos os consequentes tomados em conjunto. Portanto, pela quinta definição, assim como um dos antecedentes está para o seu consequente, também todos os antecedentes são tomados em conjunto para todos os consequentes tomados em conjunto.

$\therefore$ Se houver qualquer número de magnitudes, \&c.

\vfill\pagebreak

\starttheorem{Prop. XIII. Theor.}\label{prop:V.XIII}
\defineNewPicture{
byMagnitudeSymbolDefine.I("wedgeDown", byblue, 0);
byMagnitudeSymbolDefine.II("semicircleDown", byblue, 1);
byMagnitudeSymbolDefine.III("square", byred, 0);
byMagnitudeSymbolDefine.IV("rhombus", byyellow, 0);
byMagnitudeSymbolDefine.V("sectorUp", byblack, 1);
byMagnitudeSymbolDefine.VI("circle", byblack, 0);
}
\problem{S}{e}{o primeiro tiver para o segundo a mesma proporção que o terceiro tem para o quarto, mas o terceiro para o quarto tiver uma proporção maior do que o quinto para o sexto; o primeiro também terá para o segundo uma proporção maior do que o quinto para o sexto.}

\begin{center}
Let $\drawMagnitude{I} : \drawMagnitude{II} = \drawMagnitude{III} : \drawMagnitude{IV}$, but $\magnitudeIII : \magnitudeIV > \drawMagnitude{V} : \drawMagnitude{VI}$,\\
then $\magnitudeI : \magnitudeII > \magnitudeV : \magnitudeVI$

Pois, $\because \magnitudeIII : \magnitudeIV > \magnitudeV : \magnitudeVI$, existem alguns múltiplos ($M'$ e $m'$) de \magnitudeIII\ e \magnitudeV, e de \magnitudeIV\ e \magnitudeVI,\\ %porque
such that $M' \magnitudeIII > m' \magnitudeIV$,\\
but $M' \magnitudeV \ngtr m' \magnitudeVI$, by the seventh definition. %\mbox{ not } >

Tomemos esses múltiplos e tomemos os mesmos múltiplos de \magnitudeI\ e \magnitudeII.

$\therefore$ \byref{def:V.V} if $M' \magnitudeI >, =, \mbox{ ou } < m' \magnitudeII$;\\
then will $M' \magnitudeIII >, =, \mbox{ ou } < m' \magnitudeIV$,\\
but $M' \magnitudeIII > m' \magnitudeIV$ (construction);

$\therefore M' \magnitudeI > m' \magnitudeII$,\\
but $M' \magnitudeV \mbox{ is } \ngtr m' \magnitudeVI$ (construction);\\ %  \mbox{ is not } >
and therefore by the seventh definition,\\
$\magnitudeI : \magnitudeII > \magnitudeV : \magnitudeVI$

$\therefore$ Se o primeiro tiver relação com o segundo, \&c.
\end{center}

\vfill\pagebreak

\starttheorem{Prop. XIV. Theor.}\label{prop:V.XIV}
\defineNewPicture{
byMagnitudeSymbolDefine.I("wedgeDown", byred, 0);
byMagnitudeSymbolDefine.II("semicircleDown", byblack, 1);
byMagnitudeSymbolDefine.III("square", byyellow, 0);
byMagnitudeSymbolDefine.IV("rhombus", byblue, 0);
}
\problem{S}{e}{o primeiro tem para o segundo a mesma proporção que o terceiro tem para o quarto; então, se o primeiro for maior que o terceiro, o segundo será maior que o quarto; e se for igual, igual; e se menos, menos.}

\begin{center}
Let $\drawMagnitude{I} : \drawMagnitude{II} :: \drawMagnitude{III} : \drawMagnitude{IV}$, and first suppose\\
$\magnitudeI > \magnitudeIII$, then will $\magnitudeII > \magnitudeIV$.

Para $\magnitudeI : \magnitudeII > \magnitudeIII : \magnitudeII$ \byref{prop:V.VIII}, e pela hipótese $\magnitudeI : \magnitudeII = \magnitudeIII : \magnitudeIV$;\\
$\therefore \magnitudeIII : \magnitudeIV > \magnitudeIII : \magnitudeII$ \byref{prop:V.XIII},\\
$\therefore \magnitudeIV < \magnitudeII$ \byref{prop:V.X}, or $\magnitudeII > \magnitudeIV$.

Secondly, let $\magnitudeI = \magnitudeIII$, then will $\magnitudeII = \magnitudeIV$.

For $\magnitudeI : \magnitudeII = \magnitudeIII : \magnitudeII$ \byref{prop:V.VII},\\
and $\magnitudeI : \magnitudeII = \magnitudeIII : \magnitudeIV$ \byref{\hypref};\\
$\therefore \magnitudeIII : \magnitudeII = \magnitudeIII : \magnitudeIV$ \byref{prop:V.XI},\\
and $\therefore \magnitudeII = \magnitudeIV$ \byref{prop:V.IX}.

Thirdly, if $\magnitudeI < \magnitudeIII$, then will $\magnitudeII < \magnitudeIV$;\\
$\because \magnitudeIII > \magnitudeI$ and $\magnitudeIII : \magnitudeIV = \magnitudeI : \magnitudeII$;\\ %because
$\therefore \magnitudeIV > \magnitudeII$, by the first case,\\
that is, $\magnitudeII < \magnitudeIV$.

$\therefore$ Se o primeiro tiver a mesma proporção, \&c.
\end{center}

\vfill\pagebreak

\starttheorem{Prop. XV. Theor.}\label{prop:V.XV}
\defineNewPicture{
byMagnitudeSymbolDefine.I("circle", byred, 0);
byMagnitudeSymbolDefine.II("square", byyellow, 0);
}
\problem{A}{s}{magnitudes têm entre si a mesma proporção que seus equimúltiplos têm.}

\begin{center}
Sejam \drawMagnitude{I} e \drawMagnitude{II} duas grandezas;\\
then, $\magnitudeI : \magnitudeII :: M' \magnitudeI : M' \magnitudeII$.

$\begin{aligned}
\mbox{For } \magnitudeI : \magnitudeII &= \magnitudeI : \magnitudeII \\
&= \magnitudeI : \magnitudeII \\
&= \magnitudeI : \magnitudeII \\
\end{aligned}$\\
$\therefore \magnitudeI : \magnitudeII :: 4 \magnitudeI : 4 \magnitudeII$. \byref{prop:V.XII}.

E como o mesmo raciocínio é geralmente aplicável, temos\\
$\magnitudeI : \magnitudeII :: M' \magnitudeI : M' \magnitudeII$.

As magnitudes $\therefore$ têm a mesma proporção, \&c.
\end{center}

\vfill\pagebreak

\startdefinition[8]{Definição XIII}\label{def:V.XIII}
\defineNewPicture{
byMagnitudeSymbolDefine.I("circle", byyellow, 0);
byMagnitudeSymbolDefine.II("rhombus", byblack, 0);
byMagnitudeSymbolDefine.III("wedgeDown", byred, 0);
byMagnitudeSymbolDefine.IV("square", byblue, 0);
}
O termo técnico permutando, ou alternando, por permutação ou alternadamente, é utilizado quando há quatro proporcionais, e infere-se que a primeira tem a mesma proporção para a terceira que a segunda tem para a quarta; ou que o primeiro está para o terceiro assim como o segundo está para o quarto: como é mostrado na seguinte proposição: -

\begin{center}
Let $\drawMagnitude{I} : \drawMagnitude{II} :: \drawMagnitude{III} : \drawMagnitude{IV}$,\\
by \enquote{permutando} or \enquote{alternando} it is inferred\\
$\magnitudeI : \magnitudeIII :: \magnitudeII : \magnitudeIV$.
\end{center}

Pode ser necessário aqui observar que as grandezas \magnitudeI, \magnitudeII, \magnitudeIII, \magnitudeIV, devem ser homogêneas, isto é, da mesma natureza ou semelhança de espécie; devemos, portanto, em tais casos, comparar linhas com linhas, superfícies com superfícies, sólidos com sólidos, \&c. Portanto, o estudante perceberá prontamente que uma linha e uma superfície, uma superfície e um sólido, ou outras grandezas heterogêneas, nunca podem estar na relação de antecedente e consequente.

\vfill\pagebreak

\starttheorem{Prop. XVI. Theor.}\label{prop:V.XVI}
\defineNewPicture{
byMagnitudeSymbolDefine.I("wedgeDown", byred, 0);
byMagnitudeSymbolDefine.II("semicircleDown", byblack, 1);
byMagnitudeSymbolDefine.III("square", byyellow, 0);
byMagnitudeSymbolDefine.IV("rhombus", byblue, 0);
}
\problem{S}{e}{quatro grandezas do mesmo tipo forem proporcionais, elas também serão proporcionais quando tomadas alternadamente.}

\begin{center}
Let $\drawMagnitude{I} : \drawMagnitude{II} :: \drawMagnitude{III} : \drawMagnitude{IV}$, then $\magnitudeI : \magnitudeIII :: \magnitudeII : \magnitudeIV$.

For $M \magnitudeI : M \magnitudeII :: \magnitudeI : \magnitudeII$ \byref{prop:V.XV},\\
and $M \magnitudeI : M \magnitudeII :: \magnitudeIII : \magnitudeIV$ \byref{\hypref} and \byref{prop:V.XI};\\
also $m \magnitudeIII : m \magnitudeIV :: \magnitudeIII : \magnitudeIV$ \byref{prop:V.XV};

$\therefore M \magnitudeI : M \magnitudeII :: m \magnitudeIII : m \magnitudeIV$ \byref{prop:V.XIV},\\
and $\therefore$ if $M \magnitudeI >, = \mbox { or } < m \magnitudeIII$,\\
then will $M \magnitudeII >, =, \mbox{ ou } < m \magnitudeIV$ \byref{prop:V.XIV};

therefore, by the fifth definition,\\
 $\magnitudeI : \magnitudeIII :: \magnitudeII : \magnitudeIV$

 $\therefore$ Se houver quatro magnitudes do mesmo tipo, \&c.
\end{center}

\startdefinition[16]{Definição XVI}\label{def:V.XVI}
\def\varA{\textcolor{byred}{A}}
\def\varB{\textcolor{byblack}{B}}
\def\varC{\textcolor{byblue}{C}}
\def\varD{\textcolor{byyellow}{D}}
Dividendo, por divisão, quando há quatro proporcionais, e infere-se que o excesso da primeira sobre a segunda está para a segunda, assim como o excesso da terceira sobre a quarta, está para a quarta.

\begin{center}
Seja $\varA : \varB :: \varC : \varD$;\\
by \enquote{dividendo} it is inferred\\
$\varA - \varB : \varB :: \varC - \varD : \varD$. % duas vezes \mbox{ menos }
\end{center}

De acordo com o acima exposto, $\varA$ deve ser maior que $\varB$ e $\varC$ maior que $\varD$; se este não for o caso, mas para ter $\varB$ maior que $\varA$, e $\varD$ maior que $\varC$, $\varB$ e $\varD$ podem ser considerados antecedentes, e $\varA$ e $\varC$ como consequentes, por \enquote{inversão}

\begin{center}
$\varB : \varA :: \varD : \varC$;\\
then, by \enquote{dividendo,} we infer\\
$\varB - \varA : \varA :: \varD - \varC : \varC$ % duas vezes \mbox{ menos }
\end{center}

\vfill\pagebreak

\starttheorem{Prop. XVII. Theor.}\label{prop:V.XVII}
\defineNewPicture{
byMagnitudeSymbolDefine.I("wedgeDown", byred, 0);
byMagnitudeSymbolDefine.II("semicircleDown", byblack, 1);
byMagnitudeSymbolDefine.III("square", byyellow, 0);
byMagnitudeSymbolDefine.IV("rhombus", byblue, 0);
}
\problem{S}{e}{as grandezas, tomadas em conjunto, forem proporcionais, também serão proporcionais quando tomadas separadamente: isto é, se duas grandezas juntas têm para uma delas a mesma razão que duas outras têm para uma delas, a restante das duas primeiras terá para a outra a mesma razão que a restante das duas últimas tem para a outra.}

\begin{center}
Let $\drawMagnitude{I} + \drawMagnitude{II} : \magnitudeII :: \drawMagnitude{III} + \drawMagnitude{IV} : \magnitudeIV$,\\
then will $\magnitudeI : \magnitudeII :: \magnitudeIII : \magnitudeIV$.

Take $\magnitudeI > m \magnitudeII$ to each add $M \magnitudeII$,\\
then we have $M \magnitudeI + M \magnitudeII > m \magnitudeII + M \magnitudeII$,\\
or $M (\magnitudeI + \magnitudeII) > (m + M) \magnitudeII$:\\
but $\because \magnitudeI + \magnitudeII: \magnitudeII :: \magnitudeIII + \magnitudeIV: \magnitudeIV$ \byref{\hypref},\\ %because
and $M (\magnitudeI + \magnitudeII) > (m + M) \magnitudeII$;\\
$\therefore M(\magnitudeIII + \magnitudeIV) > (m + M) \magnitudeIV$ \byref{def:V.V};\\
$\therefore M \magnitudeIII + M \magnitudeIV > m \magnitudeIV + M \magnitudeIV$;\\
$\therefore M \magnitudeIII > m \magnitudeIV$, by taking $M \magnitudeIV$ from both sides:\\
that is, when $M \magnitudeI > m \magnitudeII$, then $M \magnitudeIII > m \magnitudeIV$.

Da mesma maneira pode ser provado que se $M \magnitudeI = \mbox{ or } < m \magnitudeII$, então $M \magnitudeIII = \mbox{ or } < m \magnitudeIV$;\\
and $\therefore \magnitudeI : \magnitudeII :: \magnitudeIII : \magnitudeIV$ \byref{def:V.V}

$\therefore$ Se as magnitudes forem tomadas em conjunto, \&c.
\end{center}

\vfill\pagebreak

\startdefinition[15]{Definição XV}\label{def:V.XV}
\def\varA{\textcolor{byred}{A}}
\def\varB{\textcolor{byblack}{B}}
\def\varC{\textcolor{byyellow}{C}}
\def\varD{\textcolor{byblue}{D}}
O termo componendo, por composição, é utilizado quando há quatro proporcionais; e infere-se que o primeiro junto com o segundo está para o segundo assim como o terceiro junto com o quarto está para o quarto.

\begin{center}
Seja $\varA : \varB :: \varC : \varD$;

then, by term \enquote{componendo,} it is inferred that\\
$\varA + \varB : \varB :: \varC + \varD: \varD$
\end{center}

Bу \enquote{invertion} $\varB$ and $\varD$ may become the first and the third, and $\varA$ and $\varC$ the second and fourth, as

\begin{center}
$\varB : \varA :: \varD : \varC$,

then, by \enquote{componendo,} we infer that\\
$\varB + \varA : \varA :: \varD + \varC : \varC$.
\end{center}

\vfill\pagebreak

\starttheorem{Prop. XVIII. Theor.}\label{prop:V.XVIII}
\defineNewPicture{
byMagnitudeSymbolDefine.I("wedgeDown", byred, 0);
byMagnitudeSymbolDefine.II("semicircleDown", byblack, 1);
byMagnitudeSymbolDefine.III("square", byyellow, 0);
byMagnitudeSymbolDefine.IV("rhombus", byblue, 0);
byMagnitudeSymbolDefine.V("circle", byblack, 0);
}
\problem{S}{e}{as grandezas, tomadas separadamente, forem proporcionais, também serão proporcionais tomadas conjuntamente: isto é, se a primeira estiver para a segunda como a terceira estiver para a quarta, a primeira e a segunda juntas estarão para a segunda assim como a terceira e a quarta juntas estão para a quarta.}

\begin{center}
Let $\drawMagnitude{I} : \drawMagnitude{II} :: \drawMagnitude{III} : \drawMagnitude{IV}$,\\
then $\magnitudeI + \magnitudeII : \magnitudeII :: \magnitudeIII + \magnitudeIV : \magnitudeIV$;

for if not, let $\magnitudeI + \magnitudeII : \magnitudeII :: \magnitudeIII + \drawMagnitude{V} : \magnitudeV$,\\
supposing $\magnitudeV \neq \magnitudeIV$; %\mbox{ not } =

$\therefore \magnitudeI : \magnitudeII :: \magnitudeIII : \magnitudeV$ \byref{prop:V.XVII}\\
but $\magnitudeI : \magnitudeII :: \magnitudeIII : \magnitudeIV$ \byref{\hypref};

$\therefore \magnitudeIII : \magnitudeV :: \magnitudeIII : \magnitudeIV$ \byref{prop:V.XI};

$\therefore \magnitudeV = \magnitudeIV$ \byref{prop:V.IX},\\
o que é contrário à suposição;

$\therefore \magnitudeV \mbox{ não é desigual } \magnitudeIV$;\\
that is $\magnitudeV = \magnitudeIV$;

$\therefore \magnitudeI + \magnitudeII : \magnitudeII :: \magnitudeIII + \magnitudeIV : \magnitudeIV$.

$\therefore$ If magnitudes, taken separately, \&c.
\end{center}

\vfill\pagebreak

\starttheorem{Prop. XIX. Theor.}\label{prop:V.XIX}
\defineNewPicture{
byMagnitudeSymbolDefine.I("wedgeDown", byred, 0);
byMagnitudeSymbolDefine.II("semicircleDown", byblack, 1);
byMagnitudeSymbolDefine.III("square", byblue, 0);
byMagnitudeSymbolDefine.IV("rhombus", byyellow, 0);
}
\problem{S}{e}{uma grandeza inteira está para um todo, assim como uma grandeza tirada da primeira está para uma grandeza tirada da outra; o restante será para o lembrete, como o todo para o todo.}

\begin{center}
Let $\drawMagnitude{I} + \drawMagnitude{II} : \drawMagnitude{III} + \drawMagnitude{IV} :: \magnitudeI : \magnitudeIII$,\\
then will $\magnitudeII : \magnitudeIV :: \magnitudeI + \magnitudeII : \magnitudeIII + \magnitudeIV$.

For $\magnitudeI + \magnitudeII : \magnitudeI :: \magnitudeIII + \magnitudeIV : \magnitudeIII$ (alter.),

$\therefore \magnitudeII : \magnitudeI :: \magnitudeIV : \magnitudeIII$ (divid.),

again $\magnitudeII : \magnitudeIV :: \magnitudeI : \magnitudeIII$ (alter.),

but $\magnitudeI + \magnitudeII : \magnitudeIII + \magnitudeIV :: \magnitudeI : \magnitudeIII$ \byref{\hypref};

therefore $\magnitudeII : \magnitudeIV :: \magnitudeI + \magnitudeII : \magnitudeIII + \magnitudeIV$ \byref{prop:V.XI}.

$\therefore$ Se uma magnitude inteira for um todo, \&c.
\end{center}

\startdefinition[17]{Definição XVII}\label{def:V.XVII}
O termo \enquote{convertendo,} por conversão, é utilizado pelos geômetras, quando há quatro proporcionais, e infere-se que a primeira está em seu excesso acima da segunda, assim como a terceira está em seu excesso acima da quarta. Veja a seguinte proposição: -

\vfill\pagebreak

\starttheoremAZ{Prop. E. Theor.}\label{prop:V.E} % Proposition with letters
\defineNewPicture{
byMagnitudeSymbolDefine.i("circle", byblue, 0);
byMagnitudeSymbolDefine.II("sectorUp", byblack, 1);
byMagnitudeSymbolDefine.iii("square", byred, 0);
byMagnitudeSymbolDefine.IV("rhombus", byyellow, 0);
byMagnitudeDefine.I(0, true)(1, 1)(i, II);
byMagnitudeDefine.III(0, true)(1, 1)(iii, IV);
}
\problem{S}{e}{quatro grandezas são proporcionais, elas também o são por conversão: isto é, a primeira está em seu excesso acima da segunda, assim como a terceira está em seu excesso acima da quarta.}

\begin{center}
Let $\drawMagnitude{I} : \drawMagnitude{II} :: \drawMagnitude{III} : \drawMagnitude{IV}$,\\
then shall $\magnitudeI : \drawMagnitude{i} :: \magnitudeIII : \drawMagnitude{iii}$.

$\because \magnitudeI : \magnitudeII :: \magnitudeIII : \magnitudeIV$;\\ %Because
$\therefore \magnitudei : \magnitudeII :: \magnitudeiii : \magnitudeIV$ (divid.), %therefore

$\therefore \magnitudeII : \magnitudei :: \magnitudeIV : \magnitudeiii$ (inver.),

$\therefore \magnitudeI : \magnitudei :: \magnitudeIII : \magnitudeiii$ (compo.).

$\therefore$ If four magnitudes, \&c.
\end{center}

\startdefinition[18]{Definição XVIII}\label{def:V.XVIII}
\enquote{Ex \ae quali} (sc.\ distanteiâ), ou ex \ae quo, de igualdade de distância: quando há qualquer número de magnitudes maior que dois, e tantos outros, tais que são proporcionais quando tomados dois mais dois de cada classificação, e infere-se que o primeiro está para o último da primeira classificação de magnitudes, como o primeiro está para o último dos outros: \enquote{disso existem dois tipos seguintes, que surgem da ordem diferente em que as magnitudes são medidas, dois e dois.}

\vfill\pagebreak

\startdefinition[19]{Definição XIX}\label{def:V.XIX}
\def\varA{\textcolor{byred}{A}}
\def\varB{\textcolor{byred}{B}}
\def\varC{\textcolor{byyellow}{C}}
\def\varD{\textcolor{byyellow}{D}}
\def\varE{\textcolor{byblue}{E}}
\def\varF{\textcolor{byblue}{F}}
\def\varL{\textcolor{byred}{L}}
\def\varM{\textcolor{byred}{M}}
\def\varN{\textcolor{byyellow}{N}}
\def\varO{\textcolor{byyellow}{O}}
\def\varP{\textcolor{byblue}{P}}
\def\varQ{\textcolor{byblue}{Q}}
\enquote{Ex \ae quali,} da igualdade. Este termo é usado simplesmente por si só, quando a primeira grandeza está para a segunda da primeira categoria, como a primeira está para a segunda da outra categoria; e assim como o segundo está para o terceiro da primeira categoria, assim é o segundo para o terceiro da outra; e assim por diante: e a inferência é como mencionada na definição anterior; daí isso é chamado de proposição ordenada. Isso é demonstrado no \byref{prop:V.XXII}.

\begin{center}
Thus, if there be two ranks of magnitudes,\\
$\varA, \varB, \varC, \varD, \varE, \varF$, the first rank,
e $\varL, \varM, \varN, \varO, \varP, \varQ$, o segundo,
tal que $\varA : \varB :: \varL : \varM$, $\varB : \varC :: \varM : \varB$, $\varC : \varD :: \varN : \varO$, $\varD : \varE :: \varO : \varP$, $\varE : \varF :: \varP : \varQ$;\\
inferimos pela expressão \enquote{ex~\ae quali} que $\varA : \varF :: \varL : \varQ$
\end{center}

\startdefinition[20]{Definição XX}\label{def:V.XX}
\def\varA{\textcolor{byred}{A}}
\def\varB{\textcolor{byred}{B}}
\def\varC{\textcolor{byblue}{C}}
\def\varD{\textcolor{byblue}{D}}
\def\varE{\textcolor{byyellow}{E}}
\def\varF{\textcolor{byyellow}{F}}
\def\varL{\textcolor{byyellow}{L}}
\def\varM{\textcolor{byyellow}{M}}
\def\varN{\textcolor{byblue}{N}}
\def\varO{\textcolor{byblue}{O}}
\def\varP{\textcolor{byred}{P}}
\def\varQ{\textcolor{byred}{Q}}
\enquote{Ex \ae quali in proporcional perturbatâ seu inordinatâ,} de igualdade em perturbada, ou proporção desordenada. Este termo é usado quando a primeira magnitude está para a segunda da primeira classificação como a última, mas uma está para a última da segunda classificação; e assim como o segundo está para o terceiro da primeira categoria, o penúltimo dois está para o penúltimo da segunda categoria; e assim como o terceiro está para o quarto da primeira categoria, assim também o terceiro está para o penúltimo penúltimo da segunda categoria; e assim por diante em ordem cruzada: e a inferência está na 18ª definição. É demonstrado em \byref{prop:V.XXIII}.

\begin{center}
Thus, if there be two ranks of magnitudes,\\
$\varA, \varB, \varC, \varD, \varE, \varF$, the first rank,
e $\varL, \varM, \varN, \varO, \varP, \varQ$, o segundo,
tal que $\varA : \varB :: \varP : \varQ$, $\varB : \varC :: \varO : \varP$, $\varC : \varD :: \varN : \varO$, $\varD : \varE :: \varM : \varN$, $\varE : \varF :: \varL : \varM$;\\
o termo \enquote{ex \ae quali in proporcional perturbatâ seu inordinatâ} infere que $\varA : \varF :: \varL : \varQ$
\end{center}

\vfill\pagebreak

\starttheorem{Prop. XX. Theor.}\label{prop:V.XX}
\defineNewPicture{
byMagnitudeSymbolDefine.I("wedgeDown", byblue, 0);
byMagnitudeSymbolDefine.II("semicircleDown", byred, 1);
byMagnitudeSymbolDefine.III("square", byyellow, 0);
byMagnitudeSymbolDefine.IV("rhombus", byblue, 0);
byMagnitudeSymbolDefine.V("sectorUp", byred, 1);
byMagnitudeSymbolDefine.VI("circle", byyellow, 0);
}
\problem{S}{e}{há três grandezas, e outras três, que são consideradas dois e dois, têm a mesma razão; então, se o primeiro for maior que o terceiro, o quarto será maior que o sexto; e se for igual, igual; e se menos, menos.}

\begin{center}
Seja \drawMagnitude{I}, \drawMagnitude{II}, \drawMagnitude{III}, as três primeiras magnitudes,\\
e \drawMagnitude{IV}, \drawMagnitude{V}, \drawMagnitude{VI}, sejam os outros três,\\
such that $\magnitudeI : \magnitudeII :: \magnitudeIV : \magnitudeV$, and $\magnitudeII : \magnitudeIII :: \magnitudeV : \magnitudeVI$.

Then, if $\magnitudeI >, =, \mbox{ ou } < \magnitudeIII$,\\
then will $\magnitudeIV >, =, \mbox{ ou } < \magnitudeVI$.

Da hipótese, por alternância, temos\\
$\magnitudeI : \magnitudeIV :: \magnitudeII : \magnitudeV$,\\
and $\magnitudeII : \magnitudeV :: \magnitudeIII : \magnitudeVI$

$\therefore \magnitudeI : \magnitudeIV :: \magnitudeIII : \magnitudeVI$ \byref{prop:V.XI};

$\therefore$ if $\magnitudeI >, =, \mbox{ ou } < \magnitudeIII$,\\
then will $\magnitudeIV >, =, \mbox{ ou } < \magnitudeVI$ \byref{prop:V.XIV}.

$\therefore$ If there be three magnitudes, \&c.
\end{center}

\vfill\pagebreak

\starttheorem{Prop. XXI. Theor.}\label{prop:V.XXI}
\defineNewPicture{
byMagnitudeSymbolDefine.I("wedgeDown", byyellow, 0);
byMagnitudeSymbolDefine.II("wedgeUp", byred, 0);
byMagnitudeSymbolDefine.III("square", byblue, 0);
byMagnitudeSymbolDefine.IV("rhombus", byblue, 0);
byMagnitudeSymbolDefine.V("sectorUp", byred, 1);
byMagnitudeSymbolDefine.VI("circle", byyellow, 0);
}
\problem{S}{e}{houver três grandezas e outras três que tenham a mesma razão, tomadas duas e duas, mas em ordem cruzada; então, se a primeira grandeza for maior que a terceira, a quarta será maior que a sexta; e se for igual, igual, e se for menor, menos.}
\unskip

\begin{center}
Let \drawMagnitude{I}, \drawMagnitude{II}, \drawMagnitude{III}, be the first three magnitudes,\\
e \drawMagnitude{IV}, \drawMagnitude{V}, \drawMagnitude{VI}, os outros três,\\
such that $\magnitudeI : \magnitudeII :: \magnitudeV : \magnitudeVI$,\\
and $\magnitudeII : \magnitudeIII :: \magnitudeIV : \magnitudeV$.

Then if $\magnitudeI >, =, \mbox{ ou } < \magnitudeIII$,\\
then will $\magnitudeIV >, =, \mbox{ ou } < \magnitudeVI$.

First, let $\magnitudeI \mbox{ be } > \magnitudeIII$:\\
then, because \magnitudeII\ is any other magnitude,\\
$\magnitudeI : \magnitudeII > \magnitudeIII : \magnitudeII$ \byref{prop:V.VIII};\\
but $\magnitudeV : \magnitudeVI :: \magnitudeI : \magnitudeII$ \byref{\hypref};\\
$\therefore \magnitudeV : \magnitudeVI > \magnitudeIII : \magnitudeII$ \byref{prop:V.XIII};\\
and $\because \magnitudeII : \magnitudeIII :: \magnitudeIV : \magnitudeV$ \byref{\hypref};\\ %because
$\therefore \magnitudeIII : \magnitudeII :: \magnitudeV : \magnitudeIV$ (inv.),\\
and it was shown that $\magnitudeV : \magnitudeVI > \magnitudeIII : \magnitudeII$,
$\therefore \magnitudeV : \magnitudeVI > \magnitudeV : \magnitudeIV$ \byref{prop:V.XIII};\\
$\therefore \magnitudeVI < \magnitudeIV$,\\
that is $\magnitudeIV > \magnitudeVI$.

Secondly, let $\magnitudeI = \magnitudeIII$; then shall $\magnitudeIV = \magnitudeVI$.\\
For $\because \magnitudeI = \magnitudeIII$,\\ %because
$\magnitudeI : \magnitudeII = \magnitudeIII : \magnitudeII$ \byref{prop:V.VII};\\
but $\magnitudeI : \magnitudeII = \magnitudeV : \magnitudeVI$ \byref{\hypref},\\
and $\magnitudeIII : \magnitudeII = \magnitudeV : \magnitudeIV$ (\hypstr\ and inv.),\\
$\therefore \magnitudeV : \magnitudeVI = \magnitudeV : \magnitudeIV$ \byref{prop:V.XI},\\
$\therefore \magnitudeIV = \magnitudeVI$ \byref{prop:V.IX}.

Next, let $\magnitudeI \mbox{ be } < \magnitudeIII$, then $\magnitudeIV \mbox{ shall be } < \magnitudeVI$;\\
for $\magnitudeIII > \magnitudeI$,\\
and it has been shown that $\magnitudeIII : \magnitudeII = \magnitudeV : \magnitudeIV$,\\
and $\magnitudeII : \magnitudeI = \magnitudeVI : \magnitudeV$;\\
$\therefore$ by the first case $\magnitudeVI \mbox{ is } > \magnitudeIV$,\\
that is, $\magnitudeIV < \magnitudeVI$.

$\therefore$ If there be three, \&c.
\end{center}

\vfill\pagebreak

\starttheorem{Prop. XXII. Theor.}\label{prop:V.XXII}
\defineNewPicture{
byMagnitudeSymbolDefine.I("wedgeDown", byred, 0);
byMagnitudeSymbolDefine.II("rhombus", byblue, 0);
byMagnitudeSymbolDefine.III("square", byyellow, 0);
byMagnitudeSymbolDefine.IV("rhombus", byred, 0);
byMagnitudeSymbolDefine.V("sectorUp", byblue, 1);
byMagnitudeSymbolDefine.VI("circle", byyellow, 0);
byMagnitudeSymbolDefine.Ia("wedgeDown", byblue, 0);
byMagnitudeSymbolDefine.IIa("rhombus", byblack, 0);
byMagnitudeSymbolDefine.IIIa("square", byyellow, 0);
byMagnitudeSymbolDefine.IVa("rhombus", byred, 0);
byMagnitudeSymbolDefine.Va("sectorUp", byblue, 1);
byMagnitudeSymbolDefine.VIa("circle", byblack, 0);
byMagnitudeSymbolDefine.VIIa("halfsquare", byyellow, 0);
byMagnitudeSymbolDefine.VIIIa("halfrhombusUp", byred, 0);
}
\problem{I}{f}{há qualquer número de grandezas, e tantas outras, que, tomadas duas e duas em ordem, têm a mesma razão; a primeira terá para a última das primeiras grandezas a mesma razão que a primeira das outras tem para a última das mesmas.
N.B.— Geralmente é citado pelas palavras \enquote{ex \ae quali,} ou \enquote{ex \ae quo.}}

\begin{center}
Primeiro, sejam as magnitudes \drawMagnitude{I}, \drawMagnitude{II}, \drawMagnitude{III},\\
and as many others \drawMagnitude{IV}, \drawMagnitude{V}, \drawMagnitude{VI},\\
such that\\
$\magnitudeI : \magnitudeII :: \magnitudeIV : \magnitudeV$,\\
and $\magnitudeII : \magnitudeIII :: \magnitudeV : \magnitudeVI$,\\
then shall $\magnitudeI : \magnitudeIII :: \magnitudeIV : \magnitudeVI$.\\
\end{center}

Seja essas grandezas, bem como quaisquer equimúltiplos quaisquer dos antecedentes e consequentes das razões, ficarem como segue: -

\begin{center}
$\magnitudeI, \magnitudeII, \magnitudeIII, \magnitudeIV, \magnitudeV, \magnitudeVI,$\\
and\\
$M \magnitudeI, m \magnitudeII, N \magnitudeIII, M \magnitudeIV, m \magnitudeV, N \magnitudeVI,$\\
$\because \magnitudeI : \magnitudeII :: \magnitudeIV : \magnitudeV$,\\ %because
$\therefore M \magnitudeI : m \magnitudeII :: M \magnitudeIV : m \magnitudeV$ \byref{prop:V.IV}.

For the same reason\\
$m \magnitudeII : N \magnitudeIII :: m \magnitudeV : N \magnitudeVI$;\\
e porque existem três magnitudes\\
$M \magnitudeI, m \magnitudeII, N \magnitudeIII$,\\
and other three, $M \magnitudeIV, m \magnitudeV, N \magnitudeVI$,\\
que, considerados dois e dois, têm a mesma proporção;

$\therefore$ if $M \magnitudeI >, = \mbox{ ou } < N \magnitudeIII$\\
then will $M \magnitudeIV >, = \mbox{ ou } < N \magnitudeVI$, by \byref{prop:V.XX},\\
and $\therefore \magnitudeI : \magnitudeIII :: \magnitudeIV : \magnitudeVI$ \byref{def:V.V}.

Next, let there be four magnitudes, \drawMagnitude{Ia}, \drawMagnitude{IIa}, \drawMagnitude{IIIa}, \drawMagnitude{IVa},\\
and other four, \drawMagnitude{Va}, \drawMagnitude{VIa}, \drawMagnitude{VIIa}, \drawMagnitude{VIIIa},\\
que, tomados dois e dois, têm a mesma proporção,\\
isto é, $\magnitudeIa : \magnitudeIIa :: \magnitudeVa : \magnitudeVIa$,\\
$\magnitudeIIa : \magnitudeIIIa :: \magnitudeVIa : \magnitudeVIIa$,\\
and $\magnitudeIIIa : \magnitudeIVa :: \magnitudeVIIa : \magnitudeVIIIa$,\\
then shall $\magnitudeIa : \magnitudeIVa :: \magnitudeVa : \magnitudeVIIIa$;\\
for, $\because \magnitudeIa, \magnitudeIIa, \magnitudeIIIa$, are three magnitudes,\\ %because
and $\magnitudeVa, \magnitudeVIa, \magnitudeVIIa$, other three,\\
que, tomados dois e dois, têm a mesma proporção;\\
therefore, by the foregoing case $\magnitudeIa : \magnitudeIIIa :: \magnitudeVa : \magnitudeVIIa$,\\
but $\magnitudeIIIa : \magnitudeIVa :: \magnitudeVIIa : \magnitudeVIIIa$;\\
therefore again, by the first case, $\magnitudeIa : \magnitudeIVa :: \magnitudeVa : \magnitudeVIIIa$;\\
e assim por diante, qualquer que seja o número de magnitudes.

$\therefore$ If there be any number, \&c.
\end{center}

\vfill\pagebreak

\starttheorem{Prop. XXIII. Theor.}\label{prop:V.XXIII}
\defineNewPicture{
byMagnitudeSymbolDefine.I("wedgeDown", byyellow, 0);
byMagnitudeSymbolDefine.II("semicircleDown", byblue, 1);
byMagnitudeSymbolDefine.III("square", byred, 0);
byMagnitudeSymbolDefine.IV("rhombus", byyellow, 0);
byMagnitudeSymbolDefine.V("sectorUp", byblue, 1);
byMagnitudeSymbolDefine.VI("circle", byred, 0);
byMagnitudeSymbolDefine.VII("halfsquare", byblack, 0);
byMagnitudeSymbolDefine.VIII("halfrhombusUp", byblack, 0);
}
\problem{I}{f}{há qualquer número de magnitudes, e tantas outras, que, tomadas duas mais duas em ordem cruzada, têm a mesma razão; a primeira terá para a última das primeiras grandezas a mesma razão que a primeira das outras tem para a última das mesmas.
N.B.— Isso geralmente é citado pelas palavras \enquote{ex \ae quali in propore perturbatâ;} ou \enquote{ex \ae quo perturbato.}}

\begin{center}
Primeiro, sejam três magnitudes \drawMagnitude{I}, \drawMagnitude{II}, \drawMagnitude{III},\\
and other three, \drawMagnitude{IV}, \drawMagnitude{V}, \drawMagnitude{VI},\\
que, considerados dois mais dois em ordem cruzada, têm a mesma proporção;

that is, $\magnitudeI : \magnitudeII :: \magnitudeV : \magnitudeVI$,\\
and $\magnitudeII : \magnitudeIII :: \magnitudeIV : \magnitudeV$,\\
then shall $\magnitudeI : \magnitudeIII :: \magnitudeIV : \magnitudeVI$.
\end{center}

Seja essas magnitudes e seus respectivos equimúltiplos serem organizados da seguinte forma: -

\begin{center}
$\magnitudeI, \magnitudeII, \magnitudeIII, \magnitudeIV, \magnitudeV, \magnitudeVI,$\\
$M \magnitudeI, M \magnitudeII, m \magnitudeIII, M \magnitudeIV, m \magnitudeV, m \magnitudeVI,$\\
then $\magnitudeI : \magnitudeII :: M \magnitudeI : M \magnitudeII$ \byref{prop:V.XV};\\
and for the same reason\\
$\magnitudeV : \magnitudeVI :: m \magnitudeV : m \magnitudeVI$;\\
but $\magnitudeI : \magnitudeII :: \magnitudeV : \magnitudeVI$ \byref{\hypref},\\
$\therefore M \magnitudeI : M \magnitudeII :: \magnitudeV : \magnitudeVI$ \byref{prop:V.XI};\\
and $\because \magnitudeII : \magnitudeIII :: \magnitudeIV : \magnitudeV$ \byref{\hypref},\\ %because
$\therefore M \magnitudeII : m \magnitudeIII :: M \magnitudeIV : m \magnitudeV$ \byref{prop:V.IV};\\
then, because there are three magnitudes,\\
$M \magnitudeI, M \magnitudeII, m \magnitudeII$,\\
and other three $M \magnitudeIV, m \magnitudeV, m \magnitudeVI$,\\
que, tomados dois mais dois em ordem cruzada, têm a mesma proporção;\\
therefore, if $M \magnitudeI >, =, \mbox{ ou } < m \magnitudeIII$,\\
then will $M \magnitudeIV >, =, \mbox{ ou } < m \magnitudeVI$ \byref{prop:V.XXI},\\
and $\therefore \magnitudeI : \magnitudeIII :: \magnitudeIV : \magnitudeIV$ \byref{def:V.V}.

Next, let there be four magnitudes,\\
\magnitudeI, \magnitudeII, \magnitudeIII, \magnitudeIV,\\
and other four,\\
\magnitudeV, \magnitudeVI, \drawMagnitude{VII}, \drawMagnitude{VIII},\\
que, quando tomados dois mais dois em ordem cruzada, têm a mesma proporção;\\
namely, $\magnitudeI : \magnitudeII :: \magnitudeVII : \magnitudeVIII$,\\
$\magnitudeII : \magnitudeIII :: \magnitudeVI : \magnitudeVII$,\\
and $\magnitudeIII : \magnitudeIV :: \magnitudeV : \magnitudeVI$,\\
then shall $\magnitudeIII : \magnitudeIV :: \magnitudeV : \magnitudeVI$.\\
For, $\because \magnitudeI, \magnitudeII, \magnitudeIII$ are three magnitudes,\\ %because
and $\magnitudeVI, \magnitudeVII, \magnitudeVIII$, other three,\\
que, tomados dois mais dois em ordem cruzada, têm a mesma proporção,\\
therefore, by the first case, $\magnitudeI : \magnitudeIII :: \magnitudeVI : \magnitudeVIII$,\\
but $\magnitudeIII : \magnitudeIV :: \magnitudeV : \magnitudeVI$,\\
therefore again, by the first case,\\
$\magnitudeIII : \magnitudeIV :: \magnitudeV : \magnitudeVI$;\\
e assim por diante, qualquer que seja o número de tais magnitudes.

$\therefore$ If there be any number, \&c.
\end{center}

\vfill\pagebreak

\starttheorem{Prop. XXIV. Theor.}\label{prop:V.XXIV}
\defineNewPicture{
byMagnitudeSymbolDefine.I("wedgeDown", byred, 0);
byMagnitudeSymbolDefine.II("semicircleDown", byblack, 1);
byMagnitudeSymbolDefine.III("square", byblue, 0);
byMagnitudeSymbolDefine.IV("rhombus", byyellow, 0);
byMagnitudeSymbolDefine.V("sectorUp", byred, 1);
byMagnitudeSymbolDefine.VI("circle", byblue, 0);
}
\problem{S}{e}{o primeiro tiver para o segundo a mesma proporção que o terceiro tem para o quarto, e o quinto para o segundo a mesma que o sexto tem para o quarto, o primeiro e o quinto juntos terão para o segundo a mesma proporção que o terceiro e o sexto juntos têm para o quarto.}

\begin{center}
%\unprotect
%\ \hfill\vbox{
%\offinterlineskip\lineskip3pt
%\halign{\hfil # \hfil & \hfil # \hfil & \hfil # \hfil & \hfil # \hfil\\
\begin{tabular}{c c c c}
{\footnotesize First} & {\footnotesize Second} & {\footnotesize Third} & {\footnotesize Fourth} \\
\drawMagnitude{I} & \drawMagnitude{II} & \drawMagnitude{III} & \drawMagnitude{IV} \\
{\footnotesize Fifth} & & {\footnotesize Sixth} & \\
\drawMagnitude{V} & & \drawMagnitude{VI} & \\
\end{tabular}
%}}\hfill\
%\protect
\end{center}

\begin{center}
Let $\magnitudeI : \magnitudeII :: \magnitudeIII : \magnitudeIV$,\\
and $\magnitudeV : \magnitudeII :: \magnitudeVI : \magnitudeIV$,\\
then $\magnitudeI + \magnitudeV : \magnitudeII :: \magnitudeIII + \magnitudeVI : \magnitudeIV$.

For $\magnitudeV : \magnitudeII :: \magnitudeVI : \magnitudeIV$ \byref{\hypref},\\
and $\magnitudeII : \magnitudeI :: \magnitudeIV : \magnitudeIII$ \byref{\hypref} and (invert.),

$\therefore \magnitudeV : \magnitudeI :: \magnitudeVI : \magnitudeIII$ \byref{prop:V.XXII};\\
e, porque essas magnitudes são proporcionais, elas são proporcionais quando tomadas em conjunto,

$\therefore \magnitudeI + \magnitudeV : \magnitudeV :: \magnitudeVI + \magnitudeIII : \magnitudeVI$ \byref{prop:V.XVIII},\\
but $\magnitudeV : \magnitudeII :: \magnitudeVI : \magnitudeIV$ \byref{\hypref},

$\therefore \magnitudeI + \magnitudeV : \magnitudeII :: \magnitudeVI + \magnitudeIII : \magnitudeIV$ \byref{prop:V.XXII}

$\therefore$ If the first, \&c.
\end{center}

\vfill\pagebreak

\starttheorem{Prop. XXV. Theor.}\label{prop:V.XXV}
\defineNewPicture{
byMagnitudeSymbolDefine.Ia("wedgeDown", byred, 0);
byMagnitudeSymbolDefine.III("semicircleDown", byblack, 1);
byMagnitudeSymbolDefine.IIa("square", byblue, 0);
byMagnitudeSymbolDefine.IV("rhombus", byyellow, 0);
}
\problem{S}{e}{quatro grandezas do mesmo tipo são proporcionais, a maior e a menor delas juntas são maiores que as outras duas juntas.}

\begin{center}
Sejam quatro grandezas, $\drawMagnitude{Ia} + \drawMagnitude{III}, \drawMagnitude{IIa} + \drawMagnitude{IV}, \magnitudeIII, \mbox{ e } \magnitudeIV$, do mesmo tipo, proporcionais, ou seja,\\
$\magnitudeIa + \magnitudeIII : \magnitudeIIa + \magnitudeIV :: \magnitudeIII : \magnitudeIV$,\\
e seja $\magnitudeIa + \magnitudeIII$ o maior dos quatro, e consequentemente por \byref{prop:V.A,prop:V.XIV}, \magnitudeIV\ é o menor;\\
then will $\magnitudeIa + \magnitudeIII + \magnitudeIV \mbox{ be } > \magnitudeIIa + \magnitudeIV + \magnitudeIII$;\\
$\because \magnitudeIa + \magnitudeIII : \magnitudeIIa + \magnitudeIV :: \magnitudeIII : \magnitudeIV$, %because

$\therefore \magnitudeIa : \magnitudeIIa :: \magnitudeIa + \magnitudeIII : \magnitudeIIa + \magnitudeIV$ \byref{prop:V.XIX},\\
but $\magnitudeIa + \magnitudeIII > \magnitudeIIa + \magnitudeIV$ \byref{\hypref},

$\therefore \magnitudeIa > \magnitudeIIa$ \byref{prop:V.A};\\
a cada um deles adicione $\magnitudeIII + \magnitudeIV$,\\
$\therefore \magnitudeIa + \magnitudeIII + \magnitudeIV > \magnitudeIIa + \magnitudeIII + \magnitudeIV$.

$\therefore$ If four magnitudes, \&c.
\end{center}

\vfill\pagebreak

\startdefinition[10]{Definição X}\label{def:V.X}
\def\varA{\textcolor{byred}{A}}
\def\varB{\textcolor{byyellow}{B}}
\def\varC{\textcolor{byblue}{C}}
\def\vararS{\textcolor{byred}{ar^2}}
\def\varar{\textcolor{byyellow}{ar}}
\def\vara{\textcolor{byblue}{a}}
Quando três grandezas são proporcionais, diz-se que a primeira tem para a terceira o razão duplicada daquele que tem para a segunda.

Por exemplo, se $\varA$, $\varB$, $\varC$, forem proporcionais continuados, ou seja, $\varA : \varB :: \varB : \varC$, $\varA$ tem que $\varC$ o razão duplicada de $\varA : \varB$;

\begin{center}
or $\dfrac{\varA}{\varC} = \mbox{ a praça de } \dfrac{\varA}{\varB}$.

Esta propriedade será mais facilmente observada nas grandezas $\vararS$, $\varar$, $\vara$, para $\vararS : \varar :: \varar : \vara$;

and $\dfrac{\vararS}{\vara} = r^2 = \mbox{ a praça de } \dfrac{\vararS}{\varar} = r$,\\
or of $\vara$, $\varar$, $\vararS$;\\
for $\dfrac{\vara}{\vararS} = \dfrac{1}{r^2} = \mbox{ o quadrado de } \dfrac{\vara}{\varar} = \dfrac{1}{r}$.
\end{center}

\vfill\pagebreak

\startdefinition[11]{Definição XI}\label{def:V.XI}
\def\varA{\textcolor{byred}{A}}
\def\varB{\textcolor{byyellow}{B}}
\def\varC{\textcolor{byblue}{C}}
\def\varD{\textcolor{byblack}{D}}
\def\vararQ{\textcolor{byred}{ar^3}}
\def\vararS{\textcolor{byyellow}{ar^2}}
\def\varar{\textcolor{byblue}{ar}}
\def\vara{\textcolor{byblack}{a}}
Quando quatro grandezas são proporcionais contínuas, diz-se que a primeira tem para a quarta o razão triplicada daquela que tem para a segunda; e assim por diante, quadruplicar, \&c.\ aumentando a denominação ainda pela unidade, em qualquer número de proporcionais.

Por exemplo, sejam $\varA$, $\varB$, $\varC$, $\varD$ quatro proporcionais contínuas, ou seja, $\varA : \varB :: \varB : \varC :: \varC : \varD$; diz-se que $\varA$ tem para $\varD$, o razão triplicada de $\varA$ a $\varB$;

\begin{center}
or $\dfrac{\varA}{\varD} = \mbox{o cubo de } \dfrac{\varA}{\varB}$.
\end{center}

Esta definição será melhor compreendida e aplicada a um número maior de magnitudes do que quatro que são proporcionais continuadas, como segue: -

\begin{center}
Sejam $\vararQ$, $\vararS$, $\varar$, $\vara$ quatro magnitudes em proporção contínua, ou seja, $\vararQ : \vararS :: \vararS : \varar :: \varar : \vara$,\\
then $\dfrac{\vararQ}{\vara} = r^3 = \mbox{ o cubo de } \dfrac{\vararQ}{\vararS} = r$.
\end{center}

Ou seja $ar^5$, $ar^4$, $ar^3$, $ar^2$, $ar$, $a$, seis magnitudes em proporção, isto é

\begin{center}
$ar^5 : ar^4 :: ar^4 : ar^3 :: ar^3 : ar^2 :: ar^2 : ar :: ar : a$,\\
then the ratio $\dfrac{ar^5}{a} = r^5 = \mbox{ a quinta potência de } \dfrac{ar^5}{ar^4} = r$
\end{center}

Or, let $a$, $ar$, $ar^2$, $ar^3$, $ar^4$, be five magnitudes in continued proportion; then $\dfrac{a}{ar^4} = \dfrac{1}{r^4} = \mbox{ the fourth power of } \dfrac{a}{ar} = \dfrac{1}{r}$.

\vfill\pagebreak

\startdefinitionAZ[1]{Definição A}\label{def:V.A}
\def\varA{\textcolor{byred}{A}}
\def\varB{\textcolor{byred}{B}}
\def\varC{\textcolor{byred}{C}}
\def\varD{\textcolor{byred}{D}}
\def\varE{\textcolor{byblue}{E}}
\def\varF{\textcolor{byblue}{F}}
\def\varG{\textcolor{byblue}{G}}
\def\varH{\textcolor{byblue}{H}}
\def\varK{\textcolor{byblue}{K}}
\def\varL{\textcolor{byblue}{L}}
\def\varM{\textcolor{byyellow}{M}}
\def\varN{\textcolor{byyellow}{N}}
Para conhecer um razão composta:—

Quando há qualquer número de grandezas da mesma espécie, diz-se que a primeira tem para a última a razão composta pela razão que a primeira tem para a segunda, e pela razão que a segunda tem para a terceira, e da razão que a terceira tem para a quarta; e assim por diante, até a última magnitude.

\marginpar{
~\hfill\fbox{\parbox{0.5\marginparwidth}{\centering
~$\varA\ \varB\ \varC\ \varD$~\\
~$\varE\ \varF\ \varG\ \varH\ \varK\ \varL$~\\
~$\varM\ \varN$~
}}\hfill~}
Por exemplo, se $\varA, \varB, \varC, \varD$ tiver quatro grandezas do mesmo tipo, diz-se que o primeiro $\varA$ tem para o último $\varD$ a razão composta pela razão de $\varA$ para $\varB$, e da razão de $\varB$ para $\varC$, e da razão de $\varC$ para $\varD$; ou, diz-se que a proporção de $\varA$ para $\varD$ é composta pelas proporções de $\varA$ para $\varB$, $\varB$ para $\varC$ e $\varC$ para $\varD$.

E se $\varA$ tem para $\varB$ a mesma proporção que $\varE$ tem para $\varF$, e $\varB$ para $\varC$ a mesma proporção que $\varG$ tem para $\varH$, e $\varC$ para $\varD$ a mesma que $\varK$ tem para $\varL$; então, por esta definição, diz-se que $\varA$ tem para $\varD$ a proporção composta de proporções que são iguais às proporções de $\varE$ para $\varF$, $\varG$ para $\varH$ e $\varK$ para $\varL$. E a mesma coisa deve ser dita, $\varA$ tem para $\varD$ a proporção composta pelas proporções de $\varE$ para $\varF$, $\varG$ para $\varH$ e $\varK$ para $\varL$.

Da mesma maneira, supondo as mesmas coisas; se $\varM$ tem para $\varN$ a mesma proporção que $\varA$ tem para $\varD$, então, por questões de abreviação, diz-se que $\varM$ tem para $\varN$ a proporção composta pelas proporções $\varE$ para $\varF$, $\varG$ para $\varH$, e $\varK$ a $\varL$.

\vskip \baselineskip

Esta definição pode ser melhor compreendida a partir de uma ilustração aritmética ou algébrica; pois, de fato, uma razão composta de várias outras razões nada mais é do que uma razão que tem como antecedente o produto contínuo de todos os antecedentes das razões compostas, e como consequente o produto contínuo de todos os consequentes das razões compostas.

\begin{center}
Assim, a razão composta pelas razões de\\
$\textcolor{byred}{2} : \textcolor{byred}{3}$, $\textcolor{byyellow}{4} : \textcolor{byyellow}{7}$, $\textcolor{byblue}{6} : \textcolor{byblue}{11}$, $\textcolor{byblack}{2} : \textcolor{byblack}{5}$,\\
é a razão de $\textcolor{byred}{2} \times \textcolor{byyellow}{4} \times \textcolor{byblue}{6} \times \textcolor{byblack}{2} : \textcolor{byred}{3} \times \textcolor{byyellow}{7} \times \textcolor{byblue}{11} \times \textcolor{byblack}{5}$,\\
ou a razão de $96 : 1155$ ou $32 : 385$.
\end{center}

\def\varA{\textcolor{byred}{A}}
\def\varB{\textcolor{byred}{B}}
\def\varC{\textcolor{byyellow}{C}}
\def\varD{\textcolor{byyellow}{D}}
\def\varE{\textcolor{byblue}{E}}
\def\varF{\textcolor{byblue}{F}}
E das grandezas $\varA, \varB, \varC, \varD, \varE, \varF$, da mesma espécie, $\varA : \varF$ é a razão composta das razões de

\begin{center}
$\varA : \varB$, $\varB : \varC$, $\varC : \varD$, $\varD : \varE$, $\varE : \varF$;\\
for $\varA \times \varB \times \varC \times \varD \times \varE : \varB \times \varC \times \varD \times \varE \times \varF$,\\
ou $\dfrac{\varA \times \varB \times \varC \times \varD \times \varE}{\varB \times \varC \times \varD \times \varE \times \varF} = \dfrac{\varA}{\varF}$, ou a razão de $\varA : \varF$.
\end{center}

\vfill\pagebreak

\starttheoremAZ{Prop. F. Theor.}\label{prop:V.F} % Proposition with letters
\def\varA{\textcolor{byred}{A}}
\def\varB{\textcolor{byyellow}{B}}
\def\varC{\textcolor{byyellow}{C}}
\def\varD{\textcolor{byyellow}{D}}
\def\varE{\textcolor{byred}{E}}
\def\varF{\textcolor{byblue}{F}}
\def\varG{\textcolor{byyellow}{G}}
\def\varH{\textcolor{byyellow}{H}}
\def\varK{\textcolor{byyellow}{K}}
\def\varL{\textcolor{byblue}{L}}
\problem{A}{s}{proporções que são compostas pelas mesmas proporções são iguais entre si.}

\begin{center}
Seja $\varA : \varB :: \varF : \varG$,\\
 $\varB : \varC :: \varG : \varH$,\\
$\varC : \varD :: \varH : \varK$,\\
and $\varD : \varE :: \varK : \varL$,
\end{center}

\marginpar{
~\hfill\fbox{\parbox{0.5\marginparwidth}{\centering
~$\varA\ \varB\ \varC\ \varD\ \varE$~\\
~$\varF\ \varG\ \varH\ \varK\ \varL$~
}}\hfill~}
Então, a proporção que é composta pelas razões de $\varA : \varB$, $\varB : \varC$, $\varC : \varD$, $\varD : \varE$ ou a razão de $\varA : \varE$ é a mesma que a proporção composta pelas proporções $\varF : \varG$, $\varG : \varH$, $\varH : \varK$, $\varK : \varL$ ou a razão de $\varF : \varL$.

\begin{center}
For $\dfrac{\varA}{\varB} = \dfrac{\varF}{\varG}$,\\
$\dfrac{\varB}{\varC} = \dfrac{\varG}{\varH}$,\\
$\dfrac{\varC}{\varD} = \dfrac{\varH}{\varK}$,\\
and $\dfrac{\varD}{\varE} = \dfrac{\varK}{\varL}$;\\
$\therefore \dfrac{\varA \times \varB \times \varC \times \varD}{\varB \times \varC \times \varD \times \varE} = \dfrac{\varF \times \varG \times \varH \times \varK}{\varG \times \varH \times \varK \times \varL}$,\\
and $\therefore \dfrac{\varA}{\varE} = \dfrac{\varF}{\varL}$,\\
ou a razão de $\varA : \varE$ é igual à razão $\varF : \varL$.
\end{center}

O mesmo pode ser demonstrado em relação a qualquer número de rácios assim circunstanciados.

\begin{center}
Em seguida, seja $\varA : \varB :: \varK : \varL$,\\
 $\varB : \varC :: \varH : \varK$,\\
$\varC : \varD :: \varG : \varH$,\\
$\varD : \varE :: \varF : \varG$,
\end{center}

Então, a proporção que é composta pelas razões de $\varA : \varB$, $\varB : \varC$, $\varC : \varD$, $\varD : \varE$ ou a razão de $\varA : \varE$ é a mesma que a proporção composta pelas razões de $\varK : \varL$, $\varH : \varK$, $\varG : \varH$, $\varF : \varG$ ou a razão de $\varF : \varL$

\begin{center}
For $\dfrac{\varA}{\varB} = \dfrac{\varK}{\varL}$,\\
$\dfrac{\varB}{\varC} = \dfrac{\varH}{\varK}$,\\
$\dfrac{\varC}{\varD} = \dfrac{\varG}{\varH}$,\\
and $\dfrac{\varD}{\varE} = \dfrac{\varF}{\varG}$;\\
$\therefore \dfrac{\varA \times \varB \times \varC \times \varD}{\varB \times \varC \times \varD \times \varE} = \dfrac{\varK \times \varH \times \varG \times \varF}{\varL \times \varK \times \varH \times \varG}$,\\
and $\therefore \dfrac{\varA}{\varE} = \dfrac{\varF}{\varL}$,\\
ou a razão de $\varA : \varE$ é igual à razão $\varF : \varL$.

$\therefore$ Ratios which are compounded, \&c.
\end{center}
\vfill\pagebreak

\starttheoremAZ{Prop. G. Theor.}\label{prop:V.G} % Proposition with letters
\def\varA{\textcolor{byblue}{A}}
\def\varB{\textcolor{byblue}{B}}
\def\varC{\textcolor{byblue}{C}}
\def\varD{\textcolor{byblue}{D}}
\def\varE{\textcolor{byblue}{E}}
\def\varF{\textcolor{byblue}{F}}
\def\varG{\textcolor{byblue}{G}}
\def\varH{\textcolor{byblue}{H}}
\def\varP{\textcolor{byred}{P}}
\def\varQ{\textcolor{byred}{Q}}
\def\varR{\textcolor{byred}{R}}
\def\varS{\textcolor{byred}{S}}
\def\varT{\textcolor{byred}{T}}
\def\varV{\textcolor{byyellow}{V}}
\def\varW{\textcolor{byyellow}{W}}
\def\varX{\textcolor{byyellow}{X}}
\def\varY{\textcolor{byyellow}{Y}}
\def\varZ{\textcolor{byyellow}{Z}}
\problem{I}{f}{several ratios be the same to several ratios, each to each, the ratio which is compounded of ratios which are the same to the first ratios, each to each, shall be the same to the ratio compounded of ratios which are the same to the other ratios, each to each.}

\begin{center}
%\unprotect
%~\hfill\vbox{\halign{\vrule height2.2ex depth1ex\ \hskip 5pt\ # & # & # & # & # & # & # & # \hskip 20pt & # & # & # & # & # \hskip 5pt \vrule \\
\begin{tabular}{c c c c c c c c c c c c c}
\noalign{\hrule}
$\varA$ & $\varB$ & $\varC$ & $\varD$ & $\varE$ & $\varF$ & $\varG$ & $\varH$ &
$\varP$ & $\varQ$ & $\varR$ & $\varS$ & $\varT$ \\
$a$ & $b$ & $c$ & $d$ & $e$ & $f$ & $g$ & $h$ &
$\varV$ & $\varW$ & $\varX$ & $\varY$ & $\varZ$ \\
\noalign{\hrule}
%}}\hfill~
%\protect
\end{tabular}
\vskip 0.5\baselineskip

%\unprotect
\setbox0\vbox{
%\halign{ # & # & # \\
\begin{tabular}{c c c}
$\begin{aligned}
\mbox{ If \ } \varA : \varB & :: a : b\\
\varC : \varD & :: c : d\\
\varE : \varF & :: e : f\\
\mbox{ and \ } \varG : \varH & :: g : h\\
\end{aligned}$ &
$\begin{aligned}
\mbox{ e \ } \varA : \varB & :: \varP : \varQ\\
\varC : \varD & :: \varQ : \varR\\
\varE : \varF & :: \varR : \varS\\
\varG : \varH & :: \varS : \varT\\
\end{aligned}$ &
$\begin{aligned}
\mbox{ and \ } a : b & :: \varV : \varW\\
c : d & :: \varW : \varX\\
e : f & :: \varX : \varY\\
f : h & :: \varY : \varZ\\
\end{aligned}$ \\
%}
\end{tabular}
}
%\protect
\box0\hskip-\wd0\hskip\textwidth\box0

then $\varP : \varT = \varV : \varZ$.

$\begin{aligned}
\mbox{ For } \frac{\varP}{\varQ} = \frac{\varA}{\varB}
	& = \frac{a}{b} = \frac{\varV}{\varW} \\
\frac{\varQ}{\varR} = \frac{\varC}{\varD}
	& = \frac{c}{d} = \frac{\varW}{\varX} \\
\frac{\varR}{\varS} = \frac{\varE}{\varF}
	& = \frac{e}{f} = \frac{\varX}{\varY} \\
\frac{\varS}{\varT} = \frac{\varG}{\varH}
	& = \frac{g}{h} = \frac{\varY}{\varZ} \\
\end{aligned}$

and $\therefore
\dfrac{\varP \times \varQ \times \varR \times \varS}
	{\varQ \times \varR \times \varS \times \varT} =
\dfrac{\varV \times \varW \times \varX \times \varY}
	{\varW \times \varX \times \varY \times \varZ}$,

and $\therefore \dfrac{\varP}{\varT} = \dfrac{\varV}{\varZ}$,

or $\varP : \varT = \varV : \varZ$.

$\therefore$ If several ratios, \&c.
\end{center}
\vfill\pagebreak

\starttheoremAZ{Prop. H. Theor.}\label{prop:V.H} % Proposition with letters
\def\varA{\textcolor{byred}{A}}
\def\varB{\textcolor{byred}{B}}
\def\varC{\textcolor{byred}{C}}
\def\varD{\textcolor{byred}{D}}
\def\varE{\textcolor{byred}{E}}
\def\varF{\textcolor{byred}{F}}
\def\varG{\textcolor{byred}{G}}
\def\varH{\textcolor{byred}{H}}
\def\varP{\textcolor{byblue}{P}}
\def\varQ{\textcolor{byblue}{Q}}
\def\varR{\textcolor{byblue}{R}}
\def\varS{\textcolor{byblue}{S}}
\def\varT{\textcolor{byblue}{T}}
\def\varX{\textcolor{byblue}{X}}
\problem{S}{e}{uma proporção que é composta por vários índices for igual a uma proporção que é composta por vários outros índices; e se uma das primeiras razões, ou a razão que é composta por várias delas, for igual a uma das últimas razões, ou à razão que é composta por várias delas; então o rácio restante do primeiro, ou, se houver mais de um, o rácio composto dos rácios restantes, será igual aos restantes rácios, será igual ao rácio restante do último, ou, se houver mais de um, ao rácio composto desses rácios restantes.}

\marginpar{
~\hfill\fbox{\parbox{0.5\marginparwidth}{\centering
~$\varA\ \varB\ \varC\ \varD\ \varE\ \varF\ \varG\ \varH$~\\
~$\varP\ \varQ\ \varR\ \varS\ \varT\ \varX$~\\
}}\hfill~}
Sejam $\varA : \varB$, $\varB : \varC$, $\varC : \varD$, $\varD : \varE$, $\varE : \varF$, $\varF : \varG$, $\varG : \varH$ as primeiras proporções, e $\varP : \varQ$, $\varQ : \varR$, $\varR : \varS$, $\varS : \varT$, $\varT : \varX$, o outros rácios; além disso, seja $\varA : \varH$, que é composto pelas primeiras relações, igual à razão de $\varP : \varX$, que é a razão composta pelas demais relações; e, deixe a proporção de $\varA : \varE$, que é composta pelas proporções de $\varA : \varB$, $\varB : \varC$, $\varC : \varD$, $\varD : \varE$, ser a mesma que a proporção de $\varP : \varR$, que é composta pelas proporções $\varP : \varQ$, $\varQ : \varR$.

Então, a razão que é composta pelas primeiras razões restantes, ou seja, a razão composta pelas razões $\varE : \varF$, $\varF : \varG$, $\varG : \varH$, ou seja, a razão de $\varE : \varH$, será a mesma que a razão de $\varR : \varX$, que é composta pelas razões de $\varR : \varS$, $\varS : \varT$, $\varT : \varX$, as demais proporções.

\begin{center}
Because $\dfrac
{\varA \times \varB \times \varC \times \varD \times \varE \times \varF \times \varG}
{\varB \times \varC \times \varD \times \varE \times \varF \times \varG \times \varH} =
\dfrac
{\varP \times \varQ \times \varR \times \varS \times \varT}
{\varQ \times \varR \times \varS \times \varT \times \varX}$,

or $\dfrac
{\varA \times \varB \times \varC \times \varD}
{\varB \times \varC \times \varD \times \varE}
\times
\dfrac
{\varE \times \varF \times \varG}
{\varF \times \varG \times \varH} =
\dfrac
{\varP \times \varQ}
{\varQ \times \varR}
\times
\dfrac
{\varR \times \varS \times \varT}
{\varS \times \varT \times \varX}$,

and $\dfrac
{\varA \times \varB \times \varC \times \varD}
{\varB \times \varC \times \varD \times \varE} =
\dfrac
{\varP \times \varQ}
{\varQ \times \varR}$

$\therefore \dfrac
{\varE \times \varF \times \varG}
{\varF \times \varG \times \varH} =
\dfrac
{\varR \times \varS \times \varT}
{\varS \times \varT \times \varX}$,

$\therefore \dfrac{\varE}{\varH} = \dfrac{\varR}{\varX}$,

$\therefore \varE : \varH = \varR : \varX$.

$\therefore$ Se uma proporção que, \&c.
\end{center}
\vfill\pagebreak

\starttheoremAZ[11]{Prop. K. Theor.}\label{prop:V.K} % Proposition with letters
\def\varA{\textcolor{byred}{A}}
\def\varB{\textcolor{byred}{B}}
\def\varC{\textcolor{byred}{C}}
\def\varD{\textcolor{byred}{D}}
\def\varE{\textcolor{byred}{E}}
\def\varF{\textcolor{byred}{F}}
\def\varG{\textcolor{byred}{G}}
\def\varH{\textcolor{byred}{H}}
\def\varK{\textcolor{byred}{K}}
\def\varL{\textcolor{byred}{L}}
\def\varM{\textcolor{byred}{M}}
\def\varN{\textcolor{byred}{N}}
\def\varO{\textcolor{byyellow}{O}}
\def\varP{\textcolor{byyellow}{P}}
\def\varQ{\textcolor{byyellow}{Q}}
\def\varR{\textcolor{byyellow}{R}}
\def\varS{\textcolor{byyellow}{S}}
\def\varT{\textcolor{byyellow}{T}}
\def\varV{\textcolor{byyellow}{V}}
\def\varW{\textcolor{byyellow}{W}}
\def\varX{\textcolor{byyellow}{X}}
\def\varY{\textcolor{byyellow}{Y}}
\def\vara{\textcolor{byblue}{a}}
\def\varb{\textcolor{byblue}{b}}
\def\varc{\textcolor{byblue}{c}}
\def\vard{\textcolor{byblue}{d}}
\def\vare{\textcolor{byblue}{e}}
\def\varf{\textcolor{byblue}{f}}
\def\varg{\textcolor{byblue}{g}}
\problem{S}{e}{houver qualquer número de razões, e qualquer número de outras razões, de modo que a razão que é composta de razões, que são iguais às primeiras razões, cada uma para cada, seja a mesma que a razão que é composta de razões, que são as mesmas, cada uma para cada, até as últimas razões - e se uma das primeiras razões, ou a razão que é composta de razões, que são iguais a várias das primeiras razões, cada uma para cada, for a mesma para uma das últimas proporções, ou à proporção que é composta de proporções, que são iguais, cada uma para cada uma, a várias das últimas proporções - então a proporção restante da primeira; ou, se houver mais de um, a proporção que é composta de razões iguais, cada uma para cada uma, às demais proporções da primeira, será igual à proporção restante da última; ou, se houver mais de um, à proporção que é composta de razões, que são iguais, cada uma para cada uma, às restantes razões.}

%\unprotect
%~\hfill\vbox{\halign{\vrule height2.2ex depth1ex\ \hskip 2pt\
%#&#& #&#& #&#& #&#& #&#& #&# & \hfil#\hfil
%\hskip 2pt \vrule \\
\begin{center}
\noindent\begin{tabular}{c c c c c c c c c c c c c}
\noalign{\hrule}
& & & & & & & $h$ & $k$ & $m$ & $n$ & $s$ & \\
$\varA$ & $\varB$, & $\varC$ & $\varD$, & $\varE$ & $\varF$, & $\varG$ & $\varH$, & $\varK$ & $\varL$, & $\varM$ & $\varN$ & $\vara\ \varb\ \varc\ \vard\ \vare\ \varf\ \varg$ \\
$\varO$ & $\varP$, & $\varQ$ & $\varR$, & $\varS$ & $\varT$, & $\varV$ & $\varW$, & $\varX$ & $\varY$ & & & $h\ k\ l\ m\ n\ p$ \\
& & $a$ & $b$ & $k$ & $m$ & & $e$ & $f$ & $g$ & & & \\
\noalign{\hrule}
\end{tabular}
\end{center}
%}}\hfill~
%\protect

Let $\varA : \varB$, $\varC : \varD$, $\varE : \varF$, $\varG : \varH$, $\varK : \varL$, $\varM : \varN$, be the first ratios, and $\varO : \varP$, $\varQ : \varR$, $\varS : \varT$, $\varV : \varW$, $\varX : \varY$, the other ratios;

\begin{center}
$\begin{aligned}
\mbox{e seja }
\varA : \varB &= \vara : \varb ,\\
\varC : \varD &= \varb : \varc ,\\
\varE : \varF &= \varc : \varf ,\\
\varG : \varH &= \vard : \vare ,\\
\varK : \varL &= \vare : \varf ,\\
\varM : \varN &= \varf : \varg .
\end{aligned}$
\end{center}

Então, pela definição de um razão composta, a proporção de $\vara : \varg$ é composta pelas proporções de $\vara : \varb$, $\varb : \varc$, $\varc : \vard$, $\vard : \vare$, $\vare : \varf$, $\varf : \varg$, que são iguais à proporção de $\varA : \varB$, $\varC : \varD$, $\varE : \varF$, $\varG : \varH$, $\varK : \varL$, $\varM : \varN$, cada um para cada.

\begin{center}
$\begin{aligned}
\mbox{Also, }
\varO : \varP &= h : k ,\\
\varQ : \varR &= k : l ,\\
\varS : \varT &= l : m ,\\
\varV : \varW &= m : n ,\\
\varX : \varY &= n : p ,\\
\end{aligned}$
\end{center}

Então a razão de $h : p$ será a proporção composta pelas razões de $h : k$, $k : l$, $l : m$, $m : n$, $n : p$, que são iguais às proporções de $\varO : \varP$, $\varQ : \varR$, $\varS : \varT$, $\varV : \varW$, $\varX : \varY$, cada um para cada.

\begin{center}
$\therefore$ by the hypothesis $\vara : \varg = h : p$.
\end{center}

Além disso, deixe a proporção que é composta pelas proporções de $\varA : \varB$, $\varC : \varD$, duas das primeiras proporções (ou as proporções de $\vara : \varc$ para $\varA : \varB = \vara : \varb$ e $\varC : \varD = \varb : \varc$), seja a mesma que a proporção de $a : d$, que é composta pelas proporções de $a : b$, $b : c$, $c : d$, que são iguais às proporções de $\varO : \varP$, $\varQ : \varR$, $\varS : \varT$, três das outras proporções.

E deixe as proporções de $h : s$, que são compostas pelas proporções de $h : k$, $k : m$, $m : n$, $n : s$, que são iguais às primeiras proporções restantes, ou seja, $\varE : \varF$, $\varG : \varH$, $\varK : \varL$, $\varM : \varN$; além disso, seja a razão de $e : g$ aquela que é composta pelas razões $e : f$, $f : g$, que são iguais, cada uma para cada, às restantes outras relações, nomeadamente $\varV : \varW$, $\varX : \varY$. Então as proporções de $h : s$ serão iguais às proporções de $e : g$; ou $h : s = e : g$.

\begin{center}
For $\dfrac
{\varA \times \varC \times \varE \times \varG \times \varK \times \varM}
{\varB \times \varD \times \varF \times \varH \times \varL \times \varN}
=
\dfrac
{\vara \times \varb \times \varc \times \vard \times \vare \times \varf}
{\varb \times \varc \times \vard \times \vare \times \varf \times \varg}$,

and $\dfrac
{\varO \times \varQ \times \varS \times \varV \times \varX}
{\varP \times \varR \times \varT \times \varW \times \varY}
=
\dfrac
{h \times k \times l \times m \times n}
{k \times l \times m \times n \times p}$,\\
pela composição dos rácios;

$\therefore \dfrac
{\vara \times \varb \times \varc \times \vard \times \vare \times \varf}
{\varb \times \varc \times \vard \times \vare \times \varf \times \varg}
=
\dfrac
{h \times k \times l \times m \times n}
{k \times l \times m \times n \times p}$ \byref{\hypref},

or $\dfrac{\vara \times \varb}{\varb \times \varc}
\times
\dfrac
{\varc \times \vard \times \vare \times \varf}
{\vard \times \vare \times \varf \times \varg}
=
\dfrac
{h \times k \times l}
{k \times l \times m}
\times
\dfrac
{m \times n}
{n \times p}$,

mas $\dfrac{\vara \times \varb}{\varb \times \varc} = \dfrac{\varA \times \varC}{\varB \times \varD}
=
\dfrac
{\varO \times \varQ \times \varS}
{\varP \times \varR \times \varT}
=
\dfrac{a \times b \times c}{b \times c \times d}
=
\dfrac{h \times k \times l}{k \times l \times m}$;

$\therefore
\dfrac
{\varc \times \vard \times \vare \times \varf}
{\vard \times \vare \times \varf \times \varg}
=
\dfrac
{m \times n}
{n \times p}$.

And $\dfrac
{\varc \times \vard \times \vare \times \varf}
{\vard \times \vare \times \varf \times \varg}
= \dfrac
{h \times k \times l \times m \times n}
{k \times l \times m \times n \times p}$ \byref{\hypref},

and $\dfrac
{m \times n}
{n \times p}
=
\dfrac
{e \times f}
{f \times g}$ \byref{\hypref},

$\therefore \dfrac
{h \times k \times l \times m \times n}
{k \times l \times m \times n \times p}
= \dfrac{e f}{f g}$,

$\therefore \dfrac{h}{s} = \dfrac{e}{g}$,\\
$\therefore h : s = e : g$.

$\therefore$ If there be any number, \&c.
\end{center}

\part{Livro VI}

\chapter*{Definições}

\startdefinition{}\label{def:VI.I}
\defineNewPicture{
pair a, b, c, A, B, C, d, e, f, g, D, E ,F ,G;
a := (0, 0);
b := (u, 0);
c := (1/2u, u);
A := (a scaled 2) shifted (u, -u);
B := (b scaled 2) shifted (u, -u);
C := (c scaled 2) shifted (u, -u);
d := (0, 0);
e := (2/3u, 0);
f := (u, 2/3u);
g := (1/3u, 2/3u);
D := d scaled 2;
E := e scaled 2;
F := f scaled 2;
G := g scaled 2;
d := d shifted (0, -5/2u);
e := e shifted (0, -5/2u);
f := f shifted (0, -5/2u);
g := g shifted (0, -5/2u);
D := D shifted (u, -5/2u);
E := E shifted (u, -5/2u);
F := F shifted (u, -5/2u);
G := G shifted (u, -5/2u);
draw byPolygon(a,b,c)(byblue);
draw byPolygon(A,B,C)(byblue);
draw byPolygon(d,e,f,g)(byred);
draw byPolygon(D,E,F,G)(byred);
}\drawCurrentPictureInMargin
Diz-se que figuras retilíneas são semelhantes quando têm seus ângulos iguais, cada um a cada um, e os lados que compreendem os ângulos iguais proporcionais.

\startdefinition{}\label{def:VI.II}
Diz-se que dois lados de uma figura são reciprocamente proporcionais a dois lados de outra figura quando um dos lados da primeira está para o segundo como o lado restante da segunda está para o lado restante da primeira.

\startdefinition{}\label{def:VI.III}
Diz-se que uma reta é dividida em extrema e média razão quando o todo está para o segmento maior assim como o segmento maior está para o menor.

\startdefinition{}\label{def:VI.IV}
\defineNewPicture{
pair d[];
d1 := (0, 0);
d2 := (5/2u, 0);
d3 := (5u, 0);
d4 := (15/2u, 0);
numeric h;
h := 6/5u;
pair A, B, C, H;
A := (0, 0) shifted d1;
B := (3/2u, 0) shifted d1;
C := (3/4u, h) shifted d1;
H := (xpart(C), 0);
draw byPolygon (A,B,C)(byyellow);
draw byLine(C, H)(byyellow, DASHED_LINE, REGULAR_WIDTH);
pair A, B, C, D, H, G;
A := (0, 0) shifted d2;
B := (1/2u, 0) shifted d2;
C := (3/2u, h) shifted d2;
D := (u, h) shifted d2;
H := (xpart(C), 0);
G := (xpart(D), 0);
draw byPolygon (A,B,C,D)(byblue);
byLineDefine(C, H, byblue, DASHED_LINE, REGULAR_WIDTH);
draw byLine(D, G, byblue, DASHED_LINE, REGULAR_WIDTH);
byLineDefine(A, H, byblue, DASHED_LINE, REGULAR_WIDTH);
draw byNamedLineSeq(0)(CH,AH);
pair A, B, C, D, H;
A := (0, 0) shifted d3;
B := (2/3u, 0) shifted d3;
C := (u, h) shifted d3;
H := (xpart(C), 0);
draw byPolygon (A,B,C)(byblack);
byLineDefine(C, H, byblack, DASHED_LINE, REGULAR_WIDTH);
byLineDefine(A, H, byblack, DASHED_LINE, REGULAR_WIDTH);
draw byNamedLineSeq(0)(CH,AH);
pair A, B, C, D, E, H;
A := (0, 0) shifted d4;
B := (1/2u, 0) shifted d4;
C := (u, 1/2h) shifted d4;
D := (1/4u, h) shifted d4;
E := (-1/2u, 1/2h) shifted d4;
H := (xpart(D), 0);
draw byPolygon (A,B,C,D,E)(byred);
draw byLine(D, H)(byred, DASHED_LINE, REGULAR_WIDTH);
}
A altura de qualquer figura é a linha reta traçada a partir de seu vértice perpendicular à sua base, ou base produzida.

\vskip\baselineskip

\noindent\ \hfill\drawCurrentPicture\hfill\

\vfill\pagebreak

\starttheorem{Prop. I. Teor.}\label{prop:VI.I}
\defineNewPicture[2/5]{
pair A, B, C, D, H, G, K, L, M, N, O, P, Q;
numeric w, h;
w := 9/2u;
h := 3u;
A := (6/11w, h);
B := (5/11w, 0);
C := (6/11w, 0);
D := (7/11w, 0);
H := (3/11w, 0);
G := (4/11w, 0);
K := (8/11w, 0);
L := (9/11w, 0);
M := (0/11w, 0);
N := (1/11w, 0);
O := (2/11w, 0);
P := (10/11w, 0);
Q := (11/11w, 0);
draw byPolygon(A,M,N)(byblack);
draw byPolygon(A,O,H)(byblack);
draw byPolygon(A,G,B)(byblack);
draw byPolygon(A,B,C)(byred);
draw byPolygon(A,C,D)(byblue);
draw byPolygon(A,K,L)(byyellow);
draw byPolygon(A,P,Q)(byyellow);
draw byLine(M, B)(byblack, SOLID_LINE, REGULAR_WIDTH);
draw byLine(B, C)(byblue, SOLID_LINE, REGULAR_WIDTH);
draw byLine(C, D)(byred, SOLID_LINE, REGULAR_WIDTH);
draw byLine(D, Q)(byblack, SOLID_LINE, REGULAR_WIDTH);
draw byLabelsOnPolygon(A, Q, D, C, B, M)(ALL_LABELS, 0);
}
\drawCurrentPictureInMargin
\problem{T}{riângulos}{e paralelogramos que têm a mesma altura estão entre si como suas bases.}

Seja os triângulos \drawPolygon[bottom]{ABC} e \drawPolygon[bottom]{ACD} terem um vértice comum, e suas bases \drawUnitLine{BC} e \drawUnitLine{CD} na mesma linha reta.

Produza \drawUnitLine{BC,CD} nos dois sentidos, pegue sucessivamente no \drawUnitLine{CD} linhas produzidas iguais a ele; e no \drawUnitLine{BC} produziu linhas sucessivamente iguais a ele; e trace linhas do vértice comum até suas extremidades.

The triangles
\drawFromCurrentPicture[bottom][trianglesAMC]{
draw byNamedPolygon(AMN,AOH,AGB,ABC);
draw byNamedLine(MB,BC);
draw byLabelsOnPolygon(A, C, M)(ALL_LABELS, 0);
}
assim formados são todos iguais entre si, pois suas bases são iguais. \byref{prop:I.XXXVIII}

$\therefore$ \trianglesAMC\ e sua base são respectivamente equimúltiplos de \polygonABC\ e da base \drawUnitLine{BC}.

In like manner
\drawFromCurrentPicture[bottom][trianglesACQ]{
draw byNamedPolygon(ACD,AKL,APQ);
draw byNamedLine(CD,DQ);
draw byLabelsOnPolygon(C, A, Q)(ALL_LABELS, 0);
}
e sua base são respectivamente equimúltiplos de \polygonACD\ e da base \drawUnitLine{CD}.

$\therefore$ se $m$ ou $6$ vezes \polygonABC\ $>, = \mbox{ or } < n$ ou $5$ vezes \polygonACD\ então $m$ ou $6$ vezes \drawUnitLine{BC} $>, = \mbox{ or } < n$ ou $5$ vezes \drawUnitLine{CD}, $m$ e $n$ representam todos os múltiplos tomados como na quinta definição do Quinto Livro. Embora tenhamos mostrado apenas que esta propriedade existe quando $m$ é igual a $6$ e $n$ é igual a $5$, ainda assim é evidente que a propriedade é válida para cada valor múltiplo que pode ser dado a $m$ e a $n$.

\begin{center}
$\therefore \polygonABC\ : \polygonACD\ :: \drawUnitLine{BC} : \drawUnitLine{CD}$ \byref{def:V.V}
\end{center}

Os paralelogramos com a mesma altitude são os duplos dos triângulos, nas suas bases, e são proporcionais a eles (Parte I), e portanto os seus duplos, os paralelogramos, são como suas bases \byref{prop:V.XV}.

\qed

\starttheorem{Prop. II. Teor.}\label{prop:VI.II}
\defineNewPicture[1/2]{
pair A, B, C, D, E, Ab, Bb, Cb, Db, Eb, Ac, Bc, Cc, Dc, Ec, Fc, d[];
d1 := (0, 0);
d2 := (u, -4u);
d3 := (5/2u, -9/2u);
A := (0, 0) shifted d1;
B := (-1/3u, -7/2u) shifted d1;
C := B shifted (5/2u, 0);
D := 1/2[A, B];
E := 1/2[A, C];
Ab := (0, 0) shifted d2;
Bb := (-u, -2u) shifted d2;
Cb := Bb shifted (2u, 0);
Db := 4/3[Ab, Bb];
Eb := 4/3[Ab, Cb];
Ac := (0, 0) shifted d3;
Bc := (-3u, -5u) shifted d3;
Cc := Bc shifted (2u, 0);
Dc := 3/5[Ac, Bc];
Ec := 3/5[Ac, Cc];
Fc = whatever[Bc, Ec] = whatever[Cc, Dc];
draw byLine(D, E, byblack, SOLID_LINE, REGULAR_WIDTH);
draw byLine(B, E, byred, SOLID_LINE, REGULAR_WIDTH);
draw byLine(C, D, byblue, SOLID_LINE, REGULAR_WIDTH);
byLineDefine(A, D, byred, DASHED_LINE, REGULAR_WIDTH);
byLineDefine(D, B, byyellow, SOLID_LINE, REGULAR_WIDTH);
byLineDefine(B, C, byblack, DASHED_LINE, REGULAR_WIDTH);
byLineDefine(A, E, byblue, DASHED_LINE, REGULAR_WIDTH);
byLineDefine(E, C, byyellow, DASHED_LINE, REGULAR_WIDTH);
draw byNamedLineSeq(0)(AD,DB,BC,EC,AE);
draw byLine(Bb, Eb, byred, SOLID_LINE, REGULAR_WIDTH);
draw byLine(Cb, Db, byblue, SOLID_LINE, REGULAR_WIDTH);
draw byLine(Bb, Cb, byblack, DASHED_LINE, REGULAR_WIDTH);
byLineDefine(Ab, Bb, byred, DASHED_LINE, REGULAR_WIDTH);
byLineDefine(Bb, Db, byyellow, SOLID_LINE, REGULAR_WIDTH);
byLineDefine(Db, Eb, byblack, SOLID_LINE, REGULAR_WIDTH);
byLineDefine(Ab, Cb, byblue, DASHED_LINE, REGULAR_WIDTH);
byLineDefine(Cb, Eb, byyellow, DASHED_LINE, REGULAR_WIDTH);
draw byNamedLineSeq(0)(AbBb,BbDb,DbEb,CbEb,AbCb);
draw byLine(Bc, Fc, byyellow, SOLID_LINE, REGULAR_WIDTH);
draw byLine(Fc, Ec, byred, DASHED_LINE, REGULAR_WIDTH);
draw byLine(Cc, Fc, byyellow, DASHED_LINE, REGULAR_WIDTH);
draw byLine(Fc, Dc, byblue, DASHED_LINE, REGULAR_WIDTH);
byLineDefine(Bc, Cc, byblack, DASHED_LINE, REGULAR_WIDTH);
byLineDefine(Bc, Dc, byred, SOLID_LINE, REGULAR_WIDTH);
byLineDefine(Dc, Ec, byblack, SOLID_LINE, REGULAR_WIDTH);
byLineDefine(Cc, Ec, byblue, SOLID_LINE, REGULAR_WIDTH);
draw byNamedLineSeq(0)(BcCc,BcDc,DcEc,CcEc);
byPointLabelDefine(Ab, "A");
byPointLabelDefine(Bb, "B");
byPointLabelDefine(Cb, "C");
byPointLabelDefine(Db, "D");
byPointLabelDefine(Eb, "E");
byPointLabelDefine(Bc, "B");
byPointLabelDefine(Cc, "C");
byPointLabelDefine(Dc, "D");
byPointLabelDefine(Ec, "E");
byPointLabelDefine(Fc, "F");
draw byLabelsOnPolygon(A, E, C, B, D)(ALL_LABELS, 0);
draw byLabelsOnPolygon(Ab, Cb, Eb, Db, Bb)(ALL_LABELS, 0);
draw byLabelsOnPolygon(Dc, Ec, Cc, Bc)(ALL_LABELS, 0);
draw byLabelsOnPolygon(Bc, Fc, Dc)(OMIT_FIRST_LABEL+OMIT_LAST_LABEL, 0);
}
\drawCurrentPictureInMargin
\problem{S}{e}{uma linha reta \drawUnitLine{DE} for traçada paralelamente a qualquer lado \drawUnitLine{BC}
E se qualquer linha reta \drawUnitLine{DE} dividir os lados de um triângulo, ou os lados produzidos, em segmentos proporcionais, ela será paralela ao lado restante \drawUnitLine{BC}.}

\startsubproposition{Parte I.}
\begin{center}
Let $\drawUnitLine{DE} \parallel \drawUnitLine{BC}$, then shall\\
$\drawUnitLine{DB} : \drawUnitLine{AD} :: \drawUnitLine{EC} : \drawUnitLine{AE}$.

Trace \drawUnitLine{BE} e \drawUnitLine{CD},\\
and
$\drawLine[middle][triangleBDE]{DE,BE,DB}
=
\drawLine[middle][triangleCDE]{EC,CD,DE}$ \byref{prop:I.XXXVII};

$\therefore \triangleBDE\ :
\drawLine[middle][triangleADE]{AD,AE,DE} :: \triangleCDE\ : \triangleADE$ \byref{prop:V.VII};\\
but $\triangleBDE\ : \triangleADE\ :: \drawUnitLine{DB} : \drawUnitLine{AD}$ \byref{prop:VI.I},\\
$\therefore \drawUnitLine{DB} : \drawUnitLine{AD} :: \drawUnitLine{EC} : \drawUnitLine{AE}$ \byref{prop:V.XI}.
\end{center}

\vfill\pagebreak

\startsubproposition{Parte II.}
\begin{center}
Let $\drawUnitLine{DB} : \drawUnitLine{AD} :: \drawUnitLine{EC} : \drawUnitLine{AE}$.\\
then $\drawUnitLine{DE} \parallel \drawUnitLine{BC}$.

Que permaneça a mesma construção,\\
$\left.
\begin{aligned}
\because \drawUnitLine{DB} : \drawUnitLine{AD} &:: \triangleBDE\ : \triangleADE \\ %\mbox{ because }
\mbox{ e } \drawUnitLine{EC} : \drawUnitLine{AE} &:: \triangleCDE\ : \triangleADE \\
\end{aligned}
\right\}\mbox{ \byref{prop:VI.I} }$;\\
but $\drawUnitLine{DB} : \drawUnitLine{AD} :: \drawUnitLine{EC} : \drawUnitLine{AE}$ \byref{\hypref}\\
$\therefore \triangleBDE\ : \triangleADE :: \triangleCDE\ : \triangleADE$ \byref{prop:V.XI}\\
$\therefore \triangleBDE\ = \triangleCDE$ \byref{prop:V.IX};

mas eles estão na mesma base \drawUnitLine{BC}, e no mesmo lado dela, e\\
$\therefore \drawUnitLine{DE} \parallel \drawUnitLine{BC}$ \byref{prop:I.XXXIX}.
\end{center}

\qed

\starttheorem{Prop. III. Teor.}\label{prop:VI.III}
\defineNewPicture{
pair A, B, C, D, E;
numeric a;
a := 80;
A := (0, 0);
B := A shifted (dir(220)*3u);
C = whatever[A, A + dir(220+a)] = whatever[B, B shifted dir(0)];
D = whatever[A, A + dir(220+1/2a)] = whatever[B, C];
E = whatever[A, B] = whatever[C, C shifted (A-D)];
byAngleDefine(B, A, D, byyellow, SOLID_SECTOR);
byAngleDefine(D, A, C, byblack, SOLID_SECTOR);
byAngleDefine(C, A, E, byblack, ARC_SECTOR);
byAngleDefine(A, E, C, byblue, SOLID_SECTOR);
byAngleDefine(E, C, A, byred, SOLID_SECTOR);
draw byNamedAngleResized();
draw byLine(C, A, byyellow, SOLID_LINE, REGULAR_WIDTH);
draw byLine(A, D, byblue, SOLID_LINE, REGULAR_WIDTH);
byLineDefine(A, E, byred, DASHED_LINE, REGULAR_WIDTH);
byLineDefine(C, E, byblue, DASHED_LINE, REGULAR_WIDTH);
byLineDefine(A, B, byred, SOLID_LINE, REGULAR_WIDTH);
byLineDefine(B, D, byblack, SOLID_LINE, REGULAR_WIDTH);
byLineDefine(D, C, byblack, DASHED_LINE, REGULAR_WIDTH);
draw byNamedLineSeq(0)(AE,CE,DC,BD,AB);
draw byLabelsOnPolygon(B, A, E, C, D)(ALL_LABELS, 0);
}
\drawCurrentPictureInMargin
\problem{U}{ma}{linha reta (\drawUnitLine{AD}) que divide o ângulo de um triângulo ao meio, divide o lado oposto em segmentos (\drawUnitLine{BD}, \drawUnitLine{DC}) proporcionais aos lados contíguos (\drawUnitLine{AB}, \drawUnitLine{CA}
E se uma linha reta (\drawUnitLine{AD}) traçada a partir de qualquer ângulo de um triângulo divide o lado oposto (\drawUnitLine{BD,DC}) em segmentos (\drawUnitLine{BD}, \drawUnitLine{DC}) proporcionais aos lados contíguos (\drawUnitLine{AB}, \drawUnitLine{CA}), ela divide o ângulo ao meio.}

\startsubproposition{Parte I.}
\begin{center}
Trace $\drawUnitLine{CE} \parallel \drawUnitLine{AD}$, para atender \drawUnitLine{AE};\\
then, $\drawAngle{BAD} = \drawAngle{E}$ \byref{prop:I.XXIX}.

$\therefore \drawAngle{DAC} = \drawAngle{E}$; but $\drawAngle{DAC} = \drawAngle{C}$, $\therefore \drawAngle{C} = \drawAngle{E}$,\\
$\therefore \drawUnitLine{AE} = \drawUnitLine{CA}$ \byref{prop:I.VI};\\
and $\because \drawUnitLine{AD} \parallel \drawUnitLine{CE}$,\\ %because
$\drawUnitLine{AE} : \drawUnitLine{AB} :: \drawUnitLine{DC} : \drawUnitLine{BD}$ \byref{prop:VI.II};

but $\drawUnitLine{AE} = \drawUnitLine{CA}$;\\
$\therefore \drawUnitLine{CA} : \drawUnitLine{AB} :: \drawUnitLine{DC} : \drawUnitLine{BD}$ \byref{prop:V.VII}.
\end{center}

\vfill\pagebreak

\startsubproposition{Parte II.}
\begin{center}
Que permaneça a mesma construção,\\
and $\drawUnitLine{AB} : \drawUnitLine{AE} :: \drawUnitLine{BD} : \drawUnitLine{DC}$ \byref{prop:VI.II};

but $\drawUnitLine{BD} : \drawUnitLine{DC} :: \drawUnitLine{AB} : \drawUnitLine{CA}$ \byref{\hypref}\\
$\therefore \drawUnitLine{AB} : \drawUnitLine{AE} :: \drawUnitLine{AB} : \drawUnitLine{CA}$ \byref{prop:V.XI}.

and $\therefore \drawUnitLine{AE} = \drawUnitLine{CA}$ \byref{prop:V.IX},\\
and $\therefore \drawAngle{E} = \drawAngle{C}$ \byref{prop:I.V};\\
but since $\drawUnitLine{AD} \parallel \drawUnitLine{CE}$; $\drawAngle{DAC} = \drawAngle{C}$,\\
and $\drawAngle{BAD} = \drawAngle{E}$ \byref{prop:I.XXIX};

$\therefore \drawAngle{C} = \drawAngle{E}$, and $\drawAngle{BAD} = \drawAngle{DAC}$,\\
and $\therefore$ \drawUnitLine{AD} bisects \drawAngle{BAD,DAC}.
\end{center}

\qed

\starttheorem{Prop. IV. Teor.}\label{prop:VI.IV}
\defineNewPicture{
pair A, B, C, D, E, F;
B := (0, 0);
A := (3/2u, u);
C := (2u, 0);
D := (A scaled 4/3) shifted (C-B);
E := (C scaled 4/3) shifted (C-B);
F = whatever[A, B] = whatever[D, E];
byAngleDefine(A, B, C, byyellow, SOLID_SECTOR);
byAngleDefine(B, C, A, byblue, SOLID_SECTOR);
byAngleDefine(C, A, B, byblack, ARC_SECTOR);
byAngleDefine(D, C, E, byred, SOLID_SECTOR);
byAngleDefine(C, E, D, byblack, SOLID_SECTOR);
byAngleDefine(E, D, C, byred, ARC_SECTOR);
draw byNamedAngleResized();
draw byLine(C, A, byred, SOLID_LINE, REGULAR_WIDTH);
draw byLine(C, D, byblue, DASHED_LINE, REGULAR_WIDTH);
byLineDefine(A, F, byyellow, DASHED_LINE, REGULAR_WIDTH);
byLineDefine(D, F, byyellow, SOLID_LINE, REGULAR_WIDTH);
byLineDefine(A, B, byblue, SOLID_LINE, REGULAR_WIDTH);
byLineDefine(D, E, byred, DASHED_LINE, REGULAR_WIDTH);
byLineDefine(E, C, byblack, DASHED_LINE, REGULAR_WIDTH);
byLineDefine(B, C, byblack, SOLID_LINE, REGULAR_WIDTH);
draw byNamedLineSeq(0)(AF,DF,DE,EC,BC,AB);
draw byLabelsOnPolygon(B, A, F, D, E, C)(ALL_LABELS, 0);
}
\drawCurrentPictureInMargin
\problem{I}{n}{equiangular triangles
(\drawLine[bottom][triangleABC]{CA,BC,AB}
and
\drawLine[bottom][triangleCDE]{CD,DE,EC})
os lados sobre os ângulos iguais são proporcionais e os lados opostos aos ângulos iguais são homólogos.
}

Sejam os triângulos equiangulares colocados de forma que dois lados \drawUnitLine[0.5cm]{BC}, \drawUnitLine[0.5cm]{EC} opostos aos ângulos iguais \drawAngle{D} e \drawAngle{A} possam ser contíguos e na mesma linha reta; e que os triângulos ficam do mesmo lado dessa linha reta; e que os triângulos situados no mesmo lado dessa linha reta podem ter ângulos iguais não contíguos, ou seja, \drawAngle{DCE} oposto a \drawAngle{B}, e \drawAngle{BCA} a \drawAngle{E}.

\begin{center}
Trace \drawUnitLine{AF} e \drawUnitLine{DF}.

Então, $\because \drawAngle{BCA} = \drawAngle{E}$, $\drawUnitLine{CA} \parallel \drawUnitLine{DF,DE}$ \byref{prop:I.XXVIII}; %because

e, por razão semelhante, $\drawUnitLine{CD} \parallel \drawUnitLine{AB,AF}$,

$\therefore$
\drawFromCurrentPicture[bottom][polygonACDF]{
startGlobalRotation(-lineAngle.CD);
startAutoLabeling;
draw byNamedLineSeq(0)(CA,AF,DF,CD);
stopAutoLabeling;
stopGlobalRotation;
}
é um paralelogramo.

Mas $\drawUnitLine{BC} : \drawUnitLine{EC} :: \drawUnitLine{DF} : \drawUnitLine{DE}$ \byref{prop:VI.II};

e, como $\drawUnitLine{DF} = \drawUnitLine{CA}$ \byref{prop:I.XXXIV},\\
$\drawUnitLine{BC} : \drawUnitLine{EC} :: \drawUnitLine{CA} : \drawUnitLine{DE}$;

e, por alternação, $\drawUnitLine{BC} : \drawUnitLine{CA} :: \drawUnitLine{EC} : \drawUnitLine{DE}$ \byref{prop:V.XVI}.

De modo semelhante pode-se mostrar que\\
$\drawUnitLine{AB} : \drawUnitLine{CD} :: \drawUnitLine{BC} : \drawUnitLine{EC}$;\\
e, por alternação, que\\
$\drawUnitLine{AB} : \drawUnitLine{BC} :: \drawUnitLine{CD} : \drawUnitLine{EC}$;\\
mas já foi provado que\\
$\drawUnitLine{BC} : \drawUnitLine{CA} :: \drawUnitLine{EC} : \drawUnitLine{DE}$\\
e, portanto, ex \ae\ quali,\\
$\drawUnitLine{AB} : \drawUnitLine{CA} :: \drawUnitLine{CD} : \drawUnitLine{DE}$ \byref{prop:V.XXII},\\
portanto, os lados em torno dos ângulos iguais são proporcionais e aqueles que são opostos aos ângulos iguais são homólogos.
\end{center}

\qed

\starttheorem{Prop. V. Teor.}\label{prop:VI.V}
\defineNewPicture{
pair A, B, C, D, E, F, G, d;
B := (0, 0);
A := (2u, 5/2u);
C := (3u, 0);
byAngleDefine(B, A, C, byyellow, SOLID_SECTOR);
byAngleDefine(A, B, C, byblue, SOLID_SECTOR);
byAngleDefine(B, C, A, byred, SOLID_SECTOR);
draw byNamedAngleResized(BAC, ABC, BCA);
byLineDefine(A, B, byred, DASHED_LINE, REGULAR_WIDTH);
byLineDefine(B, C, byblack, DASHED_LINE, REGULAR_WIDTH);
byLineDefine(C, A, byblue, DASHED_LINE, REGULAR_WIDTH);
draw byNamedLineSeq(0)(AB,BC,CA);
d := (0, -3u);
D := A shifted d;
E := B shifted d;
F := C shifted d;
G := (A yscaled -1) shifted d;
byAngleDefine(F, D, E, byblack, SOLID_SECTOR);
byAngleDefine(D, E, F, byyellow, SOLID_SECTOR);
byAngleDefine(E, F, D, byred, ARC_SECTOR);
byAngleDefine(E, G, F, byblack, SOLID_SECTOR);
byAngleDefine(G, E, F, byblue, SOLID_SECTOR);
byAngleDefine(E, F, G, byred, SOLID_SECTOR);
draw byNamedAngleResized(FDE, DEF, EFD, EGF, GEF, EFG);
draw byLine(E, F, byblack, SOLID_LINE, REGULAR_WIDTH);
byLineDefine(F, G, byyellow, DASHED_LINE, REGULAR_WIDTH);
byLineDefine(G, E, byyellow, SOLID_LINE, REGULAR_WIDTH);
byLineDefine(D, E, byred, SOLID_LINE, REGULAR_WIDTH);
byLineDefine(F, D, byblue, SOLID_LINE, REGULAR_WIDTH);
draw byNamedLineSeq(0)(FG,GE,DE,FD);
draw byLabelsOnPolygon(B, A, C)(ALL_LABELS, 0);
draw byLabelsOnPolygon(E, D, F, G)(ALL_LABELS, 0);
}
\drawCurrentPictureInMargin
\problem{S}{e}{dois triângulos têm os lados proporcionais
($\drawUnitLine{CA} : \drawUnitLine{BC} :: \drawUnitLine{FD} : \drawUnitLine{EF}$)
e
($\drawUnitLine{BC} : \drawUnitLine{AB} :: \drawUnitLine{EF} : \drawUnitLine{DE}$),
então são equiangulares, e os ângulos iguais são subtendidos pelos lados homólogos.}
\begin{center}
Das extremidades de \drawUnitLine{EF}, trace \drawUnitLine{FG} e \drawUnitLine{GE}, fazendo\\
$\drawAngle{GEF} = \drawAngle{B}$, $\drawAngle{EFG} = \drawAngle{C}$ \byref{prop:I.XXIII};\\
$\therefore \drawAngle{G} = \drawAngle{A}$ \byref{prop:I.XXXII},\\
e, como os triângulos são equiângulos,\\
$\drawUnitLine{AB} : \drawUnitLine{BC} :: \drawUnitLine{GE} : \drawUnitLine{EF}$ \byref{prop:VI.IV};\\
mas $\drawUnitLine{AB} : \drawUnitLine{BC} :: \drawUnitLine{DE} : \drawUnitLine{EF}$ \byref{\hypref};\\
$\therefore \drawUnitLine{DE} : \drawUnitLine{EF} :: \drawUnitLine{GE} : \drawUnitLine{EF}$,\\
e, portanto, $\drawUnitLine{DE} = \drawUnitLine{GE}$ \byref{prop:V.IX}.\\
De modo semelhante, pode-se mostrar que\\
$\drawUnitLine{FD} = \drawUnitLine{FG}$.

Portanto, os dois triângulos que têm base comum \drawUnitLine{EF}, e seus lados iguais, também têm ângulos iguais opostos a lados iguais, ou seja.\ e.\\
$\drawAngle{DEF} = \drawAngle{GEF}$ e $\drawAngle{EFD} = \drawAngle{EFG}$ \byref{prop:I.VIII}.\\
Mas $\drawAngle{GEF} = \drawAngle{B}$ \byref{\constref} e, portanto, $\drawAngle{DEF} = \drawAngle{B}$;\\
pela mesma razão, $\drawAngle{EFD} = \drawAngle{C}$,\\
e, consequentemente, $\drawAngle{D} = \drawAngle{A}$ \byref{prop:I.XXXII};\\
e portanto os triângulos são equiângulos, e é evidente que os lados homólogos são subtendidos pelos ângulos iguais.
\end{center}

\qed

\starttheorem{Prop. VI. Teor.}\label{prop:VI.VI}
\defineNewPicture[1/4]{
pair A, B, C, D, E, F, G, d;
A := (0, 0);
B := (2u, 5/2u);
C := (3u, 0);
byAngleDefine(A, B, C, byyellow, SOLID_SECTOR);
byAngleDefine(B, A, C, byblue, SOLID_SECTOR);
byAngleDefine(A, C, B, byred, SOLID_SECTOR);
draw byNamedAngleResized(BAC, ABC, ACB);
byLineDefine(A, B, byred, DASHED_LINE, REGULAR_WIDTH);
byLineDefine(C, A, byblack, DASHED_LINE, REGULAR_WIDTH);
byLineDefine(B, C, byblue, DASHED_LINE, REGULAR_WIDTH);
draw byNamedLineSeq(0)(AB,BC,CA);
d := (0, -3u);
D := A shifted d;
E := B shifted d;
F := C shifted d;
G := (B yscaled -1) shifted d;
byAngleDefine(F, D, E, byblue, ARC_SECTOR);
byAngleDefine(D, E, F, byyellow, ARC_SECTOR);
byAngleDefine(E, F, D, byred, ARC_SECTOR);
byAngleDefine(F, D, G, byblue, SOLID_SECTOR);
byAngleDefine(G, F, D, byred, SOLID_SECTOR);
byAngleDefine(D, G, F, byblack, SOLID_SECTOR);
draw byNamedAngleResized(FDE, DEF, EFD, FDG, GFD, DGF);
draw byLine(F, D, byblue, SOLID_LINE, REGULAR_WIDTH);
byLineDefine(F, G, byyellow, DASHED_LINE, REGULAR_WIDTH);
byLineDefine(G, D, byyellow, SOLID_LINE, REGULAR_WIDTH);
byLineDefine(D, E, byred, SOLID_LINE, REGULAR_WIDTH);
byLineDefine(E, F, byblack, SOLID_LINE, REGULAR_WIDTH);
draw byNamedLineSeq(0)(FG,GD,DE,EF);
draw byLabelsOnPolygon(A, B, C)(ALL_LABELS, 0);
draw byLabelsOnPolygon(D, E, F, G)(ALL_LABELS, 0);
}
\drawCurrentPictureInMargin
\problem{S}{e}{dois triângulos (\drawLine[bottom][triangleABC]{AB,BC,CA} e \drawLine[bottom][triangleDEF]{DE,EF,FD}) têm um ângulo (\drawAngle{C}) de um igual a um ângulo (\drawAngle{EFD}) do outro, e os lados em torno dos ângulos iguais proporcionais, então os triângulos são equiangulares e têm iguais os ângulos opostos aos lados homólogos.}

\begin{center}
Das extremidades de \drawUnitLine{FD}, um dos lados do triângulo DEF, em torno de \drawAngle{EFD},\\
trace \drawUnitLine{GD} e \drawUnitLine{FG},\\
fazendo $\drawAngle{GFD} = \drawAngle{C}$,\\
e $\drawAngle{FDG} = \drawAngle{A}$;\\
então $\drawAngle{G} = \drawAngle{B}$ \byref{prop:I.XXXII},\\
e, sendo os dois triângulos equiângulos,\\
$\drawUnitLine{BC} : \drawUnitLine{CA} :: \drawUnitLine{FG} : \drawUnitLine{FD}$ \byref{prop:VI.IV};\\
mas $\drawUnitLine{BC} : \drawUnitLine{CA} :: \drawUnitLine{EF} : \drawUnitLine{FD}$ \byref{\hypref};\\
$\therefore \drawUnitLine{FG} : \drawUnitLine{FD} :: \drawUnitLine{EF} : \drawUnitLine{FD}$ \byref{prop:V.XI},\\
e, consequentemente, $\drawUnitLine{FG} = \drawUnitLine{EF}$ \byref{prop:V.IX};\\
$\therefore \drawLine[bottom][triangleDEF]{DE,EF,FD} = \drawLine[bottom][triangleDGF]{FD,FG,GD}$ em todos os aspectos \byref{prop:I.IV}.\\
Mas $\drawAngle{GFD} = \drawAngle{A}$ \byref{\constref},\\
e $\therefore \drawAngle{EFD} = \drawAngle{A}$;\\
e, como também $\drawAngle{EFD} = \drawAngle{C}$,\\
$\drawAngle{E} = \drawAngle{B}$ \byref{prop:I.XXXII};\\
e, portanto, os triângulos ABC e DEF são equiangulares, com seus ângulos iguais opostos aos lados homólogos.
\end{center}

\qed

\starttheorem{Prop. VII. Teor.}\label{prop:VI.VII}
\defineNewPicture[1/4]{
pair A, B, C, D, E, F, G, d;
A := (0, 0);
B := (3u, -1/2u);
C := (5/2u, 3u);
G := 1/3[A, C];
d := (0, -4u);
D := (A scaled 4/5) shifted d;
E := (B scaled 4/5) shifted d;
F := (C scaled 4/5) shifted d;
byAngleDefine(B, C, A, byblue, SOLID_SECTOR);
byAngleDefine(C, A, B, byred, SOLID_SECTOR);
byAngleDefine(A, B, G, byblack, SOLID_SECTOR);
byAngleDefine(G, B, C, byblack, ARC_SECTOR);
byAngleDefine(C, G, B, byyellow, SOLID_SECTOR);
byAngleDefine(B, G, A, byred, SOLID_SECTOR);
draw byNamedAngleResized(BCA, CAB, ABG, GBC, CGB, BGA);
draw byLine(B, G, byblue, SOLID_LINE, REGULAR_WIDTH);
byLineDefine(A, B, byyellow, SOLID_LINE, REGULAR_WIDTH);
byLineDefine(B, C, byred, SOLID_LINE, REGULAR_WIDTH);
byLineDefine(C, A, byblack, SOLID_LINE, REGULAR_WIDTH);
draw byNamedLineSeq(0)(CA,BC,AB);
byAngleDefine(D, E, F, byblue, ARC_SECTOR);
byAngleDefine(E, F, D, byyellow, ARC_SECTOR);
byAngleDefine(F, D, E, byred, ARC_SECTOR);
draw byNamedAngleResized(DEF, EFD, FDE);
byLineDefine(D, E, byyellow, DASHED_LINE, REGULAR_WIDTH);
byLineDefine(E, F, byred, DASHED_LINE, REGULAR_WIDTH);
byLineDefine(F, D, byblue, DASHED_LINE, REGULAR_WIDTH);
draw byNamedLineSeq(0)(FD,EF,DE);
draw byLabelsOnPolygon(A, G, C, B)(ALL_LABELS, 0);
draw byLabelsOnPolygon(D, F, E)(ALL_LABELS, 0);
}
\drawCurrentPictureInMargin
\problem{S}{e}{dois triângulos (\drawLine[bottom][triangleABC]{CA,BC,AB} e \drawLine[bottom][triangleDEF]{FD,EF,DE}) têm um ângulo em cada igual (\drawAngle{F} igual a \drawAngle{C}), os lados em torno de dois outros ângulos proporcionais ($\drawUnitLine{BC} : \drawUnitLine{AB} :: \drawUnitLine{EF} : \drawUnitLine{DE}$), e cada um dos ângulos restantes (\drawAngle{A} e \drawAngle{D}) menor, ou não menor, que um ângulo reto, então os triângulos são equiângulos, e os ângulos entre os lados proporcionais são iguais.}

\begin{center}
Primeiro, suponhamos que os ângulos \drawAngle{A} e \drawAngle{D} sejam menores que um ângulo reto. Então, se \drawAngle{ABG,GBC} e \drawAngle{E}, compreendidos pelos lados proporcionais, não forem iguais, seja \drawAngle{ABG,GBC} o maior, e faça-se $\drawAngle{GBC} = \drawAngle{E}$.

$\because \drawAngle{C} = \drawAngle{F}$ \byref{\hypref} e $\drawAngle{GBC} = \drawAngle{E}$ \byref{\constref}\\ %Because
$\therefore \drawAngle{CGB} = \drawAngle{D}$ \byref{prop:I.XXXII};

$\therefore \drawUnitLine{BC} : \drawUnitLine{BG} :: \drawUnitLine{EF} : \drawUnitLine{DE}$ \byref{prop:VI.IV},\\
mas $\drawUnitLine{BC} : \drawUnitLine{AB} :: \drawUnitLine{EF} : \drawUnitLine{DE}$ \byref{\hypref};

$\therefore \drawUnitLine{BC} : \drawUnitLine{BG} :: \drawUnitLine{BC} : \drawUnitLine{AB}$;\\
$\therefore \drawUnitLine{BG} = \drawUnitLine{AB}$ \byref{prop:V.IX},\\
e $\therefore \drawAngle{A} = \drawAngle{BGA}$ \byref{prop:I.V}.

Mas \drawAngle{A} é menor que um ângulo reto \byref{\hypref}\\
$\therefore$ \drawAngle{BGA} é menor que um ângulo reto;\\
e $\therefore$ \drawAngle{CGB} deve ser maior que um ângulo reto \byref{prop:I.XIII}, mas foi comprovado $= \drawAngle{D}$ e portanto menor que um ângulo reto, o que é um absurdo. $\therefore$ \drawAngle{ABG,GBC} e \drawAngle{E} não são desiguais;

$\therefore$ eles são iguais; e, como $\drawAngle{C} = \drawAngle{F}$ \byref{\hypref},\\
$\therefore \drawAngle{A} = \drawAngle{D}$ \byref{prop:I.XXXII} e, portanto, os triângulos são equiângulos.
\end{center}

Mas, se \drawAngle{A} e \drawAngle{D} forem assumidos como sendo, cada um deles, não menores que um ângulo reto, pode-se provar, como antes, que os triângulos são equiangulares e têm os lados em torno dos ângulos iguais proporcionais \byref{prop:VI.IV}.

\qed

\starttheorem{Prop. VIII. Teor.}\label{prop:VI.VIII}
\defineNewPicture{
pair A, B, C, D;
A := (0, 0);
B := A shifted (dir(-145)*7/2u);
C = whatever[B, B shifted (1, 0)] = whatever[A, A shifted dir(-145 - 90)];
D := (xpart(A), ypart(B));
draw byPolygon(A,B,D)(byyellow);
draw byPolygon(A,D,C)(byred);
byAngleDefine(A, B, C, byblack, SOLID_SECTOR);
byAngleDefine(B, C, A, byblue, ARC_SECTOR);
byAngleDefine(B, D, A, byblue, SOLID_SECTOR);
byAngleDefine(D, A, B, byred, SOLID_SECTOR);
byAngleDefine(C, A, D, byyellow, SOLID_SECTOR);
draw byNamedAngleResized();
draw byNamedAngleDummySides(BCA);
draw byLine(A, D, byblack, SOLID_LINE, REGULAR_WIDTH);
draw byLabelsOnPolygon(B, A, C, D)(ALL_LABELS, 0);
}
\drawCurrentPictureInMargin
\problem{E}{m}{um triângulo retângulo (\drawPolygon[bottom][triangleABC]{ABD,ADC}), se uma perpendicular (\drawUnitLine{AD}) for traçada do ângulo reto ao lado oposto, os triângulos (\drawPolygon[bottom][triangleABD]{ABD} e \drawPolygon[bottom][triangleADC]{ADC}) de cada lado dela são semelhantes ao triângulo total e entre si.}

\begin{center}
$\because \drawAngle{DAB,CAD} = \drawAngle{D}$ \byref{ax:I.XI},\\
e \drawAngle{B} é comum aos triângulos ABC e ABD;\\ % Porque
$\drawAngle{C} = \drawAngle{DAB}$ \byref{prop:I.XXXII};

$\therefore$ os triângulos ABC e ABD são equiângulos; e, consequentemente, têm os lados em torno dos ângulos iguais proporcionais \byref{prop:VI.IV}, sendo portanto semelhantes \byref{def:VI.I}.

De modo semelhante, pode-se provar que o triângulo ADC é semelhante ao triângulo ABC; mas o triângulo ABD já foi mostrado semelhante ao triângulo ABC;\\
$\therefore$ os triângulos ABD e ADC são semelhantes ao todo e entre si.
\end{center}

\qed

\startproblem{Prop. IX. prob.}\label{prop:VI.IX}
\defineNewPicture{
pair A, B, C, D, E, F;
numeric n;
n := 5;
A := (0, 0);
B := (3u, 5u);
D := A shifted (0, u);
E := ((D-A) scaled n) shifted A;
F = whatever[A, B] = whatever[D, D shifted (B-E)];
C := 1/3[D, E];
draw byLine(D, F, byred, DASHED_LINE, REGULAR_WIDTH);
byLineDefine(B, E, byred, SOLID_LINE, REGULAR_WIDTH);
byLineDefine(A, F, byyellow, SOLID_LINE, REGULAR_WIDTH);
byLineDefine(A, D, byblue, SOLID_LINE, REGULAR_WIDTH);
byLineDefine(D, C, byblue, DASHED_LINE, REGULAR_WIDTH);
byLineDefine(C, E, byblack, DASHED_LINE, REGULAR_WIDTH);
byLineDefine(F, B, byyellow, DASHED_LINE, REGULAR_WIDTH);
byLineDefine(A, E, byblack, SOLID_LINE, REGULAR_WIDTH);
draw byNamedLineSeq(0)(BE, FB, AF, AD, DC, CE);
for i := 2 step 1 until n - 1:
draw byMarkLine(i/n, byblack)(AE);
endfor;
draw byLabelsOnPolygon(B, F, A, D, C, E)(ALL_LABELS, 0); % improvement: drawing in the original seemed to be wrong, fixed here
}
\drawCurrentPictureInMargin
\problem{A}{}{partir de uma determinada linha reta (\drawSizedLine{AF,FB}) para cortar qualquer peça necessária.}

\begin{center}
De qualquer extremidade da linha dada, trace \drawSizedLine{AD,DC} fazendo qualquer ângulo com \drawSizedLine{AF,FB};\\
e produza \drawSizedLine{AD,DC} até que toda a linha produzida \drawSizedLine{AD,DC,CE} contenha \drawSizedLine{AD} com a mesma freqüência que \drawSizedLine{AF,FB} contém a peça necessária.

Draw \drawSizedLine{BE},\\
e trace $\drawSizedLine{DF} \parallel \drawSizedLine{BE}$.\\

\drawSizedLine{AF} é a parte necessária do \drawSizedLine{AF,FB}.

For since $\drawSizedLine{DF} \parallel \drawSizedLine{BE}$\\
$\drawSizedLine{AF} : \drawSizedLine{FB} :: \drawSizedLine{AD} : \drawSizedLine{DC,CE}$ \byref{prop:VI.II},\\
and by composition \byref{prop:V.XVIII};\\
$\drawSizedLine{AF,FB} : \drawSizedLine{AF} :: \drawSizedLine{AD,DC,CE} : \drawSizedLine{AD}$;

but \drawSizedLine{AD,DC,CE} contains \drawSizedLine{AD} as often as \drawSizedLine{AF,FB} contains the required part \byref{\constref};

$\therefore$ \drawSizedLine{AF} é a peça necessária.
\end{center}

\qed

\startproblem{Prop. X. prob.}\label{prop:VI.X}
\defineNewPicture{
pair A', D', E', C', A, B, C, D, E, F, G, L, d;
numeric a[];
A' := (0, 0);
D' := (3/2u, 0);
E' := (5/2u, 0);
C' := (7/2u, 0);
draw byLine(A', D', byblue, SOLID_LINE, THIN_WIDTH);
draw byLine(D', E', byred, SOLID_LINE, THIN_WIDTH);
draw byLine(E', C', byyellow, SOLID_LINE, THIN_WIDTH);
d := (0, -7/2u);
a1 := 40;
a2 := -70;
A := (A' rotated a1) shifted d;
D := (D' rotated a1) shifted d;
E := (E' rotated a1) shifted d;
C := (C' rotated a1) shifted d;
L := 6/5[A, C];
F = whatever[A, A shifted dir(0)] = whatever[D, D shifted dir(a2)];
G = whatever[A, A shifted dir(0)] = whatever[E, E shifted dir(a2)];
B = whatever[A, A shifted dir(0)] = whatever[C, C shifted dir(a2)];
draw byLine(D, F, byblack, SOLID_LINE, REGULAR_WIDTH);
draw byLine(E, G, byblack, DASHED_LINE, REGULAR_WIDTH);
byLineDefine(C, B, byblack, SOLID_LINE, THIN_WIDTH);
byLineDefine(A, D, byblue, DASHED_LINE, REGULAR_WIDTH);
byLineDefine(D, E, byred, DASHED_LINE, REGULAR_WIDTH);
byLineDefine(E, C, byyellow, DASHED_LINE, REGULAR_WIDTH);
byLineDefine(C, L, byblack, SOLID_LINE, REGULAR_WIDTH);
byLineDefine(A, F, byblue, SOLID_LINE, REGULAR_WIDTH);
byLineDefine(F, G, byred, SOLID_LINE, REGULAR_WIDTH);
byLineDefine(G, B, byyellow, SOLID_LINE, REGULAR_WIDTH);
draw byNamedLineSeq(0)(noLine,CL,EC,DE,AD,AF,FG,GB,CB);
draw byLabelsOnPolygon(C, B, G, F, A, D, E, C, noPoint)(OMIT_FIRST_LABEL+OMIT_LAST_LABEL, 0);
draw byLabelsOnPolygon(E, C, noPoint)(OMIT_FIRST_LABEL+OMIT_LAST_LABEL, 0);
draw byLabelsOnPolygon(A', D', E', C', noPoint)(ALL_LABELS, 0);
}
\drawCurrentPictureInMargin
\problem{P}{ara}{dividir uma determinada linha reta (\drawSizedLine{AF,FG,GB}) de forma semelhante a uma determinada linha dividida (\drawSizedLine{A'D',D'E',E'C'}).}

\begin{center}
De qualquer extremidade da linha dada \drawSizedLine{AF,FG,GB} trace \drawSizedLine{AD,DE,EC} fazendo qualquer ângulo;

considere \drawSizedLine{AD}, \drawSizedLine{DE} e \drawSizedLine{EC} iguais a \drawSizedLine{A'D'}, \drawSizedLine{D'E'} e \drawSizedLine{E'C'} respectivamente \byref{prop:I.II};

trace \drawSizedLine{CB} e trace \drawSizedLine{EG} e \drawSizedLine{DF} $\parallel$ nele.

Como
$\left\{\vcenter{
\nointerlineskip\hbox{\drawSizedLine{CB}}
\nointerlineskip\hbox{\drawSizedLine{EG}}
\nointerlineskip\hbox{\drawSizedLine{DF}}}\right\}$
são $\parallel$,\\
$\drawSizedLine{GB} : \drawSizedLine{FG} :: \drawSizedLine{EC} : \drawSizedLine{DE}$ \byref{prop:VI.II},\\
or $\drawSizedLine{GB} : \drawSizedLine{FG} :: \drawSizedLine{E'C'} : \drawSizedLine{D'E'}$ \byref{\constref},\\
and $\drawSizedLine{FG} : \drawSizedLine{AF} :: \drawSizedLine{DE} : \drawSizedLine{AD}$ \byref{prop:VI.II},\\
$\drawSizedLine{FG} : \drawSizedLine{AF} :: \drawSizedLine{D'E'} : \drawSizedLine{A'D'}$ \byref{\constref},

e $\therefore$ a linha fornecida \drawSizedLine{AF,FG,GB} é dividida de forma semelhante a \drawSizedLine{A'D',D'E',E'C'}.
\end{center}

\qed

\startproblem{Prop. XI. prob.}\label{prop:VI.XI}
\defineNewPicture{
pair A', C', A, B, C, D, E, d;
numeric a[], l[];
a1 := 85;
a2 := 60;
l1 := 2u;
l2 := 5/2u;
A := (0, 0);
B := A shifted (dir(a1)*l1);
C := A shifted (dir(a2)*l2);
D := B shifted (dir(a1)*l2);
E = whatever[A, C] = whatever[D, D shifted (B-C)];
d := (2u, 0);
A' := A shifted d;
C' := A' shifted (dir(a1)*l2);
draw byLine(B, C, byyellow, SOLID_LINE, REGULAR_WIDTH);
byLineDefine(D, E, byyellow, DASHED_LINE, REGULAR_WIDTH);
byLineDefine(A, B, byblack, SOLID_LINE, REGULAR_WIDTH);
byLineDefine(B, D, byblue, DASHED_LINE, REGULAR_WIDTH);
byLineDefine(A, C, byred, DASHED_LINE, REGULAR_WIDTH);
byLineDefine(C, E, byred, SOLID_LINE, REGULAR_WIDTH);
draw byNamedLineSeq(0)(DE,CE,AC,AB,BD);
draw byLine(A', C', byblue, SOLID_LINE, REGULAR_WIDTH);
draw byLabelsOnPolygon(E, C, A, B, D)(ALL_LABELS, 0);
draw byLabelsOnPolygon(C', A', noPoint)(ALL_LABELS, 0);
}
\drawCurrentPictureInMargin
\problem{T}{o}{find a third proportional to two given straight lines (\drawSizedLine{A'C'} and \drawSizedLine{AB}).}

\begin{center}
Em qualquer extremidade da linha dada \drawSizedLine{AB} trace \drawSizedLine{AC,CE} fazendo um ângulo;\\
pegue $\drawSizedLine{AC} = \drawSizedLine{A'C'}$ e trace \drawSizedLine{BC};\\
make $\drawSizedLine{BD} = \drawSizedLine{A'C'}$,\\
e trace $\drawSizedLine{DE} \parallel \drawSizedLine{BC}$; \byref{prop:I.XXXI}\\
\drawSizedLine{CE} é o terceiro proporcional a \drawSizedLine{AB} e \drawSizedLine{A'C'}.

For since $\drawSizedLine{BC} \parallel \drawSizedLine{DE}$,\\
$\therefore \drawSizedLine{AB} : \drawSizedLine{BD} :: \drawSizedLine{AC} : \drawSizedLine{CE}$ \byref{prop:VI.II};\\
but $\drawSizedLine{BD} = \drawSizedLine{AC} = \drawSizedLine{A'C'}$ \byref{\constref};\\
$\therefore \drawSizedLine{AB} : \drawSizedLine{A'C'} :: \drawSizedLine{A'C'} : \drawSizedLine{CE}$ \byref{prop:V.VII}.
\end{center}

\qed

\startproblem{Prop. XII. prob.}\label{prop:VI.XII}
\defineNewPicture[1/2]{
pair A, Ae, B, Be, C, Ce, D, E, F, G, H;
numeric l[], a;
l1 := 4/3u;
l2 := 11/6u;
l3 := 3/2u;
A := (0, 0); Ae := (l1, 0);
B := (0, 0); Be := (l2, 0);
C := (0, 0); Ce := (l3, 0);
a := 45;
byLineDefine.A(A, Ae, byblue, DASHED_LINE, REGULAR_WIDTH);
byLineDefine.B(B, Be, byred, DASHED_LINE, REGULAR_WIDTH);
byLineDefine.C(C, Ce, byyellow, DASHED_LINE, REGULAR_WIDTH);
forsuffixes i=A, B, C:
lineUseLineLabel.i := true;
endfor;
D := (0, 0);
G := D shifted (dir(a)*l1);
E := G shifted (dir(a)*l2);
H := D shifted (l3, 0);
F = whatever[D, H] = whatever[E, E shifted (G-H)];
draw byLine(G, H, byblack, SOLID_LINE, THIN_WIDTH);
byLineDefine(E, F, byblack, DASHED_LINE, REGULAR_WIDTH);
byLineDefine(D, H, byyellow, SOLID_LINE, REGULAR_WIDTH);
byLineDefine(H, F, byblack, SOLID_LINE, REGULAR_WIDTH);
byLineDefine(D, G, byblue, SOLID_LINE, REGULAR_WIDTH);
byLineDefine(G, E, byred, SOLID_LINE, REGULAR_WIDTH);
draw byNamedLineSeq(0)(EF,HF,DH,DG,GE);
draw byLabelsOnPolygon(D, G, E, F, H)(ALL_LABELS, 0);
}
\drawCurrentPictureInMargin
\problem{T}{o}{find a fourth proportional to three given lines
$\left\{\vcenter{
\nointerlineskip\hbox{\drawSizedLine{A}}
\nointerlineskip\hbox{\drawSizedLine{B}}
\nointerlineskip\hbox{\drawSizedLine{C}}}\right\}$.}

\begin{center}
Trace \drawSizedLine{DG,GE} e \drawSizedLine{DH,HF} fazendo qualquer ângulo;\\
take $\drawUnitLine{DG} = \drawUnitLine{A}$,\\
and $\drawUnitLine{GE} = \drawUnitLine{B}$,\\
also $\drawUnitLine{DH} = \drawUnitLine{C}$,\\
draw \drawUnitLine{GH},\\
and $\drawUnitLine{EF} \parallel \drawUnitLine{GH}$ \byref{prop:I.XXXI};\\
\drawSizedLine{HF} é o quarto proporcional.

Por conta dos paralelos,\\
$\drawUnitLine{DG} : \drawUnitLine{GE} :: \drawUnitLine{DH} : \drawUnitLine{HF}$ \byref{prop:VI.II};\\
but
$\left\{\vcenter{
\nointerlineskip\hbox{\drawSizedLine{A}}
\nointerlineskip\hbox{\drawSizedLine{B}}
\nointerlineskip\hbox{\drawSizedLine{C}}}\right\}
=
\left\{\vcenter{
\nointerlineskip\hbox{\drawSizedLine{DG}}
\nointerlineskip\hbox{\drawSizedLine{GE}}
\nointerlineskip\hbox{\drawSizedLine{DH}}}\right\}$
\byref{\constref};

$\therefore \drawUnitLine{A} : \drawUnitLine{B} :: \drawUnitLine{C} : \drawUnitLine{HF}$ \byref{prop:V.VII}.
\end{center}

\qed

\startproblem{Prop. XIII. prob.}\label{prop:VI.XIII}
\defineNewPicture{
pair Ab, Bb, Cb, A, B, C, D, O;
numeric l[], r;
path q;
l1 := 3u;
l2 := 2u;
r := 1/2*(l1 + l2);
Ab := (0, 0);
Bb := (l1, 0);
Cb := (l1 + l2, 0);
byLineDefine.A(Ab, Bb, byblue, DASHED_LINE, REGULAR_WIDTH);
byLineDefine.B(Bb, Cb, byred, DASHED_LINE, REGULAR_WIDTH);
lineUseLineLabel.A := true;
lineUseLineLabel.B := true;
A := (0, 0);
B := (l1, 0);
C := (l1 + l2, 0);
O := 1/2[A, C];
q := (fullcircle scaled 2r) shifted O;
D := q intersectionpoint (B--(B shifted (0, r)));
byAngleDefine(A, D, C, byblue, SOLID_SECTOR);
draw byNamedAngleResized();
byLineDefine(A, C, byblack, SOLID_LINE, THIN_WIDTH);
byLineStylize(A, C, 0, 0, -1)(AC);
draw byMarkLine(1/2, byblue)(AC);
draw byLine(B, D, byblack, SOLID_LINE, REGULAR_WIDTH);
byLineDefine(A, B, byblue, SOLID_LINE, REGULAR_WIDTH);
byLineDefine(B, C, byred, SOLID_LINE, REGULAR_WIDTH);
byLineDefine(A, D, byyellow, SOLID_LINE, REGULAR_WIDTH);
byLineDefine(C, D, byyellow, DASHED_LINE, REGULAR_WIDTH);
draw byNamedLineSeq(1)(AB,BC,CD,AD);
draw byArc.O(O, C, A, r, byred, 0, 0, 0, 0);
draw byLabelsOnPolygon(C, B, A, noPoint)(ALL_LABELS, 0);
draw byLabelsOnCircle(D)(O);
}
\drawCurrentPictureInMargin
\problem{T}{o}{encontre um média proporcional entre duas linhas retas dadas $\left\{\vcenter{
\nointerlineskip\hbox{\drawSizedLine{A}}
\nointerlineskip\hbox{\drawSizedLine{B}}}\right\}$.}

\begin{center}
Trace qualquer linha reta \drawSizedLine{AB,BC},\\
make $\drawSizedLine{AB} = \drawSizedLine{A}$ and $\drawSizedLine{BC} = \drawSizedLine{B}$;\\
dividir \drawSizedLine{AB,BC} ao meio; e do ponto de bissecção como centro e da meia linha como raio, descreva um semicírculo
\drawFromCurrentPicture[bottom]{
draw byNamedLineFull(A, C, 0, 0,  0, -1)(AC);
draw byNamedArc(O);
draw byLabelsOnPolygon(C, A, noPoint)(ALL_LABELS, 0);
},\\
draw $\drawSizedLine{BD} \perp \drawSizedLine{AB}$:\\
\drawSizedLine{BD} é o média proporcional necessário.

Trace \drawUnitLine{AD} e \drawUnitLine{CD}.

Como \drawAngle{D} é um ângulo reto \byref{prop:III.XXXI},\\
e \drawUnitLine{BD} é $\perp$ dele no lado oposto,\\
$\therefore$ \drawUnitLine{BD} é um média proporcional entre \drawUnitLine{AB} e \drawUnitLine{BC} \byref{prop:VI.VIII},\\
and $\therefore$ between \drawUnitLine{A} and \drawUnitLine{B} \byref{\constref}.
\end{center}

\qed

\starttheorem{Prop. XIV. Teor.}\label{prop:VI.XIV}
\defineNewPicture[1/4]{
pair A, B, C, D, E, F, G, H, d[];
numeric a;
d1 := (3/2u, 0);
d2 := (1/2u, -2u);
d3 := 2/3d1;
d4 := 3/2d2;
C := (0, 0);
E := C shifted d1;
G := C shifted d2;
B := C shifted (d1 + d2);
F := B shifted d3;
D := B shifted d4;
A := B shifted (d3 + d4);
H = whatever[A, F] = whatever[C, E];
draw byPolygon(C,E,B,G)(byyellow);
draw byPolygon(B,F,A,D)(byblue);
draw byPolygon(E,H,F,B)(byred);
draw byLine(B, G, byred, SOLID_LINE, REGULAR_WIDTH);
draw byLine(B, F, byblack, SOLID_LINE, REGULAR_WIDTH);
draw byLine(B, E, byblue, SOLID_LINE, REGULAR_WIDTH);
draw byLine(B, D, byyellow, SOLID_LINE, REGULAR_WIDTH);
draw byLabelsOnPolygon(B, G, C, E, H, F, A, D)(ALL_LABELS, 0);
}
\drawCurrentPictureInMargin
\problem{P}{aralelogramos}{iguais \drawPolygon{BFAD} e \drawPolygon{CEBG}, que têm um ângulo em cada igual,
possuem os lados em torno desses ângulos reciprocamente proporcionais
($\drawUnitLine{BG} : \drawUnitLine{BF} :: \drawUnitLine{BD} : \drawUnitLine{BE}$).\\
E paralelogramos que têm um ângulo em cada igual, e os lados em torno deles reciprocamente proporcionais, são iguais.}

\begin{center}
Seja \drawUnitLine{BG} e \drawUnitLine{BF}; e \drawUnitLine{BD} e \drawUnitLine{BE}, sejam colocados de forma que \drawUnitLine{BG,BF} e \drawUnitLine{BD,BE} possam continuar em linha reta. É evidente que eles podem assumir esta posição. \byref{prop:I.XIII,prop:I.XIV,prop:I.XV}

Complete \drawPolygon{EHFB}.

Como $\drawPolygon{CEBG} = \drawPolygon{BFAD}$;

$\therefore \drawPolygon{CEBG} : \drawPolygon{EHFB} :: \drawPolygon{BFAD} : \drawPolygon{EHFB}$ \byref{prop:V.VII}\\
$\therefore \drawUnitLine{BG} : \drawUnitLine{BF} :: \drawUnitLine{BD} : \drawUnitLine{BE}$ \byref{prop:VI.I}

Mantida a mesma construção:\\
$\drawUnitLine{BG} : \drawUnitLine{BF} ::
\left\{
\begin{aligned}
	\drawPolygon{CEBG} &: \drawPolygon{EHFB} \mbox{\byref{prop:VI.I}}\\
	\drawUnitLine{BD} &: \drawUnitLine{BE} \mbox{\byref{\hypref}}\\
	\drawPolygon{BFAD} &: \drawPolygon{EHFB} \mbox{\byref{prop:VI.I}}\\
\end{aligned}
\right.$\\
$\therefore \drawPolygon{CEBG} : \drawPolygon{EHFB} :: \drawPolygon{BFAD} : \drawPolygon{EHFB}$ \byref{prop:V.XI}\\
e, portanto, $\drawPolygon{CEBG} = \drawPolygon{BFAD}$ \byref{prop:V.IX}.
\end{center}

\qed

\starttheorem{Prop. XV. Teor.}\label{prop:VI.XV}
\defineNewPicture{
pair A, B, C, D, E;
D := (0, 0);
B := (3u, 0);
A := (5/4u, -3/2u);
C := A shifted 3/2(A-D);
E := A shifted 3/2(A-B);
draw byPolygon(D,A,B)(byblue);
draw byPolygon(D,A,E)(byred);
draw byPolygon(A,B,C)(byyellow);
byAngleDefine(D, A, E, byblue, SOLID_SECTOR);
byAngleDefine(B, A, C, byred, SOLID_SECTOR);
draw byNamedAngleResized();
byLineDefine(B, D, byblack, DASHED_LINE, REGULAR_WIDTH);
byLineDefine(D, A, byyellow, SOLID_LINE, REGULAR_WIDTH);
byLineDefine(B, A, byblack, SOLID_LINE, REGULAR_WIDTH);
byLineDefine(A, C, byred, SOLID_LINE, REGULAR_WIDTH);
byLineDefine(A, E, byblue, SOLID_LINE, REGULAR_WIDTH);
draw byNamedLineSeq(0)(noLine,AE,BA,BD,DA,AC);
draw byLabelsOnPolygon(A, E, D, B, C)(ALL_LABELS, 0);
}
\drawCurrentPictureInMargin
\problem{T}{riângulos}{que têm um ângulo em cada igual ($\drawAngle{DAE} = \drawAngle{BAC}$) têm os lados em torno dos ângulos iguais reciprocamente proporcionais
($\drawUnitLine{AE} : \drawUnitLine{BA} :: \drawUnitLine{AC} : \drawUnitLine{DA}$).\\
E dois triângulos que têm um ângulo de um igual ao ângulo do outro, e os lados em torno dos ângulos iguais reciprocamente proporcionais, são iguais.}

\startsubproposition{I.}
Sejam os triângulos colocados de forma que os ângulos iguais \drawAngle{DAE} e \drawAngle{BAC} possam ser verticalmente opostos, ou seja, de modo que \drawUnitLine{AE} e \drawUnitLine{BA} possam estar na mesma linha reta. Daí também \drawUnitLine{AC} e \drawUnitLine{DA} devem estar na mesma linha reta \byref{prop:I.XIV}.

\begin{center}
Trace \drawUnitLine{BD}; então\\
$\begin{aligned}
\drawUnitLine{DA} : \drawUnitLine{BA}
&:: \drawPolygon{DAE} : \drawPolygon{DAB} \mbox{\byref{prop:VI.I}}\\
&:: \drawPolygon{ABC} : \drawPolygon{DAB} \mbox{\byref{prop:V.VII}}\\
&:: \drawUnitLine{AC} : \drawUnitLine{DA} \mbox{\byref{prop:VI.I}}\\
\end{aligned}$

$\therefore \drawUnitLine{AE} : \drawUnitLine{BA} :: \drawUnitLine{AC} : \drawUnitLine{DA}$ \byref{prop:V.XI}.
\end{center}

\vfill\pagebreak

\startsubproposition{II.}
\begin{center}
Mantida a mesma construção, e\\
$\drawPolygon{DAE} : \drawPolygon{DAB} :: \drawUnitLine{DA} : \drawUnitLine{BA}$ \byref{prop:VI.I}\\
e $\drawUnitLine{AC} : \drawUnitLine{DA} :: \drawPolygon{ABC} : \drawPolygon{DAB}$ \byref{prop:VI.I}

mas $ \drawUnitLine{AE} : \drawUnitLine{BA} :: \drawUnitLine{AC} : \drawUnitLine{DA}$, \byref{\hypref}

$\therefore \drawPolygon{DAE} : \drawPolygon{DAB} :: \drawPolygon{ABC} : \drawPolygon{DAB}$ \byref{prop:V.XI};

$\therefore \drawPolygon{DAE} = \drawPolygon{ABC}$ \byref{prop:V.IX}.
\end{center}

\qed

\starttheorem{Prop. XVI. Teor.}\label{prop:VI.XVI}
\defineNewPicture{
pair A, B, C, D, Eb, Fb, Ee, Fe, E, F, G, H, d[];
numeric l[], h[];
l1 := -3u;
l2 := -2u;
h1 := -3/4l2;
h2 := -3/4l1;
A := (0, 0);
B := (l1, 0);
G := (0, h1);
F := (l1, h1);
d1 := (0, -h2 - 1/2u);
C := (0, 0) shifted d1;
D := (l2, 0) shifted d1;
H := (0, h2) shifted d1;
E := (l2, h2) shifted d1;
d2 := (0, h1 + u);
d3 := (0, h1 + 1/2u);
Eb := (0, 0) shifted d2; Ee := (-h2, 0) shifted d2;
Fb := (0, 0) shifted d3; Fe := (-h1, 0) shifted d3;
draw byPolygon(A,B,F,G)(byred);
draw byPolygon(C,D,E,H)(byyellow);
byLineDefine(A, B, byyellow, SOLID_LINE, REGULAR_WIDTH);
byLineDefine(A, G, byblack, SOLID_LINE, REGULAR_WIDTH);
draw byNamedLineSeq(0)(AB,AG);
byLineDefine(C, D, byblue, SOLID_LINE, REGULAR_WIDTH);
byLineDefine(C, H, byred, SOLID_LINE, REGULAR_WIDTH);
draw byNamedLineSeq(0)(CD,CH);
draw byLineWithName(Ee, Eb, byred, 1, 0)(E);
draw byLineWithName(Fe, Fb, byblack, 1, 0)(F);
draw byLabelsOnPolygon(F, G, A, B)(ALL_LABELS, 0);
draw byLabelsOnPolygon(E, H, C, D)(ALL_LABELS, 0);
draw byLabelLine(0)(E, F);
}
\drawCurrentPictureInMargin
\problem{S}{e}{quatro retas forem proporcionais ($\drawUnitLine{AB} : \drawUnitLine{CD} :: \drawUnitLine{E} : \drawUnitLine{F}$), o retângulo ($\drawUnitLine{AB} \times \drawUnitLine{F}$) compreendido pelos extremos é igual ao retângulo ($\drawUnitLine{CD} \times \drawUnitLine{E}$) compreendido pelos médios.\\
E, se o retângulo compreendido pelos extremos for igual ao retângulo compreendido pelos médios, as quatro retas são proporcionais.}

\startsubproposition{Parte I.}
\begin{center}
Das extremidades de \drawUnitLine{AB} e \drawUnitLine{CD} trace \drawUnitLine{AG} e \drawUnitLine{CH} $\perp$ para eles e $= \drawUnitLine{F}$ e \drawUnitLine{E} respectivamente:\\
complete os paralelogramos \drawPolygon{ABFG} e \drawPolygon{CDEH}.

e, como,\\
$\drawUnitLine{AB} : \drawUnitLine{CD} :: \drawUnitLine{E} : \drawUnitLine{F}$ \byref{\hypref}\\
$\therefore \drawUnitLine{AB} : \drawUnitLine{CD} :: \drawUnitLine{CH} : \drawUnitLine{AG}$ \byref{\constref}\\
$\therefore \drawPolygon{ABFG} = \drawPolygon{CDEH}$ \byref{prop:VI.XIV},\\
isto é, o retângulo contido nos extremos, igual ao retângulo contido nos meios.
\end{center}

\vfill\pagebreak

\startsubproposition{Parte II.}
\begin{center}
Mantida a mesma construção;\\
$\because \drawUnitLine{F} = \drawUnitLine{AG}$, $\drawPolygon{ABFG} = \drawPolygon{CDEH}$\\ %because
e $\drawUnitLine{CH} = \drawUnitLine{E}$.\\
$\therefore \drawUnitLine{AB} : \drawUnitLine{CD} :: \drawUnitLine{CH} : \drawUnitLine{AG}$ \byref{prop:VI.XIV}

Mas $\drawUnitLine{CH} = \drawUnitLine{E}$,\\
e $\drawUnitLine{AG} = \drawUnitLine{F}$ \byref{\constref}\\
$\therefore \drawUnitLine{AB} : \drawUnitLine{CD} :: \drawUnitLine{E} : \drawUnitLine{F}$ \byref{prop:V.VII}.
\end{center}

\qed

\starttheorem{Prop. XVII. Teor.}\label{prop:VI.XVII}
\defineNewPicture{
pair Ab, Ae, Bb, Be, Cb, Ce, Db, De;
pair A, B, C, D, E, F, G, H;
pair d[];
numeric l[], r;
r := 3/4;
l1 := 3u;
l2 := r*l1;
l3 := r*l2;
d1 := (0, 4/2u); d2 := (0, 3/2u); d3 := (0, 2/2u); d4 := (0, 1/2u);
Ae := (0, 0) shifted d1; Ab := (-l1, 0) shifted d1;
Be := (0, 0) shifted d2; Bb := (-l2, 0) shifted d2;
Ce := (0, 0) shifted d3; Cb := (-l3, 0) shifted d3;
De := (0, 0) shifted d4; Db := (-l2, 0) shifted d4;
d5 := (0, -l2);
A := (0, 0) shifted d5; B := (-l2, 0) shifted d5; C := (-l2, l2) shifted d5; D := (0, l2) shifted d5;
d6 := (0, -l2-l1-1/2u);
E := (0, 0) shifted d6; F := (-l3, 0) shifted d6; G := (-l3, l1) shifted d6; H := (0, l1) shifted d6;
draw byPolygon(A,B,C,D)(byred);
byLineDefine(A, B, byyellow, SOLID_LINE, REGULAR_WIDTH);
byLineDefine(A, D, byblue, SOLID_LINE, REGULAR_WIDTH);
draw byNamedLineSeq(0)(AB,AD);
draw byPolygon(E,F,G,H)(byyellow);
byLineDefine(E, F, byblack, SOLID_LINE, REGULAR_WIDTH);
byLineDefine(E, H, byred, SOLID_LINE, REGULAR_WIDTH);
draw byNamedLineSeq(0)(EF,EH);
draw byLineWithName(Ab, Ae, byred, 0, 0)(A);
draw byLineWithName(Bb, Be, byblue, 0, 0)(B);
draw byLineWithName(Cb, Ce, byblack, 0, 0)(C);
draw byLineWithName(Db, De, byyellow, 0, 0)(D);
draw byLabelLine(0)(A, B, C, D);
draw byLabelPolygon(1)(ABCD);
draw byLabelPolygon(1)(EFGH);
}
\drawCurrentPictureInMargin
\problem[4]{S}{e}{três retas forem proporcionais ($\drawUnitLine{A} : \drawUnitLine{B} :: \drawUnitLine{B} : \drawUnitLine{C}$), o retângulo sob os extremos é igual ao quadrado da média.\\
E, se o retângulo sob os extremos for igual ao quadrado da média, as três retas são proporcionais.}

\startsubproposition{Parte I.}
\begin{center}
$\begin{aligned}
\mbox{ Suponha } \drawUnitLine{D} &= \drawUnitLine{B}, \\
\mbox{ e, como } \drawUnitLine{A} : \drawUnitLine{B} &:: \drawUnitLine{B} : \drawUnitLine{C}, \\
\mbox{ então } \drawUnitLine{A} : \drawUnitLine{B} &:: \drawUnitLine{D} : \drawUnitLine{C}, \\
\therefore \drawUnitLine{A} \times \drawUnitLine{C} &= \drawUnitLine{B} \times \drawUnitLine{D} \\
\end{aligned}$\\
\byref{prop:VI.XVI}.

Mas $\drawUnitLine{D} = \drawUnitLine{B}$,\\
$\therefore \drawUnitLine{B} \times \drawUnitLine{D} = \drawUnitLine{B} \times \drawUnitLine{B} \mbox{ ou } = \drawUnitLine{B}^2$;\\
portanto, se as três retas são proporcionais, o retângulo compreendido pelos extremos é igual ao quadrado da média.
\end{center}

\startsubproposition{Parte II.}
\begin{center}
Suponha $\drawUnitLine{D} = \drawUnitLine{B}$,\\
então $\drawUnitLine{A} \times \drawUnitLine{C} = \drawUnitLine{D} \times \drawUnitLine{B}$,\\
$\therefore \drawUnitLine{A} : \drawUnitLine{B} :: \drawUnitLine{D} : \drawUnitLine{C}$ \byref{prop:VI.XVI},\\
e $\drawUnitLine{A} : \drawUnitLine{B} :: \drawUnitLine{B} : \drawUnitLine{C}$.
\end{center}

\qed

\startproblem{Prop. XVIII. prob.}\label{prop:VI.XVIII}
\defineNewPicture[1/4]{
pair C, D, E, F, L, d;
C := (0, 0);
D := (3/2u, 0);
E := (u, 3u);
F := (-u, 2u);
L := (5/2u, u);
draw byPolygon(D,C,F)(byyellow);
draw byPolygon(D,F,E)(byblue);
draw byPolygon(D,E,L)(byred);
byAngleDefineExtended(D, F, E, byred, 1)(byblack);
byAngleDefineExtended(D, E, L, byyellow, 1)(byred);
byAngleDefine(D, C, F, byred, ARC_SECTOR);
byAngleDefine(F, D, C, byblue, ARC_SECTOR);
byAngleDefine(E, D, F, byblack, ARC_SECTOR);
byAngleDefine(L, D, E, byyellow, ARC_SECTOR);
draw byNamedAngleResized(DFE, DEL, DCF, FDC, EDF, LDE);
draw byLine(D, F, byred, DASHED_LINE, REGULAR_WIDTH);
byLineDefine(D, E, byyellow, DASHED_LINE, REGULAR_WIDTH);
byLineDefine(C, D, byblack, DASHED_LINE, REGULAR_WIDTH);
byLineDefine(C, F, byblue, SOLID_LINE, THIN_WIDTH);
byLineDefine(F, E, byyellow, SOLID_LINE, REGULAR_WIDTH);
draw byNamedLineSeq(0)(DE,FE,CF,CD);
pair A, B, H, G, K;
numeric s;
s := 5/6;
d := (0, -7/2u);
A := (C scaled s) shifted d;
B := (D scaled s) shifted d;
H := (E scaled s) shifted d;
G := (F scaled s) shifted d;
K := (L scaled s) shifted d;
draw byPolygon(B,A,G)(byyellow);
draw byPolygon(B,G,H)(byblue);
draw byPolygon(B,H,K)(byred);
byAngleDefineExtended(B, G, H, byblack, 1)(byred);
byAngleDefineExtended(B, H, K, byyellow, 1)(byred);
byAngleDefine(B, A, G, byred, SOLID_SECTOR);
byAngleDefine(G, B, A, byblue, SOLID_SECTOR);
byAngleDefine(H, B, G, byblack, SOLID_SECTOR);
byAngleDefine(K, B, H, byyellow, SOLID_SECTOR);
draw byNamedAngleResized(BGH, BHK, BAG, GBA, HBG, KBH);
byLineDefine(B, G, byred, SOLID_LINE, REGULAR_WIDTH);
byLineDefine(A, B, byblack, SOLID_LINE, REGULAR_WIDTH);
byLineDefine(A, G, byblue, SOLID_LINE, REGULAR_WIDTH);
byLineDefine(G, H, byblue, DASHED_LINE, REGULAR_WIDTH);
draw byNamedLineSeq(0)(noLine,BG,AB,AG,GH);
draw byLabelsOnPolygon(F, E, L, D, C)(ALL_LABELS, 0);
draw byLabelsOnPolygon(G, H, K, B, A)(ALL_LABELS, 0);
}
\drawCurrentPictureInMargin
\problem{E}{m}{uma dada linha reta (\drawUnitLine{AB}) para construir uma figura retilínea semelhante a uma dada (\drawPolygon[middle][polygonCD]{DCF,DFE,DEL}) e colocada de forma semelhante.}

\begin{center}
Resolva a figura dada em triângulos desenhando as linhas \drawUnitLine{DF} e \drawUnitLine{DE}.

Nas extremidades do \drawUnitLine{AB} faça\\
$\drawAngle{GBA} = \drawAngle{FDC}$ and $\drawAngle{A} = \drawAngle{C}$;\\
novamente nas extremidades do \drawUnitLine{BG} make\\
$\drawAngle{G} = \drawAngle{F}$ e $\drawAngle{HBG} = \drawAngle{EDF}$;\\
de modo semelhante faça\\
$\drawAngle{KBH} = \drawAngle{LDE}$ e $\drawAngle{H} = \drawAngle{E}$.

Então $\polygonCD = \drawPolygon[middle][polygonAB]{BAG,BGH,BHK}$.

É evidente pela construção e \byref{prop:I.XXXII} que os números são equiangulares; e como os triângulos \drawPolygon{DCF} e \drawPolygon{BAG} são equiângulos;\\
então, por \byref{prop:VI.IV},
$\drawUnitLine{AB}:\drawUnitLine{AG} :: \drawUnitLine{CD} : \drawUnitLine{CF}$\\
and $\drawUnitLine{AG}:\drawUnitLine{BG} :: \drawUnitLine{CF} : \drawUnitLine{DF}$

Novamente, como \drawPolygon{DFE} e \drawPolygon{BGH} são equiangulares,\\
$\drawUnitLine{BG}:\drawUnitLine{GH} :: \drawUnitLine{DF} : \drawUnitLine{FE}$\\
$\therefore$ ex \ae quali,\\
$\drawUnitLine{AG}:\drawUnitLine{GH} :: \drawUnitLine{CF} : \drawUnitLine{FE}$ \byref{prop:V.XXII} % improvement: Byrne had VI.XXII wrongly referenced here

Da mesma forma, pode-se mostrar que os restantes lados das duas figuras são proporcionais.

$\therefore$ by \byref{prop:VI.I}\\
\polygonAB\ é semelhante a \polygonCD\ e tem localização semelhante; e na linha fornecida \drawUnitLine{AB}.
\end{center}

\qed

\starttheorem{Prop. XIX. Teor.}\label{prop:VI.XIX}
\defineNewPicture[1/4]{
pair A, B, C, D, E, F, G, d;
numeric s;
A := (5/2u, -3u);
B := (0, 0);
C := (-u, ypart(A));
d := (0, -4u);
s := 3/4;
D := (A scaled s) shifted d;
E := (B scaled s) shifted d;
F := (C scaled s) shifted d;
G := B shifted (unitvector(C-B) scaled ((abs(E-F)/abs(B-C))*abs(E-F)));
draw byPolygon(A,B,G)(byblue);
draw byPolygon(A,G,C)(byred);
byAngleDefine(A, B, C, byred, SOLID_SECTOR);
draw byNamedAngleResized(ABC);
draw byLine(A, G, byyellow, DASHED_LINE, REGULAR_WIDTH);
byLineDefine(A, B, byyellow, SOLID_LINE, REGULAR_WIDTH);
byLineDefine(B, G, byblack, DASHED_LINE, REGULAR_WIDTH);
byLineDefine(G, C, byblack, SOLID_LINE, REGULAR_WIDTH);
draw byNamedLineSeq(0)(AB,BG,GC);
draw byPolygon(D,E,F)(byyellow);
byAngleDefine(D, E, F, byblack, SOLID_SECTOR);
draw byNamedAngleResized(DEF);
byLineDefine(D, E, byred, SOLID_LINE, REGULAR_WIDTH);
byLineDefine(E, F, byblue, SOLID_LINE, REGULAR_WIDTH);
draw byNamedLineSeq(0)(DE,EF);
draw byLabelsOnPolygon(F, E, D)(ALL_LABELS, 0);
draw byLabelsOnPolygon(C, G, B, A)(ALL_LABELS, 0);
}
\drawCurrentPictureInMargin
\problem{T}{riângulos}{semelhantes (\drawPolygon[bottom][triangleDEF]{DEF} e \drawPolygon[bottom][triangleABC]{ABG,AGC}) estão entre si na razão duplicada de seus lados homólogos.}

Sejam \drawAngle{E} e \drawAngle{B} ângulos iguais, e \drawUnitLine{BG,GC} e \drawUnitLine{EF} homólogos lados dos triângulos semelhantes \triangleDEF\ e \triangleABC\ e em \drawUnitLine{BG,GC} a maior dessas retas toma \drawUnitLine{BG} uma terceira proporcional, de modo que
\begin{center}
$\drawUnitLine{BG,GC} : \drawUnitLine{EF} :: \drawUnitLine{EF} : \drawUnitLine{BG}$;\\
draw \drawUnitLine{AG}.\\
$\drawUnitLine{BG,GC} : \drawUnitLine{AB} :: \drawUnitLine{EF} : \drawUnitLine{DE}$ \byref{prop:VI.IV};\\
$\therefore \drawUnitLine{BG,GC} : \drawUnitLine{EF} :: \drawUnitLine{AB} : \drawUnitLine{DE}$ \byref{prop:V.XVI},\\
but $\drawUnitLine{BG,GC} : \drawUnitLine{EF} :: \drawUnitLine{EF} : \drawUnitLine{BG}$ \byref{\constref},\\
$\therefore \drawUnitLine{EF} : \drawUnitLine{BG} :: \drawUnitLine{AB} : \drawUnitLine{DE}$\\
conseqüentemente $\triangleDEF = \drawPolygon[bottom][triangleABG]{ABG}$ pois eles têm os lados aproximadamente nos ângulos iguais \drawAngle{E} e \drawAngle{B} reciprocamente proporcionais \byref{prop:VI.XV};

$\therefore \triangleABC : \triangleDEF :: \triangleABC : \triangleABG$ \byref{prop:V.VII};\\
but $\triangleABC : \triangleABG :: \drawUnitLine{BG,GC} : \drawUnitLine{BG}$ \byref{prop:VI.I},

$\therefore \triangleABC : \triangleDEF :: \drawUnitLine{BG,GC} : \drawUnitLine{BG}$,\\
isto é, os triângulos estão entre si na razão duplicada de seus lados homólogos \drawUnitLine{EF} e \drawUnitLine{BG,GC} \byref{def:V.XI}.
\end{center}

\qed

\starttheorem{Prop. XX. Teor.}\label{prop:VI.XX}
\defineNewPicture{
pair A, B, C, D, E;
A := (-2u, 5/2u);
B := (-1/2u, 0);
C := (3/2u, ypart(B));
D := (2u, 3/2u);
E := (u, ypart(A));
draw byPolygon(B,E,A)(byblue);
draw byPolygon(B,D,E)(byred);
draw byPolygon(B,C,D)(byyellow);
byAngleDefine(B, D, E, byblue, SOLID_SECTOR);
byAngleDefine(B, C, D, byblack, SOLID_SECTOR);
byAngleDefine(B, D, C, byred, SOLID_SECTOR);
draw byNamedAngleResized(BDE, BCD, BDC);
draw byLine(B, D, byblack, DASHED_LINE, REGULAR_WIDTH);
byLineDefine(B, E, byblack, SOLID_LINE, REGULAR_WIDTH);
byLineDefine(B, C, byblue, SOLID_LINE, REGULAR_WIDTH);
byLineDefine(C, D, byblue, DASHED_LINE, REGULAR_WIDTH);
byLineDefine(D, E, byyellow, SOLID_LINE, REGULAR_WIDTH);
draw byNamedLineSeq(0)(BE,BC,CD,DE);
pair F, G, H, K, L, d;
numeric s;
s := 3/4;
d := (0, -3u);
F := (A scaled s) shifted d;
G := (B scaled s) shifted d;
H := (C scaled s) shifted d;
K := (D scaled s) shifted d;
L := (E scaled s) shifted d;
draw byPolygon(G,L,F)(byblue);
draw byPolygon(G,K,L)(byred);
draw byPolygon(G,H,K)(byyellow);
byAngleDefine(G, K, L, byblue, ARC_SECTOR);
byAngleDefine(G, H, K, byblack, ARC_SECTOR);
byAngleDefine(G, K, H, byred, ARC_SECTOR);
draw byNamedAngleResized(GKL, GHK, GKH);
draw byLine(G, K, byblack, DASHED_LINE, THIN_WIDTH);
byLineDefine(G, L, byblack, SOLID_LINE, THIN_WIDTH);
byLineDefine(G, H, byred, SOLID_LINE, THIN_WIDTH);
byLineDefine(H, K, byred, DASHED_LINE, REGULAR_WIDTH);
byLineDefine(K, L, byyellow, SOLID_LINE, REGULAR_WIDTH);
draw byNamedLineSeq(0)(GL,GH,HK,KL);
draw byLabelsOnPolygon(A, E, D, C, B)(ALL_LABELS, 0);
draw byLabelsOnPolygon(F, L, K, H, G)(ALL_LABELS, 0);
}
\drawCurrentPictureInMargin
\problem{P}{olígonos}{semelhantes podem ser divididos no mesmo número de triângulos semelhantes, cada par semelhante sendo proporcional aos polígonos; e os polígonos estão entre si no razão duplicada de seus lados homólogos.}

Trace \drawUnitLine{BE} e \drawUnitLine{BD} e \drawUnitLine{GL} e \drawUnitLine{GK}, resolvendo os polígonos em triângulos. Então, como os polígonos são semelhantes, $\drawAngle{C} = \drawAngle{H}$ e $\drawUnitLine{BC} : \drawUnitLine{CD} :: \drawUnitLine{GH} : \drawUnitLine{HK}$

\begin{center}
$\therefore$ \drawPolygon{BCD} e \drawPolygon{GHK} são semelhantes,\\
and $\drawAngle{BDC} = \drawAngle{GKH}$ \byref{prop:VI.VI};

mas $\drawAngle{BDE,BDC} = \drawAngle{GKL,GKH}$ porque são ângulos de polígonos semelhantes; portanto os restos \drawAngle{BDE} e \drawAngle{GKL} são iguais;\\
hence $\drawUnitLine{BD} : \drawUnitLine{CD} :: \drawUnitLine{GK} : \drawUnitLine{HK}$,\\
por conta dos triângulos semelhantes,\\
and $\drawUnitLine{CD} : \drawUnitLine{DE} :: \drawUnitLine{HK} : \drawUnitLine{KL}$,\\
por conta dos polígonos semelhantes,\\
$\therefore \drawUnitLine{BD} : \drawUnitLine{DE} :: \drawUnitLine{GK} : \drawUnitLine{KL}$,\\
ex \ae quali \byref{prop:V.XXII}, e como esses lados proporcionais contêm ângulos iguais, os triângulos \drawPolygon{BDE} e \drawPolygon{GKL} são semelhantes a \byref{prop:VI.VI}.

Da mesma forma pode ser mostrado que os triângulos \drawPolygon{BEA} e \drawPolygon{GLF} são semelhantes.

Mas \drawPolygon{BCD} é para \drawPolygon{GHK} na razão duplicada de \drawUnitLine{BD} para \drawUnitLine{GK} \byref{prop:VI.XIX},\\
e \drawPolygon{BDE} é para \drawPolygon{GKL} da mesma maneira, na razão duplicada de \drawUnitLine{BD} para \drawUnitLine{GK};\\
$\therefore \drawPolygon{BCD} : \drawPolygon{GHK} :: \drawPolygon{BDE} : \drawPolygon{GKL}$ \byref{prop:V.XI}.

Novamente \drawPolygon{BDE} é para \drawPolygon{GKL} na razão duplicada de \drawUnitLine{BE} para \drawUnitLine{GL}, e \drawPolygon{BEA} é para \drawPolygon{GLF} na razão duplicada de \drawUnitLine{BE} para \drawUnitLine{GL}.
$\begin{aligned}
\drawPolygon{BCD} : \drawPolygon{GHK} &:: \drawPolygon{BDE} : \drawPolygon{GKL} \\
&:: \drawPolygon{BEA} : \drawPolygon{GLF} \\
\end{aligned}$;\\
e assim como um dos antecedentes está para um dos consequentes, assim também a soma de todos os antecedentes está para a soma de todos os consequentes; isto é, os triângulos semelhantes têm entre si a mesma proporção que os polígonos \byref{prop:V.XII}.

Mas \drawPolygon{BCD} está para \drawPolygon{GHK} na razão duplicada de \drawUnitLine{BC} para \drawUnitLine{GH};\\
$\therefore$ \drawPolygon{BEA,BDE,BCD} é \drawPolygon{GLF,GKL,GHK} na razão duplicada de \drawUnitLine{BC} para \drawUnitLine{GH}.
\end{center}

\qed

\starttheorem{Prop. XXI. Teor.}\label{prop:VI.XXI}
\defineNewPicture{
pair Aa, Ab, Ac, Ba, Bb, Bc, Ca, Cb, Cc, d[];
numeric s[];
Aa := (0, 0);
Ab := (0, 3/2u);
Ac := (-3u, 0);
s1 := 5/6;
d1 := (0, -2u);
Ba := (Aa scaled s1) shifted d1;
Bb := (Ab scaled s1) shifted d1;
Bc := (Ac scaled s1) shifted d1;
s2 := 4/6;
d2 := d1 shifted (0, -2u * s1) ;
Ca := (Aa scaled s2) shifted d2;
Cb := (Ab scaled s2) shifted d2;
Cc := (Ac scaled s2) shifted d2;
draw byPolygon.A(Aa,Ab,Ac)(byred);
draw byPolygon.B(Ba,Bb,Bc)(byblue);
draw byPolygon.C(Ca,Cb,Cc)(byyellow);
byPointLabelRemove(Aa, Ba, Ca, Ab, Bb, Cb, Ac, Bc, Cc);
}
\drawCurrentPictureInMargin
\problem{F}{iguras}{retilíneas (\drawPolygon[bottom]{A} e \drawPolygon[bottom]{B}) que são semelhantes à mesma figura (\drawPolygon[bottom]{C}) também são semelhantes entre si.}

Como \polygonA\ e \polygonC\ são semelhantes, eles são equiangulares e têm os lados aproximadamente iguais aos ângulos proporcionais \byref{def:VI.I}; e como as figuras \polygonB\ e \polygonC\ também são semelhantes, elas são equiangulares e têm os lados aproximadamente iguais aos ângulos proporcionais; portanto \polygonA\ e \polygonB\ também são equiangulares e têm os lados sobre os ângulos iguais proporcionais \byref{prop:V.XI} e são, portanto, semelhantes.

\qed

\starttheorem{Prop. XXII. Teor.}\label{prop:VI.XXII}
\defineNewPicture[1/2]{
pair A, B, C, D, E, F, G, H, Ob, Oe, Pb, Pe;
pair K, L, Ma, Mb, Mc, Md, Na, Nb, Nc, Nd;
pair d[];
numeric s[];
s1 := 6/5;
A := (0, 0); B := (6/5u, 0); K := (u, u);
d1 := (7/4u, 0);
C := (A scaled s1) shifted d1;
D := (B scaled s1) shifted d1;
L := (K scaled s1) shifted d1;
d2 := (7/2u, 0);
Ob := (A scaled (s1*s1)) shifted d2;
Oe := (B scaled (s1*s1)) shifted d2;
E := (0, 0); F := (5/6u, 0); Ma := (4/3u, 2/3u); Mb := (u, 4/3u); Mc := (1/5u, 3/2u); Md := (-1/4u, 5/6u);
G := (E scaled s1) shifted d1;
H := (F scaled s1) shifted d1;
Na := (Ma scaled s1) shifted d1;
Nb := (Mb scaled s1) shifted d1;
Nc := (Mc scaled s1) shifted d1;
Nd := (Md scaled s1) shifted d1;
Pb := (E scaled (s1*s1)) shifted d2;
Pe := (F scaled (s1*s1)) shifted d2;
d3 := (1/4u, -5/2u);
forsuffixes i=E,F,Ma,Mb,Mc,Md,G,H,Na,Nb,Nc,Nd,Pb,Pe:
i := i shifted d3;
endfor;
draw byPolygon.AB(A,B,K)(byyellow);
draw byPolygon.CD(C,D,L)(byred);
draw byPolygon.EF(E,F,Ma,Mb,Mc,Md)(byblue);
draw byPolygon.GH(G,H,Na,Nb,Nc,Nd)(white);
draw byLine(A, B, byblack, SOLID_LINE, REGULAR_WIDTH);
draw byLine(C, D, byblue, SOLID_LINE, REGULAR_WIDTH);
draw byLineWithName(Ob, Oe, byblack, 1, 0)(O);
draw byLine(E, F, byred, SOLID_LINE, REGULAR_WIDTH);
draw byLine(G, H, byyellow, SOLID_LINE, REGULAR_WIDTH);
draw byLineWithName(Pb, Pe, byred, 1, 0)(P);
byPointLabelRemove(Ma, Mb, Mc, Md, Na, Nb, Nc, Nd, K, L);
draw byLabelLine(1)(AB,CD,EF,GH,O,P);
}
\drawCurrentPictureInMargin
\initialIndentation{9}
\problem{I}{f}{four straight lines be proportional ($\drawUnitLine{AB} : \drawUnitLine{CD} :: \drawUnitLine{EF} : \drawUnitLine{GH}$), the similar rectilinear figures similarly described on them are also proportional.\\
E se quatro figuras retilíneas semelhantes, descritas de forma semelhante em quatro linhas retas, forem proporcionais, as linhas retas também serão proporcionais.}

\startsubproposition{Parte I.}
\begin{center}
Considere \drawUnitLine{O} um terço proporcional a \drawUnitLine{AB} e \drawUnitLine{CD}, e \drawUnitLine{P} um terço proporcional a \drawUnitLine{EF} e \drawUnitLine{GH} \byref{prop:VI.XI};

since $\drawUnitLine{AB} : \drawUnitLine{CD} :: \drawUnitLine{EF} : \drawUnitLine{GH}$ \byref{\hypref}\\
$\drawUnitLine{CD} : \drawUnitLine{O} :: \drawUnitLine{GH} : \drawUnitLine{P}$ \byref{\constref}

$\therefore$ ex \ae quali,\\
$\drawUnitLine{AB} : \drawUnitLine{O} :: \drawUnitLine{EF} : \drawUnitLine{P}$;

but $\drawPolygon{AB} : \drawPolygon{CD} :: \drawUnitLine{AB} : \drawUnitLine{O}$ \byref{prop:VI.XX},\\
and $\drawPolygon{EF} : \drawPolygon{GH} :: \drawUnitLine{EF} : \drawUnitLine{P}$;

$\therefore \drawPolygon{AB} : \drawPolygon{CD} :: \drawPolygon{EF} : \drawPolygon{GH}$ \byref{prop:V.XI}.
\end{center}

\vfill\pagebreak

\startsubproposition{Parte II.}
\begin{center}
Seja a mesma construção permanecer;

$\drawPolygon{AB} : \drawPolygon{CD} :: \drawPolygon{EF} : \drawPolygon{GH}$ \byref{\hypref},

$\therefore \drawUnitLine{AB} : \drawUnitLine{O} :: \drawUnitLine{EF} : \drawUnitLine{P}$ \byref{\constref},

and $\therefore \drawUnitLine{AB} : \drawUnitLine{CD} :: \drawUnitLine{EF} : \drawUnitLine{GH}$ \byref{prop:V.XI}.
\end{center}

\qed

\starttheorem{Prop. XXIII. Teor.}\label{prop:VI.XXIII}
\defineNewPicture[1/2]{
pair A, B, C, D, E, F, G ,H, d[];
d1 := (u, 0);
d2 := d1 scaled -2;
d3 := (1/2u, 2u);
d4 := d3 scaled -2/3;
C := (0, 0);
G := C shifted d1;
B := C shifted d2;
E := C shifted d3;
D := C shifted d4;
H := C shifted d1 shifted d4;
A := C shifted d2 shifted d4;
F := C shifted d1 shifted d3;
draw byPolygon(A,B,C,D)(byyellow);
draw byPolygon(E,F,G,C)(byblue);
draw byPolygon(C,D,H,G)(byred);
byAngleDefine(B, C, D, byred, SOLID_SECTOR);
byAngleDefine(D, C, G, byyellow, SOLID_SECTOR);
byAngleDefine(G, C, E, byblack, SOLID_SECTOR);
draw byNamedAngleResized();
draw byLine(D, C, byblack, SOLID_LINE, REGULAR_WIDTH);
draw byLine(C, E, byred, SOLID_LINE, REGULAR_WIDTH);
draw byLine(B, C, byblue, SOLID_LINE, REGULAR_WIDTH);
draw byLine(C, G, byyellow, SOLID_LINE, REGULAR_WIDTH);
draw byLabelsOnPolygon(D, A, B, C, E, F, G, H)(ALL_LABELS, 0);
}
\drawCurrentPictureInMargin
\problem[3]{O}{s}{paralelogramos equiangulares (\drawPolygon[bottom]{ABCD} e \drawPolygon[bottom]{EFGC}) estão entre si em uma proporção composta pelas razões de seus lados.}

Seja dois dos lados \drawUnitLine{BC} e \drawUnitLine{CG} com ângulos aproximadamente iguais serem colocados de modo que possam formar uma linha reta.

\begin{center}
Como $\drawAngle{BCD} + \drawAngle{DCG} = \drawTwoRightAngles$,\\
e $\drawAngle{GCE} = \drawAngle{BCD}$ \byref{\hypref},\\
$\drawAngle{GCE} + \drawAngle{DCG} = \drawTwoRightAngles$,\\
e $\therefore$ \drawUnitLine{CE} e \drawUnitLine{DC} formam uma linha reta \byref{prop:I.XIV};

complete \drawPolygon[bottom]{CDHG}.

Como $\polygonABCD\ : \polygonCDHG\ :: \drawUnitLine{BC} : \drawUnitLine{CG}$ \byref{prop:VI.I},\\
e $\polygonCDHG\ : \polygonEFGC\ :: \drawUnitLine{DC} : \drawUnitLine{CE}$ \byref{prop:VI.I},\\
\polygonABCD\ tem para \polygonEFGC\ uma proporção composta pelas razões de \drawUnitLine{BC} para \drawUnitLine{CG} e de \drawUnitLine{DC} para \drawUnitLine{CE}.
\end{center}

\qed

\starttheorem{Prop. XXIV. Teor.}\label{prop:VI.XXIV}
\defineNewPicture{
pair A, B, C, D, E, F, G, H, K, d[];
numeric s;
d1 := (3u, 0);
d2 := (u, 3u);
s := 3/5;
A := (0, 0);
B := A shifted d1;
C := A shifted d1 shifted d2;
D := A shifted d2;
G := s[A, D];
E := s[A, B];
H := s[B, C];
K := s[D, C];
F = whatever[G, H] = whatever[E, K];
draw byLine(F, H, byblack, SOLID_LINE, THIN_WIDTH);
draw byLine(E, F, byblack, SOLID_LINE, THIN_WIDTH);
draw byPolygon(G,F,K,D)(byyellow);
draw byPolygon(A,G,F)(byblue);
draw byPolygon(F,C,K)(byred);
draw byLine(G, F, byred, SOLID_LINE, REGULAR_WIDTH);
byLineDefine(A, E, byblack, SOLID_LINE, THIN_WIDTH);
byLineDefine(E, B, byblack, SOLID_LINE, THIN_WIDTH);
byLineDefine(B, H, byblack, SOLID_LINE, THIN_WIDTH);
byLineDefine(H, C, byblack, SOLID_LINE, THIN_WIDTH);
byLineDefine(D, K, byblue, SOLID_LINE, REGULAR_WIDTH);
byLineDefine(K, C, byblue, DASHED_LINE, REGULAR_WIDTH);
byLineDefine(D, G, byred, DASHED_LINE, REGULAR_WIDTH);
byLineDefine(G, A, byyellow, SOLID_LINE, REGULAR_WIDTH);
draw byNamedLineSeq(0)(GA,DG,DK,KC,HC,BH,EB,AE);
draw byLabelsOnPolygon(A, G, D, K, C, H, B, E)(ALL_LABELS, 0);
draw byLabelsOnPolygon(H, F, E)(OMIT_FIRST_LABEL+OMIT_LAST_LABEL, 0);
}
\drawCurrentPictureInMargin
\problem[4]{I}{n}{any parallelogram
(\drawFromCurrentPicture[bottom][parallelogramABCD]{
draw byNamedLine(AE,EB,BH,HC,FH,EF);
draw byNamedPolygon(GFKD,AGF,FCK);
draw byLabelsOnPolygon(D, C, B, A)(ALL_LABELS, 0);
})
the parallelograms
(\drawFromCurrentPicture[bottom][parallelogramFHCK]{
draw byNamedLine(FH,HC);
draw byNamedPolygon(FCK);
draw byLabelsOnPolygon(K, C, H, F)(ALL_LABELS, 0);
}
and
\drawFromCurrentPicture[bottom][parallelogramAEFG]{
draw byNamedLine(AE,EF);
draw byNamedPolygon(AGF);
draw byLabelsOnPolygon(G, F, E, A)(ALL_LABELS, 0);
})
que estão aproximadamente na diagonal são semelhantes ao todo e entre si.}

\begin{center}
Como \parallelogramABCD\ e \parallelogramAEFG\ possuem um ângulo comum eles são equiangulares;\\
but $\because \drawUnitLine{GF} \parallel \drawUnitLine{DK,KC}$\\ %because
\drawFromCurrentPicture[bottom]{
startGlobalRotation(180-angle(A-C));
startAutoLabeling;
draw byNamedPolygon(AGF);
stopAutoLabeling;
stopGlobalRotation;
}
and
\drawFromCurrentPicture[bottom]{
startGlobalRotation(180-angle(A-C));
startAutoLabeling;
draw byNamedPolygon(GFKD,AGF,FCK);
stopAutoLabeling;
stopGlobalRotation;
}
are similar \byref{prop:VI.IV},\\
$\therefore \drawUnitLine{GA} : \drawUnitLine{GF} :: \drawUnitLine{GA,DG} : \drawUnitLine{DK,KC}$;

e os restantes lados opostos são iguais a esses,\\
$\therefore$ \parallelogramAEFG\ e \parallelogramABCD\ têm os lados aproximadamente iguais em ângulos proporcionais e, portanto, são semelhantes.

Da mesma forma pode-se demonstrar que os paralelogramos \parallelogramABCD\ e \parallelogramFHCK\ são semelhantes.\\
Como, portanto, cada um dos paralelogramos \parallelogramAEFG\ e \parallelogramFHCK\ é semelhante a \parallelogramABCD, eles são semelhantes entre si.
\end{center}

\qed

\startproblem{Prop. XXV. prob.}\label{prop:VI.XXV}
\defineNewPicture[1/2]{
pair d[];
pair A, B, C;
A := (5/3u, 2u);
B := (0, 0);
C := (9/4u, 0);
draw byPolygon(A,B,C)(byred);
pair L, E;
L := (xpart(B), -1/2ypart(A));
E := (xpart(C), ypart(L));
draw byPolygon(B,C,E,L)(byblue);
byAngleDefine(C, B, L, byred, SOLID_SECTOR);
draw byNamedAngleResized(CBL);
pair Da, Db, Dc, Dd, De;
Da := dir(0) scaled 1/2u;
Db := dir(72) scaled 2/5u;
Dc := dir(144) scaled 1/2u;
Dd := dir(-144) scaled 3/5u;
De := dir(-72) scaled 1/2u;
d1 := (3u, 3/2u);
byPointLabelRemove(Da,Db,Dc,Dd,De);
forsuffixes i=Da,Db,Dc,Dd,De:
	i := ((i rotated 36) scaled 3/2) shifted d1;
endfor;
draw byPolygon.D(Da,Db,Dc,Dd,De)(byblue);
pair F, M;
numeric a, l;
a := 	((abs(Da-Db)*distanceToLine(Dc, Da--Db))+
	(abs(Da-Dc)*distanceToLine(Dd, Da--Dc))+
	(abs(Da-Dd)*distanceToLine(De, Da--Dd)))/2;
l := a/abs(C-E);
F := C shifted (l, 0);
M := E shifted (l, 0);
draw byPolygon(C,F,M,E)(white);
byAngleDefine(F, C, E, byyellow, SOLID_SECTOR);
draw byNamedAngleResized(FCE);
draw byLine(C, E, byred, SOLID_LINE, REGULAR_WIDTH);
byLineDefine(B, C, byyellow, SOLID_LINE, REGULAR_WIDTH);
byLineDefine(C, F, byblack, DASHED_LINE, REGULAR_WIDTH);
draw byNamedLineSeq(0)(BC, CF);
pair K, G, H;
numeric s;
s := (a/(abs(B-C)*abs(B-L)))**(1/2);
d2 := (2/3u, -3u);
K := (A scaled s) shifted d2;
G := (B scaled s) shifted d2;
H := (C scaled s) shifted d2;
draw byPolygon(K,G,H)(byyellow);
draw byLine(G, H, byblue, SOLID_LINE, REGULAR_WIDTH);
draw byLabelsOnPolygon(C, F, M, E, L, B, A)(ALL_LABELS, 0);
draw byLabelsOnPolygon(G, K, H)(ALL_LABELS, 0);
}
\drawCurrentPictureInMargin
\problem[3]{D}{escrever}{uma figura retilínea semelhante a uma determinada figura retilínea (\drawPolygon[bottom][triangleABC]{ABC}) e igual a outra (\drawPolygon[middle][polygonD]{D}).}

\begin{center}
Sobre \drawUnitLine{BC} descreva $\drawPolygon[bottom][polygonBCEL]{BCEL} = \drawPolygon[bottom][triangleABC]{ABC}$,\\
e sobre \drawUnitLine{CE} descreva
$\drawFromCurrentPicture[bottom][polygonCFME]{
startTempAngleScale(angleScale*1/2);
draw byNamedAngle(C);
startAutoLabeling;
draw byNamedPolygon(CFME);
stopAutoLabeling;
stopTempAngleScale;
} = \polygonD$,\\
tendo $\drawAngle{B} = \drawAngle{C}$ \byref{prop:I.XLV},\\
então \drawUnitLine{BC} e \drawUnitLine{CF} ficarão na mesma linha reta \byref{prop:I.XXIX,prop:I.XIV}.

Entre \drawUnitLine{BC} e \drawUnitLine{CF} encontre uma média proporcional \drawUnitLine{GH} \byref{prop:VI.XIII},\\
e sobre \drawUnitLine{GH} descreva \drawPolygon[bottom][triangleKGH]{KGH}, semelhante a \drawPolygon[bottom][triangleABC]{ABC} e semelhantemente posta.

Então $\drawPolygon[bottom][triangleKGH]{KGH} = \polygonD$.

Pois, como \drawPolygon[bottom][triangleABC]{ABC} e \drawPolygon[bottom][triangleKGH]{KGH} são semelhantes, e\\
$\drawUnitLine{BC} : \drawUnitLine{GH} :: \drawUnitLine{GH} : \drawUnitLine{CF}$ \byref{\constref},\\
$\drawPolygon[bottom][triangleABC]{ABC} : \drawPolygon[bottom][triangleKGH]{KGH} :: \drawUnitLine{BC} : \drawUnitLine{CF}$ \byref{prop:VI.XX};

mas $\polygonBCEL\ : \polygonCFME\ :: \drawUnitLine{BC} : \drawUnitLine{CF}$ \byref{prop:VI.I};\\
$\therefore \drawPolygon[bottom][triangleABC]{ABC} : \drawPolygon[bottom][triangleKGH]{KGH} :: \polygonBCEL\ : \polygonCFME$ \byref{prop:V.XI};\\
mas $\drawPolygon[bottom][triangleABC]{ABC} = \polygonBCEL$ \byref{\constref},\\
e, portanto, $\drawPolygon[bottom][triangleKGH]{KGH} = \polygonCFME$ \byref{prop:V.XIV};\\
e $\polygonCFME = \polygonD$ \byref{\constref};\\
conseqüentemente, \drawPolygon[bottom][triangleKGH]{KGH}, que é semelhante a \drawPolygon[bottom][triangleABC]{ABC}, também é igual a \polygonD.
\end{center}

\qed

\starttheorem{Prop. XXVI. Teor.}\label{prop:VI.XXVI}
\defineNewPicture{
pair A, B, C, D, E, F, G, H, K, d[];
d1 := (3u, 0);
d2 := (u, 3u);
A := (0, 0);
B := A shifted d1;
C := A shifted d1 shifted d2;
D := A shifted d2;
F := 3/4[A, C];
E = whatever[A, B] = whatever[F, F shifted d2];
G = whatever[A, D] = whatever[F, F shifted d1];
H := 2/3[G, F];
K = whatever[A, B] = whatever[H, H shifted d2];
draw byPolygon(F,H,K,E)(byred);
draw byPolygon(G,D,C,B,E,F)(byblue);
byAngleDefine(B, A, D, byyellow, SOLID_SECTOR);
draw byNamedAngleResized();
draw byLine(K, H, byyellow, SOLID_LINE, REGULAR_WIDTH);
draw byLine(G, H, byyellow, DASHED_LINE, REGULAR_WIDTH);
draw byLine(A, C, byblack, SOLID_LINE, THIN_WIDTH);
byLineDefine(A, G, byred, SOLID_LINE, REGULAR_WIDTH);
byLineDefine(G, D, byred, DASHED_LINE, REGULAR_WIDTH);
byLineDefine(A, K, byblue, SOLID_LINE, REGULAR_WIDTH);
byLineDefine(K, E, byblue, DASHED_LINE, REGULAR_WIDTH);
byLineDefine(E, B, byblack, DASHED_LINE, REGULAR_WIDTH);
draw byNamedLineSeq(0)(GD,AG,AK,KE,EB);
draw byArbitraryFigure.AHC(A..H..C, byblack, 0, 0);
byLineDefine.KAt(K, A, byblack, SOLID_LINE, THIN_WIDTH);
byLineStylize(H, G, 1, 0, 1)(KAt);
byLineDefine.AGt(A, G, byblack, SOLID_LINE, THIN_WIDTH);
byLineStylize(K, H, 0, 0, 1)(AGt);
byLineDefine.GHt(G, H, byblack, SOLID_LINE, THIN_WIDTH);
byLineStylize(A, K, 0, 1, 1)(GHt);
byLineDefine.EAt(E, A, byblack, SOLID_LINE, THIN_WIDTH);
byLineStylize(F, G, 1, 0, 1)(EAt);
draw byLabelsOnPolygon(A, G, D, C, B, E, K)(ALL_LABELS, 0);
draw byLabelsOnPolygon(K, H, A)(OMIT_FIRST_LABEL+OMIT_LAST_LABEL, 0);
draw byLabelsOnPolygon(E, F, A)(OMIT_FIRST_LABEL+OMIT_LAST_LABEL, 0);
}
\drawCurrentPictureInMargin
\problem{I}{f}{similar and similarly posited parallelograms
(\drawFromCurrentPicture[bottom][parallelogramAEFG]{
draw byNamedPolygon(FHKE);
draw byNamedLine(KAt,AGt,GHt);
draw byLabelsOnPolygon(A, G, F, E)(ALL_LABELS, 0);
}
and
\drawFromCurrentPicture[bottom][parallelogramABCD]{
draw byNamedPolygon(GDCBEF);
draw byNamedLine(EAt,AC);
draw byNamedLineFull(E, F, 0, 1,  0, 1)(AGt);
draw byLabelsOnPolygon(A, D, C, B)(ALL_LABELS, 0);
})
têm um ângulo comum, eles têm aproximadamente a mesma diagonal.}

\begin{center}
For, if possible, let
\drawFromCurrentPicture[bottom][lineAHC]{
startGlobalRotation(180-angle(A-C));
draw byNamedArbitraryFigure(AHC);
draw byLabelsOnPolygon(A, H, C, noPoint)(ALL_LABELS, 0);
stopGlobalRotation;
}
seja a diagonal de \parallelogramABCD\ e trace $\drawUnitLine{KH} \parallel \drawUnitLine{AG}$ \byref{prop:I.XXXI}.

Como \drawLine[bottom][parallelogramAKHG]{AG,GH,KH,AK} e \parallelogramABCD\ têm aproximadamente a mesma diagonal \lineAHC e têm \drawAngle{A} comum, eles são \byref{prop:VI.XXIV} semelhantes;

$\therefore \drawSizedLine{AG} : \drawSizedLine{AK} :: \drawSizedLine{AG,GD} : \drawSizedLine{AK,KE,EB}$;\\
but $\drawSizedLine{AG} : \drawSizedLine{AK,KE} :: \drawSizedLine{AG,GD} : \drawSizedLine{AK,KE,EB}$ \byref{\hypref},\\
$\therefore \drawSizedLine{AG} : \drawSizedLine{AK} :: \drawSizedLine{AG} : \drawSizedLine{AK,KE}$,\\
e $\therefore \drawSizedLine{AK} = \drawSizedLine{AK,KE}$ \byref{prop:V.IX}, o que é um absurdo.

$\therefore$ \lineAHC\ não é a diagonal de \parallelogramABCD\ da mesma maneira que pode ser demonstrado que nenhuma outra linha é exceto \drawUnitLine{AC}.
\end{center}

\qed

\starttheorem{Prop. XXVII. Teor.}\label{prop:VI.XXVII}
\defineNewPicture[1/2]{
pair A, B, C, D, E, F, G, H, K;
numeric l, s;
l := 5u;
s := 1/5;
A := (0, 0);
B := (xpart(A), -l);
C := 1/2[A, B];
D := s[A, B];
E := (-1/2l, ypart(C));
F := (xpart(E), ypart(A));
G := (-s*l, ypart(D));
H := (xpart(G), ypart(B));
K = whatever[C, E] = whatever[G, H];
draw byPolygon(A,D,G,K,E,F)(byred);
draw byPolygon(C,D,G,K)(byblue);
draw byPolygon(B,C,K,H)(byyellow);
draw byLine(A, D, byyellow, SOLID_LINE, REGULAR_WIDTH);
draw byLine(D, C, byred, SOLID_LINE, REGULAR_WIDTH);
draw byLine(C, B, byblue, SOLID_LINE, REGULAR_WIDTH);
draw byLabelsOnPolygon(E, F, A, D, C, B, H)(ALL_LABELS, 0);
draw byLabelsOnPolygon(H, G, D)(OMIT_FIRST_LABEL+OMIT_LAST_LABEL, 0);
}
\drawCurrentPictureInMargin
\problem{D}{e}{todos os retângulos compreendido por segmentos de uma determinada reta, o maior é o quadrado descrito na metade da reta.}

\begin{center}
Seja \drawSizedLine{AD,DC,CB} a linha fornecida,\\
\drawSizedLine{AD} and \drawSizedLine{DC,CB} unequal segments,\\
e segmentos iguais \drawSizedLine{AD,DC} e \drawSizedLine{CB};

then $\drawPolygon[bottom]{ADGKEF,CDGK} > \drawPolygon[bottom]{BCKH,CDGK}$.
\end{center}

Pois já foi demonstrado \byref{prop:II.V}, que o quadrado da metade da reta é igual ao retângulo compreendido por quaisquer segmentos desiguais junto com o quadrado da parte intermediária entre o ponto médio e o ponto de seção desigual. O quadrado descrito na metade da linha excede, portanto, o retângulo compreendido por quaisquer segmentos desiguais da linha.

\qed

\startproblem{Prop. XXVIII. prob.}\label{prop:VI.XXVIII}
\defineNewPicture{
pair A, B, C, D, E, F, G, H, d;
path a;
numeric r;
A := (0, 0);
B := (4u, 0);
C := 1/2[A, B];
D := C shifted (0, 3/2u);
r := 2u;
a := (fullcircle scaled (2r)) shifted D;
E := a intersectionpoint (A--C);
F := a intersectionpoint (D--2[D, C]);
byLineDefine(D, E, byyellow, SOLID_LINE, REGULAR_WIDTH);
byLineDefine(D, C, byred, SOLID_LINE, REGULAR_WIDTH);
byLineDefine(C, F, byblack, DASHED_LINE, REGULAR_WIDTH);
draw byNamedLineSeq(0)(CF,DC,DE);
byLineDefine(A, E, byred, DASHED_LINE, REGULAR_WIDTH);
byLineDefine(E, C, byblue, SOLID_LINE, REGULAR_WIDTH);
byLineDefine(C, B, byblue, DASHED_LINE, REGULAR_WIDTH);
draw byNamedLineSeq(0)(AE, EC, CB);
draw byArcBE.a(D, -1/2, -4 + 1/2, r, byred, 0, 0, 0, 0);
d := (0, 2u);
G := A shifted d;
H := C shifted d;
byLineDefine.G(G, H, byyellow, DASHED_LINE, REGULAR_WIDTH);
lineUseLineLabel.G := true;
draw byLabelsOnCircle(F)(a);
draw byLabelsOnPolygon(E, D, C, B)(OMIT_FIRST_LABEL+OMIT_LAST_LABEL, 0);
draw byLabelsOnPolygon(A, B, noPoint)(ALL_LABELS, 0);
draw byLabelPoint(E, angle(E-D) + 45, 2);
}
\drawCurrentPictureInMargin
\problem{D}{ividir}{uma determinada reta (\drawSizedLine{AE,EC,CB}) de modo que o retângulo compreendido por seus segmentos sejam iguais a uma determinada área, não excedendo o quadrado da metade da reta.}

\begin{center}
Seja a área fornecida $=\drawSizedLine{G}^2$.

Bissete \drawSizedLine{AE,EC,CB}, ou faça $\drawSizedLine{AE,EC} = \drawSizedLine{CB}$;\\
e, se $\drawSizedLine{AE,EC}^2 = \drawSizedLine{G}^2$,\\
problema está resolvido.

Mas, se $\drawSizedLine{AE,EC}^2 \neq \drawSizedLine{G}^2$,\\
então deve ser $\drawSizedLine{AE,EC} > \drawSizedLine{G}$ \byref{\hypref}.

Trace $\drawSizedLine{DC} \perp \drawSizedLine{AE,EC} = \drawSizedLine{G}$;\\
faça $\drawSizedLine{DC,CF} = \drawSizedLine{AE,EC} \mbox{ ou } \drawSizedLine{CB}$;\\
com \drawSizedLine{DC,CF} como raio descreva um círculo cortando a linha fornecida; trace \drawSizedLine{DE}.

Então $\drawSizedLine{AE} \times \drawSizedLine{EC,CB} + \drawSizedLine{EC}^2 = \drawSizedLine{AE,EC}^2$ \byref{prop:II.V} $= \drawSizedLine{DE}^2$.

Mas $\drawSizedLine{DE}^2 = \drawSizedLine{DC}^2 + \drawSizedLine{EC}^2$ \byref{prop:I.XLVII};

$\therefore \drawSizedLine{AE} \times \drawSizedLine{EC,CB} + \drawSizedLine{EC}^2 = \drawSizedLine{DC}^2 + \drawSizedLine{EC}^2$,\\
retire $\drawSizedLine{EC}^2$ de ambos,\\
e $\drawSizedLine{AE} \times \drawSizedLine{EC,CB} = \drawSizedLine{DC}^2$.

Mas $\drawSizedLine{DC} = \drawSizedLine{G}$ \byref{\constref},\\
e $\therefore$ \drawSizedLine{AE,EC,CB} está tão dividido que $\drawSizedLine{AE} \times \drawSizedLine{EC,CB} = \drawSizedLine{G}^2$.
\end{center}

\qed

\startproblem{Prop. XXIX. prob.}\label{prop:VI.XXIX}
\defineNewPicture[1/2]{
pair A, B, C, D, E, F, G, H, d;
d := (0, -1/2u);
G := (0, 0) shifted d;
H := (3/2u, 0) shifted d;
A := (0, 0);
B := (2u, ypart(A));
C := 1/2[A, B];
D := (xpart(B), abs(G-H));
E := C shifted (-abs(C-D), 0);
F := C shifted (abs(C-D), 0);
draw byLine.G(H, G, byblack, SOLID_LINE, REGULAR_WIDTH);
draw byLineFull(C, D, byred, 0, 0)(B, D, 1, 0, 0);
draw byLine(B, D, byred, DASHED_LINE, REGULAR_WIDTH);
byLineDefine(E, A, byyellow, DASHED_LINE, REGULAR_WIDTH);
byLineDefine(A, C, byblue, SOLID_LINE, REGULAR_WIDTH);
byLineDefine(C, B, byblue, DASHED_LINE, REGULAR_WIDTH);
byLineDefine(B, F, byyellow, SOLID_LINE, REGULAR_WIDTH);
draw byNamedLineSeq(1)(EA, AC, CB, BF);
draw byArcBE.a(C, 0, 4, abs(C-D), byred, 0, 0, 0, 0);
draw byLabelsOnCircle(D)(a);
draw byLabelsOnPolygon(F, B, C, A, noPoint)(ALL_LABELS, 0);
draw byLabelLine(0)(G);
}
\drawCurrentPictureInMargin
\problem{P}{roduzir}{uma dada reta (\drawSizedLine{AC,CB}), de modo que o retângulo compreendido pelos segmentos entre as extremidades da dada reta e o ponto para o qual ela é produzida, possam ser iguais a uma determinada área, ou seja, igual ao quadrado de \drawSizedLine{G}.}

\begin{center}
Faça $\drawSizedLine{AC} = \drawSizedLine{CB}$,\\
e trace $\drawSizedLine{BD} \perp \drawSizedLine{CB} = \drawSizedLine{G}$;\\
draw \drawSizedLine{CD};\\
e com o raio \drawSizedLine{CD}, descreva um círculo que corresponde ao \drawSizedLine{AC,CB} produzido.

Então $\drawSizedLine{AC,CB,BF} \times \drawSizedLine{BF} + \drawSizedLine{CB}^2 = \drawSizedLine{AC,CB}^2$ \byref{prop:II.VI} $= \drawSizedLine{CD}^2$.

Mas $\drawSizedLine{CD}^2 = \drawSizedLine{BD}^2 + \drawSizedLine{CB}^2$ \byref{prop:I.XLVII}

$\therefore \drawSizedLine{AC,CB,BF} \times \drawSizedLine{BF} + \drawSizedLine{CB}^2 = \drawSizedLine{BD}^2 + \drawSizedLine{CB}^2$,\\
retire $\drawSizedLine{CB}^2$ de ambos,\\
e $\therefore \drawSizedLine{AC,CB,BF} \times \drawSizedLine{BF}= \drawSizedLine{BD}^2$\\
mas $\drawSizedLine{BD} = \drawSizedLine{G}$,\\
$\therefore \drawSizedLine{BD}^2 = \mbox{ a área dada.}$
\end{center}

\qed

\startproblem{Prop. XXX. prob.}\label{prop:VI.XXX}
\defineNewPicture{
pair A, B, C, D, E, F, G, H;
numeric w;
w := 3u;
A := (0, 0);
B := (w, 0);
C := (0, w);
H := (w, w);
G := 1/2[A, C] shifted (0, -abs((1/2[A, C]) - B));
D := G shifted (abs(G-A), 0);
E = whatever[D, D shifted (0, 1)] = whatever[A, B];
F = whatever[C, H] = whatever[D, E];
draw byPolygon(A,C,F,E)(byyellow);
draw byPolygon(E,F,H,B)(byblue);
byLineDefine(A, E, byred, SOLID_LINE, REGULAR_WIDTH);
byLineDefine(E, B, byred, DASHED_LINE, REGULAR_WIDTH);
draw byNamedLineSeq(1)(AE, EB);
byLineDefine(C, A, byblue, SOLID_LINE, REGULAR_WIDTH);
byLineDefine(A, G, byblue, DASHED_LINE, REGULAR_WIDTH);
byLineDefine(G, D, byyellow, SOLID_LINE, REGULAR_WIDTH);
byLineDefine(D, E, byblack, SOLID_LINE, REGULAR_WIDTH);
draw byNamedLineSeq(0)(CA, DE, GD, AG);
byLineDefine.AGt(A, G, byblack, SOLID_LINE, THIN_WIDTH);
byLineStylize(A, D, 1, 0, -1)(AGt);
byLineDefine.GDt(G, D, byblack, SOLID_LINE, THIN_WIDTH);
byLineStylize(A, E, 0, 0, -1)(GDt);
byLineDefine.DEt(D, E, byblack, SOLID_LINE, THIN_WIDTH);
byLineStylize(G, E, 0, 1, -1)(DEt);
draw byLabelsOnPolygon(A, C, F, H, B, E, D, G)(ALL_LABELS, 0);
}
\drawCurrentPictureInMargin
\problem{P}{ara}{cortar uma determinada linha reta finita (\drawProportionalLine{AE,EB}) em extrema e média razão.}

\begin{center}
Em \drawProportionalLine{AE,EB} descreva o quadrado \drawPolygon[middle][squareABHC]{ACFE,EFHB} \byref{prop:I.XLVI};\\
e prolongue \drawProportionalLine{CA}, de modo que $\drawProportionalLine{CA,AG} \times \drawProportionalLine{AG} = \drawProportionalLine{AE,EB}^2$ \byref{prop:VI.XXIX};\\
tome $\drawProportionalLine{AE} = \drawProportionalLine{AG}$,\\
e trace $\drawProportionalLine{DE} \parallel \drawProportionalLine{CA,AG}$,\\
encontrando $\drawProportionalLine{GD} \parallel \drawProportionalLine{AE,EB}$ \byref{prop:I.XXXI}.

Então $\drawFromCurrentPicture[middle][rectangleCFDG]{
draw byNamedPolygon(ACFE);
draw byNamedLine(AGt,GDt,DEt);
draw byLabelsOnPolygon(G, C, F, D)(ALL_LABELS, 0);
}
= \drawProportionalLine{CA,AG} \times \drawProportionalLine{AG}$, e é $\therefore\ = \squareABHC$;\\
e se de ambos estes iguais for retirada a parte comum \drawPolygon{ACFE},\\
\drawLine{DE,GD,AG,AE}, que é o quadrado de \drawProportionalLine{AE},\\
será igual a \drawPolygon{EFHB}, que é igual a $\drawProportionalLine{AE,EB} \times \drawProportionalLine{EB}$;\\
isto é, $\drawProportionalLine{AE}^2 = \drawProportionalLine{AE,EB} \times \drawProportionalLine{EB}$;\\
$\therefore \drawProportionalLine{AE,EB} : \drawProportionalLine{AE} :: \drawProportionalLine{AE} : \drawProportionalLine{EB}$,\\
e \drawProportionalLine{AE,EB} é dividido em extrema e média razão \byref{def:VI.III}.
\end{center}

\qed

\starttheorem{Prop. XXXI. Teor.}\label{prop:VI.XXXI}
\defineNewPicture[1/2]{
pair A, B, C, D, E, F, G, H, K, L;
numeric a, r, l[];
a := -125;
A := (0, 0);
B := A shifted (dir(a)*2u);
C = whatever[A, A shifted dir(a+90)] = whatever[B, B shifted dir(0)];
D = whatever[A, A shifted dir(-90)] = whatever[B, C];
l1 := abs(A-B);
l2 := abs(B-C);
l3 := abs(C-A);
r := 1/4;
E := A shifted (dir(a-90)*l1*r);
F := B shifted (dir(a-90)*l1*r);
G := B shifted (dir(-90)*l2*r);
H := C shifted (dir(-90)*l2*r);
K := C shifted (dir(a + 180)*l3*r);
L := A shifted (dir(a + 180)*l3*r);
draw byPolygon.AB(A,B,F,E)(byblue);
draw byPolygon.BC(B,C,H,G)(byred);
draw byPolygon.CA(C,A,L,K)(byyellow);
draw byLine(A, D, byblack, SOLID_LINE, REGULAR_WIDTH);
byLineDefine(A, B, byyellow, SOLID_LINE, REGULAR_WIDTH);
byLineDefine(B, D, byblue, DASHED_LINE, REGULAR_WIDTH);
byLineDefine(D, C, byblue, SOLID_LINE, REGULAR_WIDTH);
byLineDefine(C, A, byred, SOLID_LINE, REGULAR_WIDTH);
draw byNamedLineSeq(1)(AB,BD,DC,CA);
byPointLabelRemove(H, G, L, K, F, E);
draw byLabelsOnPolygon(B, F, E, A, L, K, C, H, G)(ALL_LABELS, 0);
draw byLabelsOnPolygon(A, D, C)(OMIT_FIRST_LABEL+OMIT_LAST_LABEL, 0);
}
\drawCurrentPictureInMargin
\problem[3]{S}{e}{quaisquer figuras retilíneas semelhantes forem descritas de forma semelhante nos lados de um triângulo retângulo (\drawLine[bottom]{CA,DC,BD,AB}), a figura descrita no lado (\drawProportionalLine{BD,DC}) que subentende o ângulo reto é igual à soma das figuras nos outros lados.}

\begin{center}
Do ângulo reto, trace \drawProportionalLine{AD} perpendicular a \drawProportionalLine{BD,DC};\\
então $\drawProportionalLine{BD,DC} : \drawProportionalLine{CA} :: \drawProportionalLine{CA} : \drawProportionalLine{DC}$ \byref{prop:VI.VIII}.

$\therefore
\drawFromCurrentPicture[bottom][figBC]{
startGlobalRotation(180-angle(B-C));
startAutoLabeling;
draw byNamedPolygon(BC);
stopAutoLabeling;
stopGlobalRotation;
} :
\drawFromCurrentPicture[bottom][figCA]{
startGlobalRotation(-angle(C-A));
startAutoLabeling;
draw byNamedPolygon(CA);
stopAutoLabeling;
stopGlobalRotation;
} :: \drawProportionalLine{BD,DC} : \drawProportionalLine{DC}$ \byref{prop:VI.XX}.\\
mas $\figBC\ :
\drawFromCurrentPicture[bottom][figAB]{
startGlobalRotation(-angle(A-B));
startAutoLabeling;
draw byNamedPolygon(AB);
stopAutoLabeling;
stopGlobalRotation;
} :: \drawProportionalLine{BD,DC} : \drawProportionalLine{BD}$ \byref{prop:VI.XX}.

Daí $\drawProportionalLine{DC} + \drawProportionalLine{BD} : \drawProportionalLine{BD,DC} :: \figAB\ + \figCA\ : \figBC$;\\
mas $\drawProportionalLine{DC} + \drawProportionalLine{BD} = \drawProportionalLine{BD,DC}$;\\
e, portanto, $\figAB\ + \figCA\ = \figBC$.
\end{center}

\qed

\starttheorem{Prop. XXXII. Teor.}\label{prop:VI.XXXII}
\defineNewPicture[1/2]{
pair A, B, C, D, E;
numeric s;
B := (0, 0);
C := (5/2u, 0);
A := (u, 2u);
s := 3/5;
D := (A scaled s) shifted C;
E := (C scaled s) shifted C;
byAngleDefine(C, A, B, byyellow, SOLID_SECTOR);
byAngleDefine(A, B, C, byred, SOLID_SECTOR);
byAngleDefine(B, C, A, byblue, SOLID_SECTOR);
byAngleDefine(E, D, C, byblack, SOLID_SECTOR);
byAngleDefine(D, C, E, byblack, ARC_SECTOR);
byAngleDefine(A, C, D, byyellow, SOLID_SECTOR);
draw byNamedAngleResized();
byLineDefine(A, B, byblue, SOLID_LINE, REGULAR_WIDTH);
byLineDefine(B, C, byyellow, SOLID_LINE, REGULAR_WIDTH);
byLineDefine(C, E, byyellow, DASHED_LINE, REGULAR_WIDTH);
byLineDefine(E, D, byred, DASHED_LINE, REGULAR_WIDTH);
byLineDefine(D, C, byblue, DASHED_LINE, REGULAR_WIDTH);
byLineDefine(C, A, byred, SOLID_LINE, REGULAR_WIDTH);
draw byNamedLineSeq(0)(ED,DC,noLine,CA,AB,BC,CE);
byLineDefine.ABt(A, B, byblack, SOLID_LINE, THIN_WIDTH);
byLineDefine.BCt(B, C, byblack, SOLID_LINE, THIN_WIDTH);
byLineDefine.CAt(C, A, byblack, SOLID_LINE, THIN_WIDTH);
byLineDefine.DCt(D, C, byblack, SOLID_LINE, THIN_WIDTH);
byLineDefine.CEt(C, E, byblack, SOLID_LINE, THIN_WIDTH);
byLineDefine.EDt(E, D, byblack, SOLID_LINE, THIN_WIDTH);
draw byLabelsOnPolygon(E, C, B, A, C, D)(ALL_LABELS, 0);
}
\drawCurrentPictureInMargin
\problem{I}{f}{two triangles
(\drawFromCurrentPicture[bottom]{
startTempAngleScale(angleScale*3/5);
draw byNamedAngle(A, B, BCA);
draw byNamedLineSeq(0)(CAt,BCt,ABt);
draw byLabelsOnPolygon(C, B, A)(ALL_LABELS, 0);
stopTempAngleScale;
}
and
\drawFromCurrentPicture[bottom]{
startTempAngleScale(angleScale*3/5);
draw byNamedAngle(D, DCE);
draw byNamedLineSeq(0)(EDt,CEt,DCt);
draw byLabelsOnPolygon(C, D, E)(ALL_LABELS, 0);
stopTempAngleScale;
}),
ter dois lados proporcionais ($\drawUnitLine{AB} : \drawUnitLine{CA} :: \drawUnitLine{DC} : \drawUnitLine{ED}$) e ser colocados em um ângulo que os lados homólogos fiquem paralelos, os lados restantes (\drawUnitLine{BC} e \drawUnitLine{CE}) formam uma linha reta.}

\begin{center}
Como $\drawUnitLine{AB} \parallel \drawUnitLine{DC}$,\\
$\drawAngle{A} = \drawAngle{ACD}$ \byref{prop:I.XXIX};\\
e também, como $\drawUnitLine{CA} \parallel \drawUnitLine{ED}$,\\
$\drawAngle{ACD} = \drawAngle{D}$ \byref{prop:I.XXIX};\\
$\therefore \drawAngle{A} = \drawAngle{D}$;\\
e, como $\drawUnitLine{AB} : \drawUnitLine{CA} :: \drawUnitLine{DC} : \drawUnitLine{ED}$ \byref{\hypref},\\
os triângulos são equiângulos \byref{prop:VI.VI};

$\therefore \drawAngle{B} = \drawAngle{DCE}$;\\
mas $\drawAngle{A} = \drawAngle{ACD}$;

$\therefore \drawAngle{BCA} + \drawAngle{ACD} + \drawAngle{DCE} = \drawAngle{BCA} + \drawAngle{A} + \drawAngle{B} = \drawTwoRightAngles$ \byref{prop:I.XXXII},\\
e $\therefore$ \drawUnitLine{BC} e \drawUnitLine{CE} estão na mesma linha reta \byref{prop:I.XIV}.
\end{center}

\qed

\starttheorem{Prop. XXXIII. Teor.}\label{prop:VI.XXXIII}
\defineNewPicture[1/2]{
pair A, B, C, D, E, F, G, H, K, L, M, N;
numeric r, a[], ba, aa, q;
path c[];
q := 8/360;
r := 9/4u;
aa := 125;
ba := 205;
a1 := 30;
a2 := 35;
G := (0, 0);
A := (dir(aa)*r) shifted G;
B := (dir(ba)*r) shifted G;
C := (dir(ba + a1)*r) shifted G;
K := (dir(ba + 2a1)*r) shifted G;
L := (dir(ba + 3a1)*r) shifted G;
byAngleDefine(B, A, C, byyellow, SOLID_SECTOR);
byAngleDefine.BC(B, G, C, byblack, SOLID_SECTOR);
byAngleDefine.CK(C, G, K, byred, SOLID_SECTOR);
byAngleDefine.KL(K, G, L, byblue, SOLID_SECTOR);
draw byNamedAngleResized(BAC, BC, CK, KL);
draw byLine(A, B, byblack, SOLID_LINE, THIN_WIDTH);
draw byLine(A, C, byblack, SOLID_LINE, THIN_WIDTH);
draw byLine(G, C, byblack, SOLID_LINE, THIN_WIDTH);
draw byLine(G, K, byblack, SOLID_LINE, THIN_WIDTH);
byLineDefine(G, B, byblack, SOLID_LINE, THIN_WIDTH);
byLineDefine(G, L, byblack, SOLID_LINE, THIN_WIDTH);
draw byNamedLineSeq(0)(GB,GL);
byArcDefine.LB(G, L, B, r, byred, 0, 0, 0, 0);
byArcDefine.BC(G, B, C, r, byblack, 0, 0, 0, 0);
byArcDefine.CK(G, C, K, r, byred, 0, 0, 0, 0);
byArcDefine.KL(G, K, L, r, byblue, 0, 0, 0, 0);
draw byNamedArcSeq(0)(LB, BC, CK, KL);
byCircleDefineR(G, r, byred, 0, 0, 0);
H := (0, -1/2u - 2r);
D := (dir(aa)*r) shifted H;
E := (dir(ba)*r) shifted H;
F := (dir(ba + a2)*r) shifted H;
M := (dir(ba + 2a2)*r) shifted H;
N := (dir(ba + 3a2)*r) shifted H;
byAngleDefine(E, D, F, byyellow, ARC_SECTOR);
byAngleDefine.EF(E, H, F, byblack, ARC_SECTOR);
byAngleDefine.FM(F, H, M, byred, ARC_SECTOR);
byAngleDefine.MN(M, H, N, byblue, ARC_SECTOR);
draw byNamedAngleResized(EDF, EF, FM, MN);
draw byLine(D, E, byblack, SOLID_LINE, THIN_WIDTH);
draw byLine(D, F, byblack, SOLID_LINE, THIN_WIDTH);
draw byLine(H, F, byblack, SOLID_LINE, THIN_WIDTH);
draw byLine(H, M, byblack, SOLID_LINE, THIN_WIDTH);
byLineDefine(H, E, byblack, SOLID_LINE, THIN_WIDTH);
byLineDefine(H, N, byblack, SOLID_LINE, THIN_WIDTH);
draw byNamedLineSeq(0)(HE,HN);
byArcDefine.NE(H, N, E, r, byblue, 0, 0, 0, 0);
byArcDefine.EF(H, E, F, r, byyellow, 1, 0, 0, 0);
byArcDefine.FM(H, F, M, r, byred, 1, 0, 0, 0);
byArcDefine.MN(H, M, N, r, byblue, 1, 0, 0, 0);
draw byNamedArcSeq(0)(NE, EF, FM, MN);
byCircleDefineR(H, r, byblue, 0, 0, 0);
draw byLabelsOnCircle(B, C, K, L)(G);
draw byLabelsOnCircle(E, F, M, N)(H);
draw byLabelsOnPolygon(B, G, L)(OMIT_FIRST_LABEL+OMIT_LAST_LABEL, 0);
draw byLabelsOnPolygon(E, H, N)(OMIT_FIRST_LABEL+OMIT_LAST_LABEL, 0);
}
\drawCurrentPictureInMargin
\problem[3]{C}{írculos}{desiguais (\drawCircle[middle][1/4]{G}, \drawCircle[middle][1/4]{H}), ângulos, seja no centro ou na circunferência, estão na mesma proporção entre si que os arcos nos quais estão ($\drawAngle{BC} : \drawAngle{EF}
\drawFromCurrentPicture[middle][arcBC]{
startGlobalRotation(180-angle(B-C));
startAutoLabeling;
draw byNamedArc(BC);
stopAutoLabeling;
stopGlobalRotation;
} : \drawFromCurrentPicture[middle][arcEF]{
startGlobalRotation(180-angle(E-F));
startAutoLabeling;
draw byNamedArc(EF);
stopAutoLabeling;
stopGlobalRotation;
}$); so also are sectors.}

Pegue a circunferência de \circleG\ qualquer número de arcos
\drawFromCurrentPicture[middle][arcCK]{
startGlobalRotation(180-angle(C-K));
startAutoLabeling;
draw byNamedArc(CK);
stopAutoLabeling;
stopGlobalRotation;
},
\drawFromCurrentPicture[middle][arcKL]{
startGlobalRotation(180-angle(K-L));
startAutoLabeling;
draw byNamedArc(KL);
stopAutoLabeling;
stopGlobalRotation;
}, \&c.\ cada $= \arcBC$, e também na circunferência de \circleH\ pegue qualquer número de arcos
\drawFromCurrentPicture[middle][arcFM]{
startGlobalRotation(180-angle(F-M));
startAutoLabeling;
draw byNamedArc(FM);
stopAutoLabeling;
stopGlobalRotation;
},
\drawFromCurrentPicture[middle][arcMN]{
startGlobalRotation(180-angle(M-N));
startAutoLabeling;
draw byNamedArc(MN);
stopAutoLabeling;
stopGlobalRotation;
}, \&c.\ cada $= \arcEF$, trace os raios para as extremidades dos arcos iguais.

Como os arcos \arcBC, \arcCK, \arcKL, \&c.\ são todos iguais, os ângulos \drawAngle{BC}, \drawAngle{CK}, \drawAngle{KL}, \&c.\ também são iguais \byref{prop:III.XXVII}; $\therefore$ \drawAngle{BC,CK,KL} é o mesmo múltiplo de \drawAngle{BC} que o arco \drawFromCurrentPicture[middle][arcBL]{
startGlobalRotation(180-angle(B-L));
startAutoLabeling;
draw byNamedArcSeq(0)(BC,CK,KL);
stopAutoLabeling;
stopGlobalRotation;
} é de \arcBC; e da mesma forma \drawAngle{EF,FM,MN} é o mesmo múltiplo de \drawAngle{EF}, que o arco \drawFromCurrentPicture[middle][arcEN]{
startGlobalRotation(180-angle(E-N));
startAutoLabeling;
draw byNamedArcSeq(0)(EF,FM,MN);
stopAutoLabeling;
stopGlobalRotation;
} é do arco \arcEF.

\begin{center}
Então é evidente \byref{prop:III.XXVII},\\
se \drawAngle{BC,CK,KL} (ou se $m$ vezes \drawAngle{BC}) $>, =, < \drawAngle{EF,FM,MN}$ (ou $n$ vezes \drawAngle{EF})\\
então \arcBL (ou $m$ vezes \arcBC) $>, =, < \arcEN$ (ou $n$ vezes \arcEF);
\end{center}

$\therefore \drawAngle{BC} : \drawAngle{EF} :: \arcBC\ : \arcEF$ \byref{def:V.V}, ou os ângulos no centro são como os arcos nos quais eles se apoiam; mas os ângulos na circunferência sendo metades dos ângulos no centro \byref{prop:III.XX} estão na mesma proporção \byref{prop:V.XV} e, portanto, são como os arcos nos quais eles se apoiam.

É evidente que os setores em círculos iguais e em arcos iguais são iguais \byref{prop:I.IV,prop:III.XXIV,prop:III.XXVII,def:III.IX}. Assim, se os setores substituirem os ângulos na demonstração acima, a segunda parte da proporção será estabelecida, isto é, em círculos iguais os setores têm a mesma proporção entre si que os arcos sobre os quais se apoiam.

\qed

\starttheoremAZ{Prop. A. theor.}\label{prop:VI.A} % Proposition with letters
\defineNewPicture{
pair A, B, C, D, E, F;
A := (0, 0);
B := (2u, 0);
C := (3u, 7/4u);
F := 4/3[A, C];
D = whatever[A, B] = whatever[C, C shifted 1/2[unitvector(F-C), unitvector(B-C)]];
E = whatever[A, C] = whatever[B, B shifted (C-D)];
byAngleDefine(C, E, B, byyellow, SOLID_SECTOR);
byAngleDefine(E, B, C, byblue, ARC_SECTOR);
byAngleDefine(C, B, D, byred, SOLID_SECTOR);
byAngleDefine(B, D, C, byyellow, ARC_SECTOR);
byAngleDefine(B, C, E, byblack, ARC_SECTOR);
byAngleDefine(D, C, B, byblue, SOLID_SECTOR);
byAngleDefine(F, C, D, byblack, SOLID_SECTOR);
draw byNamedAngleResized();
draw byLine(B, E, byred, SOLID_LINE, REGULAR_WIDTH);
draw byLine(B, C, byyellow, SOLID_LINE, REGULAR_WIDTH);
byLineDefine(A, B, byblue, SOLID_LINE, REGULAR_WIDTH);
byLineDefine(B, D, byblue, DASHED_LINE, REGULAR_WIDTH);
byLineDefine(D, C, byred, DASHED_LINE, REGULAR_WIDTH);
byLineDefine(A, E, byblack, SOLID_LINE, REGULAR_WIDTH);
byLineDefine(E, C, byblack, DASHED_LINE, REGULAR_WIDTH);
byLineDefine(C, F, byyellow, DASHED_LINE, REGULAR_WIDTH);
draw byNamedLineSeq(0)(noLine,DC,BD,AB,AE,EC,CF);
draw byLabelsOnPolygon(C, D, B, A, E)(OMIT_FIRST_LABEL+OMIT_LAST_LABEL, 0);
draw byLabelsOnPolygon(E, C, noPoint)(ALL_LABELS, 0);
}
\drawCurrentPictureInMargin
\problem[2]{S}{e}{a reta (\drawSizedLine{DC}) que bissecta um ângulo externo \drawAngle{DCB,FCD} do triângulo \drawFromCurrentPicture[bottom]{
startTempScale(1/2);
startAutoLabeling;
draw byNamedLineSeq(0)(AE,EC,BC,AB);
stopAutoLabeling;
stopTempScale;
} encontrar o lado oposto (\drawSizedLine{AB}) produzido, esse lado produzido por inteiro (\drawSizedLine{AB,BD}) e seu segmento externo (\drawSizedLine{BD}) serão proporcionais aos lados (\drawSizedLine{AE,EC} e \drawSizedLine{BC}) que contêm o ângulo adjacente ao ângulo externo bissectado.}

\begin{center}
Pois, se \drawSizedLine{BE} for traçada $\parallel \drawSizedLine{DC}$,\\
$\begin{aligned}
		\mbox{então } \drawAngle{DCB} &= \drawAngle{EBC} \mbox{, \byref{prop:I.XXIX};}\\
		& = \drawAngle{FCD} \mbox{, \byref{\hypref},}\\
		& = \drawAngle{E} \mbox{, \byref{prop:I.XXIX};}\\
\end{aligned}$\\
e, portanto, $\drawSizedLine{EC} = \drawSizedLine{BC}$, \byref{prop:I.VI},\\
e $\drawSizedLine{AE,EC} : \drawSizedLine{BC} :: \drawSizedLine{AE,EC} : \drawSizedLine{EC}$ \byref{prop:V.VII};\\
mas também $\drawSizedLine{AB,BD} : \drawSizedLine{BD} :: \drawSizedLine{AE,EC} : \drawSizedLine{EC}$ \byref{prop:VI.II};\\
e, portanto, $\drawSizedLine{AB,BD} : \drawSizedLine{BD} :: \drawSizedLine{AE,EC} : \drawSizedLine{BC}$ \byref{prop:V.XI}.
\end{center}

\qed
\starttheoremAZ{Prop. B. theor.}\label{prop:VI.B} % Proposition with letters
\defineNewPicture{
pair A, B, C, D, E, O;
numeric r;
path c;
r := 7/4u;
O := (0, 0);
c := (fullcircle scaled 2r) shifted O;
A := (dir(75)*r) shifted O;
B := (dir(180 + 10)*r) shifted O;
C := (dir(-10)*r) shifted O;
D = whatever[B, C] = whatever[A, A shifted 1/2[unitvector(B-A), unitvector(C-A)]];
E := (subpath (0, -4) of c) intersectionpoint (A--4[A, D]);
byAngleDefine(A, B, C, byyellow, SOLID_SECTOR);
byAngleDefine(A, E, C, byblack, SOLID_SECTOR);
byAngleDefine(C, A, E, byblue, SOLID_SECTOR);
byAngleDefine(E, A, B, byred, SOLID_SECTOR);
draw byNamedAngleResized();
draw byLine(A, B, byblue, SOLID_LINE, REGULAR_WIDTH);
draw byLine(B, D, byred, DASHED_LINE, REGULAR_WIDTH);
draw byLine(D, C, byred, SOLID_LINE, REGULAR_WIDTH);
draw byLine(C, A, byblack, SOLID_LINE, REGULAR_WIDTH);
byLineDefine(A, D, byyellow, SOLID_LINE, REGULAR_WIDTH);
byLineDefine(D, E, byyellow, DASHED_LINE, REGULAR_WIDTH);
draw byNamedLineSeq(0)(AD, DE);
draw byLine(C, E, byblue, DASHED_LINE, REGULAR_WIDTH);
draw byCircleABC.O(A, B, C, byyellow, 0, 0, 1/2);
draw byLabelsOnCircle(A, B, C, E)(O);
draw byLabelsOnPolygon(B, D, A)(OMIT_FIRST_LABEL+OMIT_LAST_LABEL, 0);
}
\drawCurrentPictureInMargin
\problem{S}{e}{um ângulo de um triângulo for dividido ao meio por uma linha reta, que também corta a base; o retângulo compreendido pelos lados do triângulo é igual ao retângulo compreendido pelos segmentos da base, junto com o quadrado da reta que divide o ângulo ao meio.}

\begin{center}
Trace \drawSizedLine{AD}, fazendo $\drawAngle{CAE} = \drawAngle{EAB}$;\\
então $\drawSizedLine{AB} \times \drawSizedLine{CA} = \drawSizedLine{BD} \times \drawSizedLine{DC} + \drawSizedLine{AD}^2$.

Em torno de \drawLine[bottom]{CA,DC,BD,AB} descreva \drawCircle[middle][1/4]{O} \byref{prop:IV.V},\\
produza \drawSizedLine{AD} para encontrar o círculo e trace \drawSizedLine{CE}.

Como $\drawAngle{CAE} = \drawAngle{EAB}$ \byref{\hypref},\\
e $\drawAngle{B} = \drawAngle{E}$ \byref{prop:III.XXI},\\
$\therefore$ \drawLine[bottom]{AD,BD,AB} e \drawLine[middle]{CA,CE,DE,AD} são \byref{prop:I.XXXII} equiangulares;\\
$\therefore \drawSizedLine{AB} : \drawSizedLine{AD} :: \drawSizedLine{AD,DE} : \drawSizedLine{CA}$ \byref{prop:VI.IV};\\
$\therefore \drawSizedLine{AB} \times \drawSizedLine{CA} = \drawSizedLine{AD} \times \drawSizedLine{AD,DE}$ \byref{prop:VI.XVI}\\
$= \drawSizedLine{DE} \times \drawSizedLine{AD} + \drawSizedLine{AD}^2$ \byref{prop:II.III};\\
mas $\drawSizedLine{DE} \times \drawSizedLine{AD} = \drawSizedLine{BD} \times \drawSizedLine{DC}$ \byref{prop:III.XXXV};\\
$\therefore \drawSizedLine{AB} \times \drawSizedLine{CA} = \drawSizedLine{BD} \times \drawSizedLine{DC} + \drawSizedLine{AD}^2$
\end{center}

\qed
\starttheoremAZ{Prop. C. theor.}\label{prop:VI.C} % Proposition with letters
\defineNewPicture{
pair A, B, C, D, E, O;
numeric r;
r := 9/4u;
O := (0, 0);
A := (dir(60)*r) shifted O;
B := (dir(180-5)*r) shifted O;
C := (dir(5)*r) shifted O;
D = whatever[B, C] = whatever[A, A shifted ((B-C) rotated 90)];
E := (dir(60 + 180)*r) shifted O;
byAngleDefine(A, B, C, byyellow, ARC_SECTOR);
byAngleDefine(A, E, C, byred, ARC_SECTOR);
byAngleDefine(B, D, A, byblue, SOLID_SECTOR);
byAngleDefine(B, C, A, byblue, ARC_SECTOR);
byAngleDefine(E, C, B, byyellow, SOLID_SECTOR);
byAngleDefine(C, A, B, byred, SOLID_SECTOR);
draw byNamedAngleResized();
draw byLine(A, E, byblue, SOLID_LINE, REGULAR_WIDTH);
draw byLine(C, E, byblack, SOLID_LINE, REGULAR_WIDTH);
draw byLine(A, D, byyellow, DASHED_LINE, REGULAR_WIDTH);
draw byLine(A, B, byblue, DASHED_LINE, REGULAR_WIDTH);
byLineDefine(B, D, byred, DASHED_LINE, REGULAR_WIDTH);
byLineDefine(D, C, byred, SOLID_LINE, REGULAR_WIDTH);
draw byNamedLineSeq(0)(BD, DC);
draw byLine(C, A, byyellow, SOLID_LINE, REGULAR_WIDTH);
draw byCircleABC.O(A, B, C, byred, 0, 0, 1/2);
draw byLabelsOnCircle(A, B, C, E)(O);
draw byLabelsOnPolygon(C, D, B)(OMIT_FIRST_LABEL+OMIT_LAST_LABEL, 0);
}
\drawCurrentPictureInMargin
\problem{S}{e}{de qualquer ângulo de um triângulo for traçada uma linha reta perpendicular à base; o retângulo compreendido pelos lados do triângulo são iguais ao retângulo compreendido por a perpendicular e o diâmetro do círculo descrito em torno do triângulo.}

\begin{center}
De \drawAngle{A} de \drawLine[bottom]{CA,DC,BD,AB}\\
draw $\drawUnitLine{AD} \perp \drawUnitLine{BD,DC}$;\\
então deve $\drawUnitLine{AB} \times \drawUnitLine{CA} = \drawUnitLine{AD} \times $ o diâmetro do círculo descrito.

Descreva \drawCircle[middle][1/4]{O} \byref{prop:IV.V}, trace seu diâmetro \drawUnitLine{AE} e trace \drawUnitLine{CE};\\
then $\because \drawAngle{D} = \drawAngle{BCA,ECB}$  \byref{\constref,prop:III.XXXI};\\ %because
and $\drawAngle{B} = \drawAngle{E}$ \byref{prop:III.XXI};\\
$\therefore$ \drawLine[bottom]{AD,BD,AB} é equiangular a \drawLine[middle]{CA,CE,AE} \byref{prop:VI.IV};

$\therefore \drawUnitLine{AB} : \drawUnitLine{AD} :: \drawUnitLine{AE} : \drawUnitLine{CA}$;\\
and $\therefore \drawUnitLine{AB} \times \drawUnitLine{CA} = \drawUnitLine{AD} \times \drawUnitLine{AE}$ \byref{prop:VI.XVI}.
\end{center}

\qed
\starttheoremAZ{Prop. D. theor.}\label{prop:VI.D} % Proposition with letters
\defineNewPicture[1/4]{
pair A, B, C, D, E, O;
numeric r, a[];
a1 := 110;
a2 := 190;
a3 := -50;
a4 := 15;
r := 9/4u;
O := (0, 0);
A := (dir(a1)*r) shifted O;
B := (dir(a2)*r) shifted O;
C := (dir(a3)*r) shifted O;
D := (dir(a4)*r) shifted O;
E = whatever[B, D] = whatever[A, A shifted dir((a1-180)-(a3-(a1-180)))];
byAngleDefine(B, C, A, byyellow, SOLID_SECTOR);
byAngleDefine(A, C, D, byred, ARC_SECTOR);
byAngleDefine(B, D, A, byblue, ARC_SECTOR);
byAngleDefine(A, B, D, byyellow, ARC_SECTOR);
byAngleDefine(E, A, B, byblue, SOLID_SECTOR);
byAngleDefine(C, A, E, byblack, SOLID_SECTOR);
byAngleDefine(D, A, C, byred, SOLID_SECTOR);
draw byNamedAngleResized();
draw byLine(A, E, byyellow, DASHED_LINE, REGULAR_WIDTH);
draw byLine(A, C, byblue, SOLID_LINE, REGULAR_WIDTH);
byLineDefine(B, E, byred, DASHED_LINE, REGULAR_WIDTH);
byLineDefine(E, D, byred, SOLID_LINE, REGULAR_WIDTH);
draw byNamedLineSeq(0)(BE, ED);
draw byLine(A, B, byblack, DASHED_LINE, REGULAR_WIDTH);
draw byLine(B, C, byblue, DASHED_LINE, REGULAR_WIDTH);
draw byLine(C, D, byblack, SOLID_LINE, REGULAR_WIDTH);
draw byLine(D, A, byyellow, SOLID_LINE, REGULAR_WIDTH);
draw byCircleABC.O(A, B, C, byred, 0, 0, 1/2);
draw byLabelsOnCircle(A, B, C, D)(O);
draw byLabelsOnPolygon(D, E, B)(OMIT_FIRST_LABEL+OMIT_LAST_LABEL, 0);
}
\drawCurrentPictureInMargin
\problem[4]{O}{}{retângulo compreendido por as diagonais de uma figura quadrilátera inscrito em um círculo, é igual a ambos os retângulos compreendido por seus lados opostos.}

\begin{center}
Seja \drawLine{DA,CD,BC,AB} qualquer figura quadrilátera inscrito em \drawCircle[middle][1/5]{O};\\
e trace \drawUnitLine{BE,ED} e \drawUnitLine{AC};\\
then $\drawUnitLine{BE,ED} \times \drawUnitLine{AC} = \drawUnitLine{AB} \times \drawUnitLine{CD} + \drawUnitLine{DA} \times \drawUnitLine{BC}$.

Make $\drawAngle{EAB} = \drawAngle{DAC}$ \byref{prop:I.XXIII},\\
$\therefore \drawAngle{EAB,CAE} = \drawAngle{CAE,DAC}$; and $\drawAngle{BCA} = \drawAngle{D}$ \byref{prop:III.XXI};

$\therefore \drawUnitLine{DA} : \drawUnitLine{ED} :: \drawUnitLine{AC} : \drawUnitLine{BC}$ \byref{prop:VI.IV};\\
and $\therefore \drawUnitLine{ED} \times \drawUnitLine{AC} = \drawUnitLine{DA} \times \drawUnitLine{BC}$ \byref{prop:VI.XVI};\\
again, $\because \drawAngle{EAB} = \drawAngle{DAC}$ \byref{\constref},\\ %because
and $\drawAngle{B} = \drawAngle{ACD}$ \byref{prop:III.XXI};\\
$\therefore \drawUnitLine{AB} : \drawUnitLine{BE} :: \drawUnitLine{AC} : \drawUnitLine{CD}$ \byref{prop:VI.IV};\\
and $\therefore \drawUnitLine{BE} \times \drawUnitLine{AC} = \drawUnitLine{AB} \times \drawUnitLine{CD}$ \byref{prop:VI.XVI};\\
but, from above, $\drawUnitLine{ED} \times \drawUnitLine{AC} = \drawUnitLine{DA} \times \drawUnitLine{BC}$;\\
$\therefore \drawUnitLine{BE,ED} \times \drawUnitLine{AC} = \drawUnitLine{AB} \times \drawUnitLine{CD} + \drawUnitLine{DA} \times \drawUnitLine{BC}$ \byref{prop:II.I}.
\end{center}

\qed

\tableofcontents

\end{document}
